\makeatletter
\def\input@path{{../}{./}}
\makeatother
\documentclass[10pt,a4paper,oneside]{ctexbook}
\RequirePackage{xcolor}
\definecolor{main}{RGB}{0,120,60}
\definecolor{light}{RGB}{235,250,240}

\usepackage{hxnotebook}

\begin{document}

% cover
\makecover
  {营养与食品卫生学}
  {华西大纲·PPT·往年资料整理}
  {复习笔记}
  {\today}
  {greedyCat}

% ===== Table of Contents =====
\tableofcontents

% ===== 使用说明 =====
\chapter*{使用说明}
\addcontentsline{toc}{chapter}{使用说明}

\begin{itemize}[leftmargin=*]
    \item 本复习笔记严格依照《营养与食品卫生学》教学大纲章节编排，以授课PPT内容为主要填充，往年笔记为补充参考。
    \item 各章末附有复习题（名词解释、简答题），覆盖该章核心知识点。
    \item \textbf{\color{keyred}红色框}标识的内容为必须重点掌握的核心知识（对应大纲双下划线内容）。
    \item \textbf{\color{main}主题框}为概念释义，帮助理解专业术语。
    \item \textbf{\color{darkgreen}绿色提示框}为要点提示与补充说明。
    \item \textbf{\color{orange}橙色考点框}为常见考点与易考内容。
    \item 大纲中\textbf{句尾"*"标识}为教学难点，笔记中已特别标注。
    \item 笔记末尾附有\textbf{英文缩写与专有名词汇总}，涵盖食品安全相关的国际组织、标准、术语等。
\end{itemize}

\vspace{1cm}
\begin{center}
\begin{tcolorbox}[colback=yellow!10, colframe=yellow!80!black, boxrule=0.8pt, arc=6pt, width=0.85\textwidth, breakable]
\textbf{考核形式提示：}
\begin{itemize}[leftmargin=*]
    \item 平时考核成绩 50\% + 期末理论考核 50\%
    \item 考试范围：以教学大纲和教师课堂讲授内容为主
    \item 大纲双下划线内容 = 掌握内容（本笔记已用红色框标注）
    \item 大纲单下划线内容 = 熟悉内容
    \item 大纲句尾"*" = 教学难点
\end{itemize}
\end{tcolorbox}
\end{center}

\newpage

% =====================================================================
%  导入各章
% =====================================================================

% =====================================================================
%  CHAPTER 0: 绪论 — 参考：0310营养参考_extracted.txt (PPT原文)
% =====================================================================
\chapter{绪论}
\chapterintro{
\keyword{掌握}营养学与食品卫生学的联系与区别；掌握食品、营养、营养素、膳食、食品安全等基本概念的准确定义；掌握六大营养素分类及宏量营养素与微量营养素的划分；掌握必需营养素的定义及条件必需营养素的概念。\keyword{熟悉}近代营养学关注的七大主要问题；熟悉DRIs七项指标体系；熟悉食品安全的两大关键要素及四大属性。\keyword{了解}食药物质（药食同源物质）的概念及纳入要求。
}

% ----- 第一节 营养学与食品卫生学的联系与区别 -----
\section{营养学与食品卫生学的联系与区别}

营养学与食品卫生学的联系在于，从广义上讲，两者有共同的研究对象：\keyword{食物（food）}和\keyword{人体}，即研究食物（饮食）与健康的关系；区别在于，从狭义上讲，两者在具体研究目标、研究目的、研究方法和理论体系等方面各不相同。\imp{营养学是研究食物中的有益成分与健康的关系}，\imp{食品卫生学则是研究食物中有害因素与健康的关系}。

% ----- 第二节 基本概念 -----
\section{基本概念}

\subsection{食品、食物与食药物质}

\begin{defbox}
\textbf{食品：}指各种供人食用或者饮用的成品和原料以及按照传统既是食品又是中药材（药品）的物品，但是不包括以治疗为目的的物品。\imp{一切食物都是食品}。
\end{defbox}

\begin{keybox}
\textbf{食物（Food）的基本概念与功能：}

\textbf{食物：}供人类食用/饮用；含有营养素；不含有害物质；不是治疗用的药物。

食物的基本功能包括：
\begin{enumerate}[leftmargin=*]
    \item \textbf{生理功能：}提供能量和营养素，维持生命活动
    \item \textbf{心理功能：}满足感官享受和心理需求
    \item \textbf{社会功能：}饮食文化、社交媒介
\end{enumerate}
\end{keybox}

\begin{defbox}
\textbf{食药物质（药食同源物质）：}传统作为食品，且列入《中华人民共和国药典》的物质。纳入食药物质目录的物质应当符合下列要求：
\begin{enumerate}[leftmargin=*]
    \item 有传统上作为食品食用的习惯；
    \item 已经列入《中华人民共和国药典》；
    \item 安全性评估未发现食品安全问题；
    \item 符合中药材资源保护、野生动植物保护、生态保护等相关法律法规规定。
\end{enumerate}
\end{defbox}

\subsection{营养与营养素}

\begin{defbox}
\textbf{营养（Nutrition）：}生物从外界摄入食物，在体内经过消化、吸收、代谢以满足其自身生理功能和从事各种活动需要的必要生物学过程。整个过程包括：\imp{消化吸收 $\rightarrow$ 中间代谢 $\rightarrow$ 排泄}。
\end{defbox}

\begin{defbox}
\textbf{膳食（Diet）：}人们日常食用的饮食，由多种食物组成的。
\end{defbox}

\begin{keybox}
\textbf{营养素（Nutrients）六大分类：}

营养素是指维持机体繁殖、生长发育和生存等一切生命活动和过程，需要从外界环境中摄取的物质。来自食物中的营养素种类繁多，人类所需大约40多种，根据其化学性质和生理作用分为六大类：
\begin{enumerate}[leftmargin=*]
    \item \textbf{蛋白质（Protein）}——提供能量，构建机体，调节代谢
    \item \textbf{脂类（Lipids）}——提供能量，构成细胞膜，溶解脂溶性维生素
    \item \textbf{碳水化合物（Carbohydrate）}——最主要的能量来源
    \item \textbf{矿物质（Mineral）}——构成机体组织，调节生理功能
    \item \textbf{维生素（Vitamin）}——调节机体代谢
    \item \textbf{水（Water）}——构成人体组织，参与物质代谢
\end{enumerate}

\imp{根据人体的需要量或体内含量多少，可将营养素分为宏量营养素和微量营养素：}

\begin{itemize}[leftmargin=*]
    \item \textbf{宏量营养素（Macronutrients）：}人体需要量较多的营养素，包括蛋白质、脂类和碳水化合物。这三种营养素可通过体内的氧化过程释放能量，又称为\imp{产能营养素（Energy-yielding Nutrients）}。
    \item \textbf{微量营养素（Micronutrients）：}相对宏量营养素来说，人体需要量较少的营养素，包括矿物质和维生素。根据在体内的含量不同，矿物质又可分为\imp{常量元素（Macroelements，$>0.01\%$体重）}和\imp{微量元素（Microelements，$<0.01\%$体重）}。根据溶解性质不同，维生素又可分为\imp{脂溶性维生素（Lipid-soluble Vitamins）}和\imp{水溶性维生素（Water-soluble Vitamins）}。
\end{itemize}
\end{keybox}

\begin{keybox}
\textbf{必需营养素与条件必需营养素：}

\textbf{必需营养素（Essential Nutrients）：}机体存活、正常生长和功能所必需，但不能由机体自身合成或合成不足，而必须从食物中获得的营养素。与其他膳食成分相比，缺乏必需营养素可造成特异性缺乏病，甚至死亡。

\textbf{条件必需营养素（Conditionally Essential Nutrients）：}特指人体正常状态下不一定需要，但对于体内不能足量合成的人群是必需供给的营养素。补充该类营养素可纠正缺乏导致的异常表现。这一概念用于全肠外营养的患者、生长发育不全、某些病理状态以及遗传缺陷等条件下人体所需的营养素。

\textbf{其他膳食成分：}食物中除必需营养素以外的具有生物活性的物质，如植物化学物（Phytochemicals），对维护健康、预防疾病具有重要意义。
\end{keybox}

% ----- 第三节 近代营养学主要关注问题 -----
\section{近代营养学主要关注问题}

\begin{keybox}
\textbf{近代营养学关注的七大领域：}
\begin{enumerate}[leftmargin=*]
    \item \textbf{膳食营养素供给量和膳食指南：}营养素的生理功能不仅仅局限于预防营养缺乏病的作用。而营养素发挥这些新功能一般都需要比以往的人体需要量或推荐的供给（\keyword{RDA：供给量}，摄入上下限。\keyword{DRI：参考量}）更高的摄入量。为了满足当前消费者预防慢性病和延缓衰老，而增加营养素摄入量的需求，美国学者首先提出：每日参考摄入量（Daily Reference Intake, DRI）的概念，在RDA的基础上，增加了\imp{适宜摄入量（Adequate Intake, AI）}和\imp{摄入量高限（Upper Limit, UL）}。
    \item \textbf{营养与慢性病：}研究膳食营养因素在肥胖、糖尿病、心血管疾病、癌症等慢性病发生发展中的作用及预防策略。
    \item \textbf{营养与基因关系研究：}营养基因组学（Nutrigenomics）探讨营养素对基因表达的调控及基因多态性对营养素需求的影响。
    \item \textbf{食物成分研究：}系统分析食物中已知和未知的营养成分及生物活性成分。
    \item \textbf{营养与抗氧化功能研究：}在正常生理状态下，人体内的\imp{氧自由基（Oxygen Free Radical, OFR）}不断产生又不断被清除，基本处于动态平衡状态。当机体内氧化和抗氧化之间的平衡严重偏向氧化状态时，称为\imp{氧化应激状态（Oxidative Stress）}。此时，过多的活性氧自由基导致细胞损伤。
    \item \textbf{膳食指南（Dietary Guidelines）：}根据营养学原理提出的一组以食物为基础的建议，指导大众合理选择与搭配食物。
    \item \textbf{功能食品（Functional Food）：}除基本营养功能外，具有调节机体功能、预防疾病作用的食品。
\end{enumerate}
\end{keybox}

% ----- 第四节 食品安全 -----
\section{食品安全的概念及与食品卫生的关系}

\begin{defbox}
\textbf{食品安全（Food Safety）：}
\begin{enumerate}[leftmargin=*]
    \item 对食品按其原定用途进行制作、食用时不会使消费者健康受到损害的一种担保（1996年世界卫生组织）
    \item 食物中有毒、有害物质对人体健康影响的公共卫生问题（WHO）
    \item 食品安全，指食品无毒、无害，符合应当有的营养要求，对人体健康不造成任何急性、亚急性或者慢性危害（《中华人民共和国食品安全法》）
\end{enumerate}
\end{defbox}

食品安全问题涉及两个关键要素：\imp{一是有毒有害物质}，\imp{一是对人体健康的不良影响}。两个要素必须同时存在，才构成食品安全问题。

\begin{keybox}
\textbf{Food Safety与Food Security的区别：}
\begin{itemize}[leftmargin=*]
    \item \textbf{Food Safety（食品安全）：}侧重食品质的方面，即食品无毒、无害，不会对消费者健康造成损害。
    \item \textbf{Food Security（食品安全或粮食安全）：}通常是指食品量的安全，即是否有能力得到或者提供足够的食物或者食品。（注：security在餐馆业是指保护消费者和雇员不受暴虐和犯罪危害。）
\end{itemize}
\end{keybox}

\begin{defbox}
\textbf{营养安全（Nutrition Security）：}人类日常饮食中有足够、平衡且含有人体发育必需营养素的供给，以达到完善的食品安全。
\end{defbox}

\begin{tipbox}
\textbf{隐性饥饿（Hidden Hunger）：}机体因营养不平衡或缺乏某些维生素/必需矿物质，同时其他营养成分过度摄入，产生隐藏性营养需求的饥饿症状。全球约20亿人遭受隐性饥饿，中国约3亿人受影响。70\%的慢性病（糖尿病、心血管病、癌症等）与营养元素摄取不均衡有关。
\end{tipbox}

\begin{keybox}
\textbf{食品安全与食品卫生的关系：}
\begin{enumerate}[leftmargin=*]
    \item 概念从属关系：\imp{食品安全是属概念（上位概念）}，\imp{食品卫生是种概念（下位概念）}
    \item 覆盖范围不同：食品安全涵盖种植、养殖、加工、包装、贮藏、运输、销售、消费全环节；食品卫生通常不包含种植、养殖环节
    \item 侧重点不同：食品安全要求结果安全与过程安全完整统一；食品卫生更加侧重过程安全
\end{enumerate}
\end{keybox}

\textbf{食品安全的四大属性（国际共识）：}
\begin{enumerate}[leftmargin=*]
    \item \textbf{综合概念：}包含食品卫生、食品质量、食品营养等，覆盖食品全流程
    \item \textbf{社会概念：}属社会治理概念，不同国家/发展阶段面临的食品安全问题不同
    \item \textbf{政治概念：}食品安全与企业、政府的责任和承诺紧密关联
    \item \textbf{法律概念：}各国从社会系统工程角度推行食品安全综合立法
\end{enumerate}

% ----- 第五节 DRIs概览 -----
\section{膳食营养素参考摄入量（DRIs）概览}

\begin{keybox}
\textbf{中国居民膳食营养素参考摄入量（Dietary Reference Intakes, DRIs）七项指标体系：}

DRIs是一组每日平均膳食营养素摄入量的参考值，用于评价膳食营养素供给量能否满足人体需要及是否存在过量摄入风险。

DRIs包括以下七项指标：
\begin{enumerate}[leftmargin=*]
    \item \textbf{平均需要量（Estimated Average Requirement, EAR）：}满足群体中50\%个体需要量的摄入水平，是制定RNI的基础。
    \item \textbf{推荐摄入量（Recommended Nutrient Intake, RNI）：}满足群体中97\%～98\%个体需要量的摄入水平，是个体每日摄入该营养素的目标值。
    \item \textbf{适宜摄入量（Adequate Intake, AI）：}当个体需要量资料不足、无法计算EAR时，通过观察或实验获得的健康人群某种营养素的摄入量。
    \item \textbf{可耐受最高摄入量（Tolerable Upper Intake Level, UL）：}平均每日可以摄入营养素的最高限量，对几乎所有个体不产生健康危害。
    \item \textbf{宏量营养素可接受范围（Acceptable Macronutrient Distribution Ranges, AMDR）：}蛋白质、脂肪和碳水化合物理想的摄入量范围（以占总能量\%表示）。
    \item \textbf{预防非传染性慢性病的建议摄入量（Proposed Intakes for Preventing Non-communicable Chronic Diseases, PI-NCD）：}以NCD一级预防为目标的必需营养素每日摄入量。
    \item \textbf{特定建议值（Specific Proposed Levels, SPL）：}以降低成年人膳食相关NCD风险为目标的其他膳食成分（如植物化学物）的每日摄入量。
\end{enumerate}

注：DRIs每4～5年修订一次。具体指标体系和应用将在第五章"公共营养"中详述。
\end{keybox}

% ----- 第六节 发展历史 -----
\section{营养与食品卫生学的发展历史}

\subsection{营养学发展简史}
\begin{enumerate}[leftmargin=*]
    \item 古代：经验性营养知识积累（如《黄帝内经》中的食养理论）
    \item 18-19世纪：拉瓦锡提出呼吸是氧化过程，能量代谢研究起步
    \item 20世纪：维生素的陆续发现、必需氨基酸概念的提出；中国营养学会成立（1945年）
    \item 21世纪：分子营养学、营养基因组学等新兴领域发展
\end{enumerate}

\subsection{食品卫生学发展历史}
\begin{enumerate}[leftmargin=*]
    \item 周朝（约3000年前）：设置"凌人"一职，专职负责食品冷藏、防腐
    \item 唐代《唐律》：变质肉食致人患病，剩余肉食必须销毁，违者杖九十
    \item 1965年：颁布《食品卫生管理试行条例》
    \item 1979年：颁布《中华人民共和国食品卫生管理条例》，正式步入法制化阶段
    \item 2015年：修订《食品安全法》正式实施
    \item 2018年：完成《食品安全法》修正，组建国家市场监督管理总局
\end{enumerate}

% ----- 第七节 研究内容与方法 -----
\section{研究内容与方法}

\textbf{研究内容：}
\begin{enumerate}[leftmargin=*]
    \item \textbf{营养学部分：}食物营养素和生物活性成分、各类食物营养、特殊人群营养、公共营养、营养与疾病
    \item \textbf{食品卫生学部分：}食品污染及预防、食品添加剂、各类食品卫生、食源性疾病、食品安全风险分析、食品安全监督管理
\end{enumerate}

\begin{tipbox}
\textbf{公众认知 vs 科学标准的食品危害性排序：}

公众排序：食品添加剂 > 环境污染 > 营养素搭配不良 > 天然毒素 > 微生物污染

科学排序：营养素搭配不良（膳食行为）> 微生物污染 > 天然毒素 > 环境污染 > 食品添加剂
\end{tipbox}

% ----- 复习题 -----
\reviewcontent{
\section{复习题}

\subsection*{一、名词解释}
\begin{enumerate}[leftmargin=*, itemsep=8pt]
    \item \textbf{营养（Nutrition）}
    \item \textbf{营养素（Nutrients）}
    \item \textbf{食品卫生学（Food Hygiene）}
    \item \textbf{食品安全（Food Safety）}
    \item \textbf{粮食安全（Food Security）}
    \item \textbf{隐性饥饿（Hidden Hunger）}
    \item \textbf{必需营养素（Essential Nutrients）}
    \item \textbf{条件必需营养素（Conditionally Essential Nutrients）}
    \item \textbf{膳食（Diet）}
    \item \textbf{食药物质（药食同源物质）}
\end{enumerate}

\subsection*{二、简答题}
\begin{enumerate}[leftmargin=*, itemsep=8pt]
    \item \textbf{简述营养学与食品卫生学的联系与区别。}
    \item \textbf{简述营养素的六大分类，并说明宏量营养素与微量营养素的划分。}
    \item \textbf{简述必需营养素与条件必需营养素的定义及区别。}
    \item \textbf{近代营养学主要关注哪些问题？（答出七个方面）}
    \item \textbf{简述食品安全与食品卫生的关系。}
    \item \textbf{简述食品安全的四大属性。}
    \item \textbf{简述DRIs的七项指标体系。}
\end{enumerate}

\subsection*{参考答案要点}

\begin{enumerate}[leftmargin=*, itemsep=5pt]
    \item 见本章各概念定义框。
    \item 联系：共同研究对象——食物和人体，即研究食物与健康的关系。区别：营养学研究食物中的有益成分与健康的关系，食品卫生学研究食物中有害因素与健康的关系。
    \item 六大类：蛋白质（Protein）、脂类（Lipids）、碳水化合物（Carbohydrate）、矿物质（Mineral）、维生素（Vitamin）、水（Water）。宏量营养素（需量多）：蛋白质、脂类、碳水化合物，三者可通过氧化释放能量，又称产能营养素。微量营养素（需量少）：矿物质（又分常量元素$>0.01\%$体重和微量元素$<0.01\%$体重）和维生素（又分脂溶性和水溶性）。
    \item 必需营养素：机体存活、正常生长和功能所必需，但不能由机体自身合成或合成不足，必须从食物中获得的营养素，缺乏可造成特异性缺乏病甚至死亡。条件必需营养素：正常状态下不一定需要，但体内不能足量合成的人群必需供给，适用于全肠外营养、生长发育不全、某些病理状态及遗传缺陷等条件。
    \item ①膳食营养素供给量和膳食指南（从RDA到DRI的转变，增加了AI和UL）；②营养与慢性病；③营养与基因关系研究；④食物成分研究；⑤营养与抗氧化功能研究（氧化应激）；⑥膳食指南（Dietary Guidelines）；⑦功能食品。
    \item 食品安全是属概念（上位概念），食品卫生是种概念（下位概念）。食品安全涵盖种植、养殖、加工、包装、贮藏、运输、销售、消费全环节，要求结果安全与过程安全完整统一；食品卫生通常不包含种植、养殖环节，侧重过程安全。
    \item 综合概念、社会概念、政治概念、法律概念。
    \item 平均需要量（EAR）、推荐摄入量（RNI）、适宜摄入量（AI）、可耐受最高摄入量（UL）、宏量营养素可接受范围（AMDR）、预防非传染性慢性病的建议摄入量（PI-NCD）、特定建议值（SPL）。
\end{enumerate}

\newpage
}

% =====================================================================
%  CHAPTER 1A: 蛋白质 — 参考：0310营养参考_extracted.txt (PPT原文)
% =====================================================================
\chapter{蛋白质}
\chapterintro{掌握蛋白质的生理作用与氮平衡（零/负/正）；掌握必需氨基酸（9种）、条件必需氨基酸和限制氨基酸；掌握氨基酸模式与蛋白质互补三原则；掌握蛋白质消化率、BV、NPU、PER、AAS/PDCAAS（难点*）；掌握人体蛋白质营养状况评价指标。}

\section{蛋白质的生理作用与氮平衡}

\begin{keybox}
\textbf{蛋白质的生理作用：}
\begin{enumerate}[leftmargin=*]
    \item 人体内细胞的主要成分：构成人体各组织器官（皮肤、肌肉、骨骼、内脏等）；补偿新陈代谢的消耗；修补组织损失
    \item 人体内重要物质的主要原料：激素、酶、抗体、血红蛋白
    \item 提供热量：每克蛋白质在体内可产生16.72千焦耳的能量，由蛋白质提供的能量占每日人需总能量的\imp{10\%—15\%}
\end{enumerate}

\textbf{氮平衡（Nitrogen Balance）：}反映机体摄入氮和排出氮的关系。
$$B = I - (U + F + S)$$
$B$：氮平衡；$I$：摄入氮；$U$：尿氮；$F$：粪氮；$S$：皮肤等氮损失
\begin{itemize}[leftmargin=*]
    \item \imp{零氮平衡}：正常成人
    \item \imp{负氮平衡}：老年人、饥饿过度节食、消耗性疾病
    \item \imp{正氮平衡}：生长发育期、孕期、恢复期、肌肉增长
\end{itemize}
\end{keybox}

\section{必需氨基酸与蛋白质分类}

\begin{keybox}
\textbf{必需氨基酸（Essential Amino Acid, EAA）：}人体不能合成或合成速度不能满足机体需要，必须从食物中获得的一类氨基酸。共有\imp{9种}：亮氨酸、异亮氨酸、赖氨酸、蛋氨酸、苯丙氨酸、苏氨酸、色氨酸、缬氨酸、组氨酸。其中\imp{组氨酸是婴儿的必需氨基酸}。

\textbf{条件必需氨基酸（Conditionally Essential AA）/ 半必需氨基酸：}半胱氨酸和酪氨酸在体内分别由蛋氨酸和苯丙氨酸转变而成。如果膳食中能直接提供这两种氨基酸，则人体对蛋氨酸和苯丙氨酸的需要可\imp{分别减少30\%和50\%}。代谢路径：蛋氨酸→同型半胱氨酸→胱硫醚→半胱氨酸。

\textbf{氨基酸模式（Pattern）：}蛋白质中各种必需氨基酸的构成比例。计算方法是将该种蛋白质中的色氨酸含量定为1，分别计算出其他必需氨基酸的相应比值。

\textbf{蛋白质分类：}完全蛋白质（优质蛋白，植物仅大豆）；半完全蛋白质；不完全蛋白质（胶原蛋白）。

\textbf{限制氨基酸（Limiting Amino Acid）：}当食物蛋白质中一种或几种必需氨基酸相对含量较低或缺乏，限制了食物蛋白质中的其它必需氨基酸被机体利用的程度，使其营养价值降低。含量最低的称为\imp{第一限制氨基酸}。\textbf{谷类缺赖氨酸，豆类缺蛋氨酸。}

\textbf{蛋白质互补三原则：}(1)食物生物学种属愈远愈好；(2)搭配种类愈多愈好；(3)食用时间愈近愈好，最好同时食用。
\end{keybox}

\section{蛋白质营养学评价（难点*）}

\begin{keybox}
\textbf{1. 蛋白质含量：}各种蛋白质含氮量稳定，约占蛋白质重量的\imp{16\%}。\imp{蛋白质含量 = 氮含量 × 6.25}。检测方法：凯氏定氮法（Kjedahl法）、双缩脲法、Folin-酚试剂法（Lowry法）、紫外吸收法、考马斯亮蓝法（Bradford法）。

\textbf{2. 蛋白质消化率：}消化率(\%) = 氮吸收量/摄入氮量 × 100
\begin{itemize}[leftmargin=*]
    \item \imp{真消化率(\%) = [食物氮 $-$ (粪氮 $-$ 粪代谢氮)] / 食物氮 × 100}
    \item \imp{表观消化率(\%) = (食物氮 $-$ 粪氮) / 食物氮 × 100}
\end{itemize}

\textbf{3. 生物价（BV）：}反映食物蛋白质消化吸收后被机体利用程度的指标。BV越高，利用程度越高。
$$\text{BV} = \frac{\text{储留氮}}{\text{吸收氮}} \times 100,\quad \text{储留氮} = \text{吸收氮} - (\text{尿氮} - \text{尿内源性氮})$$

\textbf{4. 蛋白质净利用率（NPU）：}
$$\text{NPU}(\%) = \text{消化率} \times \text{生物价} = \frac{\text{储留氮量}}{\text{食物氮}} \times 100$$

\textbf{5. 蛋白质功效比值（PER）：}测定生长发育中的幼小动物摄入1g蛋白质所增加的体重克数。一般选用刚断乳的雄性大鼠，被测蛋白质作为唯一蛋白质来源占饲料10\%，喂饲28天。PER = 动物体重增加(g)/摄入蛋白质(g)。校正PER = 待测PER/酪蛋白PER × 2.5。

\textbf{6. 氨基酸评分（AAS）/ 蛋白质化学评分：}将待评食物蛋白质必需氨基酸与人体必需氨基酸模式比较。
$$\text{AAS} = \frac{\text{待测蛋白质每克氮中某种氨基酸含量(mg)}}{\text{理想模式每克氮中该氨基酸含量(mg)}} \times 100$$
$$\text{经消化率修正的AAS (PDCAAS)} = \text{AAS} \times \text{真消化率}$$
\end{keybox}

\section{人体蛋白质营养状况评价}

\begin{defbox}
\textbf{临床检查：}身高、体重、发育；上臂围；上臂肌围、上臂肌区。

\textbf{实验室检查：}
\begin{itemize}[leftmargin=*]
    \item 血清白蛋白（正常值 \imp{35$\sim$50g/L}）
    \item 血清运铁蛋白（正常值 \imp{2.2$\sim$4.0g/L}）
    \item 血清甲状腺结合前蛋白（正常值 \imp{280$\sim$350mg/L}）
    \item 视黄醇结合蛋白（正常值 \imp{26$\sim$76mg/L}）
    \item 尿3-甲基组氨酸：营养不良时排出量减少
    \item 头发的毛干与毛根的形态改变
\end{itemize}

\textbf{蛋白质-能量营养不良（PEM）：}水肿型（Kwashiorkor）/ 干瘦型（Marasmus）。
\end{defbox}

\section{复习题}
\begin{enumerate}[leftmargin=*, itemsep=8pt]
    \item \textbf{【名词解释】}必需氨基酸；条件必需氨基酸；限制氨基酸；氨基酸模式；氮平衡
    \item \textbf{【名词解释】}生物价(BV)；NPU；PER；AAS；PDCAAS
    \item \textbf{【简答题】}简述氮平衡的三种类型并举例。
    \item \textbf{【简答题】}试述蛋白质互补三原则并举谷豆互补例。
    \item \textbf{【简答题】}简述蛋白质营养价值评价的六项指标及公式。
\end{enumerate}

\subsection*{参考答案要点}
\begin{enumerate}[leftmargin=*, itemsep=5pt]
    \item 见正文。EAA 9种（亮/异亮/赖/蛋/苯丙/苏/色/缬/组，组氨酸为婴儿必需）。
    \item 见正文公式。BV=储留氮/吸收氮×100；NPU=消化率×BV=储留氮/食物氮；PER=动物增重/摄入蛋白；AAS=待测AA/理想模式AA×100；PDCAAS=AAS×真消化率。
    \item B=I-(U+F+S)。零（正常成人）、负（老年/饥饿/消耗病）、正（发育/孕期/恢复/增肌）。
    \item 三原则：种属愈远愈好、种类愈多愈好、时间愈近愈好。谷类缺赖氨酸富蛋氨酸，豆类缺蛋氨酸富赖氨酸，同时食用互补。
    \item 含量（氮×6.25）、真消化率/表观消化率、BV、NPU、PER（雄性大鼠28d）、AAS/PDCAAS。
\end{enumerate}
\newpage

% =====================================================================
%  CHAPTER 1B: 脂类 — 参考：0613笔记 (PRIMARY SOURCE)
% =====================================================================
\chapter{脂类}
\chapterintro{掌握脂类的分类和生理功能；掌握脂肪酸分类（按碳链长度/饱和程度/空间结构）、脂肪酸命名Cx:y及必需脂肪酸（亚油酸/$\alpha$-亚麻酸）；熟悉食物脂类营养价值评价（消化率/EFA含量/脂肪酸比例/脂溶性维生素含量）。}

\section{脂类的分类和生理功能}

\begin{keybox}
\textbf{脂类（Lipids）的分类：}脂类 = 脂肪（Fat）和类脂（Lipoids）。脂肪 = 甘油三酯（Triglycerides），体内重要的储能和供能物质，\imp{占体内脂类总量的95\%}。类脂：主要包括\imp{磷脂}和\imp{固醇类}，是细胞膜、组织器官，尤其是神经组织的重要组成。

\textbf{脂类的生理功能：}
\begin{enumerate}[leftmargin=*]
    \item \imp{储存和提供能量：}1g脂肪 = 9kcal能量。脂肪细胞可无限地储存脂肪。脂肪不能给脑和神经细胞以及血细胞供能。
    \item \imp{保温及润滑。}
    \item \imp{节约蛋白质作用：}保护体内蛋白质不被用作能源物质。
    \item \imp{机体构成成分：}细胞膜。
    \item \imp{脂肪组织内分泌功能：}雌激素等。
    \item \imp{食物中脂肪的作用：}增加饱腹感、改善食物感官性状、提供脂溶性维生素。
\end{enumerate}
\end{keybox}

\section{脂肪酸的分类}

\begin{keybox}
\textbf{脂肪酸命名：}表示为Cx:y，其中x为碳原子个数、y表示不饱和双键个数。

\textbf{1. 按碳链长度分类：}
\begin{itemize}[leftmargin=*]
    \item \imp{长链脂肪酸：}含14\~{}24碳
    \item \imp{中链脂肪酸（MCFA）：}含8\~{}12个碳。可直接与甘油三酯酯化；水溶性好，能直接被小肠吸收；无需形成乳糜颗粒，直接进入肝脏；细胞内快速氧化供能，但不能用于糖尿病、酮中毒及酸中毒、肝硬化者
    \item \imp{短链脂肪酸（SCFA）：}含6碳以下（醋酸、丙酸、丁酸）。来源：膳食纤维、抗性淀粉、低聚糖、糖醇等，结肠被肠道微生物发酵。功能：供能、促进细胞膜脂类物质合成、可能预防和治疗溃疡性结肠炎、可预防结肠肿瘤、对内源性胆固醇的合成有抑制作用
\end{itemize}
食物中主要以18碳脂肪酸为主，人体只能利用偶数个数碳原子的脂肪酸。

\textbf{2. 按饱和程度分类：}
\begin{itemize}[leftmargin=*]
    \item \imp{饱和脂肪酸（SFA）：}碳链中没有不饱和双键。HDL$\downarrow$，LDL$\uparrow$
    \item \imp{不饱和脂肪酸（USFA）：}碳链中含一个或多个不饱和双键
    \item \imp{单不饱和脂肪酸（MUFA）：}含一个不饱和双键，油酸。HDL$\uparrow$、LDL$\downarrow$
    \item \imp{多不饱和脂肪酸（PUFA）：}含两个及以上不饱和双键，膳食中主要为亚麻酸、$\alpha$-亚麻酸，存在于植物油中。HDL$\downarrow$、LDL$\downarrow\downarrow$
\end{itemize}

\textbf{3. 按空间结构分类：}
\begin{itemize}[leftmargin=*]
    \item \imp{顺式脂肪酸：}大多不饱和脂肪酸
    \item \imp{反式脂肪酸：}天然的不存在，主要存在于牛奶、人造奶油中。HDL$\downarrow$、LDL$\uparrow$ $\rightarrow$ 增加心血管疾病的危险。可能诱发肿瘤、2型糖尿病。人造奶油、蛋糕、饼干、油炸食品、乳酪产品、花生酱等是反式脂肪酸的主要来源。
\end{itemize}
\end{keybox}

\section{必需脂肪酸}

\begin{keybox}
\textbf{必需脂肪酸（Essential Fatty Acid, EFA）：}人体不可缺少且自身不能合成，必须通过食物供给的脂肪酸，有\imp{亚油酸}（n-6系多不饱和）和\imp{$\alpha$-亚麻酸}（n-3系多不饱和）。

\textbf{EFA的功能：}
\begin{itemize}[leftmargin=*]
    \item 构成磷脂的组成成分 $\rightarrow$ 构成细胞膜
    \item 前列腺素PG合成的前体 $\rightarrow$ 血管舒缩、神经传导等
    \item 参与胆固醇代谢
    \item 参与类十二烷酸合成（PG、TXA等），调节血压、血脂、血栓形成及免疫反应等
\end{itemize}
\end{keybox}

\section{食物脂类营养价值评价}

\begin{defbox}
\textbf{食物脂类营养价值评价：}
\begin{enumerate}[leftmargin=*]
    \item \imp{脂肪的消化率：}含不饱和脂肪酸、短链脂肪酸越多的脂肪 $\rightarrow$ 熔点越低 $\rightarrow$ 越易消化。一般植物脂肪消化率 > 动物脂肪
    \item \imp{必需脂肪酸（EFA）的含量：}一般植物油中亚油酸和$\alpha$-亚麻酸的含量 > 动物脂肪（椰子油除外）
    \item \imp{各种脂肪酸的比例：}一般推荐SFA:MUFA:PUFA = 1:1:1。EFA摄入不少于总能量的3\%，n-3系:n-6系 = 1:4\~{}6
    \item \imp{脂溶性维生素含量：}脂溶性维生素含量高的脂类 $\rightarrow$ 营养价值高。植物油（特别谷类种子的胚油）VE含量丰富。器官脂肪（肝脏脂肪）含丰富VA、VD，海产鱼类肝脏脂肪中VA、VD更高
\end{enumerate}

\textbf{参考摄入量：}成人脂肪摄入量应占总能量的\imp{20\%\~{}30\%}。EFA摄入不少于总能量的3\%。
\end{defbox}

\section{复习题}
\begin{enumerate}[leftmargin=*, itemsep=8pt]
    \item \textbf{【名词解释】}脂类（Lipids）；甘油三酯（Triglycerides）；必需脂肪酸（EFA）；脂肪酸命名Cx:y
    \item \textbf{【名词解释】}中链脂肪酸（MCFA）；短链脂肪酸（SCFA）；反式脂肪酸
    \item \textbf{【简答题】}简述脂类的分类和生理功能。
    \item \textbf{【简答题】}按碳链长度、饱和程度、空间结构分别简述脂肪酸的分类及特点。
    \item \textbf{【简答题】}简述必需脂肪酸（EFA）的定义、种类和功能。
    \item \textbf{【简答题】}简述食物脂类营养价值评价的四个方面及成人脂肪参考摄入量。
\end{enumerate}

\subsection*{参考答案要点}
\begin{enumerate}[leftmargin=*, itemsep=5pt]
    \item \textbf{脂类（Lipids）：}脂类 = 脂肪（Fat）和类脂（Lipoids）。脂肪 = 甘油三酯（Triglycerides），占体内脂类总量的95\%，是体内重要的储能和供能物质，1g脂肪 = 9kcal能量。类脂主要包括磷脂和固醇类，是细胞膜、组织器官，尤其是神经组织的重要组成。\textbf{EFA}：人体不可缺少且自身不能合成，必须通过食物供给的脂肪酸，有亚油酸（n-6系）和$\alpha$-亚麻酸（n-3系）。\textbf{Cx:y}：x为碳原子个数，y表示不饱和双键个数。
    \item \textbf{MCFA}：含8\~{}12个碳，可直接与甘油三酯酯化，水溶性好，能直接被小肠吸收，无需形成乳糜颗粒直接进入肝脏，细胞内快速氧化供能，但不能用于糖尿病、酮中毒及酸中毒、肝硬化者。\textbf{SCFA}：含6碳以下（醋酸、丙酸、丁酸），来源于膳食纤维等在结肠被肠道微生物发酵，功能包括供能、促进细胞膜脂类物质合成、可能预防和治疗溃疡性结肠炎、可预防结肠肿瘤、对内源性胆固醇的合成有抑制作用。\textbf{反式脂肪酸}：天然的不存在，主要存在于牛奶、人造奶油中，HDL$\downarrow$/LDL$\uparrow$，增加心血管疾病的危险，可能诱发肿瘤、2型糖尿病。人造奶油、蛋糕、饼干、油炸食品、乳酪产品、花生酱等是主要来源。
    \item 分类：脂类 = 脂肪（甘油三酯，占体内脂类总量的95\%） + 类脂（磷脂和固醇类，细胞膜和神经组织的重要组成）。生理功能6项：(1)储存和提供能量（1g脂肪=9kcal，脂肪细胞可无限储存脂肪，脂肪不能给脑和神经细胞以及血细胞供能）；(2)保温及润滑；(3)节约蛋白质作用；(4)机体构成成分（细胞膜）；(5)脂肪组织内分泌功能（雌激素等）；(6)食物中脂肪的作用（增加饱腹感、改善食物感官性状、提供脂溶性维生素）。
    \item \textbf{按碳链长度：}长链（14\~{}24C）、中链MCFA（8\~{}12C，水溶性好，直接吸收进入肝脏，快速氧化供能，但糖尿病/酮中毒/酸中毒/肝硬化者不能用）、短链SCFA（$\leq$6C，来自膳食纤维发酵，供能/预防结肠肿瘤/抑制内源性胆固醇合成）。\textbf{按饱和程度：}SFA（无双键，HDL$\downarrow$/LDL$\uparrow$）、MUFA（1个双键，HDL$\uparrow$/LDL$\downarrow$）、PUFA（$\geq$2个双键，HDL$\downarrow$/LDL$\downarrow\downarrow$）。\textbf{按空间结构：}顺式脂肪酸（大多不饱和脂肪酸）、反式脂肪酸（天然不存在，HDL$\downarrow$/LDL$\uparrow$，增加心血管疾病危险）。
    \item EFA定义：人体不可缺少且自身不能合成，必须通过食物供给的脂肪酸。种类：亚油酸（n-6系多不饱和）和$\alpha$-亚麻酸（n-3系多不饱和）。功能：构成磷脂的组成成分$\rightarrow$构成细胞膜；前列腺素PG合成的前体$\rightarrow$血管舒缩、神经传导等；参与胆固醇代谢；参与类十二烷酸合成（PG、TXA等），调节血压、血脂、血栓形成及免疫反应等。
    \item 四方面：(1)脂肪的消化率（含不饱和脂肪酸/短链脂肪酸越多$\rightarrow$熔点越低$\rightarrow$越易消化，植物脂肪消化率 > 动物脂肪）；(2)EFA含量（植物油 > 动物脂肪，椰子油除外）；(3)各种脂肪酸比例（SFA:MUFA:PUFA = 1:1:1，EFA不少于总能量3\%，n-3:n-6 = 1:4\~{}6）；(4)脂溶性维生素含量（植物油胚油VE丰富，器官脂肪和海产鱼类肝脏脂肪含丰富VA/VD）。成人脂肪摄入量应占总能量的20\%\~{}30\%，EFA摄入不少于总能量的3\%。
\end{enumerate}
\newpage

% =====================================================================
%  CHAPTER 1C: 碳水化合物 — 参考：0310营养参考_extracted.txt (PPT原文)
% =====================================================================
\chapter{碳水化合物}
\chapterintro{掌握碳水化合物分类（单/双/寡/多糖）和膳食纤维（可溶/不可溶）的定义与功能；掌握碳水化合物的五大生理功能（供能/抗生酮/节约蛋白质/构成组织/解毒）；掌握血糖指数（GI）和血糖负荷（GL）的定义、计算和分类；熟悉参考摄入量与食物来源。}

\section{碳水化合物的分类与膳食纤维}

\begin{keybox}
\textbf{碳水化合物的分类：}
\begin{itemize}[leftmargin=*]
    \item \imp{单糖（Monosaccharide）}：葡萄糖（glucose）、果糖（fructose）、半乳糖（galactose）、糖醇类
    \item \imp{双糖（Disaccharide）}：蔗糖（sucrose）、麦芽糖（maltose）、乳糖（lactose）、海藻糖（trehalose）
    \item \imp{寡糖（Oligosaccharide）}：棉子糖（raffinose）、水苏糖（stachyose）
    \item \imp{多糖（Polysaccharide）}：可消化（淀粉starch/糊精dextrins/糖原glycogen）和不可消化（抗性淀粉/膳食纤维）
\end{itemize}

\textbf{膳食纤维（Dietary Fiber）：}食物中不能被消化利用的多糖和木质素。分类（按水溶性）：
\begin{itemize}[leftmargin=*]
    \item \imp{不溶性纤维}：纤维素、半纤维素、木质素
    \item \imp{可溶性纤维}：果胶、树胶、一部分半纤维素
\end{itemize}

\textbf{膳食纤维的生理功能：}减少便秘/预防大肠疾病；预防癌症；预防心血管疾病和胆石症；防治糖尿病；防治肥胖。
\end{keybox}

\section{碳水化合物的生理功能}

\begin{defbox}
\textbf{1. 提供能量}

\textbf{2. 构成身体组织}

\textbf{3. 辅助脂肪氧化——抗生酮作用：}脂肪酸在体内分解代谢时产生的乙酰基需与碳水化合物代谢产生的草酰乙酸结合才能进入三羧酸循环而最终被彻底氧化。当碳水化合物不足时，因草酰乙酸不足使得脂肪酸不能被彻底氧化分解而产生过多酮体，会发生酮症酸中毒。反之，当碳水化合物充足时可防止酮症酸中毒的发生，这种作用称为\imp{抗生酮作用}。

\textbf{4. 节约蛋白质作用：}如碳水化合物充足，人体首先利用碳水化合物供给能量，从而减少了蛋白质单纯作为供给能量的分解代谢，有利于改善和提高蛋白质的利用，并增强人体内氮的贮存量。这就是所谓的\imp{节约蛋白质作用}。

\textbf{5. 帮助肝脏解毒}
\end{defbox}

\section{血糖指数（GI）与血糖负荷（GL）}

\begin{keybox}
\textbf{血糖指数（Glycemic Index, GI）：}含50g碳水化合物的食品引起的血糖反应曲线下的面积与含等量碳水化合物的标准食品血糖反应之比。以葡萄糖或白面包为标准。

$$\text{GI} = \frac{\text{摄入含50g CHO食物后2h血糖曲线下面积}}{\text{摄入50g葡萄糖(或白面包)后2h血糖曲线下面积}}$$
\imp{低GI食物：GI $\leq$ 55}；中GI：55 $<$ GI $\leq$ 70；\imp{高GI食物：GI $>$ 70}

\textbf{血糖负荷（Glycemic Load, GL）} = GI × 可利用碳水化合物（CHO $-$ 膳食纤维）。GL同时考虑了食物的升糖速度（GI）和实际摄入的可利用碳水化合物量，能更准确反映一餐/一份食物对血糖的真实影响。
\end{keybox}

\section{参考摄入量与食物来源}

\begin{defbox}
成人CHO：EAR为120g/d，\imp{AMDR占总能量的50\%$\sim$65\%}。添加糖限制：建议不超过50g/d，最好低于25g/d。酒精摄入量不超过15g/d，任何形式的酒精对人体健康都无益处。

食物来源：粮谷类及制品（核心主食来源）、块根(茎)类蔬菜、杂豆类及制品、水果及制品、纯糖类及其制品。
\end{defbox}

\section{复习题}
\begin{enumerate}[leftmargin=*, itemsep=8pt]
    \item \textbf{【名词解释】}膳食纤维；血糖指数（GI）；血糖负荷（GL）；抗生酮作用；节约蛋白质作用
    \item \textbf{【简答题】}简述碳水化合物的分类及膳食纤维的生理功能。
    \item \textbf{【简答题】}试述碳水化合物的抗生酮作用和节约蛋白质作用的机制。
    \item \textbf{【简答题】}简述GI的定义、分类标准和GL的计算公式及意义。
\end{enumerate}

\subsection*{参考答案要点}
\begin{enumerate}[leftmargin=*, itemsep=5pt]
    \item 见正文。\textbf{膳食纤维}：不能被消化利用的多糖和木质素，分不溶性（纤维素/半纤维素/木质素）和可溶性（果胶/树胶）。\textbf{GI}：含50g CHO食物血糖曲线面积/标准食品血糖曲线面积。\textbf{GL}=GI×可利用CHO。
    \item 单糖/双糖/寡糖/多糖（可消化：淀粉/糊精/糖原；不可消化：抗性淀粉/膳食纤维）。膳食纤维功能：防便秘/防癌/防心血管病/防糖尿病/防肥胖。
    \item 抗生酮：脂肪酸分解需草酰乙酸→三羧酸循环，CHO不足→草酰乙酸不足→酮体过多→酮症酸中毒。节约蛋白质：CHO充足→优先供能→减少蛋白质分解供能→改善氮平衡。
    \item GI定义见第1题。低GI≤55，中55-70，高>70。GL=GI×可利用CHO(CHO-膳食纤维)，兼顾升糖速度与实际CHO摄入量。
\end{enumerate}
\newpage

% =====================================================================
%  CHAPTER 1D: 营养学基础 - 能量
% =====================================================================
\chapter{营养学基础（四）：能量}
\setcounter{chapter}{1}
\chapterintro{
掌握\keyword{能量单位}、\keyword{能量系数}（产能营养素热价）、\keyword{基础代谢（BEE）}的定义及影响因素、\keyword{食物热效应（TEF/SDA）*}的概念及差异、人体能量消耗的构成及\keyword{三大产能营养素的供能比}（RNI）。熟悉\keyword{体力活动水平（PAL）}的分级。了解\keyword{能量需要量测定方法*}（直接测热法、间接测热法、双标水法）和\keyword{呼吸商（RQ）}的概念。\\[4pt]
{\color{keyred}\textbf{教学难点*：}}食物热效应（TEF/SDA）、能量需要量测定方法。
}

% ----- 第一节 能量概述与能量单位 -----
\section{能量概述与能量单位}

\begin{defbox}
\textbf{能量（Energy）}是维持生命活动的基础。人体所需能量来源于食物中的\keyword{碳水化合物}、\keyword{脂肪}和\keyword{蛋白质}，三者统称为产能营养素。能量代谢遵循热力学第一定律——能量守恒定律。
\end{defbox}

\begin{keybox}
\textbf{能量单位：}
\begin{itemize}[leftmargin=*]
    \item 焦耳（Joule, J）：国际单位制（SI）的能量单位
    \item 千焦（kJ）：$1\ \text{kJ} = 10^3\ \text{J}$
    \item 兆焦（MJ）：$1\ \text{MJ} = 10^6\ \text{J}$
    \item 卡（calorie, cal）：非国际单位制，但营养学中仍广泛使用
    \item 千卡（kcal）：$1\ \text{kcal} = 10^3\ \text{cal}$
\end{itemize}

\textbf{换算关系：}
\begin{center}
\imp{$1\ \text{kJ} = 0.239\ \text{kcal}$} \hspace{2cm} \imp{$1\ \text{kcal} = 4.184\ \text{kJ}$}
\end{center}
\end{keybox}

% ----- 第二节 能量系数（产能营养素的生理有效热价）-----
\section{能量系数}

\begin{keybox}
\textbf{能量系数（生理有效热价）：}每克产能营养素在体内氧化分解后为机体供给的净能量。

\begin{table}[H]
\centering
\caption{三大产能营养素的能量系数}
\begin{tabular}{cccc}
\toprule
\textbf{产能营养素} & \textbf{kcal/g} & \textbf{kJ/g} & \textbf{说明} \\
\midrule
碳水化合物 & 4 & 16.81 & 供能主力（55\%～65\%） \\
脂肪         & 9 & 37.56 & 能量密度最高 \\
蛋白质       & 4 & 16.74 & 含氮，部分热能由尿素携带损失 \\
\bottomrule
\end{tabular}
\end{table}
\end{keybox}

\begin{tipbox}
\textbf{注意区分物理热价与生理有效热价：}
\begin{itemize}[leftmargin=*]
    \item 物理热价（弹式测热计测定）：蛋白质约5.65 kcal/g
    \item 生理有效热价（体内实际可利用）：蛋白质约4 kcal/g
    \item 差异原因：蛋白质在体内氧化不彻底，尿素、肌酐等代谢产物仍含有能量随尿排出
\end{itemize}
脂肪和碳水化合物的物理热价等于其生理有效热价；仅蛋白质两者不同。
\end{tipbox}

% ----- 第三节 人体能量消耗的构成 -----
\section{人体能量消耗的构成}

人体每日总能量消耗（Total Energy Expenditure, TEE）由三部分组成：\keyword{基础代谢}、\keyword{体力活动}和\keyword{食物热效应}。儿童、孕妇等特殊人群还包括生长发育消耗和组织储存的能量。

\subsection{基础代谢（Basal Energy Expenditure, BEE）}

\begin{keybox}
\textbf{基础代谢（BEE，教学难点*）：}

\keyword{基础代谢}是指维持人体基本生命活动的能量消耗，即维持体温、呼吸、循环、排泄、腺体分泌、肌肉张力和细胞基本功能等所需的最低能量。

\textbf{测定条件：}
\begin{enumerate}[leftmargin=*]
    \item 清晨、清醒、静卧（完全休息状态）
    \item 环境温度 18～25${}^\circ$C
    \item 空腹（禁食 12～14 h，即清晨未进餐前）
    \item 精神安宁，无紧张和焦虑情绪
\end{enumerate}

\textbf{基础代谢率（Basal Metabolic Rate, BMR）：}
单位时间内人体基础代谢所消耗的能量。常用单位为：
\begin{itemize}[leftmargin=*]
    \item kJ/(m$^2\cdot$h) —— 按体表面积计
    \item kJ/(kg$\cdot$h) —— 按体重计
    \item kcal/(kg$\cdot$h)
\end{itemize}
\end{keybox}

\subsubsection{基础代谢的计算方法}

\begin{enumerate}[leftmargin=*]
    \item \textbf{体表面积计算法：}
    \begin{itemize}[leftmargin=*]
        \item 基础代谢能量 $=$ 体表面积（m$^2$）$\times$ BMR [kJ/(m$^2\cdot$h)] $\times$ 24h
        \item 体表面积（m$^2$）$= 0.00659 \times$ 身高（cm）$+ 0.0126 \times$ 体重（kg）$- 0.1603$
    \end{itemize}
    \item \textbf{直接公式法（WHO简化公式）：}
    \begin{itemize}[leftmargin=*]
        \item 18～30岁男性：BMR $= 15.3 \times$ 体重（kg）$+ 679$
        \item 18～30岁女性：BMR $= 14.7 \times$ 体重（kg）$+ 496$
        \item 30～60岁男性：BMR $= 11.6 \times$ 体重（kg）$+ 879$
        \item 30～60岁女性：BMR $= 8.7 \times$ 体重（kg）$+ 829$
    \end{itemize}
    \item \textbf{体重计算法：}基础代谢约 $=$ 1 kcal/kg$\cdot$h。成人基础代谢约占总能量消耗的 60\%～70\%。
\end{enumerate}

\subsubsection{影响基础代谢的因素}

\begin{keybox}
\textbf{影响基础代谢的因素：}
\begin{enumerate}[leftmargin=*]
    \item \textbf{体表面积：}体表面积越大，散热量越多，BMR越高
    \item \textbf{身体成分（体组织）：}瘦体重（肌肉组织）代谢活性高于脂肪组织。瘦体重比例高者BMR高。男性BMR高于女性约 5\%～10\%
    \item \textbf{年龄：}婴幼儿、儿童青少年生长发育期BMR最高，随年龄增长BMR逐渐下降
    \item \textbf{性别：}同龄男性BMR高于女性（因瘦体重比例更高）
    \item \textbf{环境温度：}寒冷地区BMR高于热带地区；体温每升高1${}^\circ$C，BMR升高约13\%
    \item \textbf{内分泌：}甲状腺功能亢进→BMR升高（甲状腺素是调节BMR最重要的激素）；儿茶酚胺（肾上腺素、去甲肾上腺素）→BMR升高；垂体功能低下→BMR降低
    \item \textbf{其他因素：}妊娠、哺乳、发热、应激状态均升高BMR
\end{enumerate}
\end{keybox}

\subsection{体力活动能量消耗}

体力活动的能量消耗约占人体总能量消耗的 \imp{15\%～30\%}，是能量消耗中变化最大的部分。

\begin{keybox}
\textbf{体力活动水平（Physical Activity Level, PAL）分级：}

\begin{table}[H]
\centering
\caption{中国成人活动水平分级（PAL值）}
\begin{tabular}{ccc}
\toprule
\textbf{活动强度} & \textbf{男性 PAL} & \textbf{女性 PAL} \\
\midrule
轻体力活动 & 1.55 & 1.56 \\
中体力活动 & 1.78 & 1.64 \\
重体力活动 & 2.10 & 1.82 \\
\bottomrule
\end{tabular}
\end{table}

\textbf{TEE计算：} $ \text{TEE} = \text{BMR} \times \text{PAL} $
\end{keybox}

\begin{tipbox}
轻体力活动如：办公室工作、售货员、讲课等；中体力活动如：学生日常活动、机动车驾驶、电工等；重体力活动如：农业劳动、炼钢、舞蹈、体育运动等。
\end{tipbox}

\subsection{食物热效应（Thermic Effect of Food, TEF）}

\begin{keybox}
\textbf{食物热效应（TEF/SDA，教学难点*）：}

\keyword{食物热效应（TEF）}又称\keyword{特殊动力作用（Specific Dynamic Action, SDA）}，指人体在摄食过程中由于对食物中的营养素进行消化、吸收、代谢转化等引起的额外能量消耗。

\textbf{三大产能营养素的TEF差异：}
\begin{itemize}[leftmargin=*]
    \item \imp{脂肪：4\%～5\%}（最低）
    \item \imp{碳水化合物：5\%～6\%}（中等）
    \item \imp{蛋白质：30\%～40\%}（最高！相当于本身提供能量的30\%～40\%）
\end{itemize}

混合膳食的食物热效应约占总能量的 \textbf{10\%}。
\end{keybox}

\begin{exambox}
\textbf{考点提示：}蛋白质的特殊动力作用最高（30\%～40\%），是考试常考知识点。原因：蛋白质合成、尿素生成等代谢过程高度耗能。
\end{exambox}

% ----- 第四节 能量需要量的测定方法 -----
\section{能量需要量的测定方法\textsuperscript{*}}

\begin{keybox}
\textbf{能量需要量测定方法（教学难点*）：}
\begin{enumerate}[leftmargin=*]
    \item \textbf{计算法：}能量需要量 = BMR $\times$ PAL。是最常用的方法，适用于群体能量需要量评估。
    \item \textbf{膳食调查法：}通过记录食物摄入量计算能量摄入量。简单易行但依赖受试者的配合准确性。
    \item \textbf{仪器测定法：}
    \begin{enumerate}
        \item \textbf{直接测热法（Direct Calorimetry）：}在人体测热室中直接测量机体散热。原理精确但设备昂贵、操作复杂。
        \item \textbf{间接测热法（Indirect Calorimetry）：}通过测量O$_2$消耗量和CO$_2$产生量，结合\keyword{呼吸商（RQ）}计算能量消耗。是最准确的能量消耗测定方法之一。
        \item \textbf{双标水法（Doubly Labeled Water, DLW）：}受试者饮用含$^2$H$_2$$^{18}$O的双标水后，通过测定体液中$^2$H和$^{18}$O的消除速率，计算CO$_2$产生量和能量消耗。原理是$^2$H反映水代谢、$^{18}$O反映水和CO$_2$代谢，二者之差即CO$_2$生成量。该方法不影响受试者正常活动、测定结果接近真实生活状态，被认为是测定自由生活状态下总能量消耗的\imp{金标准}。
    \end{enumerate}
\end{enumerate}
\end{keybox}

\begin{defbox}
\textbf{呼吸商（Respiratory Quotient, RQ）：}机体在一定时间内CO$_2$产生量与O$_2$消耗量的比值。RQ $=$ CO$_2$产生量 / O$_2$消耗量。

不同产能营养素的RQ：碳水化合物 RQ $= 1.00$，脂肪 RQ $= 0.71$，蛋白质 RQ $= 0.80$。混合膳食的RQ通常为0.85左右。RQ可用于估测主要供能物质的类型。
\end{defbox}

% ----- 第五节 能量平衡 -----
\section{能量平衡}

\begin{keybox}
\textbf{能量平衡（Energy Balance）：}

能量摄入量与能量消耗量之间的关系决定了体重的变化：
\begin{itemize}[leftmargin=*]
    \item \textbf{能量正平衡：}摄入 > 消耗 $\rightarrow$ 剩余能量以脂肪形式储存，体重增加
    \item \textbf{能量负平衡：}摄入 < 消耗 $\rightarrow$ 动员体内脂肪和蛋白质供能，体重下降
    \item \textbf{能量平衡：}摄入 = 消耗 $\rightarrow$ 体重相对稳定
\end{itemize}

能量的摄入和支出是维持体重的关键因素。长期能量正平衡导致\keyword{超重}和\keyword{肥胖}，长期能量负平衡则是\keyword{消瘦}和\keyword{营养不良}的原因。
\end{keybox}

% ----- 第六节 能量推荐摄入量及产能营养素供能比 -----
\section{能量推荐摄入量及产能营养素供能比}

\begin{keybox}
\textbf{中国居民膳食能量推荐摄入量（RNI）——产能营养素供能比：}
\begin{itemize}[leftmargin=*]
    \item \imp{碳水化合物：占总能量的 55\%～65\%}
    \item \imp{脂肪：占总能量的 20\%～30\%}
    \item \imp{蛋白质：占总能量的 10\%～12\%}
\end{itemize}

\textbf{能量 RNI 示例（轻体力活动成人）：}
\begin{itemize}[leftmargin=*]
    \item 男性（18～49岁）：约2250 kcal/d（9.41 MJ/d）
    \item 女性（18～49岁）：约1800 kcal/d（7.53 MJ/d）
\end{itemize}
\end{keybox}

\begin{exambox}
\textbf{考点提示：}产能营养素的供能比是考试常考点。碳水化合物的供能比55\%～65\%最为核心。此为本课程第一章营养素章节的"总结归纳"指标，须牢记。
\end{exambox}

% ----- 复习题 -----
\reviewcontent{
\section{复习题}

\begin{enumerate}[leftmargin=*, itemsep=8pt]
    \item \textbf{【名词解释】}能量系数、基础代谢（BEE）、基础代谢率（BMR）、食物热效应（TEF/SDA）、呼吸商（RQ）、体力活动水平（PAL）、双标水法（DLW）
    \item \textbf{【简答题】}写出能量单位的换算关系（kJ与kcal）。
    \item \textbf{【简答题】}列出三大产能营养素的能量系数，并比较其物理热价与生理有效热价的差异。
    \item \textbf{【简答题】}简述基础代谢（BEE）的测定条件及影响因素。
    \item \textbf{【简答题】}比较三大产能营养素的食物热效应（TEF），并说明蛋白质TEF最高的原因。
    \item \textbf{【简答题】}简述能量需要量的测定方法。
    \item \textbf{【简答题】}说明中国居民膳食产能营养素的推荐供能比。
\end{enumerate}

\subsection*{参考答案要点}

\begin{enumerate}[leftmargin=*, itemsep=5pt]
    \item 见本章各概念定义框。
    \item 1 kJ = 0.239 kcal；1 kcal = 4.184 kJ。
    \item 碳水化合物4 kcal/g（16.81 kJ/g），脂肪9 kcal/g（37.56 kJ/g），蛋白质4 kcal/g（16.74 kJ/g）。蛋白质物理热价（约5.65 kcal/g）> 生理有效热价（约4 kcal/g），因体内氧化不彻底，尿素等含氮产物仍有能量排出。脂肪和碳水化合物的物理热价＝生理有效热价。
    \item 测定条件：清晨清醒静卧、18～25${}^\circ$C、空腹12～14 h、精神安宁。影响因素：体表面积、身体成分、年龄、性别、环境温度、内分泌（甲状腺素为最主要调节激素）、妊娠哺乳等。
    \item 脂肪4\%～5\% < 碳水化合物5\%～6\% < 蛋白质30\%～40\%。蛋白质TEF最高，原因为蛋白质合成、尿素生成等代谢过程高度耗能。
    \item 计算法（BMR×PAL）、膳食调查法、仪器测定法（直接测热法、间接测热法、双标水法）。双标水法是测定自由生活状态下总能量消耗的金标准。
    \item 碳水化合物 55\%～65\%，脂肪 20\%～30\%，蛋白质 10\%～12\%。
\end{enumerate}

\newpage
}

% =====================================================================
%  CHAPTER 1E: 矿物质 — 参考：0310营养参考_extracted.txt (PPT原文)
% =====================================================================
\chapter{矿物质}
\chapterintro{掌握矿物质的定义与分类（常量/微量元素）；掌握钙的分布、功能、吸收影响因素和食物来源（难点*）；掌握铁的功能、缺乏三阶段及血红素铁/非血红素铁吸收特点；掌握碘的分布与来源；熟悉锌/硒的功能。}

\section{矿物质概述}

\begin{defbox}
\textbf{矿物质（Minerals）：}除碳、氢、氧和氮主要以有机化合物形式存在外，人体组织中其余元素统称为矿物质（无机盐/灰分）。

\textbf{分类：}\imp{常量元素（Macroelement）}：占体重0.01\%以上——Ca/P/K/Na/Cl/Mg/S。\imp{微量元素（Microelement）}：占体重0.01\%以下。必需微量元素：Fe/Cu/Zn/Se/Cr/I/Co/Mo。可能必需：Mn/Si/Ni/B/V。潜在毒性：F/Pb/Cd/Hg/As/Al/Sn/Li。

\textbf{营养特点：}(1)除食物外是唯一可通过天然水获取的营养素；(2)体内分布极不均匀；(3)\imp{生理剂量与中毒剂量范围较窄}；(4)功能：构成组织/维持渗透压与酸碱平衡/维持神经肌肉兴奋性及细胞膜通透性/构成活性物质/酶激活剂。
\end{defbox}

\section{钙（Calcium）}

\begin{keybox}
\textbf{分布：}\imp{人体含量最多的无机元素}（1.5\%$\sim$2.0\%），总量约1200g，\imp{99\%在骨骼牙齿中}（羟磷灰石Ca$_{10}$[PO$_4$]$_6$[OH]$_2$），1\%为混溶钙池。

\textbf{功能：}构成骨骼牙齿；酶激活剂/参与细胞代谢；参与凝血；构成混溶钙池；维持神经肌肉应激性；维持心跳节律。

\textbf{吸收影响因素：}\imp{促进}——乳糖/氨基酸/VD；\imp{抑制}——植酸盐/膳食纤维/草酸盐/脂肪/乙醇。\textbf{来源：}小虾皮/奶与奶制品/豆与豆制品。
\end{keybox}

\section{铁（Iron）}

\begin{keybox}
\textbf{功能：}血红蛋白原料（O$_2$/CO$_2$运输）；多种酶原料（生物氧化）；促进胡萝卜素→VA；促进抗体产生/胶原合成。

\textbf{缺乏：}好发人群——婴幼儿/青少年/育龄妇女。\imp{铁缺乏三阶段：}铁减少期→缺铁性红细胞生成期→缺铁性贫血期。

\textbf{吸收：}\imp{血红素铁}：与卟啉结合，直接吸收，不受干扰。\imp{非血红素铁}：以Fe(OH)$_3$形式存在于植物性食物中，须经胃酸分离并还原为二价铁才可吸收，吸收率低且受干扰。
\end{keybox}

\section{磷与镁}

\begin{defbox}
\textbf{磷（Phosphorus）：}
\begin{itemize}[leftmargin=*]
    \item 人体含量约650g，\imp{约70\%在骨骼中}，30\%在软组织
    \item 功能：构成骨骼牙齿（羟磷灰石）；核酸/磷脂/含磷辅酶（NAD$^+$、NADP$^+$、TPP等）的组成成分；参与能量代谢（ATP）；维持酸碱平衡
    \item 吸收：受维生素D促进；植酸抑制
    \item 钙磷比（Ca:P）：成人适宜1:1$\sim$1:1.5，婴儿母乳2:1（最适宜）
\end{itemize}

\textbf{镁（Magnesium）：}
\begin{itemize}[leftmargin=*]
    \item 成人含镁约25g，\imp{60\%$\sim$65\%在骨骼中}，27\%在肌肉
    \item 功能：300多种酶的辅因子；参与蛋白质合成和能量代谢；维持神经肌肉兴奋性；心血管保护作用
    \item 缺乏：神经肌肉兴奋性增高、心律失常、低钙血症（镁缺乏同时导致低血钙）
    \item 来源：绿叶蔬菜（叶绿素含镁）、坚果、豆类、全谷物
\end{itemize}
\end{defbox}

\section{锌（Zinc）}

\begin{keybox}
\textbf{功能：}
\begin{itemize}[leftmargin=*]
    \item 多种酶的组成成分和激活剂（超200种含锌酶）：碱性磷酸酶、RNA/DNA聚合酶、碳酸酐酶
    \item 促进生长发育和组织再生：参与核酸和蛋白质合成
    \item 维持正常味觉和食欲：味觉素的组成成分
    \item 参与免疫功能：促进淋巴细胞增殖和抗体生成
    \item 抗氧化：超氧化物歧化酶（SOD）的组成成分
    \item 促进维生素A代谢：参与视黄醇结合蛋白合成
\end{itemize}

\textbf{缺乏：}婴幼儿/儿童——\imp{生长发育迟缓、性发育不全、脑发育受损、食欲不振、异食癖}；成人——味觉减退、伤口愈合延迟、免疫功能下降、脱发、皮肤损害。

\textbf{来源：}动物性食物（贝壳类海产品、红肉、动物内脏）锌含量和利用率均高；\imp{母乳锌含量>牛乳}，初乳最高，生物价值更好；植物性食物受植酸抑制吸收。

\textbf{促进吸收因素：}动物蛋白质、组氨酸、半胱氨酸、柠檬酸。\textbf{抑制因素：}植酸、膳食纤维、高钙、高铜、高镉。
\end{keybox}

\section{硒（Selenium）}

\begin{keybox}
\textbf{功能：}
\begin{itemize}[leftmargin=*]
    \item \imp{谷胱甘肽过氧化物酶（GSH-Px）}的重要组成成分——抗氧化、保护细胞膜
    \item 参与甲状腺激素代谢：碘甲腺原氨酸脱碘酶（将T$_4$转化为活性T$_3$）
    \item 维持免疫功能：促进抗体生成
    \item 重金属解毒：与汞、镉、铅等结合形成不溶性硒化物
    \item 保护心血管、抗肿瘤
\end{itemize}

\textbf{缺乏：}\imp{克山病（Keshan Disease）}——以心肌坏死为特征的地方性心肌病（中国东北至西南的克山病带，硒缺乏为主因）；\imp{大骨节病（Kaschin-Beck Disease）}——骨关节变形、软骨坏死。过量可引起\imp{硒中毒}（脱发、脱甲、神经系统损害）。

\textbf{来源：}动物肝、肾、海产品、肉类；植物性食物硒含量取决于土壤硒水平。中国成人RNI：60 $\mu$g/d。
\end{keybox}

\section{碘（Iodine）}

\begin{defbox}
人体含碘约20$\sim$50mg，\imp{70\%$\sim$80\%在甲状腺}。

\textbf{功能：}参与甲状腺激素（T$_3$、T$_4$）合成，调节基础代谢、促进生长发育（尤其是神经系统发育）。

\textbf{缺乏：}\imp{地方性甲状腺肿（Goiter）}——缺碘→T$_3$/T$_4$合成减少→TSH反馈性升高→甲状腺代偿性肿大。\imp{克汀病（Cretinism）}——孕期/婴幼儿期严重缺碘导致不可逆的智力低下和生长发育障碍。缺碘地区：食用碘盐（食盐加碘是最经济有效的预防措施）。

\textbf{来源：}海盐/海产品（海带/紫菜/海鱼/贝类等最丰富）。成人RNI：120 $\mu$g/d。
\end{defbox}

\section{其他微量元素}

\begin{defbox}
\textbf{铜（Copper）：}参与铁代谢（铜蓝蛋白促进铁的吸收和利用）；参与结缔组织中胶原蛋白和弹性蛋白的交联（赖氨酰氧化酶）；参与黑色素合成（酪氨酸酶）。缺乏：小细胞低色素性贫血（与铁代谢障碍有关）、中性粒细胞减少、骨质疏松。

\textbf{铬（Chromium）：}\imp{三价铬是葡萄糖耐量因子（GTF）的组成成分}，增强胰岛素作用，促进葡萄糖利用。缺乏：葡萄糖耐量降低、胰岛素抵抗。来源：全谷物、肉类、啤酒酵母。成人AI：30 $\mu$g/d。

\textbf{氟（Fluoride）：}促进骨骼和牙齿的矿化，预防龋齿。\imp{缺乏→龋齿发病率增高；过量→氟斑牙（Dental Fluorosis）和氟骨症（Skeletal Fluorosis）}。来源：饮用水（主要来源）、海产品、茶叶（含氟量高）。生理剂量与中毒剂量范围极窄！

\textbf{钴（Cobalt）：}维生素B$_{12}$的组成成分（唯一含金属元素的维生素），通过VitB$_{12}$发挥生理作用。缺乏即VitB$_{12}$缺乏——巨幼红细胞性贫血。
\end{defbox}

\section{复习题}
\begin{enumerate}[leftmargin=*, itemsep=8pt]
    \item \textbf{【名词解释】}矿物质；常量元素；微量元素；混溶钙池；血红素铁；非血红素铁；克山病；克汀病；葡萄糖耐量因子（GTF）
    \item \textbf{【简答题】}简述矿物质的分类和营养特点。
    \item \textbf{【简答题】}简述钙的分布、生理功能、吸收影响因素和食物来源。
    \item \textbf{【简答题】}简述铁的功能、缺乏三阶段及血红素铁与非血红素铁的区别。
    \item \textbf{【简答题】}简述锌的生理功能及缺乏表现。
    \item \textbf{【简答题】}简述硒的生理功能、缺乏病（克山病/大骨节病）及食物来源。
    \item \textbf{【简答题】}简述碘的分布、功能及缺乏病（地方性甲状腺肿/克汀病）。
    \item \textbf{【简答题】}简述钙磷比的意义及磷、镁的主要功能。
\end{enumerate}

\subsection*{参考答案要点}
\begin{enumerate}[leftmargin=*, itemsep=5pt]
    \item 见正文定义。\textbf{克山病}：以心肌坏死为特征的地方性心肌病，硒缺乏为主因。\textbf{克汀病}：孕期/婴幼儿期严重缺碘导致的不可逆智力低下和生长发育障碍。\textbf{GTF}：三价铬组成的葡萄糖耐量因子，增强胰岛素作用。
    \item 分类：常量(>0.01\%：Ca/P/K/Na/Cl/Mg/S)+微量(<0.01\%：必需Fe/Cu/Zn/Se/Cr/I/Co/Mo；可能必需Mn/Si/Ni/B/V；潜在毒性F/Pb/Cd/Hg/As/Al/Sn/Li)。特点：(1)唯一可通过天然水获取；(2)分布极不均匀；(3)生理剂量与中毒剂量范围窄；(4)构成组织/渗透压酸碱平衡/神经肌肉兴奋性/酶激活剂。
    \item 含量最多无机元素(1.5-2.0\%)，99\%骨骼牙齿(羟磷灰石)，1\%混溶钙池。功能：骨骼牙齿/酶激活/凝血/神经肌肉应激性/心跳节律。促进吸收：乳糖/氨基酸/VD。抑制：植酸盐/草酸盐/膳食纤维/脂肪/乙醇。来源：小虾皮/奶制品/豆制品。成人RNI 800mg/d。
    \item 功能：血红蛋白(O$_2$/CO$_2$运输)/生物氧化/促进胡萝卜素→VA/抗体+胶原。缺乏三阶段：铁减少期→缺铁性红细胞生成期→缺铁性贫血期。血红素铁(与卟啉结合，直接吸收不受干扰)vs非血红素铁(Fe(OH)$_3$形式，需还原为Fe$^{2+}$，吸收率低受植酸/草酸干扰)。男性RNI 12mg/d，女性20mg/d。
    \item 功能：多种酶(超200种)组成成分和激活剂/促进生长发育和组织再生/维持味觉食欲(味觉素)/免疫功能/SOD抗氧化/促进VA代谢。缺乏：生长发育迟缓/性发育不全/脑发育受损/食欲不振/异食癖/味觉减退/伤口愈合延迟/免疫功能下降。来源：贝壳类海产品/红肉/动物内脏，母乳>牛乳。成人男性RNI 12.5mg/d，女性7.5mg/d。
    \item 功能：GSH-Px组成成分(抗氧化/保护细胞膜)/参与甲状腺激素代谢(脱碘酶T$_4$→T$_3$)/维持免疫/重金属解毒。缺乏：克山病(心肌坏死，缺硒地带)、大骨节病(骨关节变形)。过量：硒中毒(脱发/脱甲/神经损害)。来源：动物肝肾/海产品/肉类(植物受土壤硒水平影响)。成人RNI 60$\mu$g/d。
    \item 70-80\%在甲状腺。功能：合成T$_3$/T$_4$，调节基础代谢/促进生长发育(尤其神经系统)。缺乏：地方性甲状腺肿(缺碘→TSH升高→甲状腺代偿性肿大)/克汀病(孕期/婴幼儿缺碘→不可逆智力低下)。来源：海产品(海带/紫菜/海鱼最丰富)。成人RNI 120$\mu$g/d。食盐加碘最经济有效。
    \item Ca:P成人1:1-1:1.5，婴儿母乳2:1(最适宜)。磷功能：骨骼牙齿/核酸磷脂/含磷辅酶/能量代谢(ATP)/酸碱平衡。镁功能：300多种酶辅因子/蛋白质合成/能量代谢/神经肌肉兴奋性/心血管保护。缺乏→神经肌肉兴奋性增高/心律失常/低钙血症。来源：绿叶蔬菜(叶绿素含镁)/坚果/豆类/全谷物。
\end{enumerate}
\newpage

% =====================================================================
%  CHAPTER 1F: 营养学基础 - 维生素与水
% =====================================================================
\chapter{营养学基础（六）：维生素与水}
\setcounter{chapter}{1}
\chapterintro{
掌握\keyword{维生素的共同特点}、\keyword{脂溶性维生素与水溶性维生素的特性对比}、\keyword{维生素缺乏发展的四个阶段}；掌握\keyword{维生素A（视黄醇）的概念、生理功能、缺乏与过量表现及食物来源}；掌握\keyword{维生素D的概念、生理功能、缺乏病及来源}；掌握\keyword{维生素C（抗坏血酸）的理化性质、生理功能、坏血病的表现及食物来源}；掌握\keyword{维生素B$_1$（硫胺素）的理化性质、TPP辅酶功能、脚气病的三型表现及食物来源}；掌握\keyword{维生素B$_2$（核黄素）的理化性质（对紫外光敏感）、口腔-生殖系统综合征（缺乏症）及食物来源}；掌握\keyword{烟酸（维生素PP/Niacin）的活性形式（NAD$^+$和NADP$^+$）、癞皮病（Pellagra）的三D症状}。熟悉维生素E、维生素K、维生素B$_6$、叶酸、维生素B$_{12}$的基本知识。了解泛酸、生物素、肉碱、牛磺酸、肌醇、柠檬素（生物类黄酮）及水的生理功能。
}

% ==================== 第一节 维生素概述 ====================
\section{维生素概述}

\begin{defbox}
\textbf{维生素（Vitamin）}是维持机体正常生理功能所必需的一类微量低分子有机化合物。它们既不是构成机体组织的原料，也不为机体提供能量，但在物质代谢中发挥不可或缺的调节作用。
\end{defbox}

\subsection{维生素共同特点}

\begin{keybox}
\textbf{维生素的共同特点：}
\begin{enumerate}[leftmargin=*]
    \item \textbf{以本体或前体形式存在于天然食物中}，必须经常由食物提供
    \item \textbf{需要量甚微}（mg或$\mu$g水平），但在调节机体代谢方面作用重要
    \item \textbf{不构成组织，也不提供能量}（有别于产能营养素）
    \item \textbf{多以辅酶或辅基的形式发挥功能}
    \item \textbf{有的具有几种结构相近、活性相同的化合物}
\end{enumerate}
\end{keybox}

\subsection{维生素的命名与分类}

\begin{keybox}
\textbf{维生素的命名：}
\begin{itemize}[leftmargin=*]
    \item 按发现顺序命名（如VitA、VitB、VitC、VitD等）
    \item 按化学结构命名（如视黄醇、硫胺素、核黄素等）
    \item 按生理功能命名（如抗坏血酸、抗干眼病维生素等）
\end{itemize}
\end{keybox}

\begin{keybox}
\textbf{维生素的分类*（按溶解性）：}

\textbf{脂溶性维生素：}
\begin{itemize}[leftmargin=*]
    \item \keyword{维生素A（视黄醇）}、\keyword{维生素D（钙化醇）}、\keyword{维生素E（生育酚）}、\keyword{维生素K（凝血维生素）}
    \item 吸收需要脂肪和胆汁的存在
    \item 可在体内储存（主要在肝脏和脂肪组织），缺乏症状出现缓慢
    \item 摄入过量可引起中毒（尤其VitA和VitD，VitK除外）
    \item 除VitK外，均有UL限制
\end{itemize}

\textbf{水溶性维生素：}
\begin{itemize}[leftmargin=*]
    \item \keyword{B族维生素}：B$_1$（硫胺素）、B$_2$（核黄素）、B$_3$（烟酸/维生素PP）、B$_5$（泛酸）、B$_6$（吡哆醇）、B$_7$（生物素）、B$_9$（叶酸）、B$_{12}$（钴胺素）
    \item \keyword{维生素C（抗坏血酸）}
\end{itemize}

\textbf{水溶性维生素的特点：}
\begin{itemize}[leftmargin=*]
    \item 体内储存量极少（B$_{12}$除外），多余者以原形或代谢产物经尿排出
    \item 缺乏症状出现较快（尤其B族维生素）
    \item 一般不引起中毒（过量摄入可通过尿液快速排泄）
    \item 饱和剂量以上几乎不增加组织浓度
\end{itemize}
\end{keybox}

\begin{exambox}
\textbf{考点提示：}两类维生素的比较是常考基础知识点。关键记忆点：脂溶性$\rightarrow$可储存、慢缺乏、易中毒（VitK除外）；水溶性$\rightarrow$储存少、快缺乏、不易中毒（尿液排泄）。
\end{exambox}

\subsection{维生素缺乏}

\begin{keybox}
\textbf{维生素缺乏的原因：}
\begin{itemize}[leftmargin=*]
    \item 维生素摄入量不足
    \item 吸收利用降低
    \item 维生素需要量相对增高
\end{itemize}

\textbf{缺乏分类：}
\begin{itemize}[leftmargin=*]
    \item 按原因分：\textbf{原发性缺乏}和\textbf{继发性缺乏}
    \item 按缺乏程度分：\textbf{临床缺乏}和\textbf{亚临床缺乏（边缘缺乏）}
\end{itemize}

\textbf{维生素缺乏发展四阶段（由轻到重）：}
\begin{enumerate}[leftmargin=*]
    \item \textbf{组织中维生素储量降低}（亚临床/边缘缺乏，无明显症状）
    \item \textbf{生化指标异常，生理功能降低}（生化检查可发现，如酶活性下降）
    \item \textbf{组织病理变化，出现临床症状和体征}（典型的缺乏病表现）
    \item \textbf{停止生命}（极端严重缺乏，可致死）
\end{enumerate}
\end{keybox}

% ==================== 第二节 维生素A ====================
\section{维生素A（视黄醇，Retinol）}

\subsection{概念与理化性质}

\begin{defbox}
\textbf{维生素A类}是指含有$\beta$-白芷酮环的多烯基结构、并具有视黄醇生物活性的一大类物质。
\end{defbox}

\begin{keybox}
\textbf{存在形式：}
\begin{itemize}[leftmargin=*]
    \item \textbf{已形成的维生素A：}动物体内具有视黄醇生物活性的维生素A称为已形成的维生素A，包括\keyword{视黄醇（Retinol）}、\keyword{视黄醛（Retinal）}、\keyword{视黄酸（Retinoic acid）}等物质
    \item \textbf{维生素A原（Provitamin A）：}黄、绿、红色植物含有类胡萝卜素，其中一部分在体内转变成维生素A的类胡萝卜素称为维生素A原，其中主要有$\alpha$-胡萝卜素、\imp{$\beta$-胡萝卜素}、$\gamma$-胡萝卜素和隐黄素4种，\imp{以$\beta$-胡萝卜素的活性最高}
\end{itemize}

\textbf{理化性质：}维生素A和胡萝卜素都对酸、碱和热稳定，一般烹调和罐头加工不易破坏，但\imp{易被氧化和受紫外线破坏}。当食物中含有磷脂、维生素E、维生素C和其它抗氧化剂时，视黄醇和胡萝卜素较为稳定，脂肪酸败可引起其严重破坏。

\textbf{单位：}
\begin{itemize}[leftmargin=*]
    \item 国际单位（IU）：1 IU 维生素A = 0.3 $\mu$g 全反式视黄醇
    \item 视黄醇当量（Retinol Equivalent, RE）
    \item 视黄醇活性当量（Retinol Activity Equivalent, RAE）
\end{itemize}
\end{keybox}

\subsection{生理功能}

\begin{keybox}
\textbf{维生素A的生理功能*：}
\begin{enumerate}[leftmargin=*]
    \item \textbf{维持正常视觉：}视网膜杆状细胞中的\keyword{视紫红质（Rhodopsin）}由视蛋白和11-顺视黄醛组成。光照$\rightarrow$视紫红质分解$\rightarrow$神经冲动$\rightarrow$视觉。暗处视紫红质再合成需要VitA持续供给。缺乏导致暗适应延长，严重者为\imp{夜盲症（Night Blindness）}
    \item \textbf{维持上皮的正常生长与分化：}维持皮肤和黏膜（呼吸道、消化道、泌尿道）的完整性
    \item \textbf{促进生长发育：}视黄酸通过核受体（RAR/RXR）调控基因表达
    \item \textbf{抑癌作用}
    \item \textbf{维持机体正常免疫功能：}促进免疫细胞分化和抗体生成
    \item \textbf{抗氧化作用：}$\beta$-胡萝卜素具有抗氧化功能，可淬灭单线态氧和清除自由基
\end{enumerate}
\end{keybox}

\subsection{维生素A缺乏与过量}

\begin{keybox}
\textbf{维生素A缺乏*：}

\textbf{长期缺乏的表现：}
\begin{itemize}[leftmargin=*]
    \item \textbf{夜盲症（Night Blindness）、干眼病（Xerophthalmia）}：暗适应能力下降$\rightarrow$夜盲症（最早出现的症状！）$\rightarrow$干眼病$\rightarrow$角膜软化$\rightarrow$角膜溃疡$\rightarrow$穿孔$\rightarrow$失明
    \item 身体多种上皮组织角化，保护内脏器官的作用减弱，抵抗力降低
    \item 影响生长发育和生殖能力
    \item 致癌
\end{itemize}

\textbf{维生素A营养状况鉴定：}
\begin{itemize}[leftmargin=*]
    \item 血清维生素A水平：参考指标
    \item 改进的相对剂量反应实验：诊断边缘状态
    \item 暗适应功能测定：适用大规模人群调查
    \item 血浆视黄醇结合蛋白：较好指标
    \item 眼结膜印迹细胞学法及眼部症状检查
\end{itemize}
\end{keybox}

\begin{tipbox}
\textbf{维生素A过量的危害：}
\begin{itemize}[leftmargin=*]
    \item 皮肤干燥、瘙痒、脱皮、脱发，出现皮疹
    \item 食欲下降，腹泻，肝脾肿大
    \item 易激动，肌肉无力，头痛，恶心，呕吐
    \item 易出血，凝血时间延长
    \item 破骨细胞活性增强，造成骨骼脆性增加，关节变形
    \item \textbf{致畸性：}孕早期过量VitA可导致胎儿畸形
    \item 注意：不包括$\beta$-胡萝卜素过量（仅引起皮肤黄染，停药后可消退）
\end{itemize}
\end{tipbox}

\subsection{视黄醇当量与参考摄入量}

\begin{keybox}
\textbf{视黄醇当量（Retinol Equivalent, RE）*：}

不同来源的维生素A折算公式：
\begin{center}
\imp{视黄醇当量 ($\mu$g RE) = 视黄醇 ($\mu$g) + $\beta$-胡萝卜素 ($\mu$g) $\times$ 0.167}
\end{center}

即：1 $\mu$g视黄醇 = 1 $\mu$g RE；1 $\mu$g $\beta$-胡萝卜素 = 0.167 $\mu$g RE = 1/6 $\mu$g RE

\textbf{膳食参考摄入量：}
\begin{itemize}[leftmargin=*]
    \item 成年男性 RNI：\imp{800 μg RE/d}
    \item 成年女性 RNI：\imp{700 μg RE/d}
    \item 孕妇早期：700 μg RE/d；中、晚期：770 μg RE/d
    \item 乳母：1300 μg RE/d
    \item UL（可耐受最高摄入量）：\imp{3000 μg RE/d}
\end{itemize}
\end{keybox}

\begin{exambox}
\textbf{维生素A的食物来源（常考）：}
\begin{itemize}[leftmargin=*]
    \item \textbf{动物来源：}动物肝脏（猪肝、牛肝、羊肝、鸡肝、鸭肝）、奶制品、鸡蛋、河蟹
    \item \textbf{植物来源（由$\beta$-胡萝卜素转化而来）：}胡萝卜、菠菜、冬苋菜、茴香、芥蓝、金针菜、芦笋、西兰花、芒果、紫菜
\end{itemize}
\end{exambox}

% ==================== 第三节 维生素D ====================
\section{维生素D（钙化醇，Calciferol）}

\subsection{概念与理化性质}

\begin{defbox}
\textbf{维生素D类}是指含环戊氢烯菲环结构、并具有钙化醇生物活性的一大类物质，包括\keyword{维生素D$_2$（麦角钙化醇，Ergocalciferol）}和\keyword{维生素D$_3$（胆钙化醇，Cholecalciferol）}。
\end{defbox}

\begin{keybox}
\textbf{理化性质：}维生素D性质较稳定，能抵抗酸、碱及氧化作用，对热也较稳定。所以通常烹调加工不会引起维生素D损失。但脂肪酸败可引起维生素D破坏。

\textbf{单位换算：}1 IU 维生素D$_3$ = 0.025 $\mu$g 维生素D$_3$；1 $\mu$g = 40 IU

\textbf{来源与活化*：}
\begin{itemize}[leftmargin=*]
    \item \textbf{VitD$_2$（麦角钙化醇）：}来源于植物（酵母、真菌经紫外线照射后麦角固醇转化）
    \item \textbf{VitD$_3$（胆钙化醇）：}\imp{皮肤中的7-脱氢胆固醇经紫外线（UVB，290$\sim$315nm）照射转化生成}（内源性来源，最重要）
\end{itemize}

\textbf{维生素D的活化过程：}
肝 $\xrightarrow{\text{25-羟化酶}}$ 25-(OH)D$_3$（肝内主要储存形式）$\xrightarrow{\text{肾1}\alpha\text{-羟化酶}}$ \imp{1,25-(OH)$_2$D$_3$（骨化三醇}，生物活性最强的形式）

1,25-(OH)$_2$D$_3$被称为"活性维生素D"，本质属于\keyword{类固醇激素}。
\end{keybox}

\subsection{生理功能}

\begin{keybox}
\textbf{维生素D的生理功能*：}

维持细胞内外钙浓度，调节钙磷代谢：
\begin{enumerate}[leftmargin=*]
    \item \textbf{促进小肠钙吸收：}诱导肠上皮细胞合成钙结合蛋白（CaBP）
    \item \textbf{促进肾小管对钙、磷的重吸收}
    \item \textbf{对骨细胞呈现多种作用：}调节骨骼的钙化与脱钙，维持血钙和血磷的正常水平
    \item \textbf{调节基因转录作用}
    \item \textbf{通过维生素D内分泌系统调节血钙平衡}
\end{enumerate}
\end{keybox}

\subsection{维生素D缺乏与过量}

\begin{keybox}
\textbf{维生素D缺乏*：}
\begin{itemize}[leftmargin=*]
    \item \textbf{儿童——佝偻病（Rickets）：}骨骼钙化不良，O形腿、X形腿、鸡胸、方颅、肋骨串珠、Harrison沟
    \item \textbf{成人——骨软化症（Osteomalacia）：}骨矿物化不良、骨痛、肌无力
    \item \textbf{老年人——骨质疏松症（Osteoporosis）：}骨量减少、骨折风险增高
\end{itemize}
\end{keybox}

\begin{tipbox}
\textbf{维生素D过多症：}
\begin{itemize}[leftmargin=*]
    \item 食欲下降，恶心腹泻
    \item 皮肤瘙痒
    \item 多尿，肾小管钙化，肾功能减退
    \item 心血管系统异常
\end{itemize}

\textbf{机体营养水平鉴定：}
\begin{itemize}[leftmargin=*]
    \item 血浆中1,25-(OH)$_2$-D$_3$的浓度
    \item 血浆中25-OH-D$_3$的浓度
\end{itemize}
\end{tipbox}

\subsection{参考摄入量与食物来源}

\begin{keybox}
\textbf{维生素D的膳食参考摄入量：}
\begin{itemize}[leftmargin=*]
    \item 成人 RNI：\imp{5 μg/d（200 IU）}
    \item 婴幼儿、儿童、孕妇、乳母、老年人 RNI：\imp{10 μg/d（400 IU）}
    \item UL（可耐受最高摄入量）：50 μg/d（2000 IU）
\end{itemize}
\end{keybox}

\begin{exambox}
\textbf{维生素D的来源：}
\begin{itemize}[leftmargin=*]
    \item \textbf{食物来源：}鱼肝油、海鱼、动物肝脏；全脂乳、蛋类、黄油、奶油、奶酪
    \item \textbf{其他来源：}\imp{适度的日光照射}（最经济、最主要的方式）
\end{itemize}
\end{exambox}

% ==================== 第四节 维生素E ====================
\section{维生素E（生育酚，Tocopherol）}

\begin{defbox}
\textbf{维生素E}是一组具有抗氧化活性的生育酚（Tocopherols）和生育三烯酚的统称，其中\keyword{$\alpha$-生育酚}的生物活性最强。
\end{defbox}

\textbf{生理功能：}最重要的功能是\keyword{抗氧化作用}——保护细胞膜多不饱和脂肪酸免受自由基攻击（链阻断型抗氧化剂），防止脂质过氧化（形成脂褐素即老年斑）；维持生殖功能；抗衰老；保护红细胞膜完整性。

\textbf{特点：}维生素E是\imp{毒性最小的脂溶性维生素}，口服大量耐受良好。

\begin{keybox}
\textbf{膳食参考摄入量：}
\begin{itemize}[leftmargin=*]
    \item 成人 AI：\imp{14 mg $\alpha$-TE/d}
    \item 乳母：17 mg $\alpha$-TE/d
\end{itemize}
\end{keybox}

\textbf{食物来源：}植物油（小麦胚芽油、葵花籽油、玉米油）、坚果、种子、绿叶蔬菜。

% ==================== 第五节 维生素K ====================
\section{维生素K（凝血维生素）}

\begin{keybox}
\textbf{维生素K的生理功能：}

维生素K是肝脏合成凝血因子所必需的物质，参与\keyword{凝血因子II（凝血酶原）、VII、IX、X}的谷氨酸残基$\gamma$-羧基化修饰，使其获得与Ca$^{2+}$结合的能力从而具有凝血活性。

缺乏导致凝血障碍、出血倾向。新生儿因肠道菌群未建立且母乳中VitK含量低，易出现\keyword{新生儿出血症}，出生后常规注射维生素K预防。
\end{keybox}

\textbf{食物来源：}绿色绿叶蔬菜（菠菜、甘蓝、西蓝花）含量最丰富；肠道细菌也能合成一部分。

% ==================== 第六节 维生素B1 ====================
\section{维生素B$_1$（硫胺素，Thiamin）}

\subsection{理化性质}

\begin{keybox}
\textbf{硫胺素的理化性质：}
\begin{itemize}[leftmargin=*]
    \item 硫胺素的盐酸盐和硝酸盐形式在干燥和酸性溶液中均稳定，但加热至其熔点（249℃）即分解
    \item 在碱性环境下易被氧化失去活性
    \item \imp{硫胺素对亚硫酸盐极为敏感}，如有亚硫酸盐存在时迅速分解成嘧啶和噻唑，并丧失其活性
\end{itemize}
\end{keybox}

\subsection{生理功能}

\begin{keybox}
\textbf{维生素B$_1$的生理功能*：}

\textbf{活性形式：}\keyword{焦磷酸硫胺素（TPP）}，主要以辅酶形式参与两类酶系统的反应：

\begin{enumerate}[leftmargin=*]
    \item \textbf{参与体内碳水化合物和能量代谢：}
    \begin{itemize}[leftmargin=*]
        \item \textbf{羧化酶}——催化$\alpha$-酮酸的氧化脱羧反应，如丙酮酸$\rightarrow$乙酰辅酶A、$\alpha$-酮戊二酸$\rightarrow$琥珀酰辅酶A（糖代谢与三羧酸循环关键环节）
        \item \textbf{转酮醇酶}——转运二碳单位的作用（磷酸戊糖途径）
    \end{itemize}
    \item \textbf{对神经系统的影响：}TPP直接激活神经细胞的氯离子通道，控制神经传导的启动
    \item \textbf{对消化系统的影响：}硫胺素可抑制胆碱酯酶，促进食欲、胃肠道的正常蠕动和消化液的分泌
\end{enumerate}
\end{keybox}

\subsection{维生素B$_1$缺乏——脚气病（Beriberi）}

\begin{keybox}
\textbf{维生素B$_1$缺乏症——"多发性神经炎/脚气病"*：}

\textbf{干性脚气病（Dry Beriberi）：}
\begin{itemize}[leftmargin=*]
    \item 神经系统：烦躁健忘、精神分散、多梦，继而出现外周神经炎、肌肉麻痹、腿麻木、有蚁走感
    \item 消化系统：胃肠蠕动减弱、便秘、食欲减退、消化不良
\end{itemize}

\textbf{湿性脚气病（Wet Beriberi）：}
\begin{itemize}[leftmargin=*]
    \item 循环系统：水肿、心动过速、心力衰竭、气喘、微血管渗透性增大
    \item 消化系统：厌食、便秘
\end{itemize}

\textbf{急性恶性脚气病（Acute Malignant Beriberi）：}
胸闷、心悸、气短、内出血、血压下降
\end{keybox}

\subsection{营养状况评价与食物来源}

\begin{keybox}
\textbf{人体营养状况评价：}
\begin{itemize}[leftmargin=*]
    \item 尿硫胺素负荷试验
    \item 硫胺素和肌酐含量比值
    \item 红细胞转酮醇酶活性系数（TPP效应）
\end{itemize}
\end{keybox}

\begin{keybox}
\textbf{参考摄入量：}
\begin{itemize}[leftmargin=*]
    \item 成年男性 RNI：\imp{1.4 mg/d}
    \item 成年女性 RNI：\imp{1.3 mg/d}
\end{itemize}
\end{keybox}

\begin{exambox}
\textbf{食物来源：}
\begin{itemize}[leftmargin=*]
    \item 粗杂粮、豆类、硬果、干酵母
    \item 瘦猪肉、动物内脏、鸡蛋黄
    \item 粮食、蔬菜、牛奶
\end{itemize}
谷类加工精度越高（碾磨越白），VB$_1$损失越严重。硫胺素集中在谷类外层和胚芽。
\end{exambox}

% ==================== 第七节 维生素B2 ====================
\section{维生素B$_2$（核黄素，Riboflavin）}

\subsection{理化性质}

\begin{keybox}
\textbf{核黄素的理化性质：}
\begin{itemize}[leftmargin=*]
    \item 核黄素的性质比较稳定，耐酸、耐热，不易氧化，但在碱性溶液中容易被破坏
    \item \imp{游离型核黄素对紫外光高度敏感}：在酸性条件下可光解为光黄素，在碱性条件下光解为光色素而丧失生物活性
    \item 黄色晶体，水溶液呈黄绿色荧光
\end{itemize}

\textbf{活性形式：}
\begin{itemize}[leftmargin=*]
    \item \keyword{黄素单核苷酸（FMN）}
    \item \keyword{黄素腺嘌呤二核苷酸（FAD）}
\end{itemize}

FMN和FAD是体内氧化还原酶类（黄素蛋白）的辅基，参与生物氧化中的氢传递，在能量代谢中起核心作用。
\end{keybox}

\subsection{维生素B$_2$缺乏}

\begin{keybox}
\textbf{维生素B$_2$缺乏与过量*：}

\textbf{典型缺乏症——"口腔-生殖系统综合征"（Oral-Genital Syndrome）：}
\begin{itemize}[leftmargin=*]
    \item \keyword{口角炎（Angular Stomatitis）}
    \item \keyword{唇炎（Cheilitis）}
    \item \keyword{舌炎（Glossitis）}：舌头红肿、乳头萎缩、地图舌
    \item \keyword{睑缘炎、结膜炎}
    \item \keyword{脂溢性皮炎（Seborrheic Dermatitis）}：鼻唇沟、眉间、耳后
    \item \keyword{阴囊皮炎}
\end{itemize}

\textbf{过量：}\imp{目前尚未见任何毒副作用。}
\end{keybox}

\begin{tipbox}
\textbf{维生素B$_2$缺乏的膳食因素：}
\begin{itemize}[leftmargin=*]
    \item 谷类加工过精（VB$_2$集中于外层和胚芽）
    \item 牛奶光照储存（VB$_2$光解破坏，故牛奶不宜用透明瓶储存）
    \item 绿色蔬菜摄入不足
    \item 动物性食物摄入不足
\end{itemize}
维生素B$_2$缺乏常与其他B族维生素缺乏同时发生。
\end{tipbox}

\subsection{营养状况评价与食物来源}

\begin{keybox}
\textbf{机体营养状况评价：}
\begin{itemize}[leftmargin=*]
    \item 红细胞谷胱甘肽还原酶活性系数
    \item 尿核黄素排出量
    \item 尿负荷试验
\end{itemize}
\end{keybox}

\begin{keybox}
\textbf{参考摄入量：}
\begin{itemize}[leftmargin=*]
    \item 成年男性 RNI：\imp{1.4 mg/d}
    \item 成年女性 RNI：\imp{1.2 mg/d}
\end{itemize}
\end{keybox}

\begin{exambox}
\textbf{食物来源：}
\begin{itemize}[leftmargin=*]
    \item 动物内脏、瘦肉、水产品（黄鳝、蟹）
    \item 奶类、蛋类
    \item 豆类、绿色蔬菜、食用菌、藻类、硬果
\end{itemize}
\end{exambox}

% ==================== 第八节 烟酸 ====================
\section{烟酸（Niacin，维生素B$_3$，维生素PP）}

\subsection{吸收与代谢}

\begin{keybox}
\textbf{烟酸的活性形式*：}
\begin{itemize}[leftmargin=*]
    \item \keyword{辅酶I——烟酰胺腺嘌呤二核苷酸（NAD$^+$）}
    \item \keyword{辅酶II——烟酰胺腺嘌呤二核苷酸磷酸（NADP$^+$）}
\end{itemize}

NAD$^+$和NADP$^+$是体内多种脱氢酶的辅酶，参与氧化还原反应和能量代谢（糖酵解、三羧酸循环、氧化磷酸化）。

\textbf{体内转化：}
色氨酸可在体内合成烟酸——\imp{约60mg色氨酸转化为1mg烟酸}。

\textbf{烟酸当量（Niacin Equivalent, NE）：}
\[
\text{NE (mg)} = \text{烟酸 (mg)} + \frac{\text{色氨酸 (mg)}}{60}
\]
\end{keybox}

\subsection{烟酸缺乏——癞皮病（Pellagra）}

\begin{keybox}
\textbf{烟酸缺乏症——癞皮病（Pellagra）*：}

烟酸缺乏症又称癞皮病（pellagra），主要损害\imp{皮肤，口、舌、胃肠道粘膜以及神经系统}。

\textbf{典型病例的"三D"症状：}
\begin{enumerate}[leftmargin=*]
    \item \keyword{Dermatitis（皮炎）：}暴露部位对称性皮炎（日照处红斑、粗糙、色素沉着）——Casal项链
    \item \keyword{Diarrhea（腹泻）：}消化不良、腹泻、舌炎
    \item \keyword{Dementia（痴呆）：}精神症状、记忆减退、焦虑、抑郁（严重时）
\end{enumerate}

\textbf{病因：}以玉米为主食且缺乏优质蛋白质的地区多发（玉米中烟酸为结合型无法吸收，且色氨酸含量低）。

\textbf{过量摄入的副作用：}皮肤发红、眼部感觉异常、高尿酸血症，偶见高血糖等。
\end{keybox}

\subsection{参考摄入量与食物来源}

\begin{keybox}
\textbf{参考摄入量：}
\begin{itemize}[leftmargin=*]
    \item 成年男性 RNI：\imp{14 mg NE/d}
    \item 成年女性 RNI：\imp{13 mg NE/d}
\end{itemize}
\end{keybox}

\begin{exambox}
\textbf{食物来源：}
\begin{itemize}[leftmargin=*]
    \item 动物肝脏、肉类、鱼类、花生、豆类
    \item 谷类中的烟酸80\%$\sim$90\%存在于麸皮中
    \item 乳类和蛋类中烟酸含量虽低但色氨酸含量丰富
\end{itemize}
\end{exambox}

% ==================== 第九节 维生素B6 ====================
\section{维生素B$_6$（吡哆醇，Pyridoxine）}

\begin{keybox}
\textbf{维生素B$_6$的生理功能与缺乏：}

\textbf{三种形式：}吡哆醇（Pyridoxine）、吡哆醛（Pyridoxal）、吡哆胺（Pyridoxamine）。

\textbf{活性形式：}\keyword{磷酸吡哆醛（PLP）}——最重要的辅酶形式。

\textbf{生理功能：}
\begin{enumerate}[leftmargin=*]
    \item 氨基酸代谢：参与转氨作用和脱羧反应（PLP是转氨酶的辅酶）
    \item 糖原代谢：参与糖原磷酸化酶作用（PLP使其活化）
    \item 参与不饱和脂肪酸代谢和神经递质（5-HT、GABA、多巴胺）的合成
\end{enumerate}

\textbf{缺乏表现：}脂溢性皮炎、舌炎、口角炎（类似于VB$_2$缺乏）；神经系统症状（易激惹、抑郁、周围神经炎）；小细胞低色素性贫血（参与血红素合成第一步$\delta$-ALA的生成）。
\end{keybox}

\begin{keybox}
\textbf{参考摄入量：}
\begin{itemize}[leftmargin=*]
    \item 成人 AI：\imp{1.2 mg/d}
\end{itemize}
\end{keybox}

\textbf{食物来源：}广泛分布于动物和植物性食物中，肉类、肝脏、全谷物、豆类、香蕉、土豆含量较高。

% ==================== 第十节 叶酸 ====================
\section{叶酸（Folate，维生素B$_9$）}

\begin{defbox}
\textbf{叶酸}属蝶酰谷氨酸，因首先在菠菜叶中分离而得名。活性形式为\keyword{四氢叶酸（Tetrahydrofolate, THF）}。
\end{defbox}

\begin{keybox}
\textbf{叶酸的生理功能与缺乏*：}

\textbf{生理功能：}
\begin{enumerate}[leftmargin=*]
    \item \textbf{"一碳单位"的载体：}THF是一碳单位（-CH$_3$、-CH$_2$、-CHO等）的载体，参与核苷酸（嘌呤、胸腺嘧啶dTMP）的生物合成
    \item 参与氨基酸代谢：同型半胱氨酸$\rightarrow$蛋氨酸的甲基化过程需要THF（N$^5$-甲基THF）作为甲基供体
    \item 参与血红蛋白的合成（与VitB$_{12}$共同作用）
\end{enumerate}

\textbf{缺乏*：}
\begin{enumerate}[leftmargin=*]
    \item \textbf{\imp{巨幼红细胞性贫血（Megaloblastic Anemia）}：}因DNA合成受阻，细胞分裂减慢，细胞体积增大但无法正常分裂。为叶酸缺乏最典型的表现
    \item \textbf{\imp{胎儿神经管畸形（Neural Tube Defects, NTD）}：}孕早期叶酸缺乏导致。包括脊柱裂（Spina Bifida）和无脑儿（Anencephaly）。\imp{围孕期（孕前至孕早期）补充400$\mu$g/d叶酸可有效预防NTD}
    \item \textbf{高同型半胱氨酸血症：}同型半胱氨酸$\rightarrow$蛋氨酸的甲基化受阻，血Hcy升高，为动脉粥样硬化的独立危险因素
\end{enumerate}
\end{keybox}

\begin{keybox}
\textbf{膳食叶酸当量（Dietary Folate Equivalent, DFE）*：}

食物中的叶酸（天然）与叶酸补充剂（合成）的生物利用率不同：
\begin{center}
\imp{DFE ($\mu$g) = 天然食物中叶酸 ($\mu$g) + 1.7 $\times$ 补充剂/强化食品中叶酸 ($\mu$g)}
\end{center}

\textbf{参考摄入量：}
\begin{itemize}[leftmargin=*]
    \item 成人 RNI：\imp{400 μg DFE/d}
    \item 孕妇：600 μg DFE/d；乳母：550 μg DFE/d
    \item UL（可耐受最高摄入量）：\imp{1000 μg DFE/d}
\end{itemize}
\end{keybox}

\textbf{食物来源：}深色绿叶蔬菜（菠菜、芥蓝）、动物肝脏、豆类、坚果、酵母。食物中叶酸易在加热和储存过程中损失。

% ==================== 第十一节 维生素B12 ====================
\section{维生素B$_{12}$（钴胺素，Cobalamin）}

\begin{keybox}
\textbf{维生素B$_{12}$的生理功能与缺乏：}

\textbf{特征：}
\begin{itemize}[leftmargin=*]
    \item 含有钴（Co）元素的维生素——唯一含金属元素的维生素
    \item \imp{仅存在于动物性食品中！}植物性食物不含维生素B$_{12}$（\keyword{严格素食者（Vegan）有缺乏风险}）
    \item 吸收需要\keyword{内因子（Intrinsic Factor, IF）}——胃壁细胞分泌
\end{itemize}

\textbf{生理功能：}
\begin{enumerate}[leftmargin=*]
    \item 作为蛋氨酸合成酶的辅酶（与叶酸共同参与同型半胱氨酸$\rightarrow$蛋氨酸）
    \item 参与神经系统的维护与髓鞘的形成
\end{enumerate}

\textbf{缺乏表现：}
\begin{enumerate}[leftmargin=*]
    \item \imp{巨幼红细胞性贫血：}与叶酸缺乏的贫血难以区分
    \item \textbf{神经系统损害：}脊髓后索和侧索退行性变（亚急性联合变性），而叶酸缺乏不引起神经系统损害——\imp{是与叶酸缺乏贫血的鉴别要点}！
\end{enumerate}
\end{keybox}

\begin{keybox}
\textbf{参考摄入量：}
\begin{itemize}[leftmargin=*]
    \item 成人 AI：\imp{2.4 μg/d}
    \item 孕妇：2.6 μg/d；乳母：2.8 μg/d
\end{itemize}
\end{keybox}

\textbf{食物来源：}动物肝脏、肉类、鱼类、蛋类、奶类。植物性食物基本不含VitB$_{12}$。

% ==================== 第十二节 维生素C ====================
\section{维生素C（抗坏血酸，Ascorbic Acid）}

\subsection{命名与理化性质}

\begin{keybox}
\textbf{维生素C的命名与理化性质：}
\begin{itemize}[leftmargin=*]
    \item 维生素C又名\keyword{抗坏血酸（Ascorbic Acid）}，为含6碳的$\alpha$-酮基内酯的弱酸，有酸味
    \item 是一种很强的还原剂
    \item \imp{与其它维生素相比，维生素C最不稳定，在有氧或碱性环境中极易氧化}
    \item L-型具有生物活性
\end{itemize}
\end{keybox}

\subsection{生理功能}

\begin{keybox}
\textbf{维生素C的生理功能*：}
\begin{enumerate}[leftmargin=*]
    \item \textbf{抗氧化作用：}保护细胞免受自由基的氧化损伤，可还原VitE的氧化形式使其重新具有活性
    \item \textbf{作为羟化过程底物和酶的辅因子：}
    \begin{itemize}[leftmargin=*]
        \item 维持\keyword{胶原蛋白的正常功能}：维持脯氨酸羟化酶和赖氨酸羟化酶的活性（羟化反应以VitC为辅因子），促进前胶原中脯氨酸和赖氨酸的羟化——是构成结缔组织、血管壁、骨骼基质的基础
        \item 参与胆固醇的羟化：促进胆固醇向胆汁酸的转化
    \end{itemize}
    \item \textbf{在铁的吸收、转运和贮备、叶酸转变为四氢叶酸等方面发挥重要作用：}
    \begin{itemize}[leftmargin=*]
        \item 将Fe$^{3+}$还原为Fe$^{2+}$，并与之形成可溶螯合物促进铁的吸收
        \item 促进叶酸转化为四氢叶酸（THF），维持叶酸的活性形式
    \end{itemize}
    \item \textbf{增强免疫功能：}促进抗体生成、增强白细胞吞噬功能
\end{enumerate}
\end{keybox}

\subsection{维生素C缺乏——坏血病（Scurvy）}

\begin{keybox}
\textbf{维生素C缺乏症——坏血病（Scurvy）*：}

因胶原蛋白合成障碍，导致结缔组织和血管壁损伤。

\textbf{早期症状：}
\begin{itemize}[leftmargin=*]
    \item 牙龈疼痛出血、鼻出血
    \item 无力、食欲减退
    \item 皮肤干燥
    \item 伤口愈合不良
    \item 血管脆性增加
\end{itemize}

\textbf{严重缺乏：}体内不断出血，导致死亡

维生素C缺乏潜伏期约3$\sim$4个月（因体内有少量储存）。
\end{keybox}

\subsection{过量、营养状况评价与食物来源}

\begin{tipbox}
\textbf{维生素C过量表现：}
\begin{itemize}[leftmargin=*]
    \item 维生素C毒性很低，但一次口服2$\sim$8g时可能会出现腹泻、腹胀
    \item 患有结石的病人，摄入量过多时可能增加尿中草酸盐的排泄，增加尿路结石的危险
\end{itemize}

\textbf{机体营养状态评价：}
\begin{itemize}[leftmargin=*]
    \item 负荷试验
    \item 血浆维生素C含量
    \item 白细胞中抗坏血酸浓度
\end{itemize}
\end{tipbox}

\begin{keybox}
\textbf{维生素C的膳食参考摄入量：}
\begin{itemize}[leftmargin=*]
    \item 成人 RNI：\imp{100 mg/d}
    \item 孕妇早期：100 mg/d；中、晚期：115 mg/d
    \item 乳母：150 mg/d
    \item UL（可耐受最高摄入量）：1000 mg/d
\end{itemize}
\end{keybox}

\begin{exambox}
\textbf{食物来源：}
\begin{itemize}[leftmargin=*]
    \item \textbf{水果类：}苜蓿、刺梨、沙棘、柚子、枣、红果、猕猴桃、酸枣、柑桔、草莓、橙子等
    \item \textbf{蔬菜类：}新鲜蔬菜（绿色和红、黄色的辣椒、菠菜、西红柿、韭菜等）
\end{itemize}
注意：维生素C极易热破坏和氧化损失，因此\imp{新鲜生食为最佳摄取方式}。
\end{exambox}

% ==================== 第十三节 其他维生素（简略） ====================
\section{其他维生素}

\subsection{泛酸（Pantothenic Acid，维生素B$_5$）}
是\keyword{辅酶A（CoA）}和酰基载体蛋白（ACP）的组成成分，在糖、脂肪和蛋白质代谢中起核心作用。广泛存在于各类食物中（Pan-thenic意为"到处存在"），缺乏罕见。成人AI：5.0 mg/d。

\subsection{生物素（Biotin，维生素B$_7$）}
是多种羧化酶（丙酮酸羧化酶、乙酰辅酶A羧化酶等）的辅酶，参与CO$_2$固定反应。缺乏极为罕见，除非长期摄入大量生鸡蛋清（含\keyword{抗生物素蛋白（Avidin）}，可与生物素紧密结合无法吸收，导致"蛋白伤害"）。成人AI：40 μg/d。

% ==================== 第十四节 类维生素 ====================
\section{类维生素}

\subsection{肉碱（Carnitine）}
在肝和肾中由赖氨酸和蛋氨酸合成。核心功能是作为长链脂肪酸进入线粒体基质进行$\beta$-氧化的载体（"脂肪酸运输穿梭系统"）。成人AI：200 mg/d。食物来源：红肉最丰富。

\subsection{牛磺酸（Taurine）}
含硫$\beta$-氨基酸，游离存在于组织中。功能：结合胆汁酸、视网膜发育、渗透压调节、抗氧化。对早产儿是条件必需营养素。

\subsection{肌醇（Inositol）}
参与细胞信号转导（肌醇磷脂途径），促进脂肪代谢。成人AI：1000 mg/d。来源：水果、豆类、谷类。

\subsection{柠檬素（Citrin / Bioflavonoids，维生素P）}
是一类具有抗氧化活性的黄酮类化合物，常与维生素C共存于自然界中。主要功能：增强毛细血管壁的韧性和通透性调节能力，协同维生素C增强其抗氧化作用；抑制透明质酸酶活性。缺乏可导致毛细血管脆性增加（易出现紫癜）。食物来源：柑橘类水果（柠檬、橙子、柚子）、葡萄、樱桃、青椒等富含维生素C的果蔬。目前尚无官方AI值，建议每日摄入量为150$\sim$200 mg。

% ==================== 第十五节 水 ====================
\section{水}

\subsection{水的分布}

人体总含水量约占体重的\imp{50\%$\sim$70\%}，随年龄增大含水量下降、随体脂增加含水量下降。
\begin{itemize}[leftmargin=*]
    \item 细胞内液：约占总体水量的 2/3
    \item 细胞外液：约占总体水量的 1/3（包括血液、组织间液、淋巴液等）
\end{itemize}

\subsection{水的生理功能}

\begin{enumerate}[leftmargin=*]
    \item \textbf{构成人体组织：}是人体含量最多的成分
    \item \textbf{作为溶剂和运载工具：}运送营养物质和代谢产物
    \item \textbf{参与物质代谢：}直接参与水解、水合等化学反应
    \item \textbf{调节体温：}通过排汗散热维持体温恒定
    \item \textbf{润滑作用：}关节液、胸腔液、腹腔液等起缓冲保护作用
\end{enumerate}

\subsection{水的平衡}

成人每日水需要量约\imp{2500 mL}。来源：饮水约1200mL、食物水约1000mL、内生水（代谢水）约300mL。排出：尿约1500mL、皮肤蒸发约500mL、肺呼出约350mL、粪便约150mL。

% ----- 复习题 -----
\reviewcontent{
\section{复习题}

\begin{enumerate}[leftmargin=*, itemsep=8pt]
    \item \textbf{【名词解释】}维生素、脂溶性维生素、水溶性维生素、视黄醇当量（RE）、$\beta$-胡萝卜素、夜盲症、佝偻病、骨化三醇1,25-(OH)$_2$D$_3$、脚气病、TPP、口腔-生殖系统综合征、烟酸当量（NE）、癞皮病（Pellagra）三D症状、膳食叶酸当量（DFE）、巨幼红细胞性贫血、坏血病（Scurvy）
    \item \textbf{【简答题】}简述维生素的共同特点。
    \item \textbf{【简答题】}比较脂溶性维生素和水溶性维生素的特点（至少4条）。
    \item \textbf{【简答题】}简述维生素缺乏发展的四个阶段。
    \item \textbf{【简答题】}简述维生素A的生理功能、缺乏的表现（由轻到重）及主要食物来源。
    \item \textbf{【简答题】}写出视黄醇当量（RE）的计算公式。
    \item \textbf{【简答题】}简述维生素D的来源、活化过程及生理功能。
    \item \textbf{【简答题】}简述维生素B$_1$的理化性质、TPP参与的酶系统和脚气病的三型表现。
    \item \textbf{【简答题】}简述维生素B$_2$的理化性质（重点：紫外光敏感性）及口腔-生殖系统综合征的临床表现。
    \item \textbf{【简答题】}简述烟酸（维生素B$_3$/PP）的活性形式及癞皮病（Pellagra）的典型"三D"症状。
    \item \textbf{【简答题】}写出烟酸当量（NE）和膳食叶酸当量（DFE）的计算公式。
    \item \textbf{【简答题】}简述维生素C的理化性质、生理功能及缺乏病（坏血病）的典型表现。
    \item \textbf{【简答题】}比较叶酸缺乏与维生素B$_{12}$缺乏引起的贫血有何异同？
\end{enumerate}

\subsection*{参考答案要点}

\begin{enumerate}[leftmargin=*, itemsep=5pt]
    \item 见本章各概念定义框。
    \item ①以本体或前体形式存在于天然食物中，必须经常由食物提供；②需要量甚微（mg或$\mu$g水平），但在调节机体代谢方面作用重要；③不构成组织，也不提供能量；④多以辅酶或辅基的形式发挥功能；⑤有的具有几种结构相近、活性相同的化合物。
    \item 脂溶性：需脂肪和胆汁协助吸收、可体内储存（肝和脂肪组织）、缺乏症状出现慢、摄入过量可引起中毒（VitK除外）、除VitK外均有UL限制。水溶性：直接吸收、体内储存极少（B$_{12}$除外）、缺乏症状出现快、一般不引起中毒（多余从尿排泄）、饱和剂量以上几乎不增加组织浓度。
    \item ①组织中维生素储量降低（亚临床/边缘缺乏）；②生化指标异常，生理功能降低；③组织病理变化，出现临床症状和体征；④停止生命（极端严重缺乏可致死）。
    \item 生理功能：①维持正常视觉（视紫红质）；②维持上皮的正常生长与分化；③促进生长发育；④抑癌作用；⑤维持机体正常免疫功能。缺乏（由轻到重）：暗适应能力下降$\rightarrow$夜盲症$\rightarrow$干眼病/上皮角化$\rightarrow$角膜软化溃疡$\rightarrow$失明。\textbf{食物来源：}动物肝脏、奶制品、鸡蛋、河蟹（已形成的VitA）；胡萝卜、菠菜、西兰花、芒果等（$\beta$-胡萝卜素）。
    \item 视黄醇当量（$\mu$g RE）= 视黄醇（$\mu$g）+ $\beta$-胡萝卜素（$\mu$g）$\times$ 0.167。即1 $\mu$g $\beta$-胡萝卜素 = 0.167 $\mu$g RE = 1/6 $\mu$g RE。
    \item 来源：VitD$_2$（植物，麦角固醇+UV）、VitD$_3$（皮肤，7-脱氢胆固醇+UVB 290$\sim$315nm）。活化：肝$\rightarrow$25-(OH)D$_3$$\rightarrow$肾$\rightarrow$1,25-(OH)$_2$D$_3$（活性形式，类固醇激素）。生理功能：促进小肠钙吸收、促进肾小管对钙磷重吸收、对骨细胞呈现多种作用（调节骨钙化与脱钙）、调节基因转录作用、维持血钙平衡。
    \item \textbf{理化性质：}盐酸盐和硝酸盐在干燥和酸性溶液中稳定，熔点249℃分解；碱性环境中易氧化失活；对亚硫酸盐极为敏感。\textbf{TPP参与：}羧化酶（$\alpha$-酮酸氧化脱羧）和转酮醇酶（二碳单位转运）。\textbf{脚气病三型：}干性（周围神经炎、感觉异常）；湿性（心血管系统损害、水肿、心衰）；急性恶性（胸闷、心悸、气短、内出血）。
    \item \textbf{理化性质：}比较稳定，耐酸耐热，不易氧化；碱性溶液中易破坏；游离型核黄素对紫外光高度敏感（酸性$\rightarrow$光黄素，碱性$\rightarrow$光色素，均丧失生物活性）。\textbf{活性形式：}FMN、FAD。\textbf{口腔-生殖系统综合征：}口角炎、唇炎、舌炎、睑缘炎、结膜炎、脂溢性皮炎、阴囊皮炎。
    \item 活性形式：辅酶I（NAD$^+$）和辅酶II（NADP$^+$），为多种脱氢酶的辅酶。体内来源：色氨酸可在体内转化为烟酸（约60mg色氨酸$\rightarrow$1mg烟酸）。癞皮病三D：Dermatitis（皮炎）——暴露部位对称性皮炎、腹泻（Diarrhea）、痴呆（Dementia）。
    \item \textbf{NE (mg)} = 烟酸 (mg) + 色氨酸 (mg)/60。\textbf{DFE ($\mu$g)} = 天然食物中叶酸 ($\mu$g) + 1.7 $\times$ 补充剂/强化食品中叶酸 ($\mu$g)。
    \item \textbf{理化性质：}又名抗坏血酸，含6碳$\alpha$-酮基内酯弱酸，强还原剂；最不稳定，有氧或碱性环境极易氧化。\textbf{生理功能：}抗氧化、作为羟化过程底物和酶的辅因子（维持胶原蛋白正常功能、参与胆固醇羟化）、促进铁吸收和叶酸活化、增强免疫功能。\textbf{坏血病：}牙龈疼痛出血、鼻出血、皮下瘀斑、伤口愈合不良、血管脆性增加；严重$\rightarrow$体内不断出血致死。
    \item \textbf{同：}均可引起巨幼红细胞性贫血。\textbf{异：}VitB$_{12}$缺乏伴有神经系统损害（脊髓后索、侧索退行性变/亚急性联合变性），而叶酸缺乏不引起神经病变；VitB$_{12}$仅存在于动物性食品中，需要内因子（IF）协助吸收。
\end{enumerate}

\newpage
}

% =====================================================================
%  CHAPTER 2: 食物中的生物活性成分
% =====================================================================
\chapter{食物中的生物活性成分}
\setcounter{chapter}{2}
\chapterintro{
\keyword{掌握}植物化学物的概念、分类和生物学作用；掌握多酚类化合物、有机硫化物、植物雌激素的生物学作用；掌握SPL（特定建议值）的概念及其与营养素DRIs的区别。\keyword{熟悉}类胡萝卜素、植物固醇的作用；生物活性物质与相关疾病。\keyword{了解}其他生物活性成分。
}

% ----- 第一节 植物化学物概述 -----
\section{植物化学物概述}

\begin{keybox}
\textbf{植物化学物（Phytochemicals）：}来自植物性食物的生物活性成分，是植物能量代谢过程中产生的多种中间或末端低分子量次级代谢产物。除个别是维生素的前体物（$\beta$-胡萝卜素）外，其余均为非传统营养素成分。

\textbf{分类（10类）：}①类胡萝卜素；②植物固醇；③皂苷类化合物；④芥子油苷；⑤多酚类化合物；⑥蛋白酶抑制剂；⑦单萜类；⑧植物雌激素；⑨有机硫化物；⑩植酸。

\textbf{生物活性（6种）：}①抑制肿瘤；②抗氧化；③免疫调节；④抑制微生物；⑤降低胆固醇；⑥其他：调节血压血糖、血小板血凝，抑制炎症等。
\end{keybox}

\begin{defbox}
\textbf{生物活性成分（Bioactive Food Components）：}食物中除了营养素以外的，对人体健康有益的微量非营养素成分。过去曾被称为"非营养素活性成分"，现已统一称为"食物中的生物活性成分"。按来源可分为植物来源（植物化学物）和动物来源两大类。
\end{defbox}

% ----- 第二节 各类植物化学物分述 -----
\section{各类植物化学物分述}

\subsection{类胡萝卜素（Carotenoids）}

\begin{defbox}
\textbf{类胡萝卜素：}广泛存在于果蔬中的脂溶性植物色素，迄今已发现700余种，约50种存在于人类膳食中。分为两大类：
\begin{itemize}[leftmargin=*]
    \item \textbf{胡萝卜素类（Carotenes）：}纯碳氢化合物，不含氧。包括$\alpha$-胡萝卜素、$\beta$-胡萝卜素、番茄红素（Lycopene）。
    \item \textbf{叶黄素类（Xanthophylls）：}含氧衍生物。包括叶黄素（Lutein）、玉米黄素（Zeaxanthin）、$\beta$-隐黄质。
\end{itemize}

\textbf{主要作用：}抗氧化（番茄红素淬灭单线态氧能力最强，是$\beta$-胡萝卜素的2倍以上）；维生素A原活性（$\beta$-胡萝卜素最高，1分子$\rightarrow$2分子视黄醇；番茄红素、叶黄素、玉米黄素无此活性）；抗肿瘤；免疫调节；视觉保护（叶黄素和玉米黄素构成视网膜黄斑色素，滤过蓝光，降低年龄相关性黄斑变性AMD风险）。
\end{defbox}

\subsection{植物固醇（Phytosterols）}

\begin{defbox}
\textbf{植物固醇：}存在于植物细胞膜中的固醇类化合物，结构与胆固醇类似（均为环戊烷多氢菲骨架，C17侧链有差异）。人体不能合成，全部来自膳食。主要种类：\keyword{$\beta$-谷固醇}（含量最丰富，约占50\%\~{}65\%）、豆固醇、菜油固醇。

\textbf{降胆固醇作用（最明确的核心功能）：}与胆固醇在肠道竞争进入混合微胶粒，减少胆固醇的溶解和吸收。每日摄入2\~{}3g可使血清LDL降低5\%\~{}15\%，是公认的辅助降脂功能性成分。\textbf{食物来源：}植物油、坚果和种子、豆类、全谷物。
\end{defbox}

\subsection{皂苷类化合物（Saponins）}

皂苷是一类具有表面活性（在水中振摇可产生持久泡沫，类似肥皂）的糖苷类化合物，由苷元（甾体或三萜类）和糖组成。

\textbf{主要来源：}大豆、鹰嘴豆、花生、菠菜、芦笋、燕麦、人参、甘草。

\textbf{主要作用：}调节脂质代谢（与胆固醇和胆汁酸结合降低肠道吸收）；抗菌/抗病毒（破坏微生物细胞膜）；抗肿瘤；抗血栓形成；免疫调节（人参皂苷可激活巨噬细胞和NK细胞）。注意：皂苷具有溶血活性，但经口摄入后胃肠道吸收极少，通常不造成全身性危害。

\subsection{芥子油苷（Glucosinolates）}

芥子油苷（硫代葡萄糖苷）主要存在于\keyword{十字花科蔬菜}（西兰花、卷心菜、白菜、甘蓝、萝卜、芥菜等）。植物细胞破碎后经黑芥子酶（Myrosinase）水解，生成具有生物活性的\keyword{异硫氰酸盐（ITCs）}，如\keyword{萝卜硫素（Sulforaphane）}。

\textbf{抗肿瘤作用——核心功能：}萝卜硫素是已知最强的天然Nrf2通路激活剂之一，强力诱导II相解毒酶（GST、UGT、NQO1）表达；流行病学证据支持降低多种癌症风险（肺癌、结肠癌、前列腺癌、乳腺癌）。此外还有调节炎症反应（抑制NF-$\kappa$B通路）和抗菌（抑制幽门螺杆菌）作用。

\subsection{多酚类化合物（Polyphenols）}

\begin{keybox}
\textbf{多酚类化合物：}植物性食物中含量最丰富、研究最深入的植物化学物类别。基本结构为2-苯基色原酮核心骨架，即一个或多个芳香环连接多个酚羟基。按碳骨架结构分为七大类：

\begin{enumerate}[leftmargin=*]
    \item \textbf{黄酮及黄酮醇类：}代表为\keyword{槲皮素}（膳食中最丰富的类黄酮之一）、芦丁、山奈酚。来源：洋葱、西兰花、蓝莓、苹果、红茶。
    \item \textbf{黄烷酮类：}代表为柚皮素、橙皮素。来源：柑橘类水果。
    \item \textbf{黄烷醇类（儿茶素类）：}代表为\keyword{儿茶素}、表儿茶素、EGCG。来源：绿茶（含量最丰富）、巧克力、苹果。
    \item \textbf{异黄酮类：}代表为大豆苷、\keyword{染料木素}、葛根素。来源：大豆及其制品。
    \item \textbf{双黄酮类：}代表为银杏黄酮。来源：银杏叶。
    \item \textbf{花色素类：}代表为矢车菊素、天竺葵素。来源：浆果、紫薯、紫甘蓝。
    \item \textbf{查尔酮类：}来源：苹果、甘草。
\end{enumerate}

\textbf{主要生物学作用：}①抗氧化——膳食最重要的抗氧化活性成分来源，直接清除自由基并通过激活Nrf2/ARE通路诱导抗氧化酶（SOD、CAT、GPx、HO-1）；②抗肿瘤——茶多酚（EGCG）和大豆异黄酮证据最充分（诱导II相解毒酶、抑制I相代谢酶、阻滞细胞周期、诱导凋亡、抑制血管生成）；③心血管保护——改善内皮功能、降低血压、抑制LDL氧化、抗血小板聚集；④抗炎——抑制NF-$\kappa$B通路和COX-2、iNOS；⑤调节肠道菌群——促进有益菌增殖，抑制致病菌。
\end{keybox}

\subsection{蛋白酶抑制剂（Protease Inhibitors）}

广泛存在于植物（特别是豆类和谷物）中的一类能与蛋白水解酶结合并抑制其活性的蛋白质或多肽。典型代表：大豆胰蛋白酶抑制剂（Bowman-Birk抑制剂BBI；Kunitz胰蛋白酶抑制剂KTI）。主要作用：抗肿瘤（抑制肿瘤相关蛋白酶活性、抑制自由基生成、调节癌基因表达）；免疫调节（BBI可抑制炎症性蛋白酶如弹性蛋白酶）。

蛋白酶抑制剂历来被视为"抗营养因子"（抑制胰蛋白酶、降低蛋白质消化率），但充分加热可使其彻底灭活，且现代研究发现在低浓度下具有抗癌防炎的健康效益。

\subsection{单萜类（Monoterpenes）}

由两个异戊二烯单元（C$_{10}$骨架）组成的萜类化合物，是植物精油的主要成分。代表性物质：紫苏醇（Perillyl Alcohol）、柠檬烯（Limonene）。主要来源：柑橘类水果果皮（柠檬烯）、薄荷油、香菜籽。

主要作用：抗肿瘤——紫苏醇和柠檬烯可抑制Ras蛋白异戊二烯化（Ras为GTP结合癌蛋白，需法尼基化才能定位到细胞膜发挥致癌效应），在动物模型中显示抑制乳腺、肝脏、胰腺肿瘤。

\subsection{植物雌激素（Phytoestrogens）}

\begin{keybox}
\textbf{植物雌激素：}植物来源的、能与哺乳动物雌激素受体（ER）结合并发挥类雌激素或抗雌激素效应的非甾体化合物。核心特征为\imp{双向调节作用}——低雌激素水平时激动ER（雌激素样作用），高雌激素水平时拮抗ER（竞争性抑制内源性雌二醇），使其成为SERM（选择性雌激素受体调节剂）的天然类似物。

\textbf{分类与来源：}
\begin{enumerate}[leftmargin=*]
    \item \textbf{异黄酮类（Isoflavones）：}大豆苷、染料木素、大豆黄素。\imp{主要来源：大豆及其制品}（豆粉、豆腐、豆浆），大豆中异黄酮含量约为0.1\%\~{}0.5\%。
    \item \textbf{木酚素类（Lignans）：}开环异落叶松脂素、罗汉松脂素。在肠道菌群作用下转化为肠内酯和肠二醇活性形式。\imp{亚麻籽是木酚素含量最高的食物}，也含于全谷物、种子类。
    \item \textbf{香豆素类（Coumestans）：}香豆雌酚。来源：发芽豆类（苜蓿芽、绿豆芽）。
    \item \textbf{芪类（Stilbenes）：}代表为\keyword{白藜芦醇}（Resveratrol）。来源：葡萄/葡萄酒、花生、虎杖。
\end{enumerate}

\textbf{主要生物学作用：}①预防骨质疏松——通过SERM-like作用促进成骨细胞活性、抑制破骨细胞分化；②改善围绝经期症状（缓解潮热、盗汗等）；③抗氧化；④心血管保护；⑤抗肿瘤——高大豆摄入人群（亚洲）的乳腺癌、前列腺癌发病率较低。
\end{keybox}

\subsection{有机硫化物（Organosulfur Compounds）}

\begin{keybox}
\textbf{有机硫化物：}主要存在于\keyword{葱属植物（Allium）}中的含硫植物化学物，以大蒜含量最为丰富，是其特征风味物质的来源。典型成分：蒜素（Allicin）——完整蒜瓣切碎或压碎后由蒜氨酸（Alliin）经蒜氨酸酶催化生成，在室温下不稳定，可进一步分解为二烯丙基二硫化物（DADS）、二烯丙基三硫化物（DATS）、阿霍烯（Ajoene）等。

\textbf{主要生物学作用：}
\begin{enumerate}[leftmargin=*]
    \item \textbf{抗菌作用：}蒜素具有广谱抗菌活性，杀菌效力约相当于1\%青霉素，对革兰阳性和阴性细菌、真菌、寄生虫均有抑制作用。细菌不易对蒜素产生耐药性。
    \item \textbf{抗氧化作用：}直接清除自由基，并诱导抗氧化酶的表达。
    \item \textbf{调节血脂：}降低血清TC、TG和LDL，升高HDL。
    \item \textbf{抗血栓形成：}抑制血小板聚集，有一定纤溶促进作用（类似阿司匹林但较弱）。
    \item \textbf{抗肿瘤作用：}阻断\imp{亚硝胺}的内源性合成——这是大蒜防癌的重要机制之一；诱导II相解毒酶表达，加速致癌物灭活和排泄。
\end{enumerate}

\textbf{食用建议：}蒜素在高温下易被破坏，建议大蒜切碎后室温放置10\~{}15分钟（让蒜氨酸酶充分作用生成蒜素），再用于凉拌或低温快炒。
\end{keybox}

\subsection{植酸（Phytic Acid, IP6）}

\begin{defbox}
\textbf{植酸（肌醇六磷酸, IP6）：}豆类、全谷物和坚果中天然存在的肌醇磷酸酯。含有六个磷酸基团，对二价矿物离子（Ca$^{2+}$、Fe$^{2+}$/Fe$^{3+}$、Zn$^{2+}$、Mg$^{2+}$）具有强烈的螯合能力。

\textbf{认识的历史演变：}
\begin{enumerate}[leftmargin=*]
    \item \textbf{传统认识（抗营养因子）：}在肠道螯合矿物质，降低矿物质的生物利用度，是发展中国家以谷物为主食人群矿物质缺乏（尤其是铁、锌）的重要影响因素。
    \item \textbf{现代认识（有益生物活性成分）：}具有抗氧化（铁螯合阻断Fenton反应产生自由基）和抗肿瘤（抑制增殖、诱导分化）的生物学作用，被认为是潜在的抗癌活性物质。
\end{enumerate}
\end{defbox}

% ----- 第三节 生物活性物质与相关疾病 -----
\section{生物活性物质与相关疾病}

\begin{defbox}
\textbf{植物化学物与慢性病防治：}大量流行病学研究和临床试验表明，富含植物化学物的膳食模式（多蔬果、全谷物、大豆、坚果）与多种慢性非传染性疾病（NCD）的风险降低显著相关。

\textbf{主要作用机制：}
\begin{itemize}[leftmargin=*]
    \item \textbf{抗肿瘤：}多途径协同——诱导II相解毒酶（萝卜硫素、蒜素）、抑制I相致癌物代谢酶（多酚类）、阻滞细胞周期（多酚类）、诱导凋亡（多酚类、皂苷）、抑制血管生成（多酚类）、阻断致癌物合成（有机硫化物阻断亚硝胺合成）、抑制Ras蛋白异戊二烯化（单萜类）。
    \item \textbf{心血管保护：}降胆固醇（植物固醇竞争吸收、皂苷结合胆汁酸、有机硫化物抑制合成酶）、抗氧化/抗炎（多酚类抑制LDL氧化和NF-$\kappa$B通路）、抗血小板聚集（有机硫化物、皂苷）、改善内皮功能。
    \item \textbf{代谢调节：}调节血糖（多酚类、硫辛酸）、改善胰岛素敏感性。
    \item \textbf{骨骼健康：}植物雌激素通过SERM样作用抑制骨丢失。
    \item \textbf{视觉保护：}叶黄素/玉米黄素保护视网膜，降低AMD风险。
\end{itemize}
\end{defbox}

% ----- 第四节 SPL -----
\section{特定建议值（SPL）}

\begin{keybox}
\textbf{特定建议值（Specific Proposed Levels, SPL）：}SPL是DRIs指标体系的第七项，专门针对\textbf{非营养素生物活性成分}（主要为植物化学物）而设立。其定义为：以降低成年人膳食相关非传染性慢性病（NCD）风险为目标的\textbf{其他膳食成分}每日摄入量。

\textbf{SPL与营养素DRIs的本质区别：}
\begin{enumerate}[leftmargin=*]
    \item \textbf{适用对象不同：}SPL针对非必需的非营养素活性成分（植物化学物），DRIs的EAR/RNI/AI/UL针对必需营养素。
    \item \textbf{底层逻辑不同：}SPL回答"摄入多少有助于降低慢病风险"（促健康导向），RNI回答"需要多少满足基本生理需求"（防缺乏导向），UL回答"多少以内安全"。
    \item \textbf{不存在EAR和RNI：}植物化学物为非必需、非缺乏性物质，因此无EAR（无"满足50\%个体"的概念）和RNI。
\end{enumerate}

\textbf{已提出SPL的植物化学物举例（中国DRIs 2023版）：}
\begin{itemize}[leftmargin=*]
    \item \textbf{大豆异黄酮：}SPL以苷元计，基于改善围绝经期症状和预防骨质疏松的证据。
    \item \textbf{植物固醇：}SPL基于降低血清LDL胆固醇的证据。
    \item \textbf{原花青素、花色苷、番茄红素、叶黄素等：}亦有相应SPL值被提出或正在论证中。
\end{itemize}

\textbf{意义：}SPL的设立标志着营养科学从"预防缺乏"向"促进健康和预防慢病"的范式转变，承认了植物化学物在膳食营养指导中的独立地位，为膳食指南中"多吃蔬果、全谷物、大豆"的建议提供了剂量化依据。
\end{keybox}

% ----- 第五节 动物来源的生物活性成分 -----
\section{动物来源的生物活性成分（了解）}

\begin{itemize}[leftmargin=*]
    \item \textbf{辅酶Q（泛醌，CoQ10）：}呼吸链电子传递体，参与ATP合成；脂溶性抗氧化剂（保护细胞膜和LDL）；心血管保护。
    \item \textbf{硫辛酸：}"万能抗氧化剂"兼有水溶性和脂溶性，能透过血脑屏障。再生GSH和维生素C/E；改善胰岛素敏感性（II型糖尿病辅助治疗）；神经保护。
    \item \textbf{褪黑素：}主要由松果体合成。调节生物节律（昼夜节律和睡眠周期）；强效抗氧化剂（穿越所有生物屏障）。
    \item \textbf{$\gamma$-氨基丁酸（GABA）：}抑制性神经递质，降血压、镇静安神。
    \item \textbf{左旋肉碱（L-Carnitine）：}脂肪酸$\beta$-氧化的转运载体（将长链脂肪酸转运至线粒体基质）。来源：红肉、奶制品。
\end{itemize}

% ----- 复习题 -----
\reviewcontent{
\section{复习题}

\begin{enumerate}[leftmargin=*, itemsep=8pt]
    \item \textbf{【名词解释】}植物化学物（Phytochemicals）；植物雌激素的双向调节作用；SPL（特定建议值）
    \item \textbf{【简答题】}列出植物化学物的10个分类。
    \item \textbf{【简答题】}列出植物化学物的6种生物活性。
    \item \textbf{【简答题】}简述多酚类化合物的七大类别及其主要生物学作用。
    \item \textbf{【简答题】}简述有机硫化物的主要生物学作用及食用建议。
    \item \textbf{【简答题】}试述植物雌激素的分类、来源及其双向调节机制。
    \item \textbf{【简答题】}简述SPL与营养素DRIs（RNI、UL）的本质区别。
    \item \textbf{【简答题】}简述植酸从"抗营养因子"到"潜在抗癌活性物质"的认识转变。
    \item \textbf{【论述题】}试述植物化学物在慢性非传染性疾病（NCD）防治中的作用及机制。
\end{enumerate}

\subsection*{参考答案要点}

\begin{enumerate}[leftmargin=*, itemsep=5pt]
    \item \textbf{植物化学物：}来自植物性食物的生物活性成分，是植物能量代谢过程中产生的多种中间或末端低分子量次级代谢产物。除个别是维生素的前体物（$\beta$-胡萝卜素）外，其余均为非传统营养素成分。\textbf{双向调节作用：}低雌激素水平时激动ER（雌激素样作用），高雌激素水平时拮抗ER（竞争性抑制内源性雌二醇）。\textbf{SPL：}DRIs第七项指标，针对非营养素生物活性成分，以降低成人膳食相关NCD风险为目标的每日摄入量。
    \item ①类胡萝卜素；②植物固醇；③皂苷类化合物；④芥子油苷；⑤多酚类化合物；⑥蛋白酶抑制剂；⑦单萜类；⑧植物雌激素；⑨有机硫化物；⑩植酸。
    \item ①抑制肿瘤；②抗氧化；③免疫调节；④抑制微生物；⑤降低胆固醇；⑥其他：调节血压血糖、血小板血凝，抑制炎症等。
    \item \textbf{七大类：}①黄酮及黄酮醇类（槲皮素）；②黄烷酮类（柚皮素）；③黄烷醇类/儿茶素类（EGCG）；④异黄酮类（染料木素）；⑤双黄酮类（银杏黄酮）；⑥花色素类（矢车菊素）；⑦查尔酮类。\textbf{主要作用：}抗氧化（激活Nrf2/ARE通路诱导抗氧化酶）、抗肿瘤（诱导II相解毒酶、诱导凋亡、抑制血管生成）、心血管保护、抗炎（抑制NF-$\kappa$B）、调节肠道菌群。
    \item \textbf{作用：}①抗菌——蒜素广谱抗菌，效力约相当于1\%青霉素，不易产生耐药性；②抗氧化；③调脂——降低TC/TG/LDL，升高HDL；④抗血栓——抑制血小板聚集；⑤抗肿瘤——阻断亚硝胺内源性合成、诱导II相解毒酶。\textbf{食用建议：}大蒜切碎后室温放置10\~{}15分钟（让蒜氨酸酶充分作用生成蒜素），避免高温长时间烹饪。
    \item \textbf{分类：}异黄酮类（大豆苷、染料木素）、木酚素类（亚麻籽含量最高，经肠道菌群转化为肠内酯/肠二醇）、香豆素类（香豆雌酚，发芽豆类）、芪类（白藜芦醇，葡萄/葡萄酒）。\textbf{机制：}低雌激素水平时结合ER发挥雌激素样作用，高雌激素水平时竞争性抑制内源性雌二醇发挥抗雌激素效应，是天然的SERM。作用包括预防骨质疏松、改善围绝经期症状、抗肿瘤。
    \item ①SPL针对非必需的非营养素（植物化学物），RNI/UL针对必需营养素；②SPL回答"摄入多少有助降低慢病风险"，RNI回答"需要多少满足生理需求"，UL回答"多少以内安全"；③SPL无EAR（无需满足50\%个体），仅基于流行病学剂量-反应关系；④SPL标志着从"预防缺乏"到"促健康防慢病"的范式转变。
    \item \textbf{传统认识：}螯合二价矿物质（Ca$^{2+}$、Fe$^{2+}$/Fe$^{3+}$、Zn$^{2+}$、Mg$^{2+}$），降低矿物质生物利用度，是发展中国家以谷物为主食人群矿物质缺乏的重要影响因素（抗营养因子）。\textbf{现代认识：}具有抗氧化（铁螯合阻断Fenton反应产生自由基）和抗肿瘤（抑制增殖、诱导分化）的生物学作用，被认为是潜在的抗癌活性物质（有益生物活性成分）。
    \item \textbf{抗肿瘤：}多途径协同——诱导II相解毒酶（萝卜硫素激活Nrf2、有机硫化物）、抑制I相致癌物代谢酶（多酚类）、阻滞细胞周期和诱导凋亡（多酚类、皂苷类）、抑制血管生成（多酚类）、阻断致癌物合成（有机硫化物阻断亚硝胺）、抑制Ras蛋白异戊二烯化（单萜类）。\textbf{心血管保护：}降胆固醇（植物固醇竞争吸收、皂苷结合胆汁酸、有机硫化物抑制合成酶）、抗氧化抗炎（多酚类抑制LDL氧化和NF-$\kappa$B）、抗血小板聚集（有机硫化物、皂苷）、改善内皮功能。\textbf{其他：}植物雌激素预防骨质疏松（SERM样作用）、叶黄素/玉米黄素保护视觉、多酚类调节肠道菌群等。
\end{enumerate}

\newpage
}

% =====================================================================
%  CHAPTER 3: 各类食物的营养价值
% =====================================================================
\chapter{各类食物的营养价值}
\setcounter{chapter}{3}
\chapterintro{
掌握各类食物的营养特点（谷类、大豆、蔬果、畜禽水产品、蛋类、乳类、坚果）。熟悉食物营养价值的定义及影响因素。了解食物成分数据库。
}

% ----- 第一节 食物营养价值的定义及影响因素 -----
\section{食物营养价值的定义及影响因素}

\begin{defbox}
\textbf{食物营养价值（Nutritional Value of Food）：}食物中所含营养素和能量满足人体营养需要的程度。营养价值的高低取决于食物中营养素种类是否齐全、数量是否充足、相互比例是否适宜，以及是否易于被消化、吸收和利用。
\end{defbox}

\subsection{食物营养价值评价的维度}

\begin{enumerate}[leftmargin=*]
    \item \textbf{营养素种类与含量：}营养素种类齐全、含量丰富者营养价值较高。
    \item \textbf{营养素质量：}主要看该营养素的构成模式与人体的接近程度（如蛋白质氨基酸模式）。
    \item \textbf{消化吸收利用率：}即使含量丰富，若消化吸收差，实际营养价值也有限。
    \item \textbf{功能成分/生物活性物质：}不仅满足基本需求，更关注维护健康、预防疾病的额外价值。
\end{enumerate}

\subsection{影响食物营养价值的因素}

\begin{keybox}
\textbf{影响食物营养素组成和含量的主要因素：}
\begin{enumerate}[leftmargin=*]
    \item \textbf{品种、产地和季节：}不同品种的同种食物营养素差异大（如不同品种大米蛋白质含量差3\%\~{}4\%）；地理环境和气候条件影响产地；收获季节决定成熟度。
    \item \textbf{栽培和饲养条件：}土壤肥沃程度、施肥种类、饲料质量均显著影响农产品营养素水平。
    \item \textbf{收获成熟度：}未充分成熟的一般营养价值较低，但过度成熟可能导致营养素流失或口感变差。
    \item \textbf{加工和烹调：}加工精度（如碾米精度对B族维生素的影响极大）、热加工的温度和时间、加工方式（水煮损失水溶性维生素，油炸增加脂肪并损失热不稳定维生素）。
    \item \textbf{储藏条件：}温度、湿度、时间、光照、氧气影响营养素（尤其是维生素C、维生素E、$\beta$-胡萝卜素等抗氧化性维生素）的保留。
\end{enumerate}
\end{keybox}

\begin{tipbox}
\textbf{食物成分数据库：}我国《中国食物成分表》由营养与食品安全所主编，是国民营养与健康状况评价、膳食指导、食物营养标签等的基础性工具。包括原料型食物和加工型食物两大板块。
\end{tipbox}

% ----- 第二节 谷类、薯类及杂豆 -----
\section{谷类、薯类及杂豆}

\begin{keybox}
\textbf{谷类（Cereals）：}我国居民膳食的\keyword{主食}，提供约50\%\~{}70\%的能量和蛋白质。主要包括稻米（大米）、小麦（面粉）、玉米、小米、燕麦、荞麦等。
\end{keybox}

\subsection{谷类的主要营养特点}

\begin{enumerate}[leftmargin=*]
    \item \textbf{碳水化合物：}含量最高，主要为淀粉（70\%\~{}80\%），是人体最经济、最主要的功能来源。淀粉分为直链淀粉和支链淀粉。含少量膳食纤维（全谷中较丰富）。
    \item \textbf{蛋白质：}含量8\%\~{}12\%，\textbf{谷类蛋白的第一限制氨基酸为\keyword{赖氨酸}}。大米蛋白质含量最低（7\%\~{}8\%），生物价值在谷类中相对最高——因其赖氨酸含量稍优于其他谷类。
    \item \textbf{脂类：}含量低（1\%\~{}4\%），主要集中在胚芽（Germ）和糊粉层（Aleurone Layer）。谷类脂类主要含不饱和脂肪酸，且含有维生素E——这也是全谷物粉贮存中酸败的主要原因。
    \item \textbf{维生素：}谷类是B族维生素（尤其是硫胺素B1、尼克酸、吡哆醇B6）的重要来源，主要分布在糊粉层和胚芽。\textbf{碾磨精度越高，B族维生素损失越严重。}
    \item \textbf{矿物质：}约1.5\%\~{}3\%，以磷为主（多为植酸磷，吸收率低），钙、铁含量均较低。
\end{enumerate}

\subsection{全谷物与精制谷物的比较}

\begin{keybox}
\textbf{全谷物 vs 精制谷物：}
\begin{itemize}[leftmargin=*]
    \item \textbf{全谷物：}保留麸皮、胚乳和胚芽。富含膳食纤维、B族维生素、维生素E、矿物质和植物化学物（多酚、植物固醇、木酚素）。
    \item \textbf{精制谷物：}仅保留胚乳（淀粉+少量蛋白质），B族维生素、膳食纤维等大量损失。
    \item \textbf{推荐：}《中国居民膳食指南（2022）》建议每日摄入全谷物和杂豆50\~{}150 g。
\end{itemize}
\end{keybox}

\subsection{薯类}

\textbf{薯类：}马铃薯、甘薯、木薯等。主要营养特点：淀粉丰富（15\%\~{}25\%），蛋白质极低（1\%\~{}2\%），富含维生素C和钾，甘薯富含$\beta$-胡萝卜素（类胡萝卜素）。

\subsection{杂豆}

\textbf{杂豆（赤豆、绿豆、芸豆等）：}淀粉含量50\%\~{}60\%，蛋白质20\%\~{}25\%（高于谷类），膳食纤维丰富，B族维生素和矿物质含量较高。

% ----- 第三节 大豆类及其制品 -----
\section{大豆类及其制品}

\begin{keybox}
\textbf{大豆（Soybean）：}含蛋白质35\%\~{}40\%，是植物性食物中蛋白质含量最高者，且为\keyword{优质植物蛋白}。含脂肪15\%\~{}20\%，主要为不饱和脂肪酸（亚油酸占50\%以上）。还含大豆异黄酮、大豆皂苷、大豆磷脂、大豆低聚糖等有益健康成分。
\end{keybox}

\subsection{大豆及制品的蛋白质}

\begin{itemize}[leftmargin=*]
    \item 大豆蛋白质氨基酸模式较好，属于优质蛋白。
    \item \textbf{第一限制氨基酸为\keyword{蛋氨酸（Met）}}——这是与谷类蛋白质（第一限制氨基酸为赖氨酸）互补的基础。
    \item \textbf{谷类+豆类互补作用：}谷类的赖氨酸不足→由大豆的蛋氨酸缺陷结构得到互补，显著提高蛋白质生物价。
    \item 大豆制成的传统豆制品：豆腐（蛋白质6\%\~{}8\%）、豆浆（蛋白质2\%\~{}3\%）、腐竹（蛋白质高达44\%\~{}50\%，脂肪含量也相对较高）。
\end{itemize}

\subsection{豆制品分类}

\begin{enumerate}[leftmargin=*]
    \item \textbf{非发酵豆制品：}豆腐、豆浆、豆腐干、腐竹、豆皮等。
    \item \textbf{发酵豆制品：}酱油、豆豉、豆瓣酱、腐乳、纳豆等。发酵可提高B族维生素含量和蛋白质消化率，部分产生具有生物活性的肽类。
\end{enumerate}

\subsection{大豆中其他有益成分}

\begin{itemize}[leftmargin=*]
    \item \textbf{大豆异黄酮：}具有植物雌激素作用和抗氧化作用（见第2章）。
    \item \textbf{大豆皂苷：}降血脂、抗肿瘤、免疫调节作用。
    \item \textbf{大豆磷脂：}降血脂、保护肝脏、促进神经系统发育。
    \item \textbf{大豆低聚糖（水苏糖、棉子糖）：}促进双歧杆菌增殖（益生元作用）。
\end{itemize}

\begin{tipbox}
\textbf{豆类的抗营养因子及处理：}大豆中含胰蛋白酶抑制剂（引起胰腺肥大）、植物血凝素（凝集红细胞）、脂肪氧化酶（引起豆腥味）。这些因子均可通过充分加热（煮沸\~{}30min以上）彻底灭活。因此，豆浆必须煮沸后食用。
\end{tipbox}

% ----- 第四节 蔬菜、水果类 -----
\section{蔬菜、水果类}

\begin{keybox}
\textbf{蔬菜和水果的共同营养贡献：}蔬菜和水果是维生素C、类胡萝卜素、叶酸、维生素K、钾、镁和膳食纤维的最重要来源。同时富含多种植物化学物（多酚类、有机硫化物、类胡萝卜素等），对健康维护具有不可替代的重要意义。
\end{keybox}

\subsection{蔬菜的营养特点}

\begin{enumerate}[leftmargin=*]
    \item \textbf{碳水化合物：}含量因种类差异大（2\%\~{}20\%），主要为可溶性糖（果糖、葡萄糖、蔗糖）和不溶性膳食纤维（纤维素、半纤维素、木质素）。
    \item \textbf{蛋白质和脂肪：}含量极低（蛋白质≤2\%，脂肪<0.5\%）。
    \item \textbf{维生素：}
    \begin{itemize}[leftmargin=*]
        \item \keyword{维生素C}：深绿色叶菜、辣椒、西兰花、番茄中含量丰富。
        \item \keyword{类胡萝卜素}（维生素A原）：深绿色和橙黄色蔬菜中丰富（胡萝卜、菠菜、南瓜）。
        \item \keyword{叶酸（Folate）}：深绿色叶菜（菠菜、西兰花）、豆类和坚果中含量丰富。
        \item \keyword{维生素K}：绿叶蔬菜是最主要的膳食来源。
    \end{itemize}
    \item \textbf{矿物质：}富含钾（有助于维持血压正常）、镁。钙在部分蔬菜（如芥菜、苋菜）中含量较高，但草酸等抗营养因子影响其吸收。
    \item \textbf{植物化学物：}有机硫化物（葱属和十字花科）、多酚类（各种有色蔬菜）、皂苷等。
\end{enumerate}

\subsection{水果的营养特点}

\begin{enumerate}[leftmargin=*]
    \item \textbf{碳水化合物：}主要为简单糖（果糖、葡萄糖、蔗糖），含量6\%\~{}25\%。含水量一般80\%以上，能量较低。
    \item \textbf{维生素：}新鲜水果是\keyword{维生素C}的最主要来源（猕猴桃、柑橘、草莓、鲜枣中含量极高）。也含类胡萝卜素（芒果、杏、木瓜）。
    \item \textbf{矿物质：}钾含量丰富（与血压控制相关）。
    \item \textbf{膳食纤维：}果胶含量较高，可降低血清胆固醇。
    \item \textbf{植物化学物：}花色素（浆果）、白藜芦醇（葡萄）、类胡萝卜素、黄酮类等。
\end{enumerate}

\begin{tipbox}
\textbf{蔬菜与水果的摄入建议（膳食指南2022）：}蔬菜摄入量300\~{}500 g/d，其中深色蔬菜应占一半（\~{}250 g/d）；水果200\~{}350 g/d，不能以果汁替代整果。蔬菜和水果在营养素谱系上有互补性，不可相互替代。
\end{tipbox}

\begin{keybox}
\textbf{蔬菜营养损失的关键因素：}
\begin{itemize}[leftmargin=*]
    \item \textbf{水溶性维生素：}维生素C和B族维生素在洗切、浸泡和煮沸中极易损失。
    \item \textbf{热敏性维生素：}维生素C、叶酸和部分B族维生素在高温下大量破坏。
    \item \textbf{氧化损失：}切碎后接触空气中的氧和光照，加速维生素氧化破坏。
    \item \textbf{合理烹调建议：}先洗后切、急火快炒、炒完即食，减少加工损失。
\end{itemize}
\end{keybox}

% ----- 第五节 畜、禽、水产品 -----
\section{畜、禽、水产品}

\begin{keybox}
\textbf{动物性食物的核心营养贡献：}畜、禽和水产品是\keyword{优质蛋白质}、\keyword{血红素铁}和\keyword{维生素B${}_{12}$}的最重要来源。B12是唯一只能从动物性食物中获得的维生素。
\end{keybox}

\subsection{畜肉（猪、牛、羊）}

\begin{itemize}[leftmargin=*]
    \item \textbf{蛋白质：}含量15\%\~{}22\%，氨基酸模式与人体接近，为优质蛋白。
    \item \textbf{脂肪：}含量因种类、部位而异（5\%\~{}30\%），以\keyword{饱和脂肪酸}为主（约占40\%\~{}60\%）。内脏脂肪含量更高，且胆固醇含量高（猪脑、猪肝尤其突出）。
    \item \textbf{碳水化合物：}极低（<1\%，主要以糖原形式存在）。
    \item \textbf{矿物质：}\textbf{血红素铁（Heme Iron）}——吸收率高达15\%\~{}25\%，是植物性非血红素铁（吸收率1\%\~{}5\%）的数倍。锌的吸收利用率也高。
    \item \textbf{维生素：}B族维生素（硫胺素B1——猪肉含量最高；吡哆醇B6——鸡肉中丰富；B12——所有肉类中均有）。维生素A在内脏（肝）中极高。
\end{itemize}

\subsection{禽肉（鸡、鸭）}

\begin{itemize}[leftmargin=*]
    \item 蛋白质含量与畜肉相当，但\keyword{不饱和脂肪酸}比例高于畜肉、胆固醇含量较低。
    \item 鸡胸肉蛋白质含量最高、脂肪最少（\~{}20\%蛋白、3\%脂肪），为健身人群的首选动物蛋白来源。
\end{itemize}

\subsection{水产品（鱼、虾、贝类）}

\begin{itemize}[leftmargin=*]
    \item \textbf{蛋白质：}与畜肉相当，但更易消化。
    \item \textbf{脂肪：}显著低（大部分<5\%），以\keyword{不饱和脂肪酸}为主——尤其是\keyword{n-3多不饱和脂肪酸}$\Rightarrow$深海鱼油含EPA（二十碳五烯酸）和DHA（二十二碳六烯酸），对心血管保护和脑发育有重要作用。
    \item \textbf{矿物质：}海产品富含碘、锌、硒（牡蛎——"海中的锌宝库"）。
    \item \textbf{维生素：}鱼肝含丰富的维生素A和维生素D。
\end{itemize}

\begin{exambox}
\textbf{考点提示——动物铁 vs 植物铁：}
\begin{itemize}[leftmargin=*]
    \item 动物食物：血红素铁，直接进入小肠上皮细胞，吸收率15\%\~{}25\%，不受膳食中植酸等抑制因素影响。
    \item 植物食物：非血红素铁，需先转化为Fe$^{2+}$，吸收率1\%\~{}5\%，受植酸、多酚、钙等抑制因素影响极大。
\end{itemize}
这是营养学中判断食物铁营养价值的关键区别。
\end{exambox}

% ----- 第六节 蛋类及蛋制品 -----
\section{蛋类及蛋制品}

\begin{keybox}
\textbf{蛋类的营养地位：}全鸡蛋蛋白质的\keyword{氨基酸模式与人体最接近}——是自然食物中生物价值最高的蛋白质来源。常被作为"参考蛋白质（Reference Protein）"用于食物蛋白质质量评价。
\end{keybox}

\subsection{蛋类的化学成分}

\textbf{全蛋：}
\begin{itemize}[leftmargin=*]
    \item 蛋白质11\%\~{}14\%，大部分为卵清蛋白（蛋清中），8种必需氨基酸齐全，生物价（BV）高达94。
    \item 脂肪8\%\~{}11\%，集中于蛋黄中，主要为甘油三酯和磷脂（卵磷脂为主），不饱和脂肪酸为主。
    \item 胆固醇含量高：蛋黄中约1500 mg/100g（每个鸡蛋约200\~{}300 mg胆固醇）。
    \item 微量营养素：维生素A、D、B2、B12、生物素含量较高；矿物质以磷、钙、铁、锌和硒为主（蛋黄中）。
\end{itemize}

\textbf{蛋清 vs 蛋黄：}
\begin{table}[h!]
\centering
\begin{tabular}{p{2.5cm}p{4cm}p{4cm}}
\hline
\textbf{营养成分} & \textbf{蛋清（约占蛋重60\%）} & \textbf{蛋黄（约占蛋重30\%）} \\
\hline
蛋白质 & 以卵清蛋白为主，约10\% & 卵黄蛋白为主，约16\% \\
脂肪 & 含量极低 & 30\%\~{}33\%，含丰富的卵磷脂 \\
维生素 & 极少 & 富含A、D、E、B族 \\
矿物质 & 极少 & 钙、磷、铁、锌、硒 \\
胆固醇 & 无 & 约1500 mg/100g \\
\hline
\end{tabular}
\end{table}

\subsection{蛋制品的营养变化}

\begin{itemize}[leftmargin=*]
    \item \textbf{咸蛋：}水分减少，钠含量显著增高，其他营养素基本保留。
    \item \textbf{皮蛋：}碱性加工条件下维生素B族大量破坏，维生素A和D可保留。
    \item \textbf{蛋粉：}脱水浓缩，营养素浓缩，但维生素A和D有一定损失。
\end{itemize}

% ----- 第七节 乳及乳制品 -----
\section{乳及乳制品}

\begin{keybox}
\textbf{乳类的核心营养地位：}乳类（特别是牛乳）是\keyword{最佳膳食钙来源}——不仅钙含量高（\~{}104 mg/100g），而且钙的吸收率高达32.1\%，远超其他钙源。乳类还是维生素B2（核黄素）和优质蛋白质的重要来源。
\end{keybox}

\subsection{牛乳的营养成分}

\begin{enumerate}[leftmargin=*]
    \item \textbf{蛋白质：}3.0\%\~{}3.5\%，主要为\keyword{酪蛋白（Casein）}（占80\%）和\keyword{乳清蛋白（Whey）}（占20\%）。乳清蛋白中的$\alpha$-乳白蛋白是母乳中的优势蛋白（见第4章）。生物价87，仅次于全鸡蛋。
    \item \textbf{脂肪：}3\%\~{}4\%，以甘油三酯为主，饱和脂肪酸为主，胆固醇含量低（10\~{}15 mg/100g）。以乳化小颗粒状存在，易于消化。
    \item \textbf{碳水化合物：}主要为\keyword{乳糖（Lactose）}（4.5\%\~{}5.0\%）。
    \begin{itemize}[leftmargin=*]
        \item 乳糖在肠道被乳糖酶水解为葡萄糖+半乳糖。
        \item \textbf{乳糖不耐症（Lactose Intolerance）：}因乳糖酶活性不足，未消化的乳糖进入结肠被细菌发酵产生气体和酸性产物，出现腹胀、腹痛和腹泻。亚洲人群发生率极高（80\%\~{}90\%）。
    \end{itemize}
    \item \textbf{钙：}含量104 mg/100g，吸收率32.1\%——高含量+高吸收率使牛乳成为不可替代的膳食钙源。
    \item \textbf{维生素：}维生素A（全脂乳丰富）、维生素B2（约0.13 mg/100g，是重要的核黄素来源）、维生素D（通常强化）。维生素C含量极低，在加热中大部分损失。
    \item \textbf{矿物质：}钙磷比为1.3:1，利于钙吸收。铁含量极低（<0.05 mg/100g，因此牛乳是缺铁性贫血的危险因素——婴幼儿单纯依靠牛乳可引致缺铁）。
\end{enumerate}

\subsection{常见乳制品}

\begin{itemize}[leftmargin=*]
    \item \textbf{酸奶（Yogurt）：}发酵后乳糖含量降低，可部分缓解乳糖不耐受；益生菌有利于肠道健康。
    \item \textbf{奶酪（Cheese）：}蛋白质和钙浓缩（钙400\~{}800 mg/100g），脂肪浓缩（25\%\~{}35\%），乳糖极低（适合乳糖不耐受者）。
    \item \textbf{奶粉（Milk Powder）：}脱水制成，营养素浓缩，适宜储藏和运输。
\end{itemize}

% ----- 第八节 坚果类 -----
\section{坚果类}

\begin{keybox}
\textbf{坚果的营养特点：}坚果富含不饱和脂肪酸、优质蛋白质、维生素E、B族维生素和多种矿物质（镁、锌、硒），是能量和营养素的浓缩食物。主要分为两类：树坚果（核桃、杏仁、腰果等）和籽类（葵花籽、南瓜籽等）。
\end{keybox}

\subsection{坚果的主要营养价值}

\begin{enumerate}[leftmargin=*]
    \item \textbf{脂肪：}含量极高（\~{}50\%，有的高达60\%+），以不饱和脂肪酸（单不饱和+多不饱和）为主。核桃的$\alpha$-亚麻酸（n-3必需脂肪酸）含量在食物中名列前茅。
    \item \textbf{蛋白质：}15\%\~{}25\%，大豆除外，是植物性食物中的高蛋白来源。
    \item \textbf{维生素：}维生素E含量丰富（重要的脂溶性抗氧化剂）、B族维生素（尤其是叶酸和吡哆醇B6）。
    \item \textbf{矿物质：}镁、钾、铜、锌、硒含量较高。
    \item \textbf{植物化学物：}植物固醇、多酚类、木酚素等。
\end{enumerate}

\subsection{食用注意事项}

\begin{itemize}[leftmargin=*]
    \item \textbf{能量密度极高：}100 g核桃约600\~{}700 kcal，过量摄入易导致能量过剩和体重增加。
    \item \textbf{中国居民膳食指南（2022）：}每日建议摄入量10\~{}15 g（约一小把），不宜贪多。
    \item \textbf{过敏风险：}花生和坚果是最常见的食物过敏原之一，在婴幼儿添加食物中须注意。
\end{itemize}

% ----- 复习题 -----
\reviewcontent{
\section{复习题}

\begin{enumerate}[leftmargin=*, itemsep=8pt]
    \item \textbf{【名词解释】}食物营养价值、优质蛋白质、血红素铁、乳糖不耐症、氨基酸模式、限制氨基酸、蛋白质互补作用
    \item \textbf{【简答题】}影响食物营养价值的因素有哪些？
    \item \textbf{【简答题】}谷类食物的主要营养特点是什么？为什么精制谷物不如全谷物？
    \item \textbf{【简答题】}简述大豆蛋白质的营养特点，及其与谷类蛋白质互补作用的原理。
    \item \textbf{【简答题】}试述畜肉、禽肉和水产品在蛋白质、脂肪和铁含量上的主要差异。
    \item \textbf{【简答题】}为什么全鸡蛋蛋白质被称为"参考蛋白质"？鸡蛋中的微量营养素集中在哪部分？
    \item \textbf{【简答题】}为什么牛乳被认为是最佳膳食钙来源？乳糖不耐受症的机制是什么？
    \item \textbf{【简答题】}简述蔬菜和水果在维生素C、类胡萝卜素和膳食纤维上的共同营养贡献，以及二者的差异。
    \item \textbf{【论述题】}从蛋白质、脂肪、矿物质和生物活性成分四个维度，比较分析畜肉类、水产品、蛋类和奶类的营养价值差异。
\end{enumerate}

\subsection*{参考答案要点}

\begin{enumerate}[leftmargin=*, itemsep=5pt]
    \item 见本章各概念定义框。
    \item 品种/产地/季节、栽培和饲养条件、收获成熟度、加工/烹调、储藏条件。
    \item 谷类：碳水化合物70\%\~{}80\%（最经济的热能来源），蛋白质8\%\~{}12\%（赖氨酸为第一限制氨基酸），低脂，B族维生素（精加工损失严重）。全谷物优于精制：多含膳食纤维、B族维生素、维生素E、矿物质和植物化学物。
    \item 大豆蛋白35\%\~{}40\%，优质植物蛋白，第一限制氨基酸为蛋氨酸；谷类第一限制氨基酸为赖氨酸。二者混合，谷类提供蛋氨酸、豆类提供赖氨酸，实现必需氨基酸互补。
    \item 蛋白质均优质。脂肪：畜肉饱和脂肪酸多，禽肉不饱和脂肪酸比例较高，水产品以n-3多不饱和脂肪酸为特色。铁：三者均含血红素铁，吸收率15\%\~{}25\%。
    \item 全鸡蛋氨基酸模式与人体最接近，BV高达94。微量营养素集中在蛋黄中（维生素A、D、E、B12、磷、铁、锌、硒、胆固醇）。
    \item 钙含量高（104 mg/100g），吸收率高（32.1\%）。乳糖不耐受机制：乳糖酶活性不足，未被消化的乳糖在结肠经细菌发酵产酸产气。
    \item 共同贡献：维生素C、类胡萝卜素、叶酸、钾、镁、膳食纤维、植物化学物。差异：蔬菜密度更低、可作为主菜食用量大，水果含糖量（简单糖）高，不宜替代蔬菜。
    \item (1) 蛋白质：均优质，蛋类BV最高——参考蛋白质；(2) 脂肪：水产品最低，n-3不饱和脂肪酸具有心血管保护作用；(3) 矿物质：畜肉等动物肉含高吸收率的血红素铁和锌，乳类的钙最高吸收率；(4) 生物活性成分：水产品含EPA/DHA，乳含益生菌/免疫球蛋白，蛋含卵磷脂。
\end{enumerate}

\newpage
}

% =====================================================================
%  CHAPTER 4: 特殊人群营养
% =====================================================================
\chapter{特殊人群营养}
\setcounter{chapter}{4}
\chapterintro{
\keyword{掌握}生命早期1000天机遇窗口及DOHaD理论；掌握孕妇、乳母、婴幼儿、学龄前儿童、学龄儿童、青少年、老年人的生理特点、营养需要及膳食指南关键推荐。\keyword{熟悉}运动员营养代谢特点及营养需要；熟悉各特殊人群常见营养问题及预防。\keyword{了解}高、低温环境人员营养需要和膳食原则。
}

% =====================================================================
\section{生命早期营养与DOHaD理论}
% =====================================================================

\begin{tipbox}
\textbf{DOHaD理论（Developmental Origins of Health and Disease，健康与疾病的发育起源）：}
生命早期（胎儿、婴幼儿期）经历不利因素（如营养不良、环境不良等），将会增加其成年后患肥胖、糖尿病、心血管疾病等慢性疾病的几率。其核心机制是"发育程序化（Developmental Programming）"——早期营养环境通过表观遗传调控（DNA甲基化、组蛋白修饰等）改变基因表达模式和组织分化方向，对个体远期疾病风险产生持久性影响。

\textbf{关键机制——"节俭表型"假说（Thrifty Phenotype Hypothesis）：}
宫内营养不良使胎儿产生"预测性适应"——调整代谢和内分泌系统以应对"匮乏"环境。若出生后实际面临"营养过剩"环境，则代谢-环境不匹配将大幅增加成年期代谢综合征风险。

\textbf{生命早期1000天（The First 1000 Days of Life）：}
生命早期1000天关键期 = 孕期270天 + 0-6月龄180天（母乳喂养）+ 7-24月龄（辅食添加）。这一阶段是生长发育的关键期和敏感期，此时营养状态对个体终身的体格发育、认知潜能、免疫功能和慢性疾病风险产生深远和不可逆的影响，是进行营养干预以获得最大健康收益的"机遇窗口期"。
\end{tipbox}

\begin{exambox}
\textbf{考点提示：} (1) 1000天的计算——孕期270d + 0-6月龄180d + 7-24月龄；注意不等于简单地"出生后730天"；(2) DOHaD的核心逻辑——早期环境 → 发育程序化（表观遗传）→ 远期健康/疾病；(3) "不匹配"是DOHaD解释代谢综合征的核心机制——宫内"预测匮乏" + 出生后"实际营养过剩" → 代谢-环境失匹配 → 代谢综合征高风险。
\end{exambox}

% =====================================================================
\section{孕妇营养}
% =====================================================================

\subsection{孕期生理特点}

\begin{keybox}
\textbf{妊娠期——生命早期1000天机遇窗口的起始阶段。}

\textbf{1. 内分泌变化：}
\begin{itemize}[leftmargin=*]
    \item \textbf{人绒毛膜促性腺激素（HCG）：}受精1周开始出现，8-9周达高峰。由胎盘合体滋养层细胞分泌，维持妊娠早期所需的孕酮分泌。
    \item \textbf{人绒毛膜生长激素（HCS）：}妊娠6周出现，36周达高峰至分娩。由胎盘合体滋养层细胞分泌，促进乳腺发育和胎儿生长，同时降低母体胰岛素敏感性。
    \item \textbf{雌激素：}分泌增加，增加子宫和胎盘的血流量，促进母体乳房发育。孕晚期可达非孕期的1000倍。
    \item \textbf{孕激素（孕酮）：}分泌增加，减少子宫收缩，维持妊娠。促进乳腺腺泡发育。
    \item \textbf{甲状腺激素：}水平升高 → 基础代谢率（BM）升高。
    \item \textbf{胰岛素敏感性下降：}由HCS、雌激素、孕激素等共同介导 → 妊娠期糖尿病（GDM）风险升高，属生理性胰岛素抵抗。
\end{itemize}

\textbf{2. 消化系统变化：}
\begin{itemize}[leftmargin=*]
    \item 贲门括约肌松弛 → 恶心呕吐（早孕反应）。
    \item 孕酮分泌增加 → 胃排空时间延长、胃肠蠕动减慢 → 胃肠胀气便秘。
    \item 对钙、铁、VitB$_{12}$、叶酸等营养素吸收增强——有利于营养储备。
\end{itemize}

\textbf{3. 循环系统变化：}
血容量增加，血液稀释：血浆容积增加幅度 $>$ 红细胞增加幅度 → \textbf{生理性贫血}。通常在孕20-30周最明显。

\textbf{WHO贫血标准（2024）：}
\begin{itemize}[leftmargin=*]
    \item 孕期：Hb $\geq$ 110 g/L
    \item 孕中期：Hb $\geq$ 105 g/L
\end{itemize}
\end{keybox}

\subsection{孕期体重增长}

\begin{table}[h!]
\centering
\caption{孕期推荐体重增长范围}
\begin{tabular}{lcc}
\hline
\textbf{孕前BMI分类} & \textbf{总增长范围（kg）} & \textbf{中晚期每周增长（kg/周）} \\
\hline
低体重（BMI $<$ 18.5） & 11.0-16.0 & 0.46（0.37-0.56） \\
正常体重（BMI 18.5-24.0） & 8.0-14.0 & 0.37（0.26-0.48） \\
超重（BMI 24.0-28.0） & 7.0-11.0 & 0.30（0.22-0.37） \\
肥胖（BMI $\geq$ 28.0） & 5.0-9.0 & 0.22（0.15-0.30） \\
\hline
\end{tabular}
\end{table}

\subsection{孕期营养需要}

\begin{keybox}
\textbf{孕期各阶段能量与宏量营养素增加量：}

\begin{center}
\begin{tabular}{lcccc}
\hline
\textbf{营养素} & \textbf{非孕期} & \textbf{孕早期} & \textbf{孕中期} & \textbf{孕晚期} \\
\hline
能量EER（kcal/d） & 1700 & +0 & +250 & +400 \\
蛋白质（g/d） & -- & +0 & +15 & +30 \\
\hline
\end{tabular}
\end{center}

\textbf{脂类：}供能比20\%-30\%不变。亚油酸占总能量4\%，$\alpha$-亚麻酸0.6\%，EPA+DHA达到250 mg/d。

\textbf{碳水化合物：}孕早期不低于130 g/d——预防酮体对胎儿中枢神经系统的毒性。

\textbf{关键微量营养素：}

\begin{center}
\begin{tabular}{lccc}
\hline
\textbf{营养素} & \textbf{非孕期} & \textbf{孕期RNI} & \textbf{备注} \\
\hline
钙（mg/d） & 800 & 800（不增加） & 孕晚期胎儿大量储存，钙吸收率代偿性升高 \\
铁（mg/d） & 18 & 孕中25，孕晚29 & 用于母体红细胞增加、胎儿胎盘需要、产时失血储备 \\
锌（mg/d） & 8.5 & 10.5 & -- \\
碘（$\mu$g/d） & 120 & 230 & 严重缺乏 → 克汀病 \\
叶酸（$\mu$g DFE/d） & -- & 400 & 备孕3个月起补充，持续整个妊娠期 \\
维生素D（$\mu$g/d） & 10 & 10（400 IU/d） & 缺乏 → 新生儿低钙血症和佝偻病风险增加 \\
\hline
\end{tabular}
\end{center}
\end{keybox}

\begin{keybox}
\textbf{叶酸——最重要的孕前/孕期微量营养素：}
\begin{itemize}[leftmargin=*]
    \item \textbf{推荐：}备孕妇女从计划怀孕前3个月开始每日补充叶酸\textbf{400 $\mu$g DFE/d}，持续整个妊娠期。
    \item \textbf{意义：}预防神经管缺陷（NTDs）——叶酸缺乏是胎儿神经管未能正常闭合（如脊柱裂、无脑儿）的确定病因。充足补充可使NTDs发生率降低70\%。
    \item \textbf{为什么必须从孕前开始？}孕后3-8周是神经管闭合关键期，此时许多妇女尚不知已怀孕。
    \item \textbf{膳食来源：}深绿色叶菜、肝脏、豆类、强化面粉。
\end{itemize}
\end{keybox}

\subsection{孕期膳食指南（2022版）}

\begin{keybox}
\textbf{《备孕和孕期妇女膳食指南（2022年）》6条关键推荐：}
\begin{enumerate}[leftmargin=*]
    \item \textbf{调整孕前体重至正常范围，保证孕期体重适宜增长。}
    \item \textbf{常吃含铁丰富的食物，选用碘盐，合理补充叶酸和维生素D。}
    \item \textbf{孕吐严重者，可少量多餐，保证摄入含必需量碳水化合物的食物。}
    \item \textbf{孕中晚期适量增加奶、鱼、禽、蛋、瘦肉的摄入。}
    \item \textbf{经常户外活动，禁烟酒，保持健康生活方式。}
    \item \textbf{愉快孕育新生命，积极准备母乳喂养。}
\end{enumerate}
\end{keybox}

\subsection{孕期常见营养问题及预防}

\begin{defbox}
\textbf{孕期常见营养问题：}

\begin{enumerate}[leftmargin=*]
    \item \textbf{早孕反应（妊娠剧吐 hyperemesis gravidarum）：}
    \begin{itemize}[leftmargin=*]
        \item 发生机制：HCG升高 + 贲门括约肌松弛 + 胃肠蠕动减慢。
        \item 营养风险：严重者进食困难 → 能量和营养素摄入不足 → 酮体产生（对胎儿中枢神经系统有毒性）。
        \item 预防处理：少量多餐（每日6-8次）；保证碳水化合物最低摄入量130 g/d（严重者需静脉补充葡萄糖，防止酮症）；晨起前吃干食物（如苏打饼干）；避免油腻和辛辣食物；补充维生素B$_{6}$有一定效果。
    \end{itemize}

    \item \textbf{孕期便秘：}
    \begin{itemize}[leftmargin=*]
        \item 发生机制：孕酮分泌增加 → 胃肠蠕动减慢；增大的子宫压迫直肠。
        \item 预防：增加膳食纤维（新鲜蔬果、全谷物）和足够饮水，适量运动。
    \end{itemize}

    \item \textbf{妊娠期糖尿病（GDM）：}
    \begin{itemize}[leftmargin=*]
        \item 发生机制：HCS、雌激素、孕激素等综合作用 → 生理性胰岛素抵抗。
        \item 发生率约10\%-20\%。高危因素：孕前超重/肥胖、有糖尿病家族史、高龄妊娠。
        \item 危害：增加巨大儿、剖宫产率、新生儿低血糖风险，母亲远期2型糖尿病风险升高。
        \item 预防：合理控制孕期体重增长，避免高GI食物和含糖饮料，适量运动。
    \end{itemize}

    \item \textbf{妊娠期高血压疾病：}
    \begin{itemize}[leftmargin=*]
        \item 与钙摄入不足可能有关，保证充足钙摄入（800 mg/d）可降低高血压风险。
        \item 控制钠盐摄入，维持适宜体重增长。
    \end{itemize}

    \item \textbf{孕期贫血（缺铁性贫血最常见）：}
    \begin{itemize}[leftmargin=*]
        \item 原因：血容量扩张 + 胎儿和胎盘铁转运 + 产时失血储备 → 铁需要量大幅增加。
        \item 孕中期铁RNI 25 mg/d，孕晚期29 mg/d（比非孕期18 mg/d显著增加）。
        \item 预防：常吃含铁丰富食物（动物肝脏、动物血、红肉），同时摄入维生素C促铁吸收。
    \end{itemize}

    \item \textbf{孕期钙营养与骨骼健康：}
    \begin{itemize}[leftmargin=*]
        \item 孕期钙RNI 800 mg/d（与孕前相同），因钙吸收率代偿性升高。
        \item 钙摄入严重不足时，胎儿动员母体骨骼钙 → 母体骨量损失。充足的奶及奶制品摄入可预防。
    \end{itemize}
\end{enumerate}
\end{defbox}

% =====================================================================

\subsection{泌乳量与母乳分期}

\begin{keybox}
\textbf{哺乳期每日泌乳量：}700-800 ml（平均750 ml）。

\textbf{母乳分期：}
\begin{enumerate}[leftmargin=*]
    \item \textbf{初乳（产后第1周）：}淡黄色，蛋白质多，富含免疫球蛋白（分泌型IgA和乳铁蛋白），乳糖和脂肪相对较少。
    \item \textbf{过渡乳（产后第2周）：}蛋白质$\downarrow$，乳糖和脂肪$\uparrow$。
    \item \textbf{成熟乳（第2周以后）：}乳白色，蛋白质相对较少，乳糖和脂肪多。
\end{enumerate}
\end{keybox}

\subsection{母乳喂养的优点}

\begin{keybox}
\textbf{母乳喂养的优点：}
\begin{enumerate}[leftmargin=*]
    \item \textbf{蛋白质：}乳清蛋白为主，必需氨基酸比例适当。
    \item \textbf{脂肪：}脂肪颗粒小 + 乳脂酶 → 易消化吸收。
    \item \textbf{糖类：}富含乳糖 → 促进乳酸杆菌、抑制大肠埃希菌生长，利于钙、铁、锌吸收。
    \item \textbf{矿物质：}钙磷比例适宜（2:1）。
    \item \textbf{含大量免疫物质：}分泌型IgA、乳铁蛋白、溶菌酶等。
    \item \textbf{不易过敏。}
    \item \textbf{经济、方便、卫生。}
    \item \textbf{促进产后恢复，增进母婴交流。}
\end{enumerate}
\end{keybox}

\subsection{乳母的营养需要}

\begin{keybox}
\textbf{乳母关键营养素较非孕妇女的增加量：}

\begin{center}
\begin{tabular}{lcl}
\hline
\textbf{营养素} & \textbf{增加量} & \textbf{备注} \\
\hline
能量EER & +400 kcal/d & 满足泌乳能量消耗 \\
蛋白质 & +25 g/d & 高品质蛋白质 \\
钙 & 800 mg/d（不增） & 乳汁中含量较稳定，每日排出300 mg \\
  & & 乳母钙供给不足会用自身骨骼中的钙满足乳汁钙含量 \\
铁 & -- & 铁不能通过乳腺输送到乳汁 \\
\hline
\end{tabular}
\end{center}

\textbf{重要说明：}
\begin{itemize}[leftmargin=*]
    \item \textbf{钙：}乳汁中钙含量较稳定（约300 mg/L），每日排出约300 mg。乳母钙供给不足时会动用自身骨骼中的钙来满足乳汁钙含量——哺乳期骨骼密度下降，断奶后可逆性恢复。
    \item \textbf{铁：}铁不能通过乳腺输送到乳汁，因此母乳中铁含量不受母体铁状态影响。
\end{itemize}
\end{keybox}

\subsection{哺乳期膳食指南}

\begin{keybox}
\textbf{哺乳期膳食指南关键推荐：}
\begin{enumerate}[leftmargin=*]
    \item 增加富含优质蛋白质及维生素A的动物性食物和海产品摄入。
    \item 保证充足的能量和液体摄入——能量在孕前基础上增加400 kcal/d。
    \item 增加钙摄入，重视哺乳期骨骼健康——每日饮奶400-500 ml，保证钙摄入800 mg/d。
    \item 保持身心愉快，保证充足睡眠——充足休息是维持泌乳量的必要条件。
    \item 坚持哺乳，适量身体活动——有助于产后体重恢复。
    \item 禁烟酒，避免浓茶和咖啡。
\end{enumerate}
\end{keybox}

\subsection{乳母常见营养问题及预防}

\begin{defbox}
\textbf{乳母常见营养问题：}

\begin{enumerate}[leftmargin=*]
    \item \textbf{能量和营养素摄入不足：}
    \begin{itemize}[leftmargin=*]
        \item 原因：产后急于恢复体型而过度控制饮食；照顾婴儿导致饮食不规律。
        \item 危害：泌乳量减少，乳汁质量降低，母体营养耗竭。
        \item 预防：能量在孕前基础上增加400 kcal/d，蛋白质增加25 g/d；保证充足液体摄入（多喝汤水，泌乳需水约1 L/d额外补充）。
    \end{itemize}

    \item \textbf{乳母骨骼健康问题：}
    \begin{itemize}[leftmargin=*]
        \item 乳汁每日排出约300 mg钙，钙供给不足时动用人骨钙 → 哺乳期骨密度暂时性下降（约3\%-5\%）。
        \item 预防：保证每日钙摄入800 mg/d（饮奶400-500 ml）；断奶后骨密度一般可恢复。
    \end{itemize}

    \item \textbf{铁营养与产后贫血：}
    \begin{itemize}[leftmargin=*]
        \item 铁不能通过乳腺进入乳汁，影响的是母亲自身的铁储备。
        \item 分娩失血 + 哺乳消耗可导致产后缺铁性贫血。预防：注意补铁（动物肝脏、红肉），产后继续补充铁剂至血象恢复。
    \end{itemize}

    \item \textbf{乳腺炎：}
    \begin{itemize}[leftmargin=*]
        \item 充分排空乳汁是最重要的预防手段之一。营养方面注意保证充足的蛋白质和维生素摄入（增强免疫力）。
    \end{itemize}

    \item \textbf{维生素缺乏：}
    \begin{itemize}[leftmargin=*]
        \item 维生素A：乳汁维生素A含量受乳母膳食维生素A影响较大，维生素A摄入不足导致乳汁含量降低。预防：摄入富含维生素A的动物性食物和海产品。
        \item 维生素D：乳母维生素D缺乏 → 乳汁维生素D含量低 → 婴儿维生素D不足。乳母需保证维生素D摄入（10 $\mu$g/d）。
        \item 维生素B$_{1}$和B$_{2}$需要量与能量摄入相关（0.5 mg/1000 kcal），需相应增加摄入。
    \end{itemize}
\end{enumerate}
\end{defbox}

% =====================================================================

\subsection{婴幼儿生理特点}

\begin{keybox}
\textbf{婴幼儿（0-3岁）主要生理特点：}

\begin{enumerate}[leftmargin=*]
    \item \textbf{生长发育旺盛：}第一生长高峰期——婴儿期体重增长为出生时的3倍（1岁时约10 kg），身高增加约25 cm。2岁至青春期前：体重每年增长约2 kg，身高5-7 cm。
    \item \textbf{代谢率高：}基础代谢率约为成人的2倍，能量需要量显著高于成人（按每千克体重计）。
    \item \textbf{消化系统不成熟：}胃容量小（新生儿约30-50 ml，1岁时250-350 ml），肠道通透性高，消化酶分泌不足——对食物的消化吸收能力有限，易发生消化紊乱。
    \item \textbf{神经系统：}大脑快速发育，3岁时脑重约为成人80\%，6岁时可达成人90\%。突触形成和髓鞘化在2岁内最为活跃，DHA和ARA充足供给对脑发育至关重要。
    \item \textbf{免疫功能：}新生儿免疫系统尚未发育完全，IgG可从母体获得（胎盘转运），sIgA主要来自母乳（初乳中含量极高）——母乳喂养对婴儿免疫保护至关重要。
    \item \textbf{肾功能：}肾小球滤过率和肾小管重吸收功能都不成熟，肾溶质负荷有限——不能处理高浓度废物（需大量水分），不宜摄入过多蛋白质和矿物质（尤其钠）。
    \item \textbf{体成分特点：}体内水分含量高（新生儿约75\%，成人约60\%），体表面积相对较大 → 水分蒸发量大 → 每日需水量150 ml/(kg$\cdot$d)。
    \item \textbf{铁贮存：}足月新生儿体内贮存约300 mg铁，可供生后4-6个月需要。6月龄后贮存铁耗竭，必须从辅食补充。
\end{enumerate}
\end{keybox}

\subsection{0-6月龄婴儿——纯母乳喂养}

\begin{keybox}
\textbf{《0-6月龄婴儿母乳喂养指南（2022）》准则：}
\begin{enumerate}[leftmargin=*]
    \item \textbf{准则1：母乳是婴儿最理想的食物，坚持6月龄内纯母乳喂养。}
    \item \textbf{准则2：生后1小时内开奶。}
    \item \textbf{准则3：回应式喂养。}
    \item \textbf{准则4：适当补充维生素D，母乳喂养无需补钙。}
    \item \textbf{准则5：任何动摇母乳喂养的想法和举动都必须咨询医生。}
    \item \textbf{准则6：定期监测体格指标。}
\end{enumerate}
\end{keybox}

\subsection{7-24月龄婴幼儿——辅食添加}

\begin{keybox}
\textbf{《7-24月龄婴幼儿喂养指南（2022）》准则：}
\begin{enumerate}[leftmargin=*]
    \item \textbf{准则1：继续母乳喂养，满6月龄起必须添加辅食，从富含铁的泥糊状食物开始。}
    \item \textbf{准则2：及时引入多样化食物，重视动物性食物的添加。}
    \item \textbf{准则3：尽量少加糖、盐，油脂适当，保持食物原味。}
    \item \textbf{准则4：提倡回应式喂养，鼓励但不强迫进食。}
    \item \textbf{准则5：注重饮食卫生和进食安全。}
    \item \textbf{准则6：定期监测体格指标，追求健康生长。}
\end{enumerate}
\end{keybox}

\begin{keybox}
\textbf{辅食添加的意义：}
\begin{enumerate}[leftmargin=*]
    \item \textbf{补充营养素：}6月龄后母乳已不能满足婴儿能量和营养素（尤其铁、锌）的全部需要——\textbf{7-12月龄99\%的铁必须从辅食中获得}。
    \item \textbf{促进消化系统发育：}辅食刺激胃肠道消化酶分泌和功能成熟。
    \item \textbf{培养咀嚼和吞咽能力：}不同质地食物的逐步引入有利于口腔运动和咀嚼功能发展。
    \item \textbf{促进味觉发育和饮食习惯形成：}辅食添加期是接受新食物和建立食物偏好的关键窗口期。
    \item \textbf{促进神经心理发育：}自主进食习惯的培养有利于精细运动和认知发育。
\end{enumerate}
\end{keybox}

\begin{keybox}
\textbf{辅食添加的原则：}
\begin{itemize}[leftmargin=*]
    \item 每次只添加一种新食物，观察3-5天以排除过敏。
    \item 由少到多，由稀到稠，由细到粗。
    \item 从泥糊状 → 碎末状 → 小块状 → 条状。
    \item 第一口辅食首选\textbf{强化铁的米粉、动物肝脏泥、瘦肉泥}等富含铁的泥糊状食物。
    \item 应在婴儿健康、消化功能正常时添加。
    \item 保持原味，不加盐（1岁前）、糖及刺激性调味品；1岁后逐渐尝试淡口味家庭膳食。
    \item 6月龄时婴儿体内储存的铁已耗尽——\textbf{7-12月龄99\%的铁必须从辅食中获得}。
\end{itemize}
\end{keybox}

\subsection{婴幼儿常见营养问题及预防}

\begin{defbox}
\textbf{婴幼儿常见营养问题：}

\begin{enumerate}[leftmargin=*]
    \item \textbf{缺铁性贫血（IDA）：}
    \begin{itemize}[leftmargin=*]
        \item 6月龄至2岁为高发期——原因：6月后贮存铁耗竭，辅食铁来源不足。
        \item 7-12月龄铁RNI = 10 mg/d，必须通过辅食获得。
        \item 危害：影响认知发育和行为能力（即使纠正后，认知缺陷可能持续）。
        \item 预防：保证6月龄及时添加富含铁辅食（强化铁米粉、肝脏泥、肉泥）；适当补充维生素C促进铁吸收。
    \end{itemize}

    \item \textbf{维生素D缺乏性佝偻病：}
    \begin{itemize}[leftmargin=*]
        \item 原因：母乳中维生素D含量极低，日照不足，未补充维生素D制剂。
        \item 预防：生后数日开始每日补充维生素D 400 IU（10 $\mu$g/d），至2岁。
    \end{itemize}

    \item \textbf{蛋白质-能量营养不良（PEM）：}
    \begin{itemize}[leftmargin=*]
        \item 辅食添加不及时、辅食质量差、断奶后喂养不当。
        \item 表现：生长迟缓（身高别年龄低于标准）、消瘦（体重别身高低于标准）。
        \item 预防：保证辅食的质和量，及时引入动物性食物。
    \end{itemize}

    \item \textbf{锌缺乏：}
    \begin{itemize}[leftmargin=*]
        \item 表现：生长发育迟缓、食欲不振、异食癖、免疫功能下降。
        \item 预防：辅食中添加富含锌的食物（红肉、肝脏、海产品）。
    \end{itemize}

    \item \textbf{碘缺乏：}
    \begin{itemize}[leftmargin=*]
        \item 严重缺乏可致克汀病（先天性碘缺乏综合征：矮小、智力严重障碍）。
        \item 预防：使用碘盐。
    \end{itemize}

    \item \textbf{维生素K缺乏：}
    \begin{itemize}[leftmargin=*]
        \item 母乳中维生素K含量很低，单纯母乳喂养新生儿易缺乏 → 出血性疾病风险。
        \item 预防：新生儿出生后立即注射维生素K（肌注1 mg）。
    \end{itemize}

    \item \textbf{龋齿：}
    \begin{itemize}[leftmargin=*]
        \item 含糖饮料（含果汁、加糖乳品）和甜食是婴幼儿龋齿的主要膳食因素。
        \item 预防：1岁前不加糖，不喂含糖饮料，1岁后减少游离糖摄入。
    \end{itemize}
\end{enumerate}
\end{defbox}

\subsection{婴幼儿营养需要}

\begin{keybox}
\textbf{婴幼儿关键营养素需要特点：}

\begin{itemize}[leftmargin=*]
    \item \textbf{蛋白质：}优质蛋白占总蛋白质1/2。
    \item \textbf{钙/磷比值：}适宜比值（2.3:1），比牛奶（1.4:1）合理。
    \item \textbf{铁：}足月新生儿有足够铁贮存，供4-6个月需要。\textbf{6m-2y易患缺铁性贫血}。7-12月龄RNI = 10 mg/d。
    \item \textbf{锌：}婴幼儿期缺锌 → 生长发育迟缓、脑发育受损、食欲不振、异食癖。
    \item \textbf{碘：}缺乏 → 克汀病（先天性碘缺乏综合征：矮小、智力严重障碍）。
    \item \textbf{维生素D：}乳类含量微量，需要额外补充（生后数日开始补充400 IU/d）。
    \item \textbf{维生素K：}母乳中很少，单纯母乳喂养易缺乏。
\end{itemize}
\end{keybox}

% =====================================================================
\section{学龄前儿童营养}
% =====================================================================

\subsection{生理特点}

\begin{defbox}
\textbf{学龄前儿童（3-6岁）主要生理特点：}

\begin{enumerate}[leftmargin=*]
    \item \textbf{生长发育：}每年体重增长约2 kg，每年身高增长约5-7 cm。生长速率较婴幼儿期明显趋缓。
    \item \textbf{神经系统：}依旧属于脑细胞生长发育关键时期。3岁脑重约为成人80\%，6岁可达成人90\%。
    \item \textbf{咀嚼和消化能力仍有限：}3岁乳牙已出齐，6岁恒牙已萌出，但咀嚼和消化能力远低于成人。
    \item \textbf{心理发育特点：}注意力分散（约15分钟），无法专心进食。好奇心强，模仿能力增强，此时期是培养良好饮食行为习惯的关键窗口期。
    \item \textbf{常见营养问题：}挑食偏食、零食摄入过多、含糖饮料摄入过量、钙和铁摄入不足。
\end{enumerate}
\end{defbox}

\subsection{学龄前儿童膳食指南（2022版）}

\begin{keybox}
\textbf{学龄前儿童膳食指南（2022）核心推荐：}
\begin{enumerate}[leftmargin=*]
    \item \textbf{食物多样，规律就餐，自主进食}——培养良好饮食习惯。
    \item \textbf{每天饮奶，足量饮水，合理选择零食。}
    \begin{itemize}[leftmargin=*]
        \item 奶类每日300-500 mL。
        \item 水每日600-800 mL（以白开水为主）。
    \end{itemize}
    \item \textbf{合理烹调，少调料少油炸}——采用蒸、煮、炖、煨等方式。每人每日食盐 $<$ 3 g。
    \item \textbf{参与食物选择制作}——增进对食物的认知与喜爱。
    \item \textbf{经常户外活动，定期体格测量。}
    \begin{itemize}[leftmargin=*]
        \item 每天至少60分钟体育活动。
        \item 每天看电视、玩平板电脑累计时间不超过2小时。
    \end{itemize}
\end{enumerate}
\end{keybox}

\subsection{学龄前儿童常见营养问题及预防}

\begin{defbox}
\textbf{学龄前儿童常见营养问题：}

\begin{enumerate}[leftmargin=*]
    \item \textbf{挑食偏食：}
    \begin{itemize}[leftmargin=*]
        \item 学龄前儿童自主意识增强，对陌生食物有"恐新症"。
        \item 预防：家长以身作则，树立良好榜样；多次尝试（8-15次）让儿童接受新食物；不强迫进食，营造愉快用餐氛围；让儿童参与食物选择和制作。
    \end{itemize}

    \item \textbf{零食摄入过多：}
    \begin{itemize}[leftmargin=*]
        \item 劣质零食（高能量、高脂肪、高盐、高糖、低营养）挤占主食和奶类摄入。
        \item 预防：合理选择零食（奶制品、水果、坚果），控制零食量和时间（两餐之间，不代替正餐）。
    \end{itemize}

    \item \textbf{含糖饮料摄入过量：}
    \begin{itemize}[leftmargin=*]
        \item 含糖饮料 → 能量过剩 → 超重/肥胖 + 龋齿风险。
        \item 预防：以白开水为主要饮品（每日600-800 ml），不提供含糖饮料。
    \end{itemize}

    \item \textbf{钙摄入不足：}
    \begin{itemize}[leftmargin=*]
        \item 奶类摄入不足是最常见原因 → 影响骨骼发育和峰值骨量积累。
        \item 预防：保证每日奶类300-500 ml。
    \end{itemize}

    \item \textbf{铁和锌摄入不足：}
    \begin{itemize}[leftmargin=*]
        \item 动物性食物摄入不足的儿童易缺铁和锌。
        \item 预防：保证肉类、肝脏、海产品适量摄入。
    \end{itemize}

    \item \textbf{超重和肥胖：}
    \begin{itemize}[leftmargin=*]
        \item 能量摄入大于消耗，零食和含糖饮料过多，体育活动不足。
        \item 预防：控制总能量和零食，保证每日户外活动至少60分钟。
    \end{itemize}
\end{enumerate}
\end{defbox}

% =====================================================================
\section{学龄儿童营养}
% =====================================================================

\subsection{生理特点}

\begin{defbox}
\textbf{学龄儿童（6-12岁）主要生理特点：}

\begin{enumerate}[leftmargin=*]
    \item \textbf{生长发育：}身高和体重较快增长但低于学龄前期和青春期。
    \item \textbf{骨骼：}骨胶多、钙盐少弹性大、硬度小，不易骨折而易弯曲变形。
    \item \textbf{换牙：}恒牙6-7岁；钙质沉积更早。
    \item \textbf{神经系统：}脑细胞分化基本完成，神经传导逐渐迅速、准确。
    \item \textbf{心理特点：}控制力仍较弱；精神易紧张，情绪易波动，过多的脑力活动易疲劳。
\end{enumerate}

\textbf{生长发育4规律：}
\begin{enumerate}[leftmargin=*]
    \item 连续性
    \item 不同步
    \item 波浪式
    \item 差异性
\end{enumerate}
\end{defbox}

\subsection{学龄儿童膳食指南（2022版）}

\begin{keybox}
\textbf{学龄儿童膳食指南（2022）5条核心推荐：}
\begin{enumerate}[leftmargin=*]
    \item \textbf{主动参与食物制作和选择，提高营养素养。}
    \item \textbf{吃好早餐，合理选择零食，培养健康饮食行为。}
    \item \textbf{天天喝奶，足量饮水，不喝含糖饮料，禁止饮酒。}
    \item \textbf{多户外活动，少视屏时间，每天60分钟以上的中高强度身体活动。}
    \item \textbf{定期监测体格发育，保持体重适宜增长。}
\end{enumerate}
\end{keybox}

\subsection{学龄儿童常见营养问题及预防}

\begin{defbox}
\textbf{学龄儿童常见营养问题：}

\begin{enumerate}[leftmargin=*]
    \item \textbf{不吃早餐或早餐质量差：}
    \begin{itemize}[leftmargin=*]
        \item 学龄儿童不吃早餐比例较高，影响上午学习效率和认知功能。
        \item 预防：培养规律早餐习惯；早餐应包含谷类、奶类/蛋类/肉类和蔬菜水果三类食物以上。
    \end{itemize}

    \item \textbf{零食和含糖饮料过度摄入：}
    \begin{itemize}[leftmargin=*]
        \item 校门口摊贩和自动售货机的零食饮料影响正餐摄入。
        \item 预防：合理选择零食，不喝含糖饮料，以白开水为日常饮水。
    \end{itemize}

    \item \textbf{钙和维生素D不足：}
    \begin{itemize}[leftmargin=*]
        \item 奶类摄入不足 + 日照减少（室内活动多） → 钙和维生素D缺乏风险。
        \item 预防：天天喝奶（300-500 ml/d），增加户外活动。
    \end{itemize}

    \item \textbf{超重肥胖或消瘦两极分化：}
    \begin{itemize}[leftmargin=*]
        \item 超重肥胖：高能量密度食物过多，身体活动不足。
        \item 消瘦：偏食、厌食、对体型的不正确认知。
        \item 预防：平衡膳食 + 每天至少60分钟中高强度身体活动 + 定期监测体格发育。
    \end{itemize}

    \item \textbf{脊柱和骨骼发育问题：}
    \begin{itemize}[leftmargin=*]
        \item 学龄儿童骨骼骨胶多、钙盐少、硬度小，易弯曲变形。
        \item 姿势不良（长时间不正确的读写姿势） → 脊柱侧弯、驼背。
        \item 预防：保证钙摄入，保持正确坐姿，适量运动。
    \end{itemize}
\end{enumerate}
\end{defbox}

% =====================================================================
\section{青少年营养}
% =====================================================================

\subsection{生理特点}

\begin{defbox}
\textbf{青少年——第二生长发育高峰期。}

\textbf{1. 体格发育：}
\begin{itemize}[leftmargin=*]
    \item \textbf{身高性别差异：}男较女晚2年；男7-9 cm/y；女5-7 cm/y。
    \item \textbf{体重：}4-5 kg/y，时间较长。
    \item \textbf{形态：}四肢比脊柱早；脚最先；从远到近。
\end{itemize}

\textbf{2. 性发育：}10岁前处于静止状态，青春期后在激素作用下迅速发育。

\textbf{3. 骨量积累：}青春期是获得峰值骨量（Peak Bone Mass）的最关键阶段——约45\%的成人骨量在青春期积累。

\textbf{4. 认知和心理：}抽象思维能力增强，独立意识增加，同伴影响力显著上升。易出现饮食行为偏离（节食减肥、不吃早餐、在外就餐增多）。
\end{defbox}

\subsection{青少年营养需要}

\begin{keybox}
\textbf{青少年关键营养素需求特点：}
\begin{itemize}[leftmargin=*]
    \item \textbf{能量：}显著增加。11-13岁：男2050 kcal/d，女1800 kcal/d；14-17岁：男2500 kcal/d，女2000 kcal/d。
    \item \textbf{蛋白质：}55-75 g/d（男 $>$ 女），优质蛋白质占1/2以上。长期摄入不足 → 青春期生长迟滞。
    \item \textbf{钙：}1000-1200 mg/d——峰值骨量积累关键窗口期。摄入不足 → 无法达到遗传决定的峰值骨量 → 老年期骨质疏松风险大幅增加。
    \item \textbf{铁：}青春期需求量大幅增加。男孩来自肌肉和血容量增加，女孩需弥补月经铁丢失。11-13岁：男15 mg/d，女18 mg/d；14-17岁：男16 mg/d，女18 mg/d。\textbf{青春期女性是缺铁性贫血高发人群。}
    \item \textbf{锌：}9-11.5 mg/d。缺乏 → 青春期生长迟缓和性发育延迟。
    \item \textbf{碘：}110-120 $\mu$g/d。缺乏 → 甲状腺肿和认知功能下降。
    \item \textbf{维生素D：}10 $\mu$g/d（400 IU/d）。
\end{itemize}
\end{keybox}

\subsection{青少年常见营养问题及预防}

\begin{defbox}
\textbf{青少年常见营养问题：}

\begin{enumerate}[leftmargin=*]
    \item \textbf{青春期女性缺铁性贫血：}
    \begin{itemize}[leftmargin=*]
        \item 原因：月经铁丢失 + 生长发育血容量扩大双重要求 → 是缺铁性贫血高发人群。
        \item 青少年女性铁RNI = 18 mg/d。预防：常吃含铁丰富食物（动物肝脏、动物血、红肉），同时摄入维生素C促铁吸收。
    \end{itemize}

    \item \textbf{钙摄入不足与峰值骨量受影响：}
    \begin{itemize}[leftmargin=*]
        \item 青春期是获得峰值骨量的最关键阶段——约45\%的成人骨量在此期积累。
        \item 钙摄入不足 → 峰值骨量无法达到遗传决定的最大值 → 老年期骨质疏松风险大幅增加。
        \item 预防：奶及奶制品为钙最佳来源，钙RNI 1000-1200 mg/d。
    \end{itemize}

    \item \textbf{不健康减肥行为：}
    \begin{itemize}[leftmargin=*]
        \item 盲目节食、不吃主食、过度运动、使用减肥药等，女性多见。
        \item 后果：能量和营养素不足，影响生长发育，月经紊乱/闭经。
    \end{itemize}

    \item \textbf{不吃早餐和在外就餐增多：}
    \begin{itemize}[leftmargin=*]
        \item 学习压力、时间紧张 → 不吃早餐率高。
        \item 在外就餐增多 → 高油高盐高能量膳食 + 蔬果摄入不足 → 超重肥胖风险增加。
    \end{itemize}

    \item \textbf{超重肥胖：}
    \begin{itemize}[leftmargin=*]
        \item 能量摄入过剩 + 缺乏运动 + 含糖饮料摄入过多。
        \item 预防：控制总能量，增加体力活动（每天至少60分钟中高强度运动），减少屏幕时间。
    \end{itemize}

    \item \textbf{锌和碘缺乏：}
    \begin{itemize}[leftmargin=*]
        \item 锌缺乏 → 青春期生长迟缓和性发育延迟（锌RNI 9-11.5 mg/d）。
        \item 碘缺乏 → 甲状腺肿和认知功能下降（碘RNI 110-120 $\mu$g/d）。
    \end{itemize}

    \item \textbf{情绪化进食与饮食行为偏离：}
    \begin{itemize}[leftmargin=*]
        \item 青春期内分泌波动 + 心理压力 → 情绪化进食（暴饮暴食）、厌食。
        \item 同伴影响力上升 → 饮食行为易受同伴影响（偏食、节食模仿等）。
    \end{itemize}
\end{enumerate}
\end{defbox}

% =====================================================================
\section{老年人营养}
% =====================================================================

\subsection{老年人生理特点与少肌症}

\begin{defbox}
\textbf{老年人主要生理变化：}

\textbf{1. 体成分变化——少肌症（Sarcopenia）：}
\begin{itemize}[leftmargin=*]
    \item 50岁以后，骨骼肌量平均每年减少1\%-2\%。
    \item 60岁以上慢性肌肉丢失估计约30\%，80岁以上为50\%。
    \item 表现为骨骼肌质量、力量和功能的进行性丧失。后果：跌倒、骨折、失能、生活质量下降和死亡率增加。
    \item 体内水分和瘦组织减少，体脂含量增加（尤其向心性脂肪分布）。
\end{itemize}

\textbf{2. 代谢变化：}
\begin{itemize}[leftmargin=*]
    \item 基础代谢率（BMR）下降：约每10年下降2\%。
    \item 分解代谢 $>$ 合成代谢：常出现负氮平衡。
\end{itemize}

\textbf{3. 消化系统功能减弱：}
\begin{itemize}[leftmargin=*]
    \item 胃酸分泌减少（萎缩性胃炎高发 → 影响VitB$_{12}$和钙吸收）。
    \item 肠蠕动减慢 → 便秘。
    \item 味觉（味蕾数量减少）和嗅觉敏感度下降 → 食欲降低。
\end{itemize}
\end{defbox}

\subsection{老年人营养需要}

\begin{keybox}
\textbf{老年人关键营养素RNI：}

\begin{center}
\begin{tabular}{lcc}
\hline
\textbf{营养素} & \textbf{50-64岁} & \textbf{65岁以上} \\
\hline
蛋白质（男/女，g/d） & 65 / 55 & 72 / 62 \\
钙（mg/d） & 800 & 800 \\
维生素D（$\mu$g/d） & 10 & 15 \\
\hline
\end{tabular}
\end{center}

\textbf{重要说明：}
\begin{enumerate}[leftmargin=*]
    \item \textbf{蛋白质需要量不减反增：}老年人蛋白质推荐摄入量为\textbf{1.0-1.5 g/(kg$\cdot$bw)/d}以预防少肌症。65岁以上：男72 g/d，女62 g/d——均高于50岁组的65 g/d和55 g/d。原因：合成抵抗（Anabolic Resistance）——衰老肌肉对较低剂量氨基酸/蛋白质的合成代谢反应下降，需要更多蛋白质刺激同等幅度的肌蛋白合成。
    \item \textbf{能量需求降低：}BMR下降和身体活动减少。
    \item \textbf{脂肪：}供能比20\%-30\%。
    \item \textbf{碳水化合物：}供能比50\%-65\%，选择低GI食物。
    \item \textbf{维生素D：}65岁以上RNI = 15 $\mu$g/d（成人10 $\mu$g/d）。因皮肤合成能力下降和日照减少。
    \item \textbf{水：}老年人渴觉敏感性下降，易发生隐性脱水。每日饮水量不少于1500 ml。
\end{enumerate}
\end{keybox}

\subsection{老年人膳食指南（2022版）}

\begin{keybox}
\textbf{《中国老年人膳食指南》条目（2022）：}
\begin{enumerate}[leftmargin=*]
    \item \textbf{食物品种丰富，动物性食物充足，常吃大豆制品。}
    \item \textbf{鼓励共同进餐，保持良好食欲。}
    \item \textbf{积极户外活动，延缓肌肉衰减，保持适宜体重。}
    \begin{itemize}[leftmargin=*]
        \item 结合抗阻运动（举轻哑铃、弹力带训练）最大程度对抗少肌症。
        \item BMI应不低于20 kg/m$^{2}$，不宜过度追求"瘦"。
    \end{itemize}
    \item \textbf{定期健康体检，测评营养状况，预防营养缺乏。}
\end{enumerate}
\end{keybox}

\begin{defbox}
\textbf{成功衰老（Successful Aging）：}在老年期尽可能延长健康寿命，维持高水平的生理和认知功能，仅于生命末期极短时间段出现功能下降和失能。三大要素（Rowe \& Kahn模型）：(1) 避免疾病和疾病相关失能；(2) 维持高水平的认知和身体功能；(3) 积极投入生活。
\end{defbox}

\subsection{老年人常见营养问题及预防}

\begin{defbox}
\textbf{老年人常见营养问题：}

\begin{enumerate}[leftmargin=*]
    \item \textbf{少肌症（Sarcopenia）：}
    \begin{itemize}[leftmargin=*]
        \item 定义：随年龄增长的骨骼肌质量、力量和功能的进行性丧失。
        \item 50岁后骨骼肌量年均减少1\%-2\%，60岁以上慢性肌肉丢失约30\%，80岁以上约50\%。
        \item 后果：跌倒、骨折、失能、生活质量下降和死亡率增加。
        \item 预防：充足优质蛋白质（1.0-1.5 g/(kg$\cdot$d)）+ 抗阻运动（举轻哑铃、弹力带训练）。
    \end{itemize}

    \item \textbf{蛋白质-能量营养不良（PEM）：}
    \begin{itemize}[leftmargin=*]
        \item 原因：食欲降低（味觉嗅觉减退）、咀嚼困难（牙齿脱落）、购买和烹饪能力下降、独居进食单调。
        \item 预防：保证动物性食物充足，常吃大豆制品，鼓励共同进餐。
    \end{itemize}

    \item \textbf{骨质疏松和骨折：}
    \begin{itemize}[leftmargin=*]
        \item 骨量随年龄下降，雌激素减少（女性绝经后骨丢失加速）。
        \item 预防：充足钙（800-1000 mg/d）和维生素D（65岁以上15 $\mu$g/d），规律负重运动。
    \end{itemize}

    \item \textbf{维生素D缺乏：}
    \begin{itemize}[leftmargin=*]
        \item 皮肤合成能力下降 + 日照减少 → 维生素D缺乏率高。
        \item 预防：多晒太阳，必要时补充制剂（65岁以上RNI 15 $\mu$g/d）。
    \end{itemize}

    \item \textbf{维生素B$_{12}$缺乏：}
    \begin{itemize}[leftmargin=*]
        \item 萎缩性胃炎高发 → 胃酸分泌减少 → 食物结合态B$_{12}$释放和吸收障碍。
        \item 预防：保证动物性食物摄入（肉类、鱼类、蛋类、奶类），严重者需补充制剂。
    \end{itemize}

    \item \textbf{便秘：}
    \begin{itemize}[leftmargin=*]
        \item 肠蠕动减慢 + 膳食纤维摄入不足 + 饮水不足 + 活动减少。
        \item 预防：增加膳食纤维和饮水（不少于1500 ml/d），适量运动。
    \end{itemize}

    \item \textbf{隐性脱水：}
    \begin{itemize}[leftmargin=*]
        \item 原因：渴觉敏感性下降 → 不自觉饮水不足。
        \item 预防：提醒喝水，每日饮水量不少于1500 ml（心肾功能正常者）。
    \end{itemize}

    \item \textbf{牙齿和口腔问题：}
    \begin{itemize}[leftmargin=*]
        \item 牙齿脱落、咀嚼困难 → 固体食物摄入受限 → 蔬菜、肉类摄入不足 → 多种营养素缺乏。
        \item 预防：食物制作应细软，保证荤素搭配，必要时制作成泥状。
    \end{itemize}
\end{enumerate}
\end{defbox}

% =====================================================================
\section{运动员营养}
% =====================================================================

\subsection{运动员营养代谢特点}

\begin{defbox}
\textbf{运动员五大系统生理改变：}

\begin{enumerate}[leftmargin=*]
    \item \textbf{心血管系统：}血容量增加，心输出量大幅提升，供氧供能需求暴涨。
    \item \textbf{神经系统：}超负荷训练易导致神经调节紊乱、疲劳、运动能力下降。
    \item \textbf{消化系统：}运动时胃肠道供血减少，营养素吸收减弱。
    \item \textbf{免疫系统：}长期高强度训练 → 免疫力下降、易呼吸道感染；中低强度有氧运动可提升免疫功能。
    \item \textbf{内分泌系统：}长期大强度训练可致激素紊乱；女性运动员多见月经不调、闭经。
\end{enumerate}

\textbf{整体代谢核心特点：}
\begin{itemize}[leftmargin=*]
    \item 能量消耗极高，易出现氧债，恢复期消耗仍大。
    \item 蛋白质分解增强、尿氮流失多，易发生运动性贫血和中枢疲劳。
    \item 高强度运动时脂肪氧化受限；糖原消耗快，是耐力短板。
    \item 大量出汗，水、钠钾钙镁等矿物质及水溶性维生素流失严重。
    \item 训练刺激微量营养素代谢，需求量高于普通人。
\end{itemize}
\end{defbox}

\subsection{运动员各类营养素需要}

\begin{defbox}
\textbf{1. 能量：}
\begin{itemize}[leftmargin=*]
    \item 总能量：3500-4400 kcal/d（210-280 kJ/kg）。
    \item 供能配比：蛋白质:脂肪:碳水化合物 = 1:(0.7-0.8):4。
    \item 影响因素：运动项目、强度、时长、体重、膳食结构。
\end{itemize}

\textbf{2. 蛋白质：}
\begin{itemize}[leftmargin=*]
    \item 适宜摄入：1.2-2.0 g/(kg$\cdot$d)，优质蛋白占1/3以上。
    \item 作用：修复肌肉、促进恢复、抗疲劳、预防运动性贫血。
    \item 摄入不足 → 耐力下降、乏力；过量加重肝肾负担。
\end{itemize}

\textbf{3. 脂肪：}
\begin{itemize}[leftmargin=*]
    \item 供能占总能量25\%-30\%；耐力项目可至50\%-60\%，极限可达70\%。
    \item 限制饱和脂肪，防血脂升高、降低耐力。
\end{itemize}

\textbf{4. 碳水化合物（核心供能物质）：}
\begin{itemize}[leftmargin=*]
    \item 高强度短时运动：碳水为主；长时间低强度才动用脂肪和蛋白质。
    \item \textbf{糖原填充法：}赛前1周低碳水耗糖原，赛前3天高碳水储备糖原（占膳食70\%）。
    \item 作用：快速供能、无代谢废物，决定耐力上限。
\end{itemize}

\textbf{5. 水与矿物质：}
\begin{itemize}[leftmargin=*]
    \item \textbf{补水：}运动大量出汗，少量多次补水，忌一次性暴饮；补液总量应大于出汗量，同步补电解质。
    \item 钠：高温训练流失多，低钠可致乏力恶心，食盐$<$5 g/d。
    \item 钾：流失致心律失常，推荐3-4 g/d，蔬果补充。
    \item 镁：多汗缺镁 → 抽筋，推荐400-500 mg/d。
    \item 钙：高强度+高温易缺钙、骨质疏松，推荐1000-1500 mg/d。
    \item 铁：训练流失多，易发运动性贫血；女性重点补铁，动物红肉补铁。
    \item 锌：大训练量锌流失、免疫力下降，推荐25 mg/d。
    \item 硒：提升免疫，推荐50-150 $\mu$g/d。
    \item 碘：150 $\mu$g/d，海产品补充。
\end{itemize}

\textbf{6. 维生素（需求量高于普通人群）：}
\begin{itemize}[leftmargin=*]
    \item 维生素B$_{1}$：3-5 mg/d，缺则丙酮酸堆积、神经疲劳；高碳水训练需增量。
    \item 维生素B$_{2}$：2-2.5 mg/d，缺乏损害有氧/无氧运动能力。
    \item 烟酸、B$_{6}$、B$_{12}$、叶酸：参与能量代谢、造血，缺乏致耐力下降。补充时机：赛前2-3周补制剂。
    \item 维生素C：训练推荐200 mg/d；促胶原修复、促铁吸收、抗氧化。
    \item 维生素A：视力需求高项目（射击、击剑）1800 $\mu$g RAE/d。
    \item 维生素D：10-12.5 $\mu$g/d；保护骨骼，缺乏易运动骨折。
    \item 维生素E：高原/耐力项目15-50 mg $\alpha$-TE/d；抗氧化、减少乳酸堆积；不可过量（蓄积中毒）。
\end{itemize}
\end{defbox}

% =====================================================================
\section{特殊环境人员营养}
% =====================================================================

\begin{tipbox}
\textbf{高温环境人员营养需要和膳食原则（了解）：}

\begin{enumerate}[leftmargin=*]
    \item \textbf{水和矿物质：}
    \begin{itemize}[leftmargin=*]
        \item 高温环境大量出汗 → 水和电解质大量丢失（钠、钾、钙、镁）。
        \item 少量多次补水，忌一次性暴饮；同时补充电解质（含钠、钾的运动饮料或淡盐水）。
        \item 食盐摄入比常温环境适当增加（通过膳食或饮水中补充）。
    \end{itemize}

    \item \textbf{能量：}高温环境基础代谢率升高（体温每升高1$^\circ$C，BMR增加约13\%），能量需要较常温增加。
    \item \textbf{蛋白质：}高温下蛋白质分解代谢增强，尿氮排出增加，蛋白质需要量相应提高（优质蛋白占1/2以上）。
    \item \textbf{维生素：}
    \begin{itemize}[leftmargin=*]
        \item 水溶性维生素（B族、C）随汗液流失显著增加。
        \item 维生素C推荐摄入提高至150-200 mg/d。
        \item 维生素B$_{1}$、B$_{2}$按能量比例相应增加（每1000 kcal：B$_{1}$ 0.5-0.6 mg，B$_{2}$ 0.55-0.65 mg）。
    \end{itemize}
    \item \textbf{膳食原则：}清淡可口，增进食欲；多提供汤类和液体食物；保证蔬菜、水果供给充足（补充水和矿物质）；适当增加食盐摄入。
\end{enumerate}

\textbf{低温环境人员营养需要和膳食原则（了解）：}

\begin{enumerate}[leftmargin=*]
    \item \textbf{能量：}低温环境下基础代谢率升高，机体需要额外产热以维持体温 → 能量需要显著增加（可比常温增加10\%-15\%以上）。
    \item \textbf{宏量营养素：}
    \begin{itemize}[leftmargin=*]
        \item 脂肪供能比可提高至35\%-40\%（增加能量密度和饱腹感）。
        \item 碳水化合物供能比维持在50\%-55\%（保证快速供能）。
        \item 蛋白质供给充足（优质蛋白占1/2以上），防止负氮平衡。
    \end{itemize}
    \item \textbf{维生素：}
    \begin{itemize}[leftmargin=*]
        \item 维生素C有助增强耐寒能力，推荐摄入量适当增加。
        \item 维生素B族需要量随能量增加而相应增加。
        \item 维生素A、D、E的摄入也应适当增加。
    \end{itemize}
    \item \textbf{矿物质：}寒冷气候下尿钠排出增加，钠需要量可适当提高；钙、铁、锌、碘的供给亦需充足。
    \item \textbf{膳食原则：}提供高能量密度食物，保证热食热饮供应；增加脂肪摄入比例；食物温热可口；保证蔬菜、水果；食盐摄入可适当增加。
\end{enumerate}
\end{tipbox}

% =====================================================================
\section{复习题}
% =====================================================================

\reviewcontent{
\section{复习题}

\subsection*{一、生命早期1000天与DOHaD}

\begin{enumerate}[leftmargin=*, itemsep=8pt]
    \item \textbf{【名词解释】}生命早期1000天
    \item \textbf{【名词解释】}DOHaD理论
    \item \textbf{【简答题】}简述DOHaD理论的核心含义及其"节俭表型"假说。
\end{enumerate}

\subsection*{二、孕妇营养}

\begin{enumerate}[leftmargin=*, itemsep=8pt, start=4]
    \item \textbf{【名词解释】}生理性贫血
    \item \textbf{【简答题】}列出孕期各阶段的能量和蛋白质增加量。
    \item \textbf{【简答题】}简述叶酸补充对孕期的重要性——为什么应从孕前3个月开始补充？
    \item \textbf{【简答题】}根据孕前BMI分类，列出各组孕期推荐体重增长范围及中晚期每周增长速率。
    \item \textbf{【简答题】}孕期碳水化合物最低摄入量为何不应低于130 g/d？
    \item \textbf{【简答题】}列出《备孕和孕期妇女膳食指南（2022年）》6条关键推荐。
    \item \textbf{【简答题】}简述孕期的内分泌变化——HCG、HCS、雌激素、孕激素各有什么作用和出现时间？
    \item \textbf{【简答题】}列举孕期常见营养问题及预防（至少4项）。
\end{enumerate}

\subsection*{三、乳母营养}

\begin{enumerate}[leftmargin=*, itemsep=8pt, start=12]
    \item \textbf{【简答题】}简述母乳三阶段（初乳、过渡乳、成熟乳）的主要成分差异。
    \item \textbf{【简答题】}列举母乳喂养的八大优点。
    \item \textbf{【简答题】}乳母的能量和蛋白质较非孕期分别增加多少？钙的需求有何特殊性？
    \item \textbf{【简答题】}列举乳母常见营养问题及预防（至少3项）。
\end{enumerate}

\subsection*{四、婴幼儿营养}

\begin{enumerate}[leftmargin=*, itemsep=8pt, start=16]
    \item \textbf{【名词解释】}回应式喂养、辅食
    \item \textbf{【简答题】}简述0-6月龄婴儿母乳喂养指南（2022）的6条准则。
    \item \textbf{【简答题】}简述7-24月龄婴幼儿喂养指南（2022）的6条准则。
    \item \textbf{【简答题】}简述辅食添加的意义（至少4点）及添加原则。
    \item \textbf{【简答题】}简述婴幼儿的铁营养特点——为何6m-2y易患缺铁性贫血？
    \item \textbf{【简答题】}母乳中乳糖的作用是什么？
    \item \textbf{【简答题】}列举婴幼儿常见营养问题及预防（至少5项）。
\end{enumerate}

\subsection*{五、学龄前儿童营养}

\begin{enumerate}[leftmargin=*, itemsep=8pt, start=23]
    \item \textbf{【简答题】}简述学龄前儿童的主要生理特点。
    \item \textbf{【简答题】}列出学龄前儿童膳食指南（2022）核心推荐。
    \item \textbf{【简答题】}学龄前儿童每日奶类、水和体育活动的推荐量分别是多少？
    \item \textbf{【简答题】}列举学龄前儿童常见营养问题及预防。
\end{enumerate}

\subsection*{六、学龄儿童与青少年营养}

\begin{enumerate}[leftmargin=*, itemsep=8pt, start=27]
    \item \textbf{【简答题】}简述学龄儿童骨骼发育的特点。
    \item \textbf{【简答题】}列出学龄儿童膳食指南（2022）5条核心推荐。
    \item \textbf{【简答题】}简述生长发育的4个规律。
    \item \textbf{【简答题】}列举学龄儿童常见营养问题及预防（至少3项）。
    \item \textbf{【简答题】}简述青少年体格发育的特点（身高、体重、形态）。
    \item \textbf{【简答题】}青春期为何是峰值骨量积累的关键期？钙摄入不足的远期后果是什么？
    \item \textbf{【简答题】}为什么青春期女性是缺铁性贫血高发人群？
    \item \textbf{【简答题】}列举青少年常见营养问题及预防（至少4项）。
\end{enumerate}

\subsection*{七、老年人营养}

\begin{enumerate}[leftmargin=*, itemsep=8pt, start=35]
    \item \textbf{【名词解释】}少肌症（Sarcopenia）、合成抵抗（Anabolic Resistance）
    \item \textbf{【简答题】}老年人蛋白质需求量"不减反增"的原因是什么？推荐摄入量是多少？
    \item \textbf{【简答题】}列出《中国老年人膳食指南》（2022）4条条目。
    \item \textbf{【简答题】}列举老年人常见营养问题及预防（至少5项）。
    \item \textbf{【论述题】}从"生命早期1000天 → 成年期 → 老年期"的时间维度，论述生命周期不同阶段营养与健康的连续联系，特别是DOHaD理论和老年人少肌症的关键节点。
\end{enumerate}

\subsection*{八、运动员营养与特殊环境人员营养}

\begin{enumerate}[leftmargin=*, itemsep=8pt, start=40]
    \item \textbf{【简答题】}简述运动员的营养代谢特点（五大系统生理改变及整体代谢特点）。
    \item \textbf{【简答题】}简述运动员各类营养素需要要点（碳水化合物、蛋白质、水与矿物质）。
    \item \textbf{【简答题】}简述高温和低温环境人员的营养需要和膳食原则。
\end{enumerate}

\subsection*{参考答案要点}

\begin{enumerate}[leftmargin=*, itemsep=5pt]
    \item \textbf{生命早期1000天：}孕期270天 + 0-6月龄180天（母乳喂养）+ 7-24月龄（辅食添加）。是生长发育的关键期和敏感期，营养状态对终身健康产生深远和不可逆的影响。
    \item \textbf{DOHaD理论（健康与疾病的发育起源）：}生命早期（胎儿、婴幼儿期）经历不利因素（如营养不良、环境不良等），将会增加其成年后患肥胖、糖尿病、心血管疾病等慢性疾病的几率。核心机制为发育程序化（表观遗传调控）。
    \item DOHaD核心含义：早期营养环境通过表观遗传机制（DNA甲基化、组蛋白修饰）对远期疾病风险产生持久性影响。节俭表型假说：宫内营养不良使胎儿产生"预测性适应"，若出生后实际面临"营养过剩"环境，则代谢-环境不匹配大幅增加成年期代谢综合征风险。
    \item \textbf{生理性贫血：}孕期血容量增加，血浆容积增加幅度 $>$ 红细胞增加幅度，导致血液稀释。通常在孕20-30周最明显。WHO标准：孕期Hb $\geq$ 110 g/L，孕中期 $\geq$ 105 g/L。
    \item 孕早期：E+0，P+0；孕中期：E+250 kcal/d，P+15 g/d；孕晚期：E+400 kcal/d，P+30 g/d。
    \item 叶酸缺乏致神经管缺陷（NTDs）。孕后3-8周是神经管闭合关键期（此时许多妇女尚不知已怀孕），故需孕前3个月开始补充。推荐量：400 $\mu$g DFE/d。
    \item 低体重（BMI$<$18.5）：总11.0-16.0 kg，每周0.46 kg；正常（18.5-24.0）：总8.0-14.0 kg，每周0.37 kg；超重（24.0-28.0）：总7.0-11.0 kg，每周0.30 kg；肥胖（$\geq$28.0）：总5.0-9.0 kg，每周0.22 kg。
    \item 防止酮体对胎儿中枢神经系统的毒性。严重早孕反应不能进食时需静脉补充葡萄糖。
    \item 6条关键推荐：(1)调整孕前体重至正常范围，保证孕期体重适宜增长；(2)常吃含铁丰富的食物，选用碘盐，合理补充叶酸和维生素D；(3)孕吐严重者，可少量多餐，保证摄入含必需量碳水化合物的食物；(4)孕中晚期适量增加奶、鱼、禽、蛋、瘦肉的摄入；(5)经常户外活动，禁烟酒，保持健康生活方式；(6)愉快孕育新生命，积极准备母乳喂养。
    \item HCG：受精1周开始出现，8-9周达高峰，维持妊娠早期孕酮分泌。HCS：妊娠6周出现，36周达高峰至分娩，促进乳腺发育和胎儿生长，降低胰岛素敏感性。雌激素：分泌增加，增加子宫和胎盘血流量，促进乳房发育。孕激素：分泌增加，减少子宫收缩，维持妊娠。
    \item 孕期常见营养问题：(1)早孕反应——HCG升高+贲门松弛，少量多餐保证碳水130g/d；(2)孕期便秘——孕酮致肠蠕动减慢+子宫压迫；(3)GDM——生理性胰岛素抵抗，控制体重和低GI饮食；(4)妊娠高血压——与钙不足可能有关，保证钙800mg/d；(5)缺铁性贫血——血容量扩张+胎铁转运，孕中25mg/d、孕晚29mg/d铁；(6)钙营养——钙吸收率代偿升高，800mg/d即可。
    \item 初乳（第1周）：淡黄色，蛋白质多，富含免疫球蛋白（sIgA和乳铁蛋白），乳糖和脂肪相对较少。过渡乳（第2周）：蛋白质$\downarrow$，乳糖和脂肪$\uparrow$。成熟乳（第2周以后）：乳白色，蛋白质相对较少，乳糖和脂肪多。
    \item 八大优点：(1)乳清蛋白为主，必需氨基酸比例适当；(2)脂肪颗粒小+乳脂酶→易消化吸收；(3)富含乳糖→促进乳酸杆菌、抑制大肠埃希菌，利于钙铁锌吸收；(4)钙磷比适宜（2:1）；(5)含大量免疫物质；(6)不易过敏；(7)经济方便卫生；(8)促进产后恢复，增进母婴交流。
    \item 能量：+400 kcal/d。蛋白质：+25 g/d。钙：RNI 800 mg/d（不增加），乳汁中钙含量较稳定（每日排出约300 mg），乳母钙供给不足时会动用自身骨骼中的钙满足乳汁钙含量——哺乳期骨骼密度下降，断奶后可逆性恢复。
    \item 乳母常见营养问题3项：(1)能量和营养素摄入不足——产后过度节食+饮食不规律→泌乳量减少，能量+400kcal/d蛋白+25g/d；(2)骨骼健康——乳汁每日排钙约300mg，钙不足动用骨钙，保证800mg/d；(3)产后贫血——分娩失血+哺乳消耗，铁不能通过乳腺入乳，母亲需注意补铁。
    \item \textbf{回应式喂养：}按婴儿饥饿信号适时喂养，而非固定时间刻板喂养。\textbf{辅食：}除母乳外，为满足婴幼儿生长发育需要而添加的其他食物。
    \item 6条准则：(1)母乳是婴儿最理想的食物，坚持6月龄内纯母乳喂养；(2)生后1小时内开奶；(3)回应式喂养；(4)适当补充维生素D，母乳喂养无需补钙；(5)任何动摇母乳喂养的想法和举动都必须咨询医生；(6)定期监测体格指标。
    \item 6条准则：(1)继续母乳喂养，满6月龄起必须添加辅食，从富含铁的泥糊状食物开始；(2)及时引入多样化食物，重视动物性食物的添加；(3)尽量少加糖、盐，油脂适当，保持食物原味；(4)提倡回应式喂养，鼓励但不强迫进食；(5)注重饮食卫生和进食安全；(6)定期监测体格指标，追求健康生长。
    \item 辅食添加意义：(1)补充营养素——6月后母乳不能满足全部需要，99\%铁需从辅食获得；(2)促进消化系统发育；(3)培养咀嚼吞咽能力；(4)促进味觉发育和饮食习惯形成；(5)促进神经心理发育。添加原则：每次一种新食物观察3-5天；由少到多由稀到稠由细到粗；健康消化功能正常时添加；不加盐糖；1岁后淡口味。
    \item 足月新生儿有足够铁贮存供4-6个月需要。6月龄时体内储存铁已耗尽，7-12月龄99\%的铁必须从辅食中获得。6m-2y为缺铁性贫血高发期。7-12月龄铁RNI = 10 mg/d。
    \item 乳糖促进乳酸杆菌生长、抑制大肠埃希菌生长；利于钙、铁、锌吸收。
    \item 婴幼儿常见营养问题5项：(1)缺铁性贫血——6m-2y高发，补充富铁辅食；(2)VD缺乏性佝偻病——生后数开始补VD 400IU/d；(3)PEM——辅食不及时质差；(4)锌缺乏——生长迟缓食欲差；(5)碘缺乏——克汀病，使用碘盐；(6)VK缺乏——新生儿肌注VK；(7)龋齿——1岁前不加糖不喝含糖饮料。
    \item (1)每年体重增长约2 kg，身高约5-7 cm；(2)神经系统仍属脑细胞生长发育关键期；(3)3岁乳牙已出齐，6岁恒牙已萌出，但咀嚼和消化能力远低于成人；(4)注意力分散（约15分钟），无法专心进食。
    \item (1)食物多样，规律就餐，自主进食；(2)每天饮奶（300-500 mL），足量饮水（600-800 mL），合理选择零食；(3)合理烹调，少调料少油炸；(4)参与食物选择制作；(5)经常户外活动，定期体格测量。
    \item 奶类每日300-500 mL；水每日600-800 mL；每天至少60分钟体育活动；每天视屏时间不超过2小时。
    \item 学龄前儿童常见营养问题：(1)挑食偏食——家长以身作则多次尝试；(2)零食过多——选健康零食两餐之间；(3)含糖饮料——白开水为主；(4)钙不足——奶类300-500ml/d；(5)铁锌不足——保证肉类肝脏；(6)超重肥胖——控能量+户外活动60min/d。
    \item 骨骼特点：骨胶多、钙盐少弹性大、硬度小，不易骨折而易弯曲变形。
    \item (1)主动参与食物制作和选择，提高营养素养；(2)吃好早餐，合理选择零食，培养健康饮食行为；(3)天天喝奶，足量饮水，不喝含糖饮料，禁止饮酒；(4)多户外活动，少视屏时间，每天60分钟以上中高强度身体活动；(5)定期监测体格发育，保持体重适宜增长。
    \item 生长发育4规律：(1)连续性；(2)不同步；(3)波浪式；(4)差异性。
    \item 学龄儿童常见营养问题：(1)不吃早餐或早餐质量差——早餐三类食物以上；(2)零食甜饮料过度——以白开水为日常饮水；(3)钙和VD不足——天天喝奶+户外活动；(4)超重肥胖或消瘦两极分化——平衡膳食+每天60min中高强度活动；(5)脊柱骨骼问题——保证钙摄入+正确坐姿+运动。
    \item 身高：性别差异——男较女晚2年，男7-9 cm/y，女5-7 cm/y。体重：4-5 kg/y，时间较长。形态：四肢比脊柱早，脚最先，从远到近。
    \item 青春期是获得峰值骨量的最关键阶段，约45\%的成人骨量在青春期积累。钙摄入不足 → 无法达到遗传决定的峰值骨量 → 老年期骨质疏松风险大幅增加。钙RNI：1000-1200 mg/d。
    \item 青春期女性需弥补月经铁的丢失，同时生长发育需要大量铁。11-13岁：18 mg/d；14-17岁：18 mg/d。是缺铁性贫血高发人群。
    \item 青少年常见营养问题：(1)青春期女性IDA——月经铁丢失+生长，铁RNI 18mg/d；(2)钙不足影响峰值骨量——1000-1200mg/d；(3)不健康减肥——盲目节食致能量和营养素不足影响发育；(4)不吃早餐和在外就餐——油盐多蔬果少；(5)超重肥胖——控能量+60min中高强度运动；(6)锌碘缺乏——锌9-11.5mg/d碘110-120$\mu$g/d；(7)情绪化进食/饮食行为偏离。
    \item \textbf{少肌症：}50岁后骨骼肌量平均每年减少1\%-2\%，60岁以上慢性肌肉丢失约30\%，80岁以上为50\%。表现为骨骼肌质量、力量和功能的进行性丧失。\textbf{合成抵抗：}衰老肌肉对较低剂量氨基酸/蛋白质的合成代谢反应下降，需要更多蛋白质刺激同等幅度的肌蛋白合成。
    \item 原因：合成抵抗（Anabolic Resistance）——老年人蛋白质合成对摄入蛋白质的敏感性下降，分解代谢亢进。推荐摄入量：1.0-1.5 g/(kg$\cdot$bw)/d。65岁以上：男72 g/d，女62 g/d（高于50岁组的65和55 g/d）。
    \item (1)食物品种丰富，动物性食物充足，常吃大豆制品；(2)鼓励共同进餐，保持良好食欲；(3)积极户外活动，延缓肌肉衰减，保持适宜体重；(4)定期健康体检，测评营养状况，预防营养缺乏。
    \item 列举老年人常见营养问题的5项以上：(1)少肌症——50岁后骨骼肌年均减1\%-2\%，蛋白质1.0-1.5 g/(kg$\cdot$d)+抗阻运动；(2)PEM——食欲降低、咀嚼困难致营养不足，保证动物性食物和共同进餐；(3)骨质疏松——钙+VD+负重运动；(4)VD缺乏——晒太阳+65岁以上RNI 15 $\mu$g/d；(5)便秘——增加膳食纤维和饮水；(6)隐性脱水——渴觉减退，每日饮水不少于1500 ml；(7)VB$_{12}$缺乏——萎缩性胃炎致吸收障碍；(8)牙齿口腔问题——食物细软。
    \item \textbf{论述题答题要点：}
    \begin{itemize}[leftmargin=*]
        \item \textbf{生命早期1000天 → DOHaD：}孕期270d + 0-6月龄180d + 7-24月龄为机遇窗口。早期营养不良经表观遗传（发育程序化）导致远期代谢综合征高风险。节俭表型假说——宫内"预测匮乏"+出生后"营养过剩"→ 代谢失匹配。
        \item \textbf{成年期：}膳食行为和生活方式继续强化或修正早期程序化效应。膳食指南2022八准则贯穿全生命周期。
        \item \textbf{老年期 → 少肌症与合成抵抗：}蛋白质需要量不减反增（1.0-1.5 g/(kg$\cdot$bw)/d），能量需求下降但蛋白质需求上升形成"老年人营养悖论"。需配合抗阻运动以实现成功衰老。
        \item \textbf{核心结论：}营养影响不是点状而是连续的——早期干预效费比最高。从生命早期1000天到老年期，营养贯彻生命周期理念。
    \end{itemize}

    \item 运动员五大系统生理改变：心血管（血容量和心输出量增加）、神经（超负荷易疲劳紊乱）、消化（运动时供血减少影响吸收）、免疫（高强度训练致免疫下降，中低强度有氧提升免疫）、内分泌（大强度可致激素紊乱，女性月经不调闭经）。整体代谢特点：能量消耗极高，蛋白分解增强尿氮流失多，糖原消耗快为耐力短板，大量出汗致水和矿物质维生素流失，微量营养素需求高于常人。
    \item 运动员营养素需要要点——碳水化合物：高强度短时运动核心供能，糖原填充法（赛前1周低碳水耗糖原，赛前3天高碳水储备占膳食70\%）；蛋白质：1.2-2.0 g/(kg$\cdot$d)，优质蛋白1/3以上；水与矿物质：少量多次补液，总补液量大于出汗量同步补电解质；钠钾钙镁铁锌硒碘均需充分补充。
    \item 高温环境人员：大量出汗致水电解质流失，少量多次补水+补电解质；BMR升高能量增加；蛋白分解增强；水溶性维生素随汗液流失显著增加（VC 150-200 mg/d）；清淡可口膳食多汤类多蔬果。低温环境人员：BMR升高能量显著增加（10\%-15\%以上）；脂肪供能比可提高至35\%-40\%；碳水50\%-55\%；蛋白充足；VC、B族增加；食盐可适当增加；保证热食热饮，食物温热可口。
\end{enumerate}

\newpage
}

% =====================================================================
%  CHAPTER 5: 公共营养
% =====================================================================
\chapter{公共营养}
\chapterintro{
\keyword{掌握}DRIs七项指标体系（EAR/RNI/AI/UL/AMDR/PI-NCD/SPL）的定义、制定逻辑和应用；掌握膳食模式与膳食指南（中国居民膳食指南2022八条核心准则与膳食宝塔2022构成及应用）；掌握营养调查膳食调查方法及膳食营养评价方法和依据；掌握BMI分级标准及NRV概念与应用。\keyword{熟悉}营养调查与评价（调查目的、内容、综合营养状况评价方法、营养监测概念及指标）；熟悉公共营养改善（营养改善措施、营养教育、营养干预的方法和应用）；熟悉食品营养标签定义内容及管理规范；熟悉功能食品概念分类及管理；熟悉食谱编制基本原则方法；熟悉全国营养调查概况。\keyword{了解}公共营养的定义与地位；了解我国营养相关政策与法规。
}

% =====================================================================
\section{概述}
% =====================================================================

\begin{defbox}
\textbf{公共营养的定义：}从医学状况+人群社会视角的群体概念。定义范围涵盖：病人、老人、慢性病患者、营养不良患者。人群分类：高危人群、亚健康人群、一般人群和特殊人群。

\textbf{合理营养与平衡膳食：}
\begin{itemize}[leftmargin=*]
    \item \textbf{合理营养（rational nutrition）：}指通过膳食提供给机体种类齐全、数量充足、比例适宜的能量和各种营养素，并与机体需要相平衡，从而达到合理营养、促进健康、预防疾病的膳食。又称为平衡膳食（balanced diet）。
    \item \textbf{平衡膳食的要求：}1. 提供种类齐全、数量充足、比例适宜的营养素；2. 保证食物安全；3. 科学烹调加工；4. 合理的进餐制度和良好的饮食习惯。
\end{itemize}

\textbf{营养不良（malnutrition）：}指由于一种或一种以上营养素的缺乏或过剩所造成的机体健康异常或疾病状态。营养不良包括两种表现：营养缺乏（nutrition deficiency）和营养过剩（nutrition excess）。
\end{defbox}

\begin{defbox}
\textbf{合理膳食与公共营养的对应关系：}
\begin{itemize}[leftmargin=*]
    \item 营养素 $\rightarrow$ 营养标准（DRIs）
    \item 食物 $\rightarrow$ 膳食指南/食物指导
    \item 膳食 $\rightarrow$ 膳食模式
\end{itemize}

\textbf{公共营养的地位与作用：}
\begin{enumerate}[leftmargin=*]
    \item 第一，促进全人群的营养与健康。
    \item 第二，增强国民体质，增强综合国力。
    \item 第三，为国家经济社会发展和民族复兴提供建议，事关国家发展的战略性问题。
\end{enumerate}
\end{defbox}

\begin{defbox}
\textbf{公共营养的工作内容：}
\begin{enumerate}[leftmargin=*]
    \item 营养标准的制定（DRIs）
    \item 居民膳食指南的制定
    \item 各类人员营养培训与指导
    \item 营养调查与评价
    \item 营养监测
    \item 营养干预
    \item 食物营养规划与营养立法
    \item 其他营养公众指导（以食物为研究对象、食物营养政策和法规）
\end{enumerate}

\textbf{公共营养的管理机构：}
\begin{itemize}[leftmargin=*]
    \item 全国营养工作的组织协调：中国营养学会
    \item 制定营养素标准：国家卫生健康委营养标准委员会
    \item 制定营养相关政策及法规：食品安全法、营养改善工作管理办法、全国营养调查工作规范、食品营养标签管理规范、《2025-2030年中国食物与营养发展纲要》、2017年7月《国民营养计划（2017-2030）》
\end{itemize}
\end{defbox}

% =====================================================================
\section{膳食营养素参考摄入量（DRIs）}\difficulty
% =====================================================================

% [补充自往届笔记] - DRIs制定基础说明
每个人每日从膳食中获取各种营养素及能量。对某种营养素的需要量因年龄、性别、生理状况而不同（如儿童、孕妇、老人、哺乳期妇女）。人们所需要的各种营养素都从每日膳食中获得。科学安排每日膳食才能满足人体的营养素需要。某种营养素长期摄入不足或摄入过多都会产生危害。因此针对各年龄、性别和劳动强度状态人群，制定了膳食营养素参考摄入量（dietary reference intakes, DRIs），用于指导膳食设计和评价膳食计划。

\begin{defbox}
\textbf{DRIs（Dietary Reference Intakes）：}一组每日平均膳食营养素摄入量的参考值。是评价膳食营养素供给量能否满足人体需要及是否存在过量摄入风险的一组参考值。

\textbf{从RDAs到DRIs的发展：}
\begin{itemize}[leftmargin=*]
    \item RDAs：营养目标——预防营养缺乏病
    \item DRIs：预防慢性病——帮助个体合理安排膳食、获得良好健康状态的新概念。用途：健康个体的膳食计划、膳食评估；并用于制定政策计划。
\end{itemize}
\end{defbox}

\begin{keybox}
\textbf{DRIs七项指标体系（DRIs 2013版）：}

\begin{table}[H]
\centering
\begin{tabularx}{\textwidth}{lXl}
\toprule
\textbf{指标} & \textbf{定义} & \textbf{制定目标} \\
\midrule
\imp{平均需要量（EAR）} & 某一特定性别、年龄及生理状况群体中对某营养素需要量的平均值。摄入量达到EAR水平时可以满足群体中50\%个体的需要，不能满足另外50\%个体对该营养素的需要。 & 预防缺乏 \\
\hline
\imp{推荐摄入量（RNI）} & 相当于RDA，可以满足某一特定性别、年龄及生理状况群体中绝大多数（97\%-98\%）个体需要量的摄入水平。长期摄入RNI水平，可以满足个体对该营养素的需要，维持组织中有适当的营养素储备和机体健康。RNI是以EAR为基础制定的。用途：作为个体每日摄入该营养素的目标值。 & 预防缺乏 \\
\hline
\imp{适宜摄入量（AI）} & 当某种营养素个体需要量的资料不足，没有办法计算出EAR，因而不能求得RNI时，可设定适宜摄入量来代替RNI。AI是通过观察或实验获得的健康人群某种营养素的摄入量，是该营养素的一个安全摄入水平，而不是通过研究营养素的个体需要量得到的。用途：作为个体每日营养素摄入的目标值。 & 预防缺乏 \\
\hline
\imp{可耐受最高摄入量（UL）} & 平均每日可以摄入营养素的最高限量。摄入量达到UL对一般人群中的几乎所有个体都不至于损害健康。在规划膳食时应使营养素摄入量低于UL，以避免造成不良反应的潜在风险。 & 防止过量 \\
\hline
\imp{宏量营养素可接受范围（AMDR）} & 指碳水化合物、脂肪和蛋白质的理想摄入量范围，以占能量摄入量的百分比(\%E)表示。该范围可提供人体对这些必需营养素的需要，并且有利于降低慢性病的发生危险。其下限为预防营养缺乏，上限为降低慢性非传染性疾病风险。如果一个个体的摄入量高于或低于推荐的范围，可能增加罹患慢性病的风险，或者导致必需营养素缺乏的可能性增加。 & 预防慢性病 \\
\hline
\imp{预防非传染性慢性病的建议摄入量（PI-NCD）} & 膳食营养素摄入过高导致的慢性病（肥胖、糖尿病、高血压、血脂异常、脑卒中、冠心病以及某些癌症等）。PI-NCD以非传染性慢性病的一级预防为目标，提出的必需营养素的每日摄入量。当NCD易感人群某些营养素的摄入量接近或达到PI值时，可以降低发生NCD的风险。某些营养素PI可能高于RNI或AI（如维生素C、钾）；另一些营养素可能低于AI（如钠）。 & 预防慢性病 \\
\hline
\imp{特定建议值（SPL）} & 研究证明某些食物成分具有改善生理功能、预防营养相关慢性病的生物学作用，其中多数为植物化学物。SPL以降低人群中膳食相关NCD风险为目标，提出其他食物成分的每日摄入水平，当该成分的摄入量达到SPL时，可能有利于降低疾病的发生风险或改善健康状况。 & 预防慢性病 \\
\bottomrule
\end{tabularx}
\end{table}
\end{keybox}

\begin{keybox}
\textbf{RNI的制定逻辑：}
RNI是以EAR为基础制定的。如果人群需要量呈正态分布，已知EAR的标准差，则RNI = EAR + 2SD表示。如果营养素需要量资料不符合正态分布，不能计算SD时，一般采用EAR乘以变异系数（10\%），即RNI = 1.2 × EAR。

\textbf{RNI设定注意事项：}RNI是根据某一特定人群中体重在正常范围内的个体的需要量设定的。对于个别身高、体重超出参考范围较多的个体，可能需要按每千克体重的需要量调整其RNI。
\end{keybox}

\begin{keybox}
\textbf{AI与RNI的异同：}
\begin{itemize}[leftmargin=*]
    \item \textbf{相同点：}都可以作为群体中个体营养素摄入量的目标值，可以满足群体中几乎所有个体的需要。
    \item \textbf{不同点：}AI的准确性不及RNI，主要用于个体营养素摄入目标；群体中营养素平均摄入量达到或超过AI，则该人群中摄入不足者的危险性很小。
\end{itemize}
\end{keybox}

\begin{keybox}
\textbf{UL的制定注意事项：}
\begin{itemize}[leftmargin=*]
    \item 随着营养素强化食品和膳食补充剂的广泛使用，需要制定ULs以指导安全消费。
    \item 如果某营养素的毒副作用与摄入总量有关，则该营养素的UL值根据食物、饮水及补充剂提供的总量来制定。
    \item 如果毒副作用仅与强化食物和补充剂有关，则UL根据这些来源的总摄入量来制定。
    \item 对于某些营养素而言，没有足够资料来制定其UL，但未定UL并不意味着过多摄入没有潜在的危害。
\end{itemize}
\end{keybox}

\begin{tipbox}
\textbf{DRIs的制定与修订程序：}
人群健康与营养状况调查和监测 $\rightarrow$ 制定DRIs $\rightarrow$ 膳食指导 $\rightarrow$ 食物指导 $\rightarrow$ 预包装食品营养标签 $\rightarrow$ 人群营养改善和健康状况监测反馈 $\rightarrow$ 营养科学与有关科学研究（实验室、临床观察、流行病学调查）$\rightarrow$ 修订DRIs。注意：每4-5年修订一次。

\textbf{DRIs的用途与层次：}
\begin{itemize}[leftmargin=*]
    \item \textbf{个体用途：}膳食规划、制定食谱、食品配方研发、食物保障/供给计划。
    \item \textbf{群体用途：}膳食评价（膳食评估和改进、营养咨询、营养风险）。
    \item \textbf{安全作用：}作为标准/依据，衔接个体膳食规划和群体膳食评价。
\end{itemize}

\textbf{评价逻辑：}个体营养素需要量 VS 膳食营养素供给量/摄入量 $\leftrightarrow$ 膳食营养素参考摄入量（DRIs）
\end{tipbox}

\begin{defbox}
\textbf{能量需要量（EER, Estimated Energy Requirement）：}
EER是指能长期保持良好健康状态、维持体重和机体构成以及体力活动水平的个体及人群，达到能量平衡时所需要的膳食能量摄入量。

群体能量推荐摄入量直接等同于该群体的能量EAR，而不同于蛋白质等其他营养素采用EAR加上2个标准差的方法。

EER的制定需考虑年龄、性别、体重、身高和体力活动的不同：
\begin{itemize}[leftmargin=*]
    \item 成人EER：一定年龄、性别、体重、身高和体力活动的健康人群，维持能量平衡所需的膳食能量。
    \item 儿童EER：一定年龄、性别3岁以上儿童正常体重、身高的个体，维持能量平衡和正常生长所需膳食能量。
    \item 孕妇EER：在成人EER基础上，再加上胎儿组织生长所需的膳食能量。
    \item 乳母EER：在成人EER基础上，再加上乳汁分泌所需的膳食能量。
\end{itemize}
\end{defbox}

% =====================================================================
\section{膳食模式与膳食指南}
% =====================================================================

\subsection{膳食结构}

\begin{defbox}
\textbf{膳食结构（Dietary Structure）的定义：}
膳食中各类食物的种类、数量及其在膳食中所占的比例。影响因素：历史时期、民族/地区、社会阶层。特性：可变性，可预见性和相对稳定性及代表性。核心构成：粮食、动物、其他。

\textbf{膳食结构研究的意义：}
\begin{itemize}[leftmargin=*]
    \item 研究一个国家或一个地区人群膳食特点和营养状况的基础。
    \item 比较差别和变化，了解发展、饮食文化、习惯的变化规律、可能的健康风险。
    \item 为制定相应的膳食指导和营养指导提供科学依据。
\end{itemize}
\end{defbox}

\begin{defbox}
\textbf{膳食模式 $\approx$ 膳食结构：}
将日常膳食中成分组合而构成的独特食物和营养素摄入的方式，即膳食的总体构成。

\textbf{素食（VU膳食）定义：}指膳食中不含畜肉、禽肉、鱼类及其制品，且不含动物性制品和动物油脂的饮食习惯。其膳食构成不同于一般膳食，能量密度较低，富含DF（膳食纤维）、镁、叶酸、维生素C、E、n-6 PUFA（多不饱和脂肪酸）、植物化学物和较低的胆固醇、SFA（饱和脂肪酸）。
\end{defbox}

\begin{keybox}
\textbf{膳食结构类型及与健康的关系：}

\begin{enumerate}[leftmargin=*]
    \item \textbf{动物性食物为主的膳食结构（西方模式/经济发达国家模式）}——欧美经济发达国家。
    \begin{itemize}
        \item 营养过剩型，"三高"一低（高能量、高脂肪、高蛋白、低膳食纤维）。
        \item 健康风险：肥胖、心脑血管疾病等代谢性疾病、结肠癌、直肠癌等。
    \end{itemize}

    \item \textbf{植物性食物为主的膳食结构（东方模式/发展中国家模式）}——亚洲、非洲部分国家/地区。
    \begin{itemize}
        \item 营养缺乏型，"一高三低"（低蛋白质、低脂肪、植物性食物为主）。
        \item 健康风险：感染性疾病增加、营养缺乏症与发达国家类似的慢性病同时并存。
    \end{itemize}

    \item \textbf{动植物食物平衡的膳食结构}——日本为代表，集合东西方膳食特点。
    \begin{itemize}
        \item 动植物消费均衡、水产品食物量大，三大宏量营养素供能比合理。
        \item 近40年也在改变中，水产品消费量依然在增加，动物性食物消费量增加。
    \end{itemize}

    \item \textbf{地中海膳食结构}——希腊、西班牙、法国和意大利南部等。
    \begin{itemize}
        \item 橄榄油为主要脂肪来源（不饱和脂肪高），大量蔬菜水果、全谷物、豆类，适量鱼类和红酒，饱和脂肪低。
        \item 大部分成年人有饮用红酒的习惯。
    \end{itemize}
\end{enumerate}
\end{keybox}

\begin{keybox}
\textbf{我国膳食结构：}
\begin{itemize}[leftmargin=*]
    \item 地域偏差：营养缺乏与某些营养素过剩并存。
    \begin{enumerate}[leftmargin=*]
        \item 少部分地区仍以贫困型膳食结构为主。
        \item 沿海城市等地区已经呈现富裕型膳食结构。
        \item 大部分城市和地区正朝富裕型膳食结构的方向转变。
        \item 营养缺乏与营养过剩双重负担并存。
    \end{enumerate}
    \item \textbf{营养失衡表现：}营养缺乏/不足（缺乏症）；营养过剩/过量（肥胖、慢性病）。
\end{itemize}
\end{keybox}

\subsection{膳食指南}

\begin{defbox}
\textbf{膳食指南（Dietary Guidelines）的定义：}
根据营养学原理，制定的一组以食物为基础的膳食行为建议，用于指导大众合理选择与搭配食物；将营养素优化与化学行为相结合，倡导合理膳食维持健康，预防膳食相关疾病。% [补充自往届笔记] DG是根据营养科学原理，结合当地食物供应和大众需求，针对实际存在的营养问题，由营养健康权威机构制定的食物选择和身体活动的指导性意见。

\textbf{膳食指南对群体的核心要求（八大特征）：}
\begin{enumerate}[leftmargin=*]
    \item 群体性：针对市民/大众，解决群体性问题，具有针对性
    \item 科学性：倡导平衡膳食，改善营养状况，预防膳食相关疾病
    \item 普及性：采用通俗易懂的大众语言，强调可操作性
    \item 有效性：核心膳食建议经科学验证有效
    \item 通俗性：使用通俗易懂的大众语言
    \item 权威性：由权威机构制定并指导
    \item 保障性：确保食物供应
    \item 预防性
\end{enumerate}
\end{defbox}

\begin{keybox}
\textbf{中国居民膳食指南（2022）——8条核心准则（针对2岁以上健康人群）：}

\begin{enumerate}[leftmargin=*]
    \item \textbf{准则一：食物多样，合理搭配。}谷类为主，平均每天摄入12种以上食物，每周25种以上。
    \item \textbf{准则二：吃动平衡，健康体重。}
    \item \textbf{准则三：多吃蔬果、奶类、全谷、大豆。}奶及奶制品300-500 ml（液态奶）。
    \item \textbf{准则四：适量吃鱼、禽、蛋、瘦肉。}平均每天120-200 g，每周最好吃鱼2次。
    \item \textbf{准则五：少盐少油，控糖限酒。}每天食盐不超过5 g，糖不超过50 g（最好25 g以下）。
    \item \textbf{准则六：规律进餐，足量饮水。}女性1500 ml，男性1700 ml（8杯），提倡喝白开水或茶水。
    \item \textbf{准则七：会烹会选，会看标签。}
    \item \textbf{准则八：公筷分餐，杜绝浪费。}
\end{enumerate}
\end{keybox}

\begin{keybox}
\textbf{中国居民平衡膳食宝塔（2022）：}

\begin{table}[H]
\centering
\begin{tabularx}{\textwidth}{clX}
\toprule
\textbf{层级} & \textbf{食物类别} & \textbf{每日推荐量} \\
\midrule
第一层 & 谷薯类 & 250-400 g/d（全谷物和杂豆50-150 g，薯类50-100 g） \\
\midrule
第二层 & 蔬菜类 & 300-500 g/d（深色蔬菜占1/2） \\
       & 水果类 & 200-350 g/d \\
\midrule
第三层 & 鱼、禽、肉、蛋 & 120-200 g/d（鱼40-75 g，畜禽肉40-75 g，蛋类1个约50 g） \\
\midrule
第四层 & 奶及奶制品 & 300-500 g/d \\
       & 大豆及坚果 & 25-35 g/d \\
\midrule
第五层 & 盐 & $<$ 5 g/d \\
       & 油 & 25-30 g/d \\
\midrule
\multicolumn{2}{l}{\textbf{其他推荐：}} & 饮水1500-1700 ml/d，活动6000步/d \\
\bottomrule
\end{tabularx}
\end{table}
\end{keybox}

\begin{keybox}
\textbf{特定人群膳食指南主要原则（熟悉）：}

\begin{enumerate}[leftmargin=*]
    \item \textbf{备孕妇女：}调整孕前体重至适宜水平；常吃含铁丰富食物，选用碘盐；孕前3个月开始补充叶酸；禁烟酒，保持健康生活方式。
    \item \textbf{孕期妇女：}补充叶酸，常吃含铁丰富食物，选用碘盐；孕吐严重者少量多餐，保证必要量碳水化合物摄入；孕中晚期适量增加奶、鱼、禽、蛋、瘦肉摄入；适量体力活动，维持孕期适宜体重。
    \item \textbf{哺乳期妇女：}增加富含优质蛋白质及维生素A的动物性食物和海产品摄入；适量增加奶类，多喝汤水；愉悦心情，充足睡眠，坚持哺乳。
    \item \textbf{6月龄内婴儿：}坚持纯母乳喂养；生后1小时内开奶，重视尽早吸吮；回应式喂养；补充维生素D，无需补钙。
    \item \textbf{7-24月龄婴幼儿：}继续母乳喂养，满6月龄添加辅食，从富铁泥糊状食物开始；引入多样化食物，重视动物性食物；保持食物原味，不加盐糖；回应式喂养。
    \item \textbf{学龄前儿童：}规律就餐，自主进食；每天饮奶，足量饮水；合理烹调，少调料少油炸；保证户外活动时间。
    \item \textbf{学龄儿童：}主动参与食物选择与制作；吃好早餐，合理选择零食；天天喝奶，足量饮水，不喝含糖饮料，禁止饮酒；多户外活动。
    \item \textbf{老年人：}食物品种丰富，动物性食物充足，常吃大豆制品；鼓励共同进餐；积极户外活动，延缓肌肉衰减；保持适宜体重，定期体检。
    \item \textbf{素食人群：}谷类为主，食物多样；增加大豆及其制品的摄入；常吃坚果、海藻和菌菇类；蔬菜水果充足；合理选择烹调油。
\end{enumerate}
\end{keybox}

% =====================================================================
\section{营养调查与评价}
% =====================================================================

% [补充自往届笔记]
营养调查（Nutrition Survey）是通过膳食调查、体格测量、营养水平生化检验和营养缺乏症临床检查等方法，全面了解个体或群体的营养状况的方法。营养状况综合评价包括四个层面：营养调查、体格测量、生化检验和临床检查。营养调查是对大规模人群膳食摄入量、临床检查、生化检测信息的综合解释。人体状况受各种营养素摄入量、需要量的综合影响，营养状况评价可以从整体上确定相关人群的健康状况。

\begin{defbox}
\textbf{营养调查与相关概念区分：}
\begin{itemize}[leftmargin=*]
    \item \textbf{营养调查（Nutrition Survey）：}指运用各种手段准确了解某一人群特定时间的各种营养指标水平，并判断其营养和健康状况。涉及一次性大规模横断面调查。
    \item \textbf{营养监测（Nutrition Surveillance）：}通过长期纵向动态监测目标人群营养状态。
    \item \textbf{营养筛查（Nutrition Screening）：}通过筛查方法快速识别群体中已患营养不良的个体。
\end{itemize}
\end{defbox}

\subsection{营养调查的目的与内容}

\begin{keybox}
\textbf{营养调查的目的：}
\begin{enumerate}[leftmargin=*]
    \item 了解不同地区、年龄、性别人群的能量和营养素摄入现状。
    \item 了解与能量和营养素摄入不足和过剩有关的营养问题的分布和严重程度。
    \item 探索营养相关疾病的病因和影响因素。
    \item 了解人群膳食结构变迁及变化趋势。
    \item 为某些营养相关国家性科学研究提供基础营养和健康数据。
    \item 为国家或地方制定食物营养政策、法规、标准、干预措施以及发展规划提供科学依据。
\end{enumerate}
\end{keybox}

% [补充自往届笔记]
为了解国民的营养和健康状况，许多国家定期有计划地开展国民营养调查工作。为研究不同经济发展时期人群的膳食营养和健康状况的变化提供基础资料。为食物生产加工、群众食物消费和制定干预措施提供依据。

\subsection{膳食调查}

\begin{defbox}
\textbf{膳食调查的定义：}
收集调查对象在一定时期内食物消费数据，通过食物消费信息了解其膳食摄入情况（能量和各种营养素的摄入量、食物来源、供能比例、加工方式和进餐制度等），并评价其能量和营养素摄入量满足机体需要的程度。

\textbf{膳食调查方法的分类维度：}
\begin{itemize}[leftmargin=*]
    \item 称重法、记账法、询问法、食物频率法、化学分析法。
    \item \textbf{方法选择依据（5W1H）：}
    \begin{itemize}
        \item Who——研究对象（个体/群体）
        \item What——研究关注点（食物/营养素）
        \item When——关注当前/通常膳食模式/长期（一天/一周/一年）
        \item Where——在家/在外食用
        \item Why——研究目的
    \end{itemize}
    \item 依据：研究目的、研究人群大小、精确度要求、经费及研究时间。调查方法各有其优缺点，确定后统一使用，不能混合使用。
\end{itemize}
\end{defbox}

\begin{defbox}
\textbf{膳食调查方法的基本要求：}
\begin{itemize}[leftmargin=*]
    \item 观察单位：个体、家庭、群体
    \item 资料收集方式：记录法、询问法（面访、电话、自动应答、电子化技术、双份法）
    \item 研究时限：既往的、目前的
    \item 食物量估计：重量、体积、尺寸估计，图片和模型或工具
    \item 食物转换为营养素方法：营养素数据库（食物成分表）、直接的化学分析
\end{itemize}
\end{defbox}

\subsection{膳食调查的主要方法}

\begin{tipbox}
\textbf{营养调查方法分类总览（了解）：}

\textbf{一、膳食调查方法：}
\begin{itemize}[leftmargin=*]
    \item \textbf{个人：}24小时回顾法、食物频率法、食物称重记录法。
    \item \textbf{家庭：}家庭食物量登记和记账法、家庭食物称重法。
    \item \textbf{群体：}食物供应调查、食物平衡表、市场供应调查。
\end{itemize}
\end{tipbox}

\begin{defbox}
\textbf{膳食调查五大方法详解：}

\textbf{（一）24小时膳食回顾法（24h History Recalls）}

又称询问法，是获得个体食物摄入量最常用的方法。要求每个调查对象回顾和描述在调查时刻以前24h内摄入的所有食物（包括饮料）的种类和数量。一般选择3天连续调查方法，每天入户回顾24小时摄入情况，连续进行3天（含2个工作日+1个周末）。

\begin{itemize}[leftmargin=*]
    \item \textbf{优点：}不改变被调查者的饮食习惯；食物量能估计比较准确；调查时间短；一年中多次回顾可提供食物和营养素通常摄入量的估计值；应答者负担轻，应答率较高；操作简单、费用低。
    \item \textbf{缺点：}食物份额量大小不易准确判断；对某些人群可能有难度（如80岁以上老年人），可能不适用。
    \item \textbf{应用：}适用于人群群体的膳食调查和评估，以及个体患者的调查（如病人、散居儿童、老人、咨询门诊等）。
\end{itemize}

\textbf{（二）食物频率问卷法（Food Frequency Questionnaires, FFQ）}

是估计被调查者在规定的一段时期内摄入某些食物频率的方法，用于评估长期膳食营养状况。调查期可从几天、1周、1个月到3个月或1年以上。食物频率问卷由调查对象的基本信息和食物频率两部分组成。

\begin{itemize}[leftmargin=*]
    \item \textbf{FFQ分类：}
    \begin{itemize}
        \item \textbf{定性FFQ：}通常只得到每种食物在特定时期内（如过去1个月）摄入的次数，不收集食物份额量大小资料。
        \item \textbf{定量FFQ：}通常要求受访者提供所摄入食物的平均数量，一般借助于测量辅助物。
        \item \textbf{半定量FFQ：}通常要求研究对象提供标准（或平均）食物份额量大小的种类，通常提供选项选择。
    \end{itemize}
    \item \textbf{优点：}能得到通常膳食摄入量；不需要训练有素的调查员；调查方法简单、费用低；习惯性膳食模式不受影响；应答者负担轻，应答率高；能获得通常膳食摄入的数据（某些食物和营养素的量），更有助于研究膳食与疾病之间的关系。
    \item \textbf{缺点：}需要对过去膳食模式的记忆，较长时期可能不准确；当前的膳食模式可能影响对过去膳食的回顾，导致偏倚；食物份额量标准大小的估计不准；食物摄入量的估计可能不准确；是一种复杂的方法，对非营养专业人士而言，自行管理的调查表错误率高；应答者负担取决于调查食物种类的数量、复杂性和频次；通常不能测量食物摄入量的每日变异，难以确定个别调查结果是否准确。
    \item \textbf{说明：}24h回顾法和FFQ联合使用可以较全面地反映人群膳食摄入的结果，弥补询问调查法的局限。在膳食与健康关系的流行病学研究中应用最广泛。
\end{itemize}

\textbf{（三）称重法（Weighed Food Records）}

指运用各种称量工具对某一饮食单位（个人或家庭）所消耗的各种食物的生重、熟重及餐后剩余量进行称重，了解食物消耗情况，计算出每人每日各种营养素的平均摄入量。调查时间一般为3-7天。

\begin{itemize}[leftmargin=*]
    \item \textbf{优点：}能测定食物份额大小或称重；能准确反映被调查对象的食物摄入量，可调查每日每餐膳食的变动情况，能准确计算每人每日各种营养素的摄入量，可以作为膳食调查的"金标准"用于验证其他方法的准确性。
    \item \textbf{缺点：}需要较大人力物力，对调查员要求高，对应答者文化水平有一定要求，额外的负担可能造成应答率下降，不能保证样本的代表性。
    \item \textbf{应用：}适用于家庭、个人及特殊工作人员的膳食调查，不太适用于大规模普查工作（如流行病学调查）。
\end{itemize}

\textbf{（四）记账法（Account Method）}

通过记录一定时期内某一饮食单位（机关、学校、部队、单位食堂等）的食物消耗总量和进餐人数，来计算每人每日各种营养素的平均摄入量。可调查较长时期膳食（一个月或更长）。

\begin{itemize}[leftmargin=*]
    \item \textbf{优点：}调查目的准确，每餐用餐人数统计较确实的情况下，使调查结果较准确；可调查较长时期的膳食。
    \item \textbf{缺点：}调查结果的准确性取决于账目记录和人数统计的准确度；不能得出个体间营养摄入差异，不能用于分析个体膳食摄入状况。
    \item \textbf{应用：}适用于有较详细账目的集体单位和家庭。
\end{itemize}

\textbf{（五）化学分析法（Chemical Analysis）}

收集受试者一日全部食物，在实验室测定营养素含量。样品收集方法：
\begin{itemize}[leftmargin=*]
    \item 双份法：是最准确的方法，制作两份完全相同的食物，一份供食用，另一份作为分析样品。
    \item 收集研究期间消耗的各种未加工食物，从当地市场上购买同样食物作为样品。
\end{itemize}
\begin{itemize}[leftmargin=*]
    \item \textbf{优点：}结果非常准确，能够准确地得出个体营养素准确摄入量。
    \item \textbf{缺点：}过程复杂、成本高，仅适于较小规模调查，如营养代谢试验或某些成分的功能评价研究需要时使用。
\end{itemize}

\textbf{食物量估计工具：}
\begin{enumerate}[leftmargin=*]
    \item 量具（碗、杯、勺等）
    \item 食物图片
    \item 食物模型
    \item 标准餐具
    \item 计算机图像
\end{enumerate}
\end{defbox}

\subsection{体格测量}

\begin{keybox}
\textbf{人体测量常用指标与评价：}

\textbf{1. Broca改良公式（理想体重/标准体重）：}
\[
\text{标准体重(kg)} = \text{身高(cm)} - 105
\]
\[
\text{标准体重} \pm 10\% \text{为正常范围}
\]
\begin{itemize}[leftmargin=*]
    \item 胖瘦 $<$ 10\% 正常
    \item $\pm$ 10\%-20\% 为瘦弱/超重
    \item $>$ 20\% 为严重瘦弱/肥胖
    \item +20\%-30\% 轻度肥胖
    \item +30\%-50\% 中度肥胖
    \item +50\% 重度肥胖
\end{itemize}

\textbf{2. 体质指数（BMI, Body Mass Index）：}
\[
\text{BMI} = \frac{\text{体重(kg)}}{[\text{身高(m)}]^{2}}
\]
我国标准：
\begin{itemize}[leftmargin=*]
    \item BMI $<$ 18.5：消瘦
    \item BMI 18.5-23.9：正常
    \item BMI 24.0-27.9：超重
    \item BMI $\ge$ 28：肥胖
\end{itemize}

\textbf{3. 儿童专用指标：}
\begin{itemize}[leftmargin=*]
    \item 年龄别体重（weight for age）：评价儿童营养状况
    \item 身高别体重（weight for height）：评价儿童营养状况
    \item 年龄别体重（weight for age）：婴幼儿专用
\end{itemize}
\end{keybox}

\begin{defbox}
\textbf{体格测量方法的原理与应用（熟悉）：}
\begin{itemize}[leftmargin=*]
    \item 体格测量用于对大体结构和皮下脂肪的测量，其大小和比例随营养程度的改变而变化。
    \item 对可衡量长期蛋白质和能量摄入不平衡的个体有较高敏感性。学龄儿童的体格测量结果可作为评价群体营养状况的指标。
    \item \textbf{常用指标：}身高（身长）、体重、上臂围、腰围、皮褶厚度。
    \item \textbf{专门指标：}坐高、头围、体成分、骨密度。
    \item 可借助设备和活动辅助测量。
\end{itemize}
\end{defbox}

\subsection{营养水平生化检验与临床检查}

\begin{tipbox}
\textbf{营养水平生化检验（实验室检查方法）（了解）：}
用于评价个体营养素缺乏的风险水平。
\begin{itemize}[leftmargin=*]
    \item 当机体组织中营养素逐渐耗竭，与营养素有关的酶水平和活性或代谢产物开始改变。
    \item 通过测定实验室/生化某些特定营养素的变化，血液学功能改变先于临床症状——出现营养素缺乏。
    \item 常用标本和测定指标：血清（血浆）——如血清铁蛋白、血浆视黄醇、血清25-(OH)D、血清钙等。
\end{itemize}

\textbf{常见营养水平鉴定生化指标（了解）：}
\begin{enumerate}[leftmargin=*]
    \item \textbf{蛋白质：}血浆总蛋白、白蛋白、前白蛋白、运铁蛋白、视黄醇结合蛋白——白蛋白半衰期较长（约20天），前白蛋白和运铁蛋白半衰期短（约2-8天），反映近期蛋白质营养状况更敏感。
    \item \textbf{铁：}血清铁蛋白（反映贮存铁）、血清铁、总铁结合力、运铁蛋白饱和度、血红蛋白。
    \item \textbf{维生素A：}血浆视黄醇水平。
    \item \textbf{维生素D：}血清25-(OH)D水平。
    \item \textbf{维生素B_{1}：}红细胞转酮醇酶活性系数（ETK-AC）。
    \item \textbf{维生素B_{2}：}红细胞谷胱甘肽还原酶活性系数（EGR-AC）。
    \item \textbf{维生素C：}血浆维生素C浓度、白细胞维生素C含量。
    \item \textbf{碘：}尿碘中位数。
    \item \textbf{钙：}血清钙（受甲状旁腺激素精密调控，不敏感）、骨密度测定。
    \item \textbf{锌：}血清/血浆锌浓度。
\end{enumerate}

\textbf{临床检查方法（熟悉）：}
\begin{enumerate}[leftmargin=*]
    \item \textbf{病史采集：}一般情况、目前健康状况、病史、既往病史（吸烟、膳食情况、饮酒和体力活动等、药物、日常体力活动）。
    \item \textbf{与营养素缺乏有关的检查：}各系统——皮肤、肌肉、骨骼、心血管、消化和神经系统的检查。检查部位：头、面部、眼、唇、舌、牙齿和齿龈、指甲。
\end{enumerate}

\textbf{常见营养缺乏症临床表现（了解）：}

\begin{table}[H]
\centering
\begin{tabularx}{\textwidth}{lX}
\toprule
\textbf{缺乏营养素} & \textbf{主要临床表现} \\
\midrule
蛋白质-能量（PEM） & 消瘦（marasmus）、水肿型（kwashiorkor）——生长迟缓、皮下脂肪减少、肌肉萎缩 \\
维生素A & 夜盲症、干眼病、毕脱氏斑（Bitot spots）、毛囊角化过度 \\
维生素D/钙 & 儿童佝偻病（颅骨软化、肋串珠、O/X型腿）、成人骨软化症、骨质疏松 \\
维生素B_{1} & 脚气病（干性：末梢神经炎；湿性：心衰水肿；婴儿型：紫绀心衰） \\
维生素B_{2} & 口角炎、唇炎、舌炎、阴囊皮炎、脂溢性皮炎 \\
烟酸 & 癞皮病（皮炎、腹泻、痴呆——"3D"：Dermatitis, Diarrhea, Dementia） \\
维生素C & 坏血病（牙龈出血、毛囊周围出血、瘀斑、伤口愈合不良） \\
叶酸/维生素B_{12} & 巨幼红细胞性贫血（舌炎、面色苍白、乏力） \\
铁 & 缺铁性贫血（面色苍白、疲乏、口角炎、反甲、异食癖） \\
碘 & 甲状腺肿、克汀病（先天性：严重智力障碍+矮小） \\
锌 & 生长迟缓、性发育延迟、食欲减退、异食癖、皮肤损害 \\
\bottomrule
\end{tabularx}
\end{table}
\end{tipbox}

\subsection{膳食营养评价}

\begin{keybox}
\textbf{膳食评价的依据：}
\begin{enumerate}[leftmargin=*]
    \item \textbf{DRIs：}评价膳食能量和营养素摄入量是否满足需要。
    \item \textbf{膳食指南：}规范膳食行为。
    \item \textbf{食物成分表：}将食物成分转换为营养素成分。
    \item \textbf{烹调加工方法合理性：}评价营养素在加工过程中的损失。
\end{enumerate}

\textbf{膳食营养素摄入量评价标准：}
\begin{itemize}[leftmargin=*]
    \item \textbf{摄入充足：}$>$ 90\% 标准
    \item \textbf{摄入不足：}$<$ 80\% 标准（长期摄入不足导致营养不良）
    \item \textbf{严重缺乏：}$<$ 60\% 标准（严重缺乏可致严重健康损害）
    \item 人群分布评估：关注高危人群比例和缺乏严重程度
\end{itemize}
\end{keybox}

\begin{keybox}
\textbf{能量及营养素评价要点（掌握）：}

\textbf{能量：}
\begin{itemize}[leftmargin=*]
    \item 摄入量与EER比较，总量 $\le$ EER $\pm$ 10\%（对成年人）。
    \item 产能营养素供能比：
    \begin{itemize}
        \item 蛋白质：10\%-15\%（儿童、老人12\%-15\%）
        \item 脂肪：20\%-30\%（儿童25\%-35\%）
        \item 碳水化合物：50\%-65\%
    \end{itemize}
    \item 蛋白质供能比（\%）= 蛋白质摄入量(g) × 4 ÷ 能量总摄入量(kcal) × 100
    \item 碳水化合物供能比（\%）= 碳水化合物摄入量(g) × 4 ÷ 能量总摄入量(kcal) × 100
    \item 脂肪供能比（\%）= 脂肪摄入量(g) × 9 ÷ 能量总摄入量(kcal) × 100
\end{itemize}

\textbf{三餐分配（能量比）：}
\begin{itemize}[leftmargin=*]
    \item 早餐:中餐:晚餐 = 3:4:3
    \item 早餐供能比 = 早餐能量摄入量 ÷ 能量总摄入量 × 100
    \item 中餐供能比 = 中餐能量摄入量 ÷ 能量总摄入量 × 100
    \item 晚餐供能比 = 晚餐能量摄入量 ÷ 能量总摄入量 × 100
\end{itemize}

\textbf{蛋白质：}
\begin{itemize}[leftmargin=*]
    \item 总量 $\ge$ 90\% RNI
    \item 优质蛋白质占 1/3-1/2（动物性蛋白+大豆蛋白）
    \item 来源：动物性食物（肉类、奶类、蛋类等）和豆类及其制品
\end{itemize}

\textbf{脂肪：}
\begin{itemize}[leftmargin=*]
    \item 动物性食物和植物性食物提供的脂肪量
    \item 脂肪酸来源：动物性食物提供的脂肪占总脂肪的百分比（E\%）
    \item 脂肪酸构成：SFA:MUFA:PUFA = 1:1:1
    \item 胆固醇和反式脂肪
\end{itemize}

\textbf{碳水化合物：}
\begin{itemize}[leftmargin=*]
    \item 占总能量百分比，单糖、双糖来源
    \item 复杂碳水化合物来源
    \item 膳食纤维来源
    \item 单糖/双糖与复合糖比例，GI和GL
\end{itemize}

\textbf{矿物质：}
\begin{itemize}[leftmargin=*]
    \item 铁总量 $\ge$ 90\% RNI，来源于动物性食物应占1/2
    \item 钙总量 $\ge$ 90\% RNI，来源于乳制品占2/3
\end{itemize}

\textbf{维生素：}
\begin{itemize}[leftmargin=*]
    \item 维生素A——食物来源与营养素评价
    \item 维生素D——晒太阳和食物来源与营养素评价
    \item 维生素B_{1}、B_{2}、叶酸、维生素C
\end{itemize}
\end{keybox}

\begin{tipbox}
\textbf{一般意义的营养缺乏：}
\begin{itemize}[leftmargin=*]
    \item 儿童生长迟缓：身高低于同年龄标准
    \item 儿童消瘦：体重低于同身高儿童
    \item 儿童超重：体重高于同身高儿童
    \item 成人超重：身体脂肪增多，BMI > 25
    \item 微量营养素缺乏症：如维生素A、锌、碘、叶酸处于较低水平
    \item 成人肥胖症：身体脂肪增多，BMI $\ge$ 30
    \item 非传染性疾病：糖尿病、心脏病、某些癌症
\end{itemize}
\end{tipbox}

% [补充自往届笔记]
平衡膳食五要求（食谱编制的核心原则）：
\begin{enumerate}[leftmargin=*]
    \item 提供种类齐全、数量充足、比例适宜的营养素
    \item 提供数量适当的能量和各种营养素
    \item 食物应新鲜卫生
    \item 正确烹调加工
    \item 良好的进餐制度和环境
\end{enumerate}

% =====================================================================
\section{营养监测与食谱编制}
% =====================================================================

\subsection{营养监测}

\begin{defbox}
\textbf{营养监测（Nutrition Surveillance）：}通过纵向持续动态监测目标人群营养状况，掌握营养状况变化趋势，为制定政策和干预措施提供依据。

\textbf{营养监测与营养调查的区别：}
\begin{itemize}[leftmargin=*]
    \item 营养调查：横断面调查，在某时间点评估目标人群营养状况。
    \item 营养监测：纵向动态监测，持续追踪营养状况变化趋势。
\end{itemize}

\textbf{常用监测指标：}
\begin{enumerate}[leftmargin=*]
    \item \textbf{健康指标：}身高、体重、贫血患病率、低出生体重率、5岁以下儿童死亡率等。
    \item \textbf{食物消费指标：}人均食物消费量、膳食能量和营养素摄入量。
    \item \textbf{社会经济指标：}恩格尔系数（Engel's Coefficient）——食物支出占家庭或个人总消费支出的百分比。
    \begin{itemize}[leftmargin=*]
        \item 恩格尔系数 $>$ 60\%：贫困
        \item 50\%-59\%：温饱
        \item 40\%-49\%：小康
        \item 30\%-39\%：富裕
        \item $<$ 30\%：最富裕
    \end{itemize}
\end{enumerate}
\end{defbox}

\subsection{食谱编制}

\begin{defbox}
\textbf{食谱编制的基本原则、依据和方法（熟悉）：}

\textbf{基本原则：}
\begin{enumerate}[leftmargin=*]
    \item \textbf{参照DRIs，保证营养均衡：}满足用膳者能量和各种营养素需要。
    \item \textbf{食物多样，合理搭配：}充分发挥营养素互补作用。
    \item \textbf{考虑用膳者特点和需求：}年龄、性别、生理状况、劳动强度、疾病状况。
    \item \textbf{注意食品安全和加工合理：}科学烹调，减少营养素损失。
    \item \textbf{考虑经济条件和食物供应：}因地制宜，合理选择食物。
\end{enumerate}

\textbf{编制依据：}
\begin{itemize}[leftmargin=*]
    \item DRIs（作为营养素摄入目标）
    \item 膳食指南和膳食宝塔（作为食物种类和数量参考）
    \item 食物成分表（计算营养素实际含量）
    \item 烹调加工方法的影响
\end{itemize}

\textbf{基本方法：}
\begin{enumerate}[leftmargin=*]
    \item 确定用膳者人数、年龄、性别、生理状况和劳动强度。
    \item 计算能量和营养素需要量（参照DRIs）。
    \item 确定主食（谷类）、副食（蔬菜、肉、蛋、奶、豆等）的种类和数量（参照膳食宝塔）。
    \item 分配一日三餐（早餐30\%、中餐40\%、晚餐30\%能量）。
    \item 编制一日食谱，计算营养素含量，与DRIs比较进行调整。
    \item 编制周食谱，避免单调重复，力求食物多样。
\end{enumerate}

\textbf{食谱编制的工作程序（"日需食谱周排"）：}% [补充自往届笔记]
建议在编制好一日食谱基础上进一步编制周食谱，避免单调重复，力求食物多样。
\end{defbox}

% =====================================================================
\section{公共营养改善}
% =====================================================================

\begin{defbox}
\textbf{公共营养改善（Nutrition Improvements）的定义：}
指采用膳食营养干预措施，改善营养不良人群的营养水平，以促使其恢复营养健康状况的活动。

\textbf{公共营养改善工作：}指为改善居民营养状况而开展的预防和控制营养缺乏、营养过剩和营养相关疾病等工作。

\textbf{公共营养改善常用方法：}
\begin{enumerate}[leftmargin=*]
    \item 营养教育与咨询
    \item 营养配餐与食谱编制
    \item 食品营养标签
    \item 食品营养强化（食品原料、新资源食品等）
\end{enumerate}

\textbf{营养干预（Nutrition Intervention）：}健康促进活动，目的是改变目标人群营养状况，包括选择干预策略、研究目标人群、研究消费者行为和实施评价研究。
\end{defbox}

\subsection{营养教育}

\begin{defbox}
\textbf{营养教育（Nutrition Education）的定义：}
通过改变人们的饮食行为达到改善营养目的的一种有计划的活动。意义：多途径、低成本、覆盖面广；有计划、有组织、有系统。

\textbf{营养教育的作用：}
\begin{enumerate}[leftmargin=*]
    \item 提供膳食行为改变所需的知识、技能和服务。
    \item 帮助个体和群体树立食品与营养的健康意识。
    \item 培养良好膳食行为与生活方式，使人们在面临营养与食品安全问题时有能力做出有益于健康的选择。
\end{enumerate}

\textbf{营养教育的主要内容：}% [补充自往届笔记]
\begin{enumerate}[leftmargin=*]
    \item 营养基础知识
    \item 健康生活方式
    \item 《中国居民膳食指南》和《中国居民平衡膳食宝塔》
    \item 我国人群膳食营养与健康状况及变化趋势
    \item 膳食营养相关的慢性疾病的膳食预防措施
    \item 营养相关的法律、法规与政策
\end{enumerate}

\textbf{营养教育的主要实施方式（了解）：}% [补充自往届笔记]
\begin{enumerate}[leftmargin=*]
    \item 各类人员知识培训：有计划地对餐饮业、农业、商业、轻工、医疗卫生、疾病控制等部门有关人员进行营养知识培训。
    \item 营养科学纳入中小学教育：将营养知识纳入中小学的教育内容和教学计划，安排一定课时的营养知识教育，使学生懂得平衡膳食的原则，培养学生良好的饮食习惯，提高自我保护意识。
    \item 社区教育网络体系：提高社区基层卫生人员和居民的膳食营养意识，培养良好的饮食行为，改善营养状况。
    \item 全社会宣传教育：利用多种宣传媒介，广泛开展群众性营养宣传活动，倡导健康膳食模式和健康生活方式，纠正不良饮食习惯等。
\end{enumerate}
\end{defbox}

\subsection{营养信息传播与营养行为干预}

\begin{tipbox}
\textbf{营养信息传播（了解）：}
\begin{itemize}[leftmargin=*]
    \item 准确信息是改变行为的基石。
    \item 营养信息传播是健康教育的一个重要组成部分，是营养教育和营养改善行动的重要手段和策略。
\end{itemize}

\textbf{营养行为干预（熟悉）：}
\begin{itemize}[leftmargin=*]
    \item 行为改变是实现营养教育计划目标的重要手段。
    \item 通过行为指导、技能训练和辅导咨询，使受教育者实现特定的饮食行为的改变。
    \item 主要方法：模拟、示范、实际操作、个别指导、小组讨论等均属干预手段。此外，还包括一些行为矫正技术。
\end{itemize}
\textbf{营养干预应用（了解）：}纳入初级卫生保健服务体系，利用多种宣传媒介开展科普宣传，将营养知识纳入中小学教育。
\end{tipbox}

\subsection{食品营养强化}

\begin{defbox}
\textbf{食品营养强化 / 食品强化（Food Fortification）：}
根据不同人群营养需要，向食品中添加一种或多种营养素（或某些天然食物成分），以提高食品营养价值的过程称为食品营养强化，或简称食品强化。

\textbf{食品强化剂：}指为增强营养成分而加入食品中的天然或者人工合成的属于天然营养素范围的食品添加剂。主要包括必需氨基酸、维生素、矿物质等。

\textbf{食品营养强化的意义：}
\begin{enumerate}[leftmargin=*]
    \item 弥补天然食物的营养缺陷，补充其不足（如谷类加赖氨酸）。
    \item 补充加工过程中人为损失的营养素（B族维生素、精白米面等）。
    \item 满足特定人群的营养需要（配方奶粉、要素膳等）。
    \item 预防营养不良（食盐加碘）。
    \item 改善食物感官性状和营养价值：婴儿配方食品。
\end{enumerate}
\end{defbox}

% [补充自往届笔记]
我国食物营养相关政策及规划（了解）：
\begin{itemize}[leftmargin=*]
    \item 食品安全法
    \item 营养改善工作管理办法
    \item 全国营养调查工作规范
    \item 食品营养标签管理规范
    \item 《2025-2030年中国食物与营养发展纲要》
    \item 2017年7月《国民营养计划（2017-2030）》
\end{itemize}

% =====================================================================
\section{食品营养标签、NRV与功能食品}
% =====================================================================

\subsection{食品营养标签}

\begin{defbox}
\textbf{食品营养标签（Food Nutrition Label）（熟悉）：}
食品营养标签是向消费者提供食品营养信息和特性的说明，包括营养成分表、营养声称和营养成分功能声称。食品营养标签是预包装食品标签的一部分，管理遵循相关法规（如《GB 28050-2011 食品安全国家标准 预包装食品营养标签通则》）。% [补充自往届笔记] 强制标示内容包括：能量、蛋白质、脂肪、碳水化合物和钠（"1+4"核心营养素）。

\textbf{营养成分表基本内容：}
\begin{itemize}[leftmargin=*]
    \item 能量和核心营养素（1+4）：能量 + 蛋白质、脂肪、碳水化合物、钠（强制标示）。
    \item 营养声称（如"高钙"、"低脂"、"富含膳食纤维"等）。
    \item 营养成分功能声称（如"钙有助于骨骼和牙齿的健康"）。
\end{itemize}
\end{defbox}

\subsection{NRV（营养素参考值）}

\begin{keybox}
\textbf{营养素参考值（NRV——Nutrient Reference Values）（掌握）：}
\begin{itemize}[leftmargin=*]
    \item \textbf{定义：}专用于食品营养标签上比较食品营养素含量的参考值。
    \item \textbf{用途：}NRV是食品营养标签上的"标尺"，用NRV\%表示食品中某营养素含量占NRV的百分比，帮助消费者快速判断该食品在每日膳食中的贡献比例，便于比较不同食品营养价值，做出合理食物选择。
    \item \textbf{示例：}某食品每份含钙200 mg，钙NRV为800 mg，则钙NRV\% = 200/800 $\times$ 100\% = 25\%。
    \item \textbf{NRV与DRIs的区别：}DRIs是一组针对不同年龄、性别、生理状况人群的参考值体系；NRV是用于营养标签的单一参考值，通常采用成人RNI或AI。NRV用于食品标签，DRIs用于膳食评价和规划。
\end{itemize}
\end{keybox}

\subsection{功能食品}

\begin{defbox}
\textbf{功能食品（Functional Food）的概念（熟悉）、分类（了解）、功能评价（了解）及管理（熟悉）：}

\textbf{概念：}功能食品是指具有特定营养保健功能的食品，即适宜于特定人群食用，具有调节机体功能，不以治疗疾病为目的的食品。在我国与"保健食品"概念基本对应。

\textbf{分类（了解）：}
\begin{enumerate}[leftmargin=*]
    \item 营养素补充剂——补充维生素、矿物质等。
    \item 具有特定保健功能的食品——按功能声称分类（如增强免疫力、辅助降血脂、辅助降血糖、抗氧化、辅助改善记忆、缓解视疲劳、促进排铅、清咽、辅助降血压、改善睡眠、促进泌乳、缓解体力疲劳、提高缺氧耐受力、对辐射危害有辅助保护功能、减肥、改善生长发育、增加骨密度、改善营养性贫血、对化学性肝损伤有辅助保护功能、祛痤疮、祛黄褐斑、改善皮肤水分、改善皮肤油分、调节肠道菌群、促进消化、通便、对胃黏膜有辅助保护功能等27类保健功能）。
\end{enumerate}

\textbf{功能评价（了解）：}功能食品需进行功能学实验验证和安全性评价。功能学实验包括动物功能实验和/或人体试食试验，证明其具有声称的保健功能。安全性毒理学评价根据原料不同进行相应的毒理学试验。

\textbf{管理（熟悉）：}我国对保健食品实行注册与备案双轨制管理。需经国家市场监督管理总局审批，获得"蓝帽子"标志（保健食品标志）后方可上市销售。产品标签和说明书必须标注保健功能、适宜人群、不适宜人群、食用方法和用量等，须标明"本品不能替代药物"。
\end{defbox}

% =====================================================================
\reviewcontent{
\section{复习题}
% =====================================================================

\begin{enumerate}[leftmargin=*, itemsep=8pt]
    \item \textbf{【名词解释】}DRIs；EAR；RNI；AI；UL；AMDR；PI-NCD；SPL；BMI；NRV；恩格尔系数；EER；公共营养；营养调查；营养监测；营养教育
    \item \textbf{【简答题】}简述公共营养的定义、地位与工作内容。
    \item \textbf{【简答题】}简述DRIs七项指标体系的定义、制定目标和相互关系。
    \item \textbf{【简答题】}比较AI与RNI的异同以及UL的制定注意事项。
    \item \textbf{【简答题】}简述DRIs的制定修订程序及其用途层次。
    \item \textbf{【简答题】}列出《中国居民膳食指南(2022)》8条核心准则及膳食指南对群体的核心要求。
    \item \textbf{【简答题】}简述膳食结构的四种类型、各自特点及其与健康的关系，以及我国膳食结构现状。
    \item \textbf{【简答题】}简述营养调查的目的、内容和营养状况综合评价的四个方面。
    \item \textbf{【简答题】}比较膳食调查五种方法（称重法、24h回顾法、FFQ、化学分析法、记账法）的特点、优缺点及应用。
    \item \textbf{【简答题】}简述膳食调查方法选择依据（5W1H原则）和基本要求。
    \item \textbf{【简答题】}简述膳食营养评价的依据、摄入量评价标准和主要营养素评价要点。
    \item \textbf{【简答题】}简述BMI分级标准、Broca改良公式及其应用。
    \item \textbf{【简答题】}简述营养监测与营养调查的概念区别，以及恩格尔系数的含义和分级标准。
    \item \textbf{【简答题】}简述食谱编制的基本原则、依据和方法。
    \item \textbf{【简答题】}简述公共营养改善常用方法、营养教育的概念和主要内容、营养干预的主要方法。
    \item \textbf{【简答题】}简述食品营养强化的概念、食品强化剂的定义及营养强化的意义。
    \item \textbf{【简答题】}简述食品营养标签的基本内容和NRV的概念、用途及其与DRIs的区别。
    \item \textbf{【简答题】}简述功能食品（保健食品）的概念、分类、功能评价及管理要求。
\end{enumerate}

\subsection*{参考答案要点}

\begin{enumerate}[leftmargin=*, itemsep=5pt]
    \item 见正文相关定义。
    \begin{itemize}
        \item \textbf{公共营养：}从医学状况+人群社会视角的群体概念。通过膳食提供给机体种类齐全、数量充足、比例适宜的能量和各种营养素，并与机体需要相平衡。
        \item \textbf{DRIs：}一组每日平均膳食营养素摄入量的参考值，用于评价膳食营养素供给量能否满足人体需要及是否存在过量摄入风险。
        \item \textbf{EAR：}满足群体中50\%个体需要量，是制订RNI的基础。
        \item \textbf{RNI：}满足97\%-98\%个体需要量，RNI = EAR + 2SD 或 RNI = 1.2 × EAR。
        \item \textbf{AI：}资料不足无法计算EAR时，通过观察或实验获得。AI主要用于个体营养素摄入目标，不能用于人群营养评价；准确度不如RNI。
        \item \textbf{UL：}平均每日摄入营养素的最高限量。摄入量 < UL对健康人群无副作用，> UL不良反应风险增加。
        \item \textbf{AMDR：}宏量营养素可接受范围——脂肪20\%-30\%，碳水化合物50\%-65\%，蛋白质10\%-15\%。
        \item \textbf{PI-NCD：}预防非传染性慢性病的建议摄入量，以NCD一级预防为目标。PI可能高于RNI/AI（如维生素C、钾）或低于AI（如钠）。
        \item \textbf{SPL：}特定建议值，专用于预防NCD，适用于其他食物成分（如植物化学物）的预防NCD建议量。
        \item \textbf{EER：}能长期保持良好健康状态、维持体重和机体构成以及体力活动水平的个体及人群，达到能量平衡时所需要的膳食能量摄入量。
        \item \textbf{BMI：}体质指数，体重(kg) / [身高(m)]^{2}。标准：<18.5消瘦，18.5-23.9正常，24.0-27.9超重，≥28肥胖。
        \item \textbf{NRV：}营养素参考值，专用于食品营养标签上比较食品营养素含量的参考值。
        \item \textbf{恩格尔系数：}食物支出占总消费支出百分比。>60\%贫困，50\%-59\%温饱，40\%-49\%小康，30\%-39\%富裕，<30\%最富裕。
        \item \textbf{营养调查：}运用各种手段准确了解某一人群特定时间的各种营养指标水平，判断其营养和健康状况。
        \item \textbf{营养监测：}通过纵向持续动态监测目标人群营养状况，掌握变化趋势。
        \item \textbf{营养教育：}通过改变人们的饮食行为达到改善营养目的的一种有计划的活动。
    \end{itemize}

    \item 公共营养定义、地位与工作内容：
    \begin{itemize}
        \item \textbf{定义：}从医学状况+人群社会视角的群体概念。合理营养（平衡膳食）：提供种类齐全、数量充足、比例适宜的能量和各种营养素。
        \item \textbf{地位：}第一，促进全人群营养与健康；第二，增强国民体质和综合国力；第三，事关国家发展战略问题。
        \item \textbf{工作内容：}营养标准制定（DRIs）、膳食指南制定、人员培训、营养调查与评价、营养监测、营养干预、食物营养规划与立法、其他营养公众指导。
        \item \textbf{营养不良：}营养缺乏（nutrition deficiency）和营养过剩（nutrition excess）。
    \end{itemize}

    \item DRIs七项指标体系：
    \begin{itemize}
        \item \textbf{EAR（平均需要量）：}满足50\%个体需要，是制订RNI的基础。目标：预防缺乏。
        \item \textbf{RNI（推荐摄入量）：}满足97\%-98\%个体需要，RNI = EAR + 2SD 或 1.2 × EAR。目标：预防缺乏。
        \item \textbf{AI（适宜摄入量）：}EAR无法计算时使用，通过观察或实验获得，准确度不如RNI。目标：预防缺乏。
        \item \textbf{UL（可耐受最高摄入量）：}安全上限。目标：防止过量。
        \item \textbf{AMDR（宏量营养素可接受范围）：}脂肪20\%-30\%，碳水化合物50\%-65\%，蛋白质10\%-15\%。目标：预防慢性病。
        \item \textbf{PI-NCD（预防慢病建议摄入量）：}以NCD一级预防为目标，与RNI/AI比较：PI>RNI/AI时以PI为准（维生素C、钾），PI<AI时以AI为准（钠）。目标：预防慢性病。
        \item \textbf{SPL（特定建议值）：}针对其他膳食成分（如植物化学物）的预防NCD建议量。目标：预防慢性病。
    \end{itemize}

    \item AI与RNI异同：
    \begin{itemize}
        \item \textbf{相同点：}都可以作为群体中个体营养素摄入量的目标值，满足群体中几乎所有个体的需要。
        \item \textbf{不同点：}AI准确度不如RNI，主要用于个体营养素摄入目标；群体摄入量≥AI时，摄入不足者危险性很小。
    \end{itemize}
    UL制定注意事项：随强化食品和使用补充剂的发展需制定ULs；毒副作用与摄入总量有关的，按食物+饮水+补充剂总量制定；毒副作用仅与强化食物和补充剂有关的，按这些来源制定；未定UL不意味着过多摄入无潜在危害。

    \item DRIs制定修订程序：人群调查监测→制定DRIs→膳食指导→食物指导→食品营养标签→人群营养改善反馈→营养科学研究→修订DRIs。每4-5年修订一次。用途层次：个体（膳食规划、食谱、配方研发）、群体（膳食评价、营养咨询）、安全作用（标准/依据）。

    \item 膳食指南8条核心准则（2022）：
    \begin{enumerate}[leftmargin=*]
        \item 食物多样，合理搭配（每天>12种，每周>25种）
        \item 吃动平衡，健康体重
        \item 多吃蔬果、奶类、全谷、大豆（奶300-500 ml/d）
        \item 适量吃鱼、禽、蛋、瘦肉（120-200 g/d，每周鱼2次）
        \item 少盐少油，控糖限酒（盐<5 g/d，糖<50 g/d最好<25 g/d）
        \item 规律进餐，足量饮水（女1500 ml，男1700 ml）
        \item 会烹会选，会看标签
        \item 公筷分餐，杜绝浪费
    \end{enumerate}
    膳食宝塔五层：第一层谷薯类250-400 g/d；第二层蔬菜300-500 g/d、水果200-350 g/d；第三层鱼禽肉蛋120-200 g/d；第四层奶及奶制品300-500 g/d、大豆及坚果25-35 g/d；第五层盐<5 g/d、油25-30 g/d。饮水1500-1700 ml/d，活动6000步/d。
    膳食指南对群体的核心要求：群体性、科学性、普及性、有效性、通俗性、权威性、保障性、预防性。

    \item 四种膳食模式：
    \begin{itemize}
        \item \textbf{动物性食物为主型（西方模式）：}"三高一低"（高能量、高脂肪、高蛋白、低膳食纤维），与肥胖、心脑血管疾病、结肠癌直肠癌等慢性病风险增加相关。
        \item \textbf{植物性食物为主型（东方模式）：}"一高三低"（低蛋白质、低脂肪、植物性食物为主），感染性疾病+慢性病并存。
        \item \textbf{动植物平衡型（日本模式）：}动植物消费均衡，水产品消费量大，三大宏量营养素供能比合理。
        \item \textbf{地中海模式：}橄榄油为主，不饱和脂肪高，大量蔬菜水果全谷物豆类，适量鱼类和红酒。
    \end{itemize}
    我国膳食结构：营养缺乏与过剩并存，地域偏差明显，正朝富裕型膳食结构转变。

    \item 营养调查目的6项（见正文）；内容4项：膳食调查、体格测量、生化检验、临床检查。

    \item 五种方法比较：
    \begin{itemize}
        \item \textbf{称重法：}准确可靠，3-7天，可作为"金标准"，费时费力，不适于大规模调查。
        \item \textbf{24h回顾法：}最常用，不改变饮食习惯，简便快速，连续3天（2工作日+1周末），依赖记忆，对某些人群不适用。
        \item \textbf{FFQ：}可调查长期（几天到1年以上）膳食模式，方法简单费用低，精确度较低，分定性/定量/半定量三种。与24h回顾法联用更全面。
        \item \textbf{化学分析法：}最准确的方法（双份法），过程复杂成本高，仅适用于小规模。
        \item \textbf{记账法：}可调查较长时期（一个月或更长），费用低，不能分析个体摄入状况。
    \end{itemize}

    \item 膳食调查方法选择依据（5W1H）：Who（个体/群体）、What（食物/营养素）、When（当前/通常模式/长期）、Where（在家/在外）、Why（研究目的），以及研究人群大小、精确度要求、经费及时间。基本要求：观察单位、资料收集方式、研究时限、食物量估计、食物-营养素转换方法。

    \item 膳食营养评价依据：DRIs（摄入充足>90\%标准、摄入不足<80\%标准、严重缺乏<60\%标准）、膳食指南、食物成分表、烹调加工方法合理性。评价要点：能量EER±10\%；蛋白质总量≥90\%RNI优质蛋白1/3-1/2供能比10\%-15\%；脂肪供能比20\%-30\%、SFA:MUFA:PUFA=1:1:1；碳水化合物50\%-65\%；三餐分配3:4:3；铁≥90\%RNI动物来源1/2；钙≥90\%RNI乳制品来源2/3。

    \item BMI分级标准：<18.5消瘦，18.5-23.9正常，24.0-27.9超重，≥28肥胖。Broca改良公式：标准体重(kg)=身高(cm)-105。±10\%正常，±10\%-20\%瘦弱/超重，>20\%严重瘦弱/肥胖，+20\%-30\%轻度肥胖，+30\%-50\%中度肥胖，+50\%重度肥胖。

    \item 营养调查为横断面调查，营养监测为纵向动态监测。恩格尔系数：食物支出/总消费支出×100\%。>60\%贫困，50\%-59\%温饱，40\%-49\%小康，30\%-39\%富裕，<30\%最富裕。

    \item 食谱编制基本原则：满足DRIs、食物多样、考虑用膳者特点和需求、注意食品安全科学加工、考虑经济条件和食物供应。依据：DRIs、膳食指南和宝塔、食物成分表、烹调加工方法。方法：确定人群特点→计算营养素需要量→确定主副食种类数量→分配三餐（3:4:3）→编制食谱计算营养素→与DRIs比较调整→编制周食谱。

    \item 公共营养改善常用方法：营养教育与咨询、营养配餐与食谱编制、食品营养标签、食品营养强化。营养教育：通过改变饮食行为改善营养的有计划活动（多途径、低成本、覆盖面广）。主要内容：营养基础知识、健康生活方式、膳食指南和宝塔、人群营养与健康状况、慢性病膳食预防、营养法律法规。营养干预方法：模拟、示范、实际操作、个别指导、小组讨论、行为矫正技术等。

    \item 食品营养强化：向食品中添加营养素以提高食品营养价值的过程。食品强化剂：天然或人工合成的属于天然营养素范围的食品添加剂（必需氨基酸、维生素、矿物质等）。意义5项：弥补天然食物营养缺陷；补充加工损失；满足特定人群需要；预防营养不良；改善食物感官性状和营养价值。

    \item NRV：专用于食品营养标签的参考值。NRV\%表示营养素含量占NRV的百分比，帮助消费者判断该食品在每日膳食中的贡献比例。NRV与DRIs区别：DRIs是针对不同人群的参考值体系，NRV是用于标签的单一参考值（取成人RNI/AI）。营养成分表强制标示"1+4"：能量+蛋白质、脂肪、碳水化合物、钠。

    \item 功能食品（保健食品）：具有特定保健功能、适宜特定人群食用、不以治疗疾病为目的的食品。分类：营养素补充剂和27类特定保健功能食品。管理：注册与备案双轨制，国家市场监督管理总局审批，获"蓝帽子"标志，标签须标明保健功能、适宜人群、不适宜人群、食用方法，须标明"本品不能替代药物"。功能评价：功能学实验（动物实验和/或人体试食试验）和安全性毒理学评价。
\end{enumerate}

}
\newpage

% =====================================================================
%  CHAPTER 6: 营养与疾病 — 参考：0613笔记 (PRIMARY SOURCE)
% =====================================================================
\chapter{营养与疾病}
\chapterintro{
\keyword{掌握}超重肥胖饮食防治策略原则和方法；掌握心血管疾病、糖尿病、癌症和痛风的营养防治原则。\keyword{熟悉}体重控制意义及方法；熟悉体力活动指南；熟悉临床营养筛查与营养评价基本方法；熟悉病人膳食分类；熟悉肠内营养和肠外营养概念、方法、适应证和并发症；熟悉FSMP概念。\keyword{了解}膳食营养与慢性病的关系；了解临床营养治疗目的和主要方法；了解肠内营养制剂、肠外营养制剂及FSMP分类。
}

% =====================================================================
%  膳食营养与慢性病的关系
% =====================================================================
\section{膳食营养与慢性病的关系}

\begin{tipbox}
\textbf{膳食营养与慢性非传染性疾病（NCD）的关系（了解）：}

慢性非传染性疾病（主要包括心血管疾病、糖尿病、癌症和慢性呼吸系统疾病）是全球最主要的死亡原因。不合理膳食是NCD最重要的可改变危险因素之一。

\textbf{膳食因素与主要慢性病的关联：}
\begin{enumerate}[leftmargin=*]
    \item \textbf{超重/肥胖：}高能量密度食物、含糖饮料、大份量和久坐 → 能量正平衡 → 肥胖 → 是心血管疾病、糖尿病、多种癌症的共同危险因素。
    \item \textbf{心血管疾病：}高钠（高血压）、高饱和脂肪和反式脂肪（动脉粥样硬化）、低蔬果和全谷物（膳食纤维和抗氧化物质不足）。
    \item \textbf{2型糖尿病：}高GI/GL膳食、含糖饮料、红肉和加工肉类摄入过多，全谷物和膳食纤维摄入不足。
    \item \textbf{癌症：}加工肉类（IARC Ⅰ类致癌物）、红肉（ⅡA类）、酒精（多种癌症）、黄曲霉毒素污染食物；保护因素：蔬果、全谷物和膳食纤维。
\end{enumerate}

\textbf{公共健康膳食策略的核心理念：}
合理膳食是NCD一级预防的基石，膳食因素可干预、可改变，通过人群层面的营养政策和个体层面的营养教育可实现NCD的有效预防。膳食营养与慢性病的关系为制定膳食指南和营养政策提供了重要的科学依据。
\end{tipbox}

% =====================================================================
%  营养与肥胖
% =====================================================================
\section{营养与肥胖}

\begin{defbox}
\textbf{肥胖的定义与判定标准：}
\begin{itemize}[leftmargin=*]
    \item \textbf{肥胖：}体内脂肪过多积累，能量摄入长期大于能量消耗所致。
    \item \textbf{体质指数（BMI）：}\imp{BMI = 体重(kg) / [身高(m)]$^2$}。我国标准：BMI $<$ 18.5 消瘦，18.5$\sim$23.9 正常，24.0$\sim$27.9 超重，\imp{BMI $\geq$ 28 肥胖}。
    \item \textbf{中心性肥胖：}腰围\imp{男 $\geq$ 85cm，女 $\geq$ 80cm}，或腰臀比男 $>$ 0.9、女 $>$ 0.85。
\end{itemize}
\end{defbox}

\begin{keybox}
\textbf{肥胖的膳食因素与预防控制：}

\textbf{膳食因素：}
\begin{itemize}[leftmargin=*]
    \item 高能量密度食物（油炸食品、甜点、快餐）
    \item 含糖饮料（添加糖占总能量 $>$ 10\%）
    \item 进食量大、进餐速度快、不吃早餐、夜间进食
    \item 膳食纤维摄入不足，蔬果和全谷物摄入不足
\end{itemize}

\textbf{能量平衡原理：}肥胖的本质是\imp{能量摄入长期大于能量消耗}。能量消耗包括基础代谢（约60\%$\sim$70\%）、身体活动（约15\%$\sim$30\%）和食物热效应（约10\%）。

\textbf{预防控制措施：}
\begin{enumerate}[leftmargin=*]
    \item \textbf{控制总能量摄入：}每天减少500$\sim$1000kcal能量摄入，每月减重1$\sim$2kg为宜。每日能量摄入不宜低于1000$\sim$1200kcal（女性）/ 1200$\sim$1600kcal（男性）。
    \item \textbf{调整宏量营养素供能比：}蛋白质\imp{20\%$\sim$25\%}，脂肪\imp{20\%$\sim$30\%}，碳水化合物\imp{45\%$\sim$60\%}。
    \item \textbf{增加膳食纤维：}每日\imp{25$\sim$30g}，增加饱腹感，延缓胃排空。
    \item \textbf{限制添加糖和饱和脂肪：}添加糖$<$总能量10\%，饱和脂肪酸$<$总能量10\%。
    \item \textbf{增加体力活动：}每周至少150分钟中等强度有氧运动 + 每周2$\sim$3次抗阻训练。推荐更高水平（每周200$\sim$300min）以维持体重下降和防止反弹。
    \item \textbf{行为矫正：}自我监测体重、记录饮食日记、规律进餐、三餐合理分配。
\end{enumerate}

\textbf{减重目标：}初始目标为减少原体重的5\%$\sim$10\%，即使未达到理想体重，减轻5\%已可显著改善血糖、血脂和血压等代谢指标。
\end{keybox}

% =====================================================================
%  体重控制与体力活动指南
% =====================================================================
\section{体重控制与体力活动指南}

\subsection{体重控制的意义及方法}

\begin{defbox}
\textbf{体重控制的意义（熟悉）：}
超重和肥胖是心血管疾病、2型糖尿病、多种癌症、骨关节炎、睡眠呼吸暂停等慢性病的重要危险因素。体重控制不仅改善外观，更重要的是降低慢性病风险、改善代谢指标、提高生活质量。即使减轻原体重的5\%-10\%，也可显著改善血糖、血脂和血压等代谢指标。

\textbf{体重控制的方法（熟悉）：}
\begin{enumerate}[leftmargin=*]
    \item \textbf{饮食控制：}控制总能量摄入——每日减少500-1000 kcal，每月减重1-2 kg为宜。女性不低于1000-1200 kcal/d，男性不低于1200-1600 kcal/d。
    \item \textbf{增加体力活动：}有氧运动结合抗阻训练（见体力活动指南）。
    \item \textbf{行为矫正：}自我监测（体重、饮食日记）、规律进餐（三餐合理分配）、控制进食速度、减少外出就餐。
    \item \textbf{综合干预：}饮食+运动+行为矫正联合干预效果优于单一措施。
    \item \textbf{药物和手术治疗：}BMI $\ge$ 28或BMI $\ge$ 24且伴有合并症者，在上述生活方式干预无效时可考虑。
\end{enumerate}
\end{defbox}

\subsection{体力活动指南}

\begin{defbox}
\textbf{体力活动指南（熟悉）：}

\textbf{一般健康人群的身体活动建议：}
\begin{enumerate}[leftmargin=*]
    \item \textbf{有氧运动：}每周至少150分钟中等强度有氧运动（如快走、骑车、游泳），或每周至少75分钟高强度有氧运动（如跑步）。每次持续至少10分钟。
    \item \textbf{抗阻训练：}每周至少2次（非连续日），涉及主要肌肉群（腿、臀、背、腹、胸、肩、臂）。
    \item \textbf{避免久坐：}任何强度的身体活动都有益，减少久坐时间。
\end{enumerate}

\textbf{控制体重的身体活动建议（更高水平）：}
\begin{itemize}[leftmargin=*]
    \item 建议每周增加有氧运动至\textbf{200-300分钟}中等强度。
    \item 维持体重下降和防止减重后反弹即需要达此水平。
\end{itemize}

\textbf{儿童青少年身体活动指南：}
\begin{itemize}[leftmargin=*]
    \item 每天至少\textbf{60分钟}中高强度身体活动。
    \item 每周至少3次高强度活动和力量训练。
    \item 视屏时间每天不超过2小时。
\end{itemize}
\end{defbox}

% =====================================================================
%  营养与痛风
% =====================================================================
\section{营养与痛风}

\begin{defbox}
\textbf{痛风（Gout）：}
痛风是由嘌呤代谢紊乱和/或尿酸排泄减少所致，以高尿酸血症为生化基础，以急性关节炎反复发作、痛风石形成、关节畸形、慢性间质性肾炎和尿酸性肾结石形成为临床特征的代谢性疾病。

\textbf{膳食因素：}
\begin{itemize}[leftmargin=*]
    \item 高嘌呤食物摄入过多（动物内脏、红肉、海鲜尤其贝类和沙丁鱼、肉汤）。
    \item 酒精（尤其啤酒和烈酒）——既增加尿酸生成又抑制尿酸排泄。
    \item 含糖饮料和果糖——促进尿酸生成。
    \item 水摄入不足——影响尿酸从肾脏排泄。
\end{itemize}
\end{defbox}

\begin{keybox}
\textbf{痛风的营养防治原则（掌握）：}

\begin{enumerate}[leftmargin=*]
    \item \textbf{限制总能量摄入：}超重/肥胖者需控制能量，减轻体重可降低血尿酸水平。但减重不宜过快（避免酮体升高抑制尿酸排泄）。
    \item \textbf{限制脂肪和蛋白质摄入：}脂肪供能比控制在20\%-25\%，过量脂肪可阻碍尿酸排泄。蛋白质不超过1.0 g/(kg$\cdot$d)，急性期0.8 g/(kg$\cdot$d)。
    \item \textbf{限制钠盐摄入：}低盐饮食（食盐$<$5 g/d）——有利于尿酸排泄。
    \item \textbf{\imp{限制嘌呤摄入：}}\\ 急性发作期：嘌呤摄入 < 150 mg/d（严格限制动物内脏、海鲜、肉汤、浓汤等）；\\ 缓解期：可适当放宽，但仍应控制高嘌呤食物。
    \item \textbf{增加蔬菜摄入：}多吃新鲜蔬菜水果，提供碱性矿物质和维生素C利于尿酸排泄。菠菜、花菜等虽含中等嘌呤但属植物性嘌呤，对血尿酸影响远小于动物性嘌呤。
    \item \textbf{保证充足饮水：}\imp{每日饮水2000-3000 ml}——稀释尿液，促进尿酸排泄，预防肾结石。
    \item \textbf{\imp{限制饮酒：}}尤其啤酒和烈酒，急性期应禁酒。
    \item \textbf{限制含果糖饮料：}果糖饮料可增加尿酸生成。
\end{enumerate}
\end{keybox}

% =====================================================================
%  第二节 营养与糖尿病
% =====================================================================
\section{营养与糖尿病}

\begin{defbox}
\textbf{糖尿病（Diabetes Mellitus）分类：}
\begin{itemize}[leftmargin=*]
    \item \textbf{1型糖尿病（T1DM）：}胰岛$\beta$细胞破坏，胰岛素绝对缺乏，多发生于儿童和青少年，需终身胰岛素治疗。
    \item \textbf{2型糖尿病（T2DM）：}胰岛素抵抗为主伴相对胰岛素缺乏，或胰岛素分泌缺陷为主伴胰岛素抵抗。占糖尿病患者90\%以上，与肥胖、不合理膳食和缺乏运动密切相关。
\end{itemize}
\end{defbox}

\begin{keybox}
\textbf{GI和GL在糖尿病饮食管理中的应用：}

\textbf{血糖生成指数（Glycemic Index, GI）：}反映食物摄入后升高血糖的速度和能力。
\begin{itemize}[leftmargin=*]
    \item \imp{低GI食物（GI $\leq$ 55）：}全谷物、豆类、牛奶
    \item \imp{中GI食物（GI 56$\sim$69）：}全麦面包、糙米
    \item \imp{高GI食物（GI $\geq$ 70）：}白面包、白米饭、含糖饮料
\end{itemize}
GI越低，血糖升高速度越慢。脂肪和蛋白质可延缓消化速度，因此脂肪和蛋白质含量高的食物GI较低。

\textbf{血糖负荷（Glycemic Load, GL）：}评价某种食物摄入量对人体血糖影响的幅度。\imp{GL $=$ GI $\times$ 食物中碳水化合物含量(g) / 100}。
\begin{itemize}[leftmargin=*]
    \item 低GL：GL $<$ 10；中GL：10$\sim$20；高GL：GL $>$ 20
    \item GI相同的食物，其GL可能不同。GL越高，餐后血糖总体水平越高。
\end{itemize}
\end{keybox}

\begin{keybox}
\textbf{糖尿病的饮食控制原则：}

\begin{enumerate}[leftmargin=*]
    \item \textbf{控制总能量：}以达到或维持理想体重为目标，根据BMI和活动水平个体化制定。
    \item \textbf{合理分配三大营养素：}
    \begin{itemize}[leftmargin=*]
        \item 碳水化合物：\imp{占总能量50\%$\sim$60\%}。优先选择低GI食物（全谷物、豆类），严格限制单糖和双糖。
        \item 蛋白质：\imp{占总能量15\%$\sim$20\%}。优质蛋白占1/3以上。有糖尿病肾病者需控制蛋白质摄入量。
        \item 脂肪：\imp{占总能量20\%$\sim$30\%}。饱和脂肪酸$<$7\%总能量，限制胆固醇$<$300mg/d。
    \end{itemize}
    \item \textbf{定时定量进餐：}三餐能量分配：早1/5、中2/5、晚2/5，或早中晚各1/3。少量多餐，避免血糖大幅波动。
    \item \textbf{增加膳食纤维：}\imp{每日25$\sim$30g}，可溶性纤维有助于降低餐后血糖。
    \item \textbf{综合管理：}饮食治疗是糖尿病\imp{最基本、最有效}的治疗措施，须与运动、药物治疗和血糖监测有机结合。
\end{enumerate}
\end{keybox}

% =====================================================================
%  第三节 营养与心血管疾病
% =====================================================================
\section{营养与心血管疾病}

\begin{keybox}
\textbf{膳食脂肪酸与血脂的关系：}

\begin{itemize}[leftmargin=*]
    \item \imp{饱和脂肪酸（SFA）：}HDL$\downarrow$，LDL$\uparrow$（LDL为心血管危险因素）→ 增加心血管疾病危险。主要来源：动物脂肪、棕榈油、椰子油。
    \item \imp{单不饱和脂肪酸（MUFA）：}HDL$\uparrow$，LDL$\downarrow$。代表：油酸（橄榄油、茶油）。对心血管有保护作用。
    \item \imp{多不饱和脂肪酸（PUFA）：}HDL$\downarrow$，LDL$\downarrow\downarrow$。代表：亚油酸（n-6）、$\alpha$-亚麻酸（n-3），存在于植物油中。虽强力降低LDL，但同时也降低HDL。
    \item \imp{反式脂肪酸：}HDL$\downarrow$，LDL$\uparrow$→ 增加心血管疾病危险。主要来源：氢化植物油、糕点、油炸食品。
\end{itemize}

\textbf{脂肪酸供能比例建议：}SFA:MUFA:PUFA $=$ 1:1:1。SFA $<$ 总能量7\%$\sim$10\%，反式脂肪酸摄入越少越好。
\end{keybox}

\begin{keybox}
\textbf{钠盐与高血压：}
\begin{itemize}[leftmargin=*]
    \item 高钠摄入是\imp{高血压最重要的膳食危险因素}。钠摄入增加→血容量增加→血压升高。
    \item \imp{每日食盐 $<$ 5g（约2000mg钠）}。减少加工食品、咸菜、调味品。
    \item \textbf{钾的降压作用：}钾可促进钠的排出，舒张血管，降低血压。多食蔬菜、水果、豆类和全谷物增加钾摄入。
    \item \textbf{钙、镁的作用：}钙和镁摄入充足也有助于血压控制。
    \item \textbf{膳食胆固醇：}\imp{$<$ 300mg/日}。减少动物内脏、脑、蛋黄等高胆固醇食物。
\end{itemize}
\end{keybox}

\begin{defbox}
\textbf{DASH饮食（Dietary Approaches to Stop Hypertension）：}高钾、高钙、高镁、高纤维、高蛋白，低饱和脂肪、低总脂肪、低胆固醇、低钠。强调蔬菜、水果、低脂奶制品、全谷物、鱼、禽肉和坚果，限制钠盐、红肉、甜食和含糖饮料。临床试验证实可\imp{降低收缩压8$\sim$14mmHg}。
\end{defbox}

\begin{keybox}
\textbf{动脉粥样硬化性心脏病的营养防治：}
\begin{enumerate}[leftmargin=*]
    \item 限制总能量摄入，维持理想体重
    \item 限制脂肪和胆固醇摄入，SFA $<$ 7\%$\sim$10\%总能量
    \item 提高植物性蛋白摄入，少吃甜食
    \item 摄入充足膳食纤维（25$\sim$30g/d）、矿物质和维生素
    \item 适当多吃富含植物化学物的食物
    \item 饮食清淡，少盐限酒
    \item 戒烟，保持积极生活方式
\end{enumerate}
\end{keybox}

% =====================================================================
%  第四节 营养与癌症
% =====================================================================
\section{营养与癌症}

\begin{keybox}
\textbf{膳食致癌因素：}

\begin{enumerate}[leftmargin=*]
    \item \textbf{高温烹调产生的致癌物：}
    \begin{itemize}[leftmargin=*]
        \item \imp{杂环胺（Heterocyclic Amines, HCAs）：}蛋白质含量较高的食物（肉类、鱼类）在高温烹调（烧、烤、煎、炸）时产生。加热温度愈高、时间愈长、水分含量愈少，产生的杂环胺愈多。美拉德反应在杂环胺形成中可能起催化作用。主要靶器官为肝脏，其次为血管、肠道、乳腺等。需经细胞色素P450代谢活化后才具有致突变性。
        \item \imp{多环芳烃（Polycyclic Aromatic Hydrocarbons, PAHs）：}代表物质为苯并(a)芘（B(a)P）。烘烤、烟熏食品中含量较高。B(a)P属于前致癌物，在体内经芳烃羟化酶作用代谢活化为多环芳烃环氧化物，与DNA、RNA和蛋白质等生物大分子结合而诱发肿瘤。
    \end{itemize}
    \item \textbf{腌制食品中的亚硝酸盐：}硝酸盐和亚硝酸盐用作食品防腐剂和护色剂，在体内可与胺类化合物反应生成\imp{亚硝胺（Nitrosamines）}——强致癌物，与食管癌、胃癌等有关。蔬菜从土壤中吸收硝酸盐也是体内硝酸盐的重要来源。
    \item \textbf{霉变食品中的黄曲霉毒素（Aflatoxin, AF）：}黄曲霉和寄生曲霉产生的代谢产物，\imp{黄曲霉毒素B$_1$为已知最强的化学致癌物之一}，主要诱发肝癌。最适产毒温度25$\sim$33$^\circ$C。花生最易受污染，其次为棉籽和油菜籽。
\end{enumerate}
\end{keybox}

\begin{keybox}
\textbf{膳食保护因素：}

\begin{enumerate}[leftmargin=*]
    \item \textbf{膳食纤维：}可吸附杂环胺并降低其活性，降低结直肠癌风险。增加粪便体积，稀释致癌物，缩短肠道转运时间。
    \item \textbf{抗氧化维生素：}维生素C、维生素E、$\beta$-胡萝卜素等可清除自由基，阻断亚硝胺体内合成，保护DNA免受氧化损伤。
    \item \textbf{植物化学物：}蔬菜水果中的酚类、黄酮类、有机硫化合物等成分有抑制杂环胺致突变性和致癌性的作用，可诱导解毒酶活性，抑制肿瘤细胞增殖。
\end{enumerate}
\end{keybox}

\begin{tipbox}
\textbf{世界癌症研究基金会（WCRF）核心建议：}保持健康体重；积极运动（每周至少150分钟）；摄入丰富的全谷物、蔬菜、水果和豆类；限制快餐和高脂/高淀粉/高糖加工食品；限制红肉（每周不超过350$\sim$500g）和加工肉类（加工肉类被IARC列为Ⅰ类致癌物）；限制含糖饮料；限制饮酒；不依赖膳食补充剂预防癌症。
\end{tipbox}

% =====================================================================
%  第五节 营养与骨质疏松
% =====================================================================
\section{营养与骨质疏松}

\begin{defbox}
\textbf{骨质疏松症（Osteoporosis）：}以骨量减少、骨组织微结构破坏、骨脆性增加和易于骨折为特征的代谢性骨病。\imp{绝经后骨质疏松症（PMOP）}是最常见的类型，与雌激素水平下降密切相关。
\end{defbox}

\begin{defbox}
\textbf{钙、维生素D、蛋白质对骨健康的影响：}

\textbf{1. 钙（Calcium）：}
\begin{itemize}[leftmargin=*]
    \item 钙是骨骼最主要的矿物成分，占人体钙总量的99\%。
    \item 长期钙摄入不足→骨钙动员→骨量丢失→骨质疏松。
    \item 膳食因素影响钙吸收：\imp{促进因素：}维生素D、乳糖、适量蛋白质。\imp{抑制因素：}草酸（菠菜）、植酸（谷物）、膳食纤维过多、高钠（钠$\uparrow$→尿钙$\uparrow$）、咖啡因。
    \item 奶制品是钙的最佳来源。
\end{itemize}

\textbf{2. 维生素D：}
\begin{itemize}[leftmargin=*]
    \item 促进肠道钙磷吸收，促进肾小管钙重吸收，调节骨钙动员和沉积。
    \item 缺乏→儿童佝偻病、成人骨软化症、老年人骨质疏松症。
    \item 来源：皮肤经紫外线照射合成（主要来源）、鱼肝油、蛋黄、强化食品。
\end{itemize}

\textbf{3. 蛋白质：}
\begin{itemize}[leftmargin=*]
    \item 适量的蛋白质是维持骨基质（胶原蛋白）合成所必需。
    \item \imp{蛋白质过剩（尤其动物蛋白过多）：}含硫氨基酸摄入过多→代谢产酸→骨钙动员缓冲→加速骨骼中钙的丢失→骨质疏松症风险增加。
\end{itemize}
\end{defbox}

\begin{defbox}
\textbf{峰值骨量（Peak Bone Mass）：}
\begin{itemize}[leftmargin=*]
    \item 人在生命早期（儿童青少年期至30岁左右）骨量达到的最高值。
    \item \imp{峰值骨量越高，晚年发生骨质疏松的风险越低。}提升峰值骨量是预防骨质疏松最重要的策略之一。
    \item 影响峰值骨量的主要因素：遗传（最主要）、钙摄入、维生素D水平、体力活动（尤其是负重运动）、性激素水平。
\end{itemize}

\textbf{骨质疏松的预防措施：}
\begin{enumerate}[leftmargin=*]
    \item 保证充足钙摄入：成人\imp{800$\sim$1000mg/d}，老年人/绝经后女性\imp{1000$\sim$1200mg/d}。
    \item 充足维生素D：多晒太阳，必要时补充维生素D制剂。
    \item 适量优质蛋白质，但避免过量动物蛋白摄入。
    \item 限制高钠饮食（减少尿钙排出）。
    \item 规律负重运动（步行、跑步、跳跃、力量训练），促进骨形成。
    \item 避免吸烟和过量饮酒。
    \item 儿童青少年期最大化峰值骨量，成年后维持骨量，老年期减缓骨丢失。
\end{enumerate}
\end{defbox}

% =====================================================================
%  临床营养
% =====================================================================
\section{临床营养筛查与营养评价}

\subsection{基本概念与营养风险筛查}

\begin{defbox}
\textbf{基本概念（熟悉）：}

\begin{enumerate}[leftmargin=*]
    \item \textbf{营养风险（Nutritional Risk）：}现存或潜在的与营养因素相关的导致病人出现不利临床结局的风险。
    \item \textbf{营养风险筛查（Nutritional Risk Screening）：}发现病人是否存在营养问题和是否需要进一步进行全面营养评估的过程。目的是发现个体是否存在营养不足和有营养不足的危险。
    \item \textbf{NRS 2002（Nutrition Risk Screening 2002）：}
    \begin{itemize}[leftmargin=*]
        \item 地位：目前以临床结局改善与否为目标的唯一风险筛查工具。
        \item 推荐：ESPEN、ASPEN、CSPEN等均给予推荐（回顾性循证研究、前瞻性队列研究验证）。
        \item 适用人群：18岁以上且住院时间超过24小时的患者，不推荐用于未成年人（2018年CSPEN专家共识）。
    \end{itemize}
    \item \textbf{营养评估（Nutritional Assessment）：}临床营养专业人员通过人体测量、生化检查、临床检查及综合营养评价方法等手段，对患者营养代谢和机体功能等进行检查和评定，以确定营养不良的类型及程度。
    \item \textbf{评估金标准：SGA（Subjective Global Assessment，主观全面营养评价）。}
\end{enumerate}

\textbf{膳食营养评价的基本方法：}
\begin{itemize}[leftmargin=*]
    \item 膳食调查（入院前/后摄食情况）、人体测量、生化检验、临床检查。
\end{itemize}

\textbf{营养评定的生化指标（了解）：}
\begin{itemize}[leftmargin=*]
    \item 维生素和微量元素：应激状态（手术、烧伤、败血症等）对维生素和微量元素需要量显著增加。
    \item 血浆氨基酸谱：重度PEM时血浆总氨基酸值明显下降；必需氨基酸（EAA）下降较非必需氨基酸（NEAA）更明显——Val、Leu、Ile和Met下降最多。
    \item 谷氨酰胺（Gln）是人体含量最丰富的氨基酸，精氨酸（Arg）可提高免疫活性、促进蛋白质合成。
\end{itemize}
\end{defbox}

% =====================================================================
\section{临床营养治疗与病人膳食}
% =====================================================================

\subsection{临床营养治疗目的和方法}

\begin{defbox}
\textbf{规范化营养支持治疗4步骤（熟悉）：}
\[
\text{营养筛查（Nutritional Screening）} \rightarrow \text{营养评定（Nutritional Assessment）} \rightarrow \text{营养干预（Nutritional Intervention）} \rightarrow \text{监测（Monitoring）}
\]

\textbf{营养治疗的五步治疗模式——临床营养干预阶梯（了解）：}
\begin{enumerate}[leftmargin=*]
    \item \textbf{首选}饮食营养治疗（膳食调整和营养教育）
    \item 进食不足 → 饮食营养治疗 + \textbf{口服营养补充（ONS）}
    \item 无法进食或只能进食流质 → \textbf{全肠内营养（TEN）}
    \item 消化功能或耐受性差 → \textbf{肠内营养 + 肠外营养（EN+PN）}
    \item 完全无法进食或无法通过胃肠提供营养 → \textbf{全肠外营养（TPN）}
\end{enumerate}
\end{defbox}

\subsection{病人膳食分类}

\begin{defbox}
\textbf{医院膳食分类（熟悉）：}

医院膳食可分为：\textbf{基本膳食（医院常规膳食）、治疗膳食、诊断试验膳食。}

\textbf{一、基本膳食（4种）：}

\begin{enumerate}[leftmargin=*]
    \item \textbf{普通膳食：}与健康人膳食基本相同，能量及营养素供应充足，符合平衡膳食原则。\\ 适用：体温正常或接近正常，无咀嚼或消化吸收功能障碍，无特殊膳食要求者。\\ 配膳原则：品种多样化，合理分配——早25\%-30\%，中40\%，晚30\%-35\%。

    \item \textbf{软食：}质地软、易咀嚼、少渣、易消化。\\ 适用：低热、消化不良或咀嚼困难者，老年患者或3-4岁患儿。

    \item \textbf{半流质膳食：}半流体、细软、更易咀嚼和消化，少食多餐。\\ 适用：发热较高、消化道疾病、耳鼻喉术后、身体虚弱者。\\ 配膳原则：能量低于软食，膳食纤维更少，少食多餐。

    \item \textbf{流质膳食：}极易消化，含渣很少，流体状态或在口腔内融化为液体。\\ 适用：高热、极度虚弱、无咀嚼能力、急性传染病、病情危重或术后，肠道手术术前准备。\\ 配膳原则：不平衡膳食（能量和营养素均不足），每日6-7餐。\\
    分类：清流质（食道/胃肠道大手术后）、浓流质（口腔手术后）、冷流质（扁桃体摘除后、胃肠道出血期）、不胀气流质（腹部手术后）。
\end{enumerate}
\end{defbox}

% =====================================================================
\section{肠内营养、肠外营养与FSMP}
% =====================================================================

\subsection{肠内营养（Enteral Nutrition, EN）}

\begin{defbox}
\textbf{肠内营养（熟悉）：}

\textbf{概念：}有胃肠道消化吸收功能的病人，因机体病理、生理改变或治疗特殊要求，需利用口服等方式给予要素膳制剂，经胃肠道消化吸收，提供能量和营养素，满足机体代谢需要的营养支持疗法。

\textbf{核心地位：}EN是最符合生理要求的途径，也是营养支持的首选途径。

\textbf{优势：}
\begin{itemize}[leftmargin=*]
    \item 营养素经胃肠消化吸收到肝脏，有利于内脏蛋白质合成和代谢调节。
    \item 实施简单、并发症少、促进胃肠功能恢复、促进胃肠激素释放。
    \item 改善门静脉循环、防止肠粘膜萎缩和细菌移位。
    \item 对技术和设备要求低，费用较低。
\end{itemize}

\textbf{局限性：}
\begin{itemize}[leftmargin=*]
    \item 受胃肠功能限制（肠梗阻、高位肠瘘、高流量肠瘘等）。
    \item 受消化功能限制（胃酸缺乏、胆汁分泌障碍、胰酶分泌障碍）。
    \item 受消化和吸收功能限制（吸收不良综合症、短肠综合症、先天性巨结肠等）。
\end{itemize}

\textbf{输入途径和方法：}
\begin{enumerate}[leftmargin=*]
    \item \textbf{经口营养 / ONS：}口服特殊制备的营养物质。当患者经口还能进食但进食量不足或结构不合理时采用。
    \item \textbf{管饲营养：}经鼻-胃管、鼻-十二指肠/空肠管、胃造瘘（PEG）、空肠造瘘等途径输入。
    \item \textbf{输注方式：}一次投给（每次250-400 ml，4-6次/天）、间歇重力滴注、经泵连续输注（12-24小时）。
\end{enumerate}

\textbf{适应证：}①无法经口摄食、摄食不足或有摄食禁忌者；②胃肠道疾病；③胃肠道外疾病。

\textbf{禁忌证：}肠道梗阻。

\textbf{肠内营养制剂（了解）：}
\begin{itemize}[leftmargin=*]
    \item 要素膳（单体膳）：氮源为氨基酸或短肽，无需消化可直接吸收，适用于消化功能严重受损者。
    \item 非要素膳（聚合膳）：氮源为整蛋白，需消化酶消化，适用于消化功能基本正常者。
    \item 组件膳：含单一或某类营养素，可根据病情需要调配。
    \item 特殊用途制剂：针对特定疾病状态（肝病、肾病、糖尿病、呼吸系统疾病等）。
\end{itemize}
\end{defbox}

\subsection{肠外营养（Parenteral Nutrition, PN）}

\begin{defbox}
\textbf{肠外营养（熟悉）：}

\textbf{概念：}通过肠道外的通路即静脉途径输注能量和各种营养素，达到纠正或预防营养不良、维持营养平衡目的的营养补充方式。

\textbf{输入途径：}
\begin{enumerate}[leftmargin=*]
    \item \textbf{周围静脉置管：}方便，可在普通病房实施。适用于短期肠外营养支持或作为肠内营养补充。局限性：需较多液体稀释以保证充足能量，不适合需限制液体的患者。（使用浅表静脉，多为上肢末梢静脉。）
    \item \textbf{中心静脉置管：}如输液港（PORT），尖端位于上腔静脉，适用长期肠外营养。
\end{enumerate}

\textbf{输注形式与方法：}
\begin{enumerate}[leftmargin=*]
    \item \textbf{持续输注：}24h内均匀输注，优点为血糖稳定，缺点为持续高胰岛素水平可能致脂肪肝和肝酶升高。
    \item \textbf{循环输注：}12-18h输注全天营养液，可预防肝毒性，改善生活质量。局限性：需心功能耐受短时间内大量液体。
\end{enumerate}

\textbf{适应证（了解）：}胃肠道功能障碍或衰竭的病人，存在营养不良或预计2周内无法正常饮食者。

\textbf{并发症（熟悉）：}
\begin{enumerate}[leftmargin=*]
    \item \textbf{置管并发症：}气胸、血胸、血肿、损伤胸导管/动脉/神经、空气栓塞等。
    \item \textbf{感染并发症：}导管性败血症是肠外营养最常见的严重并发症。
    \item \textbf{代谢并发症：}液体超负荷、糖代谢紊乱、肝功能损害、酸碱平衡失调、电解质紊乱、代谢性骨病。
    \item \textbf{肠道并发症：}肠道黏膜萎缩（长期不使用肠道）。
\end{enumerate}

\textbf{禁忌证（了解）：}严重循环/呼吸功能衰竭、水电解质平衡紊乱、肝肾衰竭等。

\textbf{肠外营养制剂（了解）：}
\begin{itemize}[leftmargin=*]
    \item 氨基酸注射液（氮源）、脂肪乳剂（能量+必需脂肪酸）、葡萄糖注射液（主要能源）、维生素制剂、微量元素制剂、电解质溶液。
    \item "全营养混合液（TNA）"即"全合一（All-in-One）"是将所有营养成分混合于一个容器中输注。
\end{itemize}
\end{defbox}

\subsection{特殊医学用途配方食品（FSMP）}

\begin{defbox}
\textbf{FSMP概念（熟悉）及分类（了解）：}

\textbf{概念：}特殊医学用途配方食品（Food for Special Medical Purposes, FSMP）是为满足进食受限、消化吸收障碍、代谢紊乱或特定疾病状态人群对营养素或膳食的特殊需要，专门加工配制而成的配方食品。必须在医生或临床营养师指导下，单独食用或与其他食品配合食用。

\textbf{FSMP不同于普通食品和保健食品：}FSMP不是药品，不能替代药物的治疗作用，但它是疾病治疗和康复过程中不可或缺的营养支持手段，属于肠内营养的范畴。

\textbf{分类（了解）：}
\begin{enumerate}[leftmargin=*]
    \item \textbf{全营养配方食品：}可作为单一营养来源满足目标人群营养需求。
    \item \textbf{特定全营养配方食品：}可作为单一营养来源满足目标人群在特定疾病或医学状态下的营养需求（如糖尿病全营养配方、肾病全营养配方、肿瘤全营养配方等）。
    \item \textbf{非全营养配方食品：}可满足目标人群部分营养需求，不适用于作为单一营养来源（如蛋白质组件、脂肪组件、碳水化合物组件、电解质配方、增稠组件等）。
\end{enumerate}

\textbf{管理：}我国对FSMP实行注册制管理，需经国家市场监督管理总局注册批准后方可生产销售。
\end{defbox}

% =====================================================================
%  复习题
% =====================================================================
\reviewcontent{
\section{复习题}
\begin{enumerate}[leftmargin=*, itemsep=8pt]
    \item \textbf{【名词解释】}BMI；中心性肥胖；GI；GL；DASH饮食；峰值骨量；痛风；NRS 2002；SGA；FSMP
    \item \textbf{【简答题】}简述膳食营养与慢性病（NCD）的关系。
    \item \textbf{【简答题】}简述体重控制的意义及方法。
    \item \textbf{【简答题】}简述体力活动指南对一般人群和体重维持的建议。
    \item \textbf{【简答题】}简述肥胖的判定标准和膳食营养防治策略。
    \item \textbf{【简答题】}简述GI和GL在糖尿病饮食管理中的应用。
    \item \textbf{【简答题】}简述糖尿病的饮食控制原则（至少五点）。
    \item \textbf{【简答题】}试述膳食脂肪酸（SFA、MUFA、PUFA、反式脂肪酸）与血脂的关系。
    \item \textbf{【简答题】}简述钠盐与高血压的关系及心血管疾病的膳食预防要点。
    \item \textbf{【简答题】}列举膳食致癌因素及其食物来源。
    \item \textbf{【简答题】}列举膳食保护因素及其防癌机制。
    \item \textbf{【简答题】}简述钙、维生素D、蛋白质对骨健康的影响。
    \item \textbf{【简答题】}简述峰值骨量的概念及骨质疏松的营养预防措施。
    \item \textbf{【简答题】}简述痛风的营养防治原则。
    \item \textbf{【简答题】}简述临床营养支持治疗的四步规范化流程和五步营养干预阶梯。
    \item \textbf{【简答题】}简述医院基本膳食的分类及各自适用对象。
    \item \textbf{【简答题】}简述肠内营养和肠外营养的概念、适应证、并发症，以及FSMP的概念和分类。
\end{enumerate}

\subsection*{参考答案要点}
\begin{enumerate}[leftmargin=*, itemsep=5pt]
    \item \textbf{BMI} = 体重(kg) / [身高(m)]$^2$，我国标准：BMI $\geq$ 28为肥胖，24.0$\sim$27.9为超重，18.5$\sim$23.9为正常。\textbf{中心性肥胖}：腰围男$\geq$85cm或女$\geq$80cm，或腰臀比男$>$0.9、女$>$0.85。\textbf{GI（血糖生成指数）：}反映食物摄入后升高血糖的速度和能力，低GI $\leq$55，中GI 56$\sim$69，高GI $\geq$70。\textbf{GL（血糖负荷）：}$=$ GI $\times$ 食物中碳水化合物含量(g) / 100，评价食物摄入量对血糖影响的幅度。\textbf{DASH饮食：}高钾/高钙/高镁/高纤维/高蛋白，低饱和脂肪/低总脂肪/低胆固醇/低钠的膳食模式，可降低收缩压8$\sim$14mmHg。\textbf{峰值骨量：}生命早期骨量达到的最高值（约30岁），峰值骨量越高，晚年骨质疏松风险越低。\textbf{痛风：}嘌呤代谢紊乱和/或尿酸排泄减少所致代谢病，以高尿酸血症为基础。\textbf{NRS 2002：}营养风险筛查2002，目前以临床结局改善为目标的唯一风险筛查工具。\textbf{SGA：}主观全面营养评价，是营养评估金标准。\textbf{FSMP：}特殊医学用途配方食品，为满足特定疾病人群特殊营养需要专门配制的配方食品。

    \item 膳食营养与NCD的关系：不合理膳食是NCD最重要可改变危险因素之一。高能量密度食物+含糖饮料等→肥胖→心血管病/糖尿病/多种癌症共同危险因素；高钠高饱和脂肪→心血管疾病；高GI/GL膳食→2型糖尿病；加工肉类（Ⅰ类致癌物）+酒精→癌症。合理膳食是NCD一级预防的基石。

    \item 体重控制意义：超重肥胖是心血管病、T2DM、多种癌症等的重要危险因素，减轻5\%-10\%体重即可显著改善代谢指标。方法：(1)饮食控制（日减500-1000kcal）；(2)增加体力活动（有氧+抗阻）；(3)行为矫正（自我监测、规律进餐）；(4)综合干预效果最佳；(5)严重肥胖者可考虑药物或手术。

    \item 体力活动指南——一般人群：每周$\ge$150min中等强度有氧或75min高强度+每周2次抗阻训练；体重维持：建议每周200-300min中等强度；儿童青少年：每天$\ge$60min中高强度活动+视屏时间$<$2h/d。

    \item \textbf{判定标准：}BMI $\geq$ 28为肥胖，24.0$\sim$27.9为超重；腰围男$\geq$85cm、女$\geq$80cm为中心性肥胖。\textbf{防治策略：}控制总能量（每日减少500$\sim$1000kcal），调整宏量营养素供能比（蛋白质20\%$\sim$25\%，脂肪20\%$\sim$30\%，碳水45\%$\sim$60\%），增加膳食纤维（25$\sim$30g/d），限制添加糖和饱和脂肪，配合有氧+抗阻运动（每周$\geq$150min中等强度），行为矫正（自我监测、规律进餐），初始减重目标为原体重的5\%$\sim$10\%，每月减重1$\sim$2kg。

    \item \textbf{GI的应用：}指导糖尿病患者选择低GI食物（全谷物、豆类、牛奶，GI $\leq$ 55），低GI有助于改善糖代谢，有效降低血糖水平；同时也可用于控制体重和慢性病。\textbf{GL的应用：}GL = GI $\times$ 碳水含量 / 100，GL $<$ 10为低GL、10$\sim$20为中GL、$>$ 20为高GL。GI相同的食物GL可能不同，GL越高餐后血糖总体水平越高，因此不仅要看GI，也要控制碳水化合物的摄入量。

    \item 五点原则：(1)控制总能量，以达到或维持理想体重为目标；(2)合理分配宏量营养素——碳水50\%$\sim$60\%（优先低GI食物）、蛋白质15\%$\sim$20\%（优质蛋白$>$1/3）、脂肪20\%$\sim$30\%（SFA $<$ 7\%）；(3)定时定量进餐，少量多餐，三餐能量分配早1/5、中2/5、晚2/5；(4)增加膳食纤维25$\sim$30g/d，可溶性纤维降低餐后血糖；(5)饮食治疗须与运动、药物治疗和血糖监测有机结合，饮食治疗是糖尿病最基本最有效的治疗。

    \item \textbf{SFA（饱和脂肪酸）：}HDL$\downarrow$、LDL$\uparrow$→增加心血管疾病危险，来源为动物脂肪、棕榈油。\textbf{MUFA（单不饱和脂肪酸）：}HDL$\uparrow$、LDL$\downarrow$→对心血管有保护作用，代表为油酸（橄榄油）。\textbf{PUFA（多不饱和脂肪酸）：}HDL$\downarrow$、LDL$\downarrow\downarrow$，强力降低LDL但也降低HDL，来源为植物油。\textbf{反式脂肪酸：}HDL$\downarrow$、LDL$\uparrow$→增加心血管疾病危险，来源为氢化植物油、糕点。建议SFA:MUFA:PUFA = 1:1:1。

    \item \textbf{钠盐与高血压：}高钠摄入是高血压最重要的膳食危险因素，钠$\uparrow$→血容量$\uparrow$→血压$\uparrow$，每日食盐$<$5g（约2000mg钠）。钾可促进钠排出、舒张血管而降低血压。\textbf{膳食预防要点：}控制能量维持理想体重，限制SFA（$<$7\%$\sim$10\%能量）和反式脂肪，限制胆固醇$<$300mg/d，控制钠摄入$<$5g/d盐，增加钾/钙/镁摄入，增加膳食纤维25$\sim$30g/d，限制饮酒，戒烟。

    \item \textbf{三类膳食致癌因素：}(1)\textbf{杂环胺：}高温烹调（烧/烤/煎/炸）蛋白质含量高的肉类产生，温度越高、时间越长、水分越少，产生越多，靶器官主要为肝脏。(2)\textbf{多环芳烃：}代表为苯并(a)芘，烟熏和烘烤食品中含量高，属前致癌物，经代谢活化后与DNA结合诱发肿瘤。(3)\textbf{亚硝酸盐：}腌制食品中的防腐剂和护色剂，与胺类反应生成亚硝胺（强致癌物），与食管癌、胃癌有关。(4)\textbf{黄曲霉毒素：}霉变食品，黄曲霉毒素B$_1$为最强化学致癌物之一，主要诱发肝癌，花生最易污染。

    \item \textbf{三大膳食保护因素：}(1)\textbf{膳食纤维：}吸附杂环胺降低其活性，稀释致癌物，缩短肠道转运时间，降低结直肠癌风险。(2)\textbf{抗氧化维生素（VC、VE、$\beta$-胡萝卜素）：}清除自由基，阻断亚硝胺体内合成，保护DNA。(3)\textbf{植物化学物（酚类、黄酮类等）：}抑制杂环胺致突变性和致癌性，诱导解毒酶活性，抑制肿瘤细胞增殖。

    \item \textbf{钙：}骨骼最主要矿物成分（占99\%），摄入不足→骨钙动员→骨质疏松。促进吸收因素：维生素D、乳糖、适量蛋白质；抑制因素：草酸、植酸、高钠、咖啡因。\textbf{维生素D：}促进肠道钙磷吸收和肾小管钙重吸收，调节骨代谢，缺乏→佝偻病/骨软化症/骨质疏松症。\textbf{蛋白质：}适量蛋白质维持骨基质合成所必需；蛋白质过剩（尤其动物蛋白过多）→含硫氨基酸产酸→骨钙动员缓冲→加速骨钙丢失。

    \item \textbf{峰值骨量：}生命早期（至30岁左右）骨量达到的最高值。峰值骨量越高，晚年骨质疏松风险越低。遗传为最主要影响因素。\textbf{预防措施：}保证充足钙摄入（成人800$\sim$1000mg/d，老年人1000$\sim$1200mg/d），充足维生素D，适量优质蛋白但避免过量，限制高钠饮食，规律负重运动，避免吸烟和过量饮酒。分三期策略：儿童青少年期最大化峰值骨量，成年后维持骨量，老年期减缓骨丢失。

    \item 痛风营养防治原则8条：(1)限制总能量摄入——减轻体重降低血尿酸，减重不宜过快；(2)限制脂肪——供能比20\%-25\%，过多阻碍尿酸排泄；(3)限制钠盐——$<$5g/d；(4)\textbf{限制嘌呤摄摄入——核心}：急性期$<$150mg/d（禁内脏海鲜肉汤），缓解期可放宽；(5)增加蔬菜摄入——提供碱性矿物质和VC；(6)\textbf{保证充足饮水——2000-3000ml/d}促尿酸排泄；(7)限酒——尤其啤酒烈酒，急性期禁酒；(8)限制含果糖饮料。

    \item 规范化营养支持治疗4步骤：营养筛查→营养评定→营养干预→监测。五步干预阶梯：(1)饮食治疗（首选）；(2)饮食+口服营养补充（ONS）；(3)全肠内营养（TEN）；(4)肠内+肠外（EN+PN）；(5)全肠外营养（TPN）。

    \item 医院基本膳食4类：普通膳食（体温正常无特殊要求）、软食（低热/咀嚼困难/老年患者）、半流质（发热较高/消化道疾病/术后）、流质（高热/极度虚弱/危重/术前，分为清流质、浓流质、冷流质、不胀气流质）。

    \item EN概念：通过胃肠途径给予要素膳制剂的营养支持方法，是营养支持首选途径（最符合生理）。适应证：无法经口摄食/摄食不足、胃肠道疾病、胃肠道外疾病。禁忌证：肠道梗阻。并发症：腹泻、腹胀、误吸。PN概念：通过静脉途径输注能量和营养素的营养补充方式。适应证：胃肠功能障碍或衰竭，存在营养不良或预计2周内无法正常饮食。并发症：置管并发症（气胸血胸）、感染（导管性败血症最常见严重并发症）、代谢并发症、肠道黏膜萎缩。FSMP：特殊医学用途配方食品，为满足特定疾病人群特殊营养需要专门配制的配方食品，须在医生/临床营养师指导下使用。分类：全营养配方、特定全营养配方、非全营养配方。
\end{enumerate}

}
\newpage

% =====================================================================
%  CHAPTER 8: 食品污染及其预防 (Food Contamination and Prevention)
%  参考：食品主要参考.txt 第1837-2110行 (PRIMARY SOURCE)
% =====================================================================
\chapter{食品污染及其预防}
\chapterintro{
\keyword{掌握}食品污染的概念、来源、分类及对健康的危害；水分活性（Aw）的定义及与微生物生长的关系；常见食品细菌菌属及特征；细菌菌相的概念及食品卫生学意义；菌落总数和大肠菌群的食品卫生学意义及MPN法；黄曲霉毒素的体内代谢、产毒条件、毒性及预防；食品腐败变质的概念、原因、鉴定指标及各类营养素的腐败产物；D值及各种食品保藏方法；$N$-亚硝基化合物、多环芳烃、杂环胺、氯丙醇、丙烯酰胺的污染及预防；食品容器包装材料的卫生问题。\keyword{熟悉}霉菌产毒条件和特点；霉菌菌相；农药和兽药残留。\keyword{了解}霉菌毒素的共同特点；其他霉菌毒素（镰刀菌毒素等）；放射性污染及物理性污染。
}

% =====================================================================
%  第一节 食品污染概述
% =====================================================================
\section{食品污染概述}

\subsection{食品污染的概念与分类}

\begin{keybox}
\textbf{食品污染（food pollution/contamination）：}食品中含有可能对身体健康造成危害或影响其食用价值及商品价值的生物性、化学性及物理性物质，称为食品污染。在食品种植、养殖、生产、加工、储存、运输、销售、烹饪、食用等各个环节中都可能进入或产生于食品中的有害因素，这些过程均可以造成食品污染，而这些有害因素则称为食品污染物（food pollutants）。

\textbf{食品危害的来源（四类）：}
\begin{enumerate}[leftmargin=*]
    \item 食品本身天然含有或自身变化产生的有毒有害物质；
    \item 食品加工储运过程中污染或形成的有害物质；
    \item 外界污染：致病性微生物、寄生虫、农药、重金属等有害化学物及其他污染物污染等；
    \item 食品加工保藏过程中故意加入的成分、食品添加剂、非法添加的有毒有害非食用物质。
\end{enumerate}
\end{keybox}

\subsection{食品污染的分类}

\begin{keybox}
\textbf{食品污染按性质分类（三大类*）：}

\textbf{（1）生物性污染：}微生物、寄生虫和昆虫污染。微生物污染主要是细菌、细菌毒素、霉菌和霉菌毒素以及病毒等的污染，其中细菌和细菌毒素对食品的污染最常见、最重要。病毒污染主要包括甲肝病毒、轮状病毒、诺如病毒和禽流感病毒等。寄生虫和昆虫污染主要有蛔虫、绦虫、囊虫及蝇蛆等。

\textbf{（2）化学性污染：}
\begin{enumerate}[leftmargin=*]
    \item 农药、兽药不合理使用，残留于食品中；
    \item 工业三废、城市生活垃圾未处理随意排放，汞、镉、铅、砷、多环芳烃等有毒有害有机污染物排放到环境中，通过环境转移到食品中；
    \item 食品接触材料、运输工具等接触食品时，食品中的有害物质；
    \item 滥用食品添加剂；
    \item 食品加工储存过程中可能产生的物质，如腌制、烟熏、烘烤等食品过程中产生的$N$-亚硝基化合物、多环芳烃、杂环胺、丙烯酰胺，以及酒中有害的醇类、醛类等；
    \item 掺假、制假过程中加入的物质，如苏丹红、三聚氰胺。
\end{enumerate}

\textbf{（3）物理性污染：}
\begin{enumerate}[leftmargin=*]
    \item 食品的杂物污染：食品生产、加工、储藏、运输、销售等过程中的污染物，包括粮食收割时混入的草籽、液体食品容器中的杂物、食品运输过程中的灰尘等；
    \item 食品的放射性污染：主要来自核爆炸、核废料排放、放射性物质的开采、冶炼、应用及核事故造成的食品污染。转移途径：水—水生生物—植物、动物—人的转移。
\end{enumerate}
\end{keybox}

\subsection{食品污染造成的危害}

\begin{keybox}
\textbf{食品污染造成的危害*：}
\begin{enumerate}[leftmargin=*]
    \item 影响食品的感官性状与营养价值，影响食品生产；
    \item 对身体健康的不良影响，包括急性中毒、慢性危害以及致畸、致突变和致癌作用等。
\end{enumerate}
\end{keybox}

% [补充自往届笔记]
\begin{tipbox}
\textbf{食品安全危害的优先控制顺序（WHO建议）：}微生物污染 $>$ 天然毒素 $>$ 化学性污染 $>$ 食品添加剂。这与公众认知恰好相反——公众往往过分担忧食品添加剂，而忽视微生物污染这一最主要的风险来源。
\end{tipbox}

% =====================================================================
%  第二节 微生物污染及其预防
% =====================================================================
\section{微生物污染及其预防}

\subsection{污染食品的微生物分类}

\begin{keybox}
\textbf{污染食品的微生物按致病能力分类（三大类*）：}

\begin{enumerate}[leftmargin=*]
    \item \textbf{致病性微生物：}能直接对人体致病并造成危害，包括致病性细菌、细菌毒素、人畜共患传染病病原菌和病毒、产毒霉菌和霉菌毒素。
    \item \textbf{条件致病性微生物（相对致病菌）：}通常情况下不致病，但在一定特殊条件下才有效病力的一些微生物。
    \item \textbf{非致病性微生物：}在自然界分布非常广泛，是引起食品腐败变质和食品卫生质量下降的主要原因。
\end{enumerate}
\end{keybox}

\subsection{水分活性（Water Activity, Aw）}

\begin{keybox}
\textbf{水分活性（Water Activity, Aw）的定义*：}在同一条件（温度、湿度、压力等）下，食品水分的蒸汽压（P）与纯水蒸汽压（$P_0$）之比，即 $\text{Aw} = P/P_0$。纯水Aw为1。由于食品中溶质能降低水分的蒸汽压，故Aw值范围为0$\sim$1。Aw反映的是溶液中可被微生物利用的实际含水量，可用来表示食品中可被微生物利用的实际含水量。每个微生物只能在一定的Aw值范围内生长。

\textbf{Aw值与食品中微生物生长繁殖的关系*：}
\begin{enumerate}[leftmargin=*]
    \item 每一种微生物在食品中生长繁殖都有其最低的Aw要求，低于这一要求，微生物的生长繁殖就会受到抑制；
    \item Aw低于0.60时，绝大多数微生物无法生长，Aw小的食品较少出现腐败变质；
    \item 一般说来，细菌生长需要的Aw $>$ 0.9，酵母菌 $>$ 0.87，霉菌 $>$ 0.8；
    \item 同属不同种的微生物对Aw的要求也可不一样；
    \item 细菌形成芽孢时需要比生长时更高的Aw值。
\end{enumerate}
\end{keybox}

% [补充自往届笔记] — supplementary details on Aw
\begin{tipbox}
往届笔记补充：水分活性（water activity, Aw）反映食品中可被微生物利用的水分。每一类微生物在食品中生长繁殖都有其最低的Aw要求，当食品Aw低于这一要求，微生物的生长繁殖就会受到抑制。Aw $<$ 0.6时，绝大多数微生物无法生长，Aw小的食品较少出现腐败变质。细菌 Aw $>$ 0.9，酵母菌 Aw $>$ 0.87，霉菌 Aw $>$ 0.8。
\end{tipbox}

\subsection{常见食品细菌}

\begin{keybox}
\textbf{引起食品腐败变质的常见细菌*：}

\begin{enumerate}[leftmargin=*]
    \item \textbf{假单胞菌属：}需氧性无芽孢杆菌，是冷藏食品（蛋白质变质、冷冻食品）腐败的主要细菌。
    \item \textbf{微球菌属：}为需氧性无芽孢杆菌，是葡萄球菌中常见的腐败菌，常见于未经处理的水和蔬菜中引起腐败。
    \item \textbf{微球菌和葡萄球菌属：}为革兰阳性菌，能产生过氧化氢酶、色素并分解食品中碳水化合物并产生色素。
    \item \textbf{芽孢杆菌属与梭状芽孢杆菌属：}为革兰阳性芽孢杆菌，是罐头食品中常见的腐败菌。
    \item \textbf{肠杆菌科：}为需氧性无芽孢杆菌，与水产制品、肉及蛋的腐败有关。肠杆菌科中除志贺菌属、沙门菌属外，也是常见食品腐败菌。大肠埃希菌是食品中常见的腐败菌，也是食品被粪便污染指示菌之一，变形杆菌属是腐败菌的代表。
\end{enumerate}
\end{keybox}

\subsection{细菌菌相}

\begin{defbox}
\textbf{细菌菌相（Bacterial Flora）：}共存于食品中的细菌种类及其相对数量的构成。其中相对数量较大的细菌称为\textbf{优势菌种（属、株）}。

\textbf{食品卫生学意义：}
\begin{enumerate}[leftmargin=*]
    \item 影响食品腐败变质的特征：使食品的食用价值降低，甚至不能食用。不同优势菌种产生的腐败产物和感官变化各不相同。
    \item 影响食品腐败变质的速度和程度：优势菌种的生长速率和分解能力直接决定腐败进程。在腐败过程中可出现霉菌繁殖，并产生霉菌毒素引起中毒。
    \item 通过分析细菌菌相可判断食品的污染来源和储藏条件。
\end{enumerate}
\end{defbox}

\subsection{菌落总数（Aerobic Plate Count）}

\begin{keybox}
\textbf{菌落总数（Aerobic Plate Count）的定义*：}指单位质量（g）、容积（mL）或表面积（$\text{cm}^2$）的被检样品，在规定条件下（培养基的pH、培养温度与时间、需氧性质等）培养后，所生成的细菌菌落总数。以菌落形成单位（colony forming unit, CFU）表示。

\textbf{食品卫生学意义*：}
\begin{enumerate}[leftmargin=*]
    \item 可作为食品被细菌污染程度及卫生状态的标志；
    \item 是可用于预测食品的存放期限（耐储性）：食品中细菌数量越多，食品腐败变质的速度就越快。
\end{enumerate}
\end{keybox}

\begin{tipbox}
\textbf{注意：}菌落总数低不等于完全无致病菌——因为致病菌可能在低菌落总数条件下存在（如少量沙门菌即可致病）；菌落总数高也不一定代表有致病菌。因此菌落总数是卫生质量的指示指标，而非安全性的直接指标，必须与致病菌检测结合使用。
\end{tipbox}

\subsection{大肠菌群（Coliform Group）与MPN法}

\begin{keybox}
\textbf{大肠菌群（Coliform）的定义*：}指在一定培养条件下能发酵乳糖、产酸产气、需氧和兼性厌氧的革兰阴性无芽孢杆菌。主要包括大肠杆菌科的埃希菌属、柠檬酸杆菌属、肠杆菌属和克雷伯菌属。

\textbf{食品卫生学意义*：}
\begin{enumerate}[leftmargin=*]
    \item \textbf{作为食品受到人与温血动物粪便污染指示菌：}因为大肠菌群直接来源是人与温血动物粪便。典型大肠杆菌的检出说明检品在近期被粪便污染，非典型大肠杆菌的检出说明检品曾经受粪便污染过一段时间（新旧粪便污染）。
    \item \textbf{作为肠道致病菌污染食品的指示菌：}因为大肠菌群与肠道致病菌来源相同，且在一般条件下大肠菌群在外界存活时间与主要肠道致病菌是一致的，食品中大肠杆菌的检出说明食品中有被肠道致病菌污染的可能。
    \item 检测食品加工者、加工设备、存储条件等的污染情况。
\end{enumerate}

\textbf{大肠杆菌检测方法：}
\begin{enumerate}[leftmargin=*]
    \item \textbf{最可能数法（Most Probable Number, MPN）：}当食品中大肠菌群数量较低时采用，相当于每克（mL）食品中大肠菌群的最可能数（MPN）来表示。MPN法是基于泊松分布的一种间接计数方法，是按照一定方法应用统计学原理推算的大肠菌群MPN值。
    \item \textbf{平板计数法：}当食品中大肠菌群数量较高时采用，用平板计数法计大肠菌群的菌落数量表示，表示为每克（mL）检品中大肠菌群的菌落数量CFU/g（mL）。
\end{enumerate}

\textbf{大肠菌群数量单位MPN（most probable number）的定义：}
食品中大肠菌群数量的表示方法，当食品中大肠菌群数量较低时以MPN表示。MPN相当于1g或1mL食品中大肠菌群的最可能数，是按照泊松分布的统计值（GB/T 4789.3-2016），以"MPN/克"或"MPN/毫升"为单位。
\end{keybox}

\subsection{粪便污染指示菌应具备的条件}

\begin{defbox}
\textbf{构成粪便污染指示菌需具备的条件：}
\begin{enumerate}[leftmargin=*]
    \item 仅来自肠道；
    \item 在肠道中数量较多，易于检出；
    \item 在外环境中能生存一定期间；
    \item 食物细菌检验方法敏感简易。
\end{enumerate}

\textbf{肠球菌（粪链球菌）作为粪便污染指示菌的意义：}
肠球菌反映的是低温类食品被粪便污染的指示菌。大肠菌群是常温类食品粪便污染的良好指示菌，但对于低温冷藏/冷冻食品，大肠菌群在低温下易死亡，不宜作为指示菌，而肠球菌对低温抵抗力较强，更适合作为低温类食品的粪便污染指示菌。% [补充自往届笔记]
\end{defbox}

% =====================================================================
%  第三节 霉菌及其毒素对食品的污染
% =====================================================================
\section{霉菌及其毒素对食品的污染}

\subsection{霉菌与霉菌毒素基本概念}

\begin{defbox}
\textbf{霉菌（fungi）：}一类不具叶绿素、无根茎叶分化、具有细胞壁的真核细胞微生物。凡是生长在营养基质上形成绒毛状、蛛网状或絮状菌丝体的真菌，统称为霉菌。

\textbf{霉菌毒素（mycotoxins）：}在霉菌检测中主要是指霉菌污染食品后所产生的有毒代谢产物。

\textbf{霉菌毒素致病特点（与细菌毒素相比）：}
\begin{enumerate}[leftmargin=*]
    \item 毒素较稳定；
    \item 通常通过天然食品中毒；
    \item 中毒以实质器官受损为主；
    \item 没有传染性和免疫性；
    \item 有明显的季节性和地区性特点。
\end{enumerate}
\end{defbox}

\subsection{霉菌产毒特点与产毒条件}

\begin{defbox}
\textbf{霉菌产毒特点*：}
\begin{enumerate}[leftmargin=*]
    \item 霉菌产毒只限于少数的产毒霉菌，而且同一菌种有不同的菌株，产毒能力也不同；
    \item 同一产毒菌株的产毒能力可具有可变性和易变性；
    \item 产毒菌种所产生的霉菌毒素不具有严格的专一性，即一种菌种或菌株可以产生多种不同的毒素，而同一霉菌毒素也可由多种霉菌产生；
    \item 产毒霉菌产生毒素也需要一定的时间。
\end{enumerate}

\textbf{霉菌产毒条件：}
\begin{enumerate}[leftmargin=*]
    \item \textbf{基质：}营养丰富的食品，霉菌生长的可能性大，天然食品上比人工合成培养基更易繁殖。
    \item \textbf{水分：}粮食水分以17\%$\sim$18\%为霉菌生长繁殖的最适条件，粮食Aw值降至0.7以下，一般霉菌均不能生长。
    \item \textbf{湿度：}在不同的相对湿度中，生长繁殖的霉菌种类不同，相对湿度为70\%时霉菌即不能产毒。
    \item \textbf{温度：}霉菌生长繁殖最适宜的温度为25$\sim$30$^\circ$C，在0$^\circ$C以下或30$^\circ$C以上时，霉菌的产毒能力明显减弱或丧失。
    \item \textbf{通风：}大部分霉菌生长繁殖和产毒需要氧气，但毛霉、厌绿曲霉、扩展青霉产毒可在厌氧条件下，能耐受高浓度的CO$_2$。
\end{enumerate}
\end{defbox}

% [补充自往届笔记] — AFB1代谢途径的完整描述
% 往届笔记中有"环氧反应：AFB1-2,3-环氧化物，一部分为解毒过程，一部分发挥毒性、致癌性和致突变性"，这是对代谢途径的重要补充

\subsection{主要产毒霉菌与霉菌毒素}

\begin{defbox}
\textbf{主要产毒霉菌：}
\begin{enumerate}[leftmargin=*]
    \item \textbf{曲霉属：}如黄曲霉、赭曲霉、杂色曲霉；
    \item \textbf{青霉属：}展青霉；
    \item \textbf{镰刀菌属：}禾谷镰刀菌；
    \item \textbf{其他：}黑绿色木霉等。
\end{enumerate}

\textbf{主要霉菌毒素：}黄曲霉毒素、赭曲霉毒素、杂色曲霉毒素、展青霉素、橘青霉素、脱氧雪腐镰刀菌烯醇、玉米赤霉烯酮、伏马菌素、曲酸、单端孢霉烯族化合物、玉米赤霉烯酮、丁烯酸内酯等。
\end{defbox}

\subsection{霉菌污染食品的食品卫生学意义与霉菌菌相}

\begin{defbox}
\textbf{霉菌污染食品的卫生学意义：}
霉菌污染食品的程度以及污染食品的霉菌种类，可从霉菌污染度和霉菌菌相构成两个方面来评定。

\textbf{霉菌菌相（Mold Flora）：}
\begin{enumerate}[leftmargin=*]
    \item \textbf{田野霉}（如气连孢霉、芽枝霉、头孢霉等）占优势：表示粮食新鲜；
    \item \textbf{曲霉和青霉}检出：不表明粮食已霉变，但提示储藏条件可能不良；
    \item \textbf{毛霉、根霉、木霉}检出：说明粮食已开始霉变——这些是储藏后期腐败菌。
\end{enumerate}
\end{defbox}

\subsection{黄曲霉毒素（Aflatoxin, AF）}

\subsubsection{化学结构与理化性质}

\begin{keybox}
\textbf{黄曲霉毒素（Aflatoxin, AF）：}是黄曲霉和寄生曲霉产生的一类有机物。并非所有黄曲霉菌株都能产生AF，但某些寄生曲霉所有菌株几乎全部会产生AF。我国黄曲霉产毒占60\%，寄生曲霉产毒占40\%，是我国长江沿岸和长江以南地区粮食和食品中常见的污染。

\textbf{化学结构：}二呋喃环 + 香豆素（氧杂萘邻酮），基本结构为二呋喃氧杂萘邻酮。AF的毒性与结构有关，含二呋喃环末端有双键者毒性较强并有致癌性。AF的毒性排序：B1 $>$ M1 $>$ G1 $>$ B2 $>$ M2。在天然污染的食品中以AFB1污染最多见，而且其毒性和致癌性也最强，因此在食品卫生监测中常以AFB1作为污染指标。

\textbf{理化性质：}
\begin{enumerate}[leftmargin=*]
    \item \textbf{耐热性强：}一般烹调温度（100$\sim$120$^\circ$C）几乎不被破坏，需\textbf{268$\sim$269$^\circ$C}才完全分解；
    \item \textbf{脂溶性：}难溶于水（溶解度仅10$\sim$20 mg/L），易溶于油脂和多种有机溶剂（氯仿、甲醇等）；
    \item \textbf{在中性或弱酸性溶液中较稳定：}在强碱性条件下（pH 9$\sim$10），内酯环被打开，可被破坏。此原理被用于碱炼去毒；
    \item \textbf{紫外线敏感：}在紫外线下产生荧光（B族发蓝紫色荧光，G族发黄绿色荧光），可用于检测。
\end{enumerate}
\end{keybox}

\subsubsection{产毒条件与污染食品}

\begin{keybox}
\textbf{产毒条件*：}
\begin{itemize}[leftmargin=*]
    \item \textbf{温度：}黄曲霉生长繁殖的温度范围12$\sim$42$^\circ$C，最适产毒温度为25$\sim$33$^\circ$C；
    \item \textbf{水分：}最适Aw值为0.93$\sim$0.98；
    \item \textbf{基质：}花生、玉米（产毒量最高），其次为大米、小麦、大麦、豆类等。
\end{itemize}

\textbf{易污染食品：}
\begin{itemize}[leftmargin=*]
    \item AF主要污染粮油及其制品，其中以花生和玉米污染最为严重，其次为棉籽油，再次为大米、小麦、大麦、豆类等；
    \item 干果类食品（核桃、杏仁、榛子）；
    \item 动物性食品（乳及乳制品、肝脏、干咸鱼）；
    \item 家庭自制发酵食品；
    \item 奶牛食入被AFB1污染的饲料后，牛乳中排出AFB1的羟基化代谢产物AFM1。牛乳及乳制品中的AFM1比较稳定，且目前一般的家庭加热及加工过程均对AFM1水平无显著影响。
\end{itemize}
\end{keybox}

\subsubsection{黄曲霉毒素的体内代谢（重点掌握）*}

\begin{keybox}
\textbf{AFB1的体内代谢过程*\difficulty：}

AFB1进入机体后主要在肝脏进行代谢，在肝细胞微粒体混合功能氧化酶系（CYP450）催化下，发生羟化、脱甲基和环氧化反应，其中主要的羟基化产物为终致癌物。

\textbf{代谢途径：}
\begin{enumerate}[leftmargin=*]
    \item \textbf{AFB1 $\rightarrow$ AFM1（羟化反应）：}AFB1在肝微粒体酶催化下的羟化产物，存在于牛乳和乳制品中。奶牛摄入AFB1污染饲料后，乳汁中排出AFM1。
    \item \textbf{AFB1 $\rightarrow$ AFQ1（羟化反应）：}AFB1的另一种羟化产物，将其转化为AFQ1是一种解毒过程。
    \item \textbf{AFB1 $\rightarrow$ AFH1（还原反应）：}AFQ1的还原产物，黄曲霉毒素醇转化为AFH1可能也是一种解毒过程。
    \item \textbf{AFB1 $\rightarrow$ AFB1-8,9-环氧化物（环氧化反应）：}由CYP450催化的环氧化反应生成，是AFB1的终致癌物。可与DNA结合形成AFB1-N$^7$-鸟嘌呤加合物，导致DNA损伤和突变。% [补充自往届笔记] AFB1-2,3-环氧化物，一部分为解毒过程，一部分发挥毒性、致癌性和致突变性
\end{enumerate}

AF有很强的急性毒性，也有明显的慢性毒性与致癌性。AF对肝脏有特殊亲和性（嗜肝毒），具有较强的肝脏毒性并有致癌作用。
\end{keybox}

\subsubsection{黄曲霉毒素的毒性}

\begin{keybox}
\textbf{黄曲霉毒素的毒性表现（三方面*）：}
\begin{enumerate}[leftmargin=*]
    \item \textbf{急性毒性：}
    \begin{itemize}[leftmargin=*]
        \item 主要靶器官：\textbf{肝脏}——引起肝细胞坏死、胆管增生、肝出血（急性中毒性肝炎）。
        \item 中毒临床表现以黄疸为主，出现发热、呕吐和食欲缺乏，严重者出现腹水、下肢水肿、肝脾大及肝硬化。
        \item AFB1经口LD$_{50}$：鸭雏 0.24$\sim$0.36 mg/kg（最敏感物种之一），大鼠 5.5$\sim$17.9 mg/kg。
    \end{itemize}
    \item \textbf{慢性毒性：}
    \begin{itemize}[leftmargin=*]
        \item 主要表现为生长发育障碍，肝脏出现亚急性或慢性损伤，肝功能降低，肝实质细胞变性、坏死，形成再生结节。
        \item 生长发育迟缓，影响蛋白质合成和营养利用。
        \item 免疫功能下降，抑制细胞免疫和体液免疫。
    \end{itemize}
    \item \textbf{致癌性（已知最强的天然致癌物！）：}
    \begin{itemize}[leftmargin=*]
        \item 主要是肝癌，诱发肝细胞型肝癌。
        \item 诱癌强度是二甲基亚硝胺的\textbf{75倍}。
        \item 可诱发多种实验动物（大鼠、小鼠、鸭、鱼、猴等）的多器官癌症，无动物种属对AF致癌完全抵抗。
        \item 灵长类动物亦可诱癌——猴经AFB1处理可发生肝癌。
        \item \textbf{与人类肝癌密切相关：}膳食AF暴露量与原发性肝癌发病率呈正相关。\textbf{尤为重要的是AF与HBV（乙肝病毒）的协同致癌作用！}HBV感染者暴露于AF时，肝癌风险急剧增高。
        \item IARC分类：AFB1为\textbf{Group 1（确认人类致癌物）}。
    \end{itemize}
\end{enumerate}
\end{keybox}

\subsubsection{黄曲霉毒素的预防措施}

\begin{keybox}
\textbf{黄曲霉毒素的预防措施*：}

\textbf{（1）食物防霉（最根本的措施）：}控制粮食水分含量在安全水分以下，注意低温保藏，注意通风。

\textbf{（2）去除毒素（7种方法）：}
\begin{enumerate}[leftmargin=*]
    \item \textbf{挑选霉粒法：}对花生、玉米去毒效果好；
    \item \textbf{碾轧加工法：}将受污染的大米加工成精米，可降低毒素含量；
    \item \textbf{加水搓洗法：}反复搓洗去除表面毒素；
    \item \textbf{植物油加碱去毒：}AF与NaOH反应，内酯结构中的酯键被水解破坏，形成香豆素钠盐。将碱炼后的水洗可去除毒素。本方法用于植物油去毒效果最好。
    \item \textbf{物理吸附法：}含毒素的植物油加入活性白陶土或活性炭等吸附剂，适用于液体食品（如植物油），效果较好，对固体食品效果不明显。
    \item \textbf{氨气处理法：}在18kg压力、72$\sim$82$^\circ$C状态下，氨气处理可使AF被去除98\%$\sim$100\%，但可使粮食中蛋白质含量增加，同时可能破坏维生素。
    \item \textbf{日光或紫外线照射法：}破坏AF的化学结构。
\end{enumerate}

\textbf{（3）制定食品中AF限量标准：}具体限量标准执行GB 2761。
\end{keybox}

\subsection{镰刀菌毒素}

\begin{defbox}
\textbf{镰刀菌毒素（Fusarium Toxins）：}由镰刀菌属（Fusarium spp.）产生的有毒次级代谢产物，根据其化学结构可分为单端孢霉烯族化合物、玉米赤霉烯酮、伏马菌素和丁烯酸内酯等。

\textbf{1. 单端孢霉烯族化合物：}
\begin{itemize}[leftmargin=*]
    \item 主要有T-2毒素、二乙酸藨草镰刀菌烯醇、雪腐镰刀菌烯醇和脱氧雪腐镰刀菌烯醇。T-2毒素是三线镰刀菌和拟枝孢镰刀菌的代谢产物，是食物中毒性白细胞缺乏症（ATA）的病原物质。
    \item 化学性质：化学性质稳定，在通常条件下不产生荧光，耐热，在中性和酸性环境中不被破坏。
    \item 毒性特点与食品卫生学意义：具有强细胞毒性、免疫抑制及致畸作用，对人体有潜在的致癌性，急性毒性强，可引起动物呕吐。主要污染麦类、玉米及麦类制品。
\end{itemize}

\textbf{2. 玉米赤霉烯酮（Zearalenone, ZEN）：}
\begin{itemize}[leftmargin=*]
    \item \textbf{来源：}又称F-2毒素，由禾谷镰刀菌、黄色镰刀菌、雪腐镰刀菌等产生。
    \item \textbf{毒性特点：}具有雌激素样作用，可表现出对生殖系统的毒性作用。对动物主要污染玉米，其次为小麦、大麦、大米等粮食作物。
\end{itemize}

\textbf{3. 伏马菌素（Fumonisins）：}
\begin{itemize}[leftmargin=*]
    \item \textbf{来源：}主要由串珠镰刀菌产生，可分为伏马菌素B1（FB1）和伏马菌素B2（FB2）两大类。食品中以FB1污染为主，主要污染玉米及玉米制品。FB1通过干扰神经鞘磷脂和鞘氨醇的合成，影响鞘脂类代谢。
    \item \textbf{主要危害：}具有神经毒性作用，可引起马的脑白质软化症；对多种动物具有肾脏毒性；伏马菌素是促癌物，主要诱发原发性肝癌。
\end{itemize}
\end{defbox}

\subsection{其他重要霉菌毒素}

\begin{defbox}
主要霉菌毒素及其毒性（了解）：
\begin{itemize}[leftmargin=*]
    \item \textbf{赭曲霉毒素（Ochratoxin A, OTA）：}主要由赭曲霉和纯绿青霉产生，主要靶器官为肾脏，可引起巴尔干肾病。IARC Group 2B。
    \item \textbf{展青霉素（Patulin）：}扩展青霉产生，污染水果及制品（苹果汁、苹果酱最突出），具有神经毒性、致突变性和免疫抑制性。
    \item \textbf{杂色曲霉毒素：}可引起肝癌等。
\end{itemize}
\end{defbox}

% =====================================================================
%  第四节 食品腐败变质
% =====================================================================
\section{食品腐败变质}

\subsection{食品腐败变质的概念与原因}

\begin{keybox}
\textbf{食品腐败变质（food spoilage）：}指食品在以微生物为主的各种环境因素作用下，食品的成分被分解、破坏，失去或降低食用价值的一切变化。包括肉、鱼、禽、蛋的腐败，粮食的霉变、蔬菜水果的溃烂、油脂酸败、牛奶的变质等。

\textbf{食品腐败变质的原因（三类因素*）：}

\textbf{（1）微生物因素：}主要是细菌、霉菌和酵母菌。
\begin{itemize}[leftmargin=*]
    \item 分解蛋白质的微生物：主要是细菌，其次是霉菌和酵母菌；
    \item 分解碳水化合物的微生物：主要有细菌和酵母菌，能分解某些碳水化合物的细菌；
    \item 分解脂肪的微生物：有能产生脂肪酶，能分解脂肪的革兰阴性菌，食品中常见的有霉菌（如黄曲霉、腊叶芽枝霉等）和假丝酵母菌等；其中酵母菌主要是解脂假丝酵母，对糖类不发酵，分解脂肪和蛋白质的能力却很强。
\end{itemize}

\textbf{（2）食物本身的组成和性质：}
\begin{itemize}[leftmargin=*]
    \item 食品中的酶：食品经酶作用，引起食品组成成分的分解，加速食品的腐败变质；
    \item 食品的营养成分和水分：食品中含有丰富的营养成分，是微生物的良好培养基，富含蛋白质的动物性食品主要以蛋白质腐败为基本特征；碳水化合物含量高的食品，主要在细菌和酵母菌的作用下，以产酸发酵为基本特征；食品中Aw值越小，微生物越不易繁殖，食品越不易腐败变质；
    \item 食品的物理性质：食品pH高低也是制约微生物生长影响食品腐败变质的重要因素之一，一般在酸性或碱性环境中都不利于细菌生长。
\end{itemize}

\textbf{（3）环境因素：}温度、湿度、氧气、光照（紫外线）等。
\end{keybox}

\subsection{食品腐败变质过程中营养成分的变化*}

\begin{keybox}
\textbf{（1）食品中蛋白质的分解：}食品中的蛋白质在微生物的蛋白酶、肽链内切酶等作用下可分解形成氨基酸，氨基酸及胺等在低分子量下可通过脱羧基、脱氨基等作用，形成各种腐败产物。
\begin{itemize}[leftmargin=*]
    \item 细菌脱羧酶作用下：酪氨酸、组氨酸、赖氨酸可分别生成酪胺、组胺和尸胺，具有恶臭味。
    \item 微生物脱氨基酶作用下：生成丁酸、己酸等。
    \item 色氨酸脱羧基色胺、甲基色胺，具有粪臭味。
    \item 含硫氨基酸在脱硫醇酶作用下生成具有恶臭味的硫醇。
\end{itemize}

\textbf{（2）食品中脂肪的酸败（rancidity）：}指食品中脂肪的腐败变质现象。脂肪腐败的化学过程复杂，主要是经水解和氧化反应产生的分解产物。
\begin{itemize}[leftmargin=*]
    \item 中性脂肪分解为甘油和脂肪酸，后者可进一步氧化为低级脂肪酸、醛、酮等；
    \item 不饱和脂肪酸中双键被氧化形成过氧化物，进一步分解为醛、酮、酸，含不饱和脂肪酸越高的食品越易氧化。
    \item 酸败过程中形成的醛、酮和某些羧酸使酸败的脂肪带有特殊的刺激性臭味，即所谓的"哈喇"味（油脂气味）。含不饱和脂肪酸越高的食品越易被氧化。
    \item 脂肪分解的早期，酸败尚不明显，即出现酸度增高，脂肪酸酸度与醛酮含量的反应很敏感。
\end{itemize}

\textbf{酸败的鉴定：}
\begin{enumerate}[leftmargin=*]
    \item 油脂气味；
    \item 过氧化值上升——脂肪酸败的最早期指标；
    \item 酸价和羰基价反应阳性；
    \item 碘、醇等食用脂肪氧化。
\end{enumerate}

\textbf{（3）碳水化合物的分解：}碳水化合物在微生物和植物组织酶的作用下，可分解成单糖、醇、醛、酮、CO$_2$和水等。该过程的主要变化是食品酸度升高，产生酸、醇、醛等气味。
\end{keybox}

\subsection{食品腐败变质的鉴定指标}

\begin{keybox}
\textbf{食品腐败变质的鉴定指标（四类*）：}

\textbf{（1）感官鉴定（最敏感的方法）：}色、香、味、形，主要通过视觉、嗅觉、味觉和触觉来鉴定食品腐败变质。例如：食品腐败常出现不良气味、颜色出现褐色、绿色或失去光泽，组织软化或发黏等。

\textbf{（2）物理指标：}常用的有油脂比重、折光率、黏度等。

\textbf{（3）化学鉴定：}通过直接测定微生物代谢产物的腐败产物。
\begin{itemize}[leftmargin=*]
    \item \textbf{挥发性盐基总氮（Total Volatile Basic Nitrogen, TVBN）：}指蛋白质丰富食品的水浸液在弱碱性条件下与水蒸气一起蒸馏出来的总氮量，在此过程中蒸馏出来的氨和形成的胺的含氮物。TVBN与食品腐败变质程度之间有明确的对应关系，为鱼、肉类蛋白腐败鉴定的化学指标。
    \item \textbf{三甲胺：}主要用于测定鱼、虾等水产品的新鲜程度。
    \item \textbf{组胺：}细菌脱羧酶作用于组氨酸生成组胺，当鱼中组胺超过200 mg/100g就可引起人类过敏性食物中毒。
    \item \textbf{K值（K value）：}指ATP分解的低级产物肌苷（HxR）和次黄嘌呤（Hx）占ATP系列分解产物ATP$+$ADP$+$AMP$+$IMP$+$HxR$+$Hx的百分比。主要用于鉴定鱼类早期腐败。$K < 20\%$说明鱼体新鲜，$K \geq 40\%$说明鱼体开始有腐败迹象。
    \item \textbf{pH值：}一般腐败开始时微生物生长受限，氨基酸脱羧后生成胺类物质，pH值出现V字形变动。
    \item \textbf{过氧化值与酸价：}过氧化值是脂肪酸败鉴定早期指标，酸价反映脂肪酸败的程度。
\end{itemize}

\textbf{（4）微生物检验：}常用指标为菌落总数、大肠菌群和细菌菌相。
\end{keybox}

\subsection{食品腐败变质的食品卫生学意义及处理原则}

\begin{keybox}
\textbf{食品卫生学意义：}
食品腐败变质可使感官性状发生改变，造成食品营养成分分解，营养价值严重降低；进而，腐败变质食品必然受到致病微生物等严重污染，并且可能带有致病菌和产毒霉菌，食用后易引起不良反应和食物中毒。

\textbf{腐败变质食品的处理原则：}
\begin{enumerate}[leftmargin=*]
    \item \textbf{严重腐败变质：}销毁或改作工业用；
    \item \textbf{轻度腐败变质的食品：}经适当的加工处理，将变质的主要指标降至安全限量以下，可以限期食用；
    \item \textbf{已经变质的食品：}禁止食用；
    \item \textbf{局部变质食品：}挑选去除变质部分，利用其他完好部位。
\end{enumerate}
\end{keybox}

% =====================================================================
%  第五节 食品保藏方法
% =====================================================================
\section{食品保藏方法}

\begin{defbox}
\textbf{食品保藏（Food Preservation）：}指为了防止食品腐败变质，延长食品可食用的期限，对食品中的微生物进行杀灭或抑制其生长繁殖，并对食品进行加工处理，通过改变食品的温度、水分、氢离子浓度、渗透压以及采用其他抑菌杀菌措施，以达到防止食品腐败变质的目的。
\end{defbox}

\subsection{化学保藏法}

\begin{keybox}
\textbf{化学保藏法：}
\begin{enumerate}[leftmargin=*]
    \item \textbf{盐腌法：}利用高渗透压抑制微生物生长，食盐浓度达10\%即可抑制大多数细菌，食品含糖量达到60\%$\sim$65\%即可。
    \item \textbf{酸渍法：}多数微生物在pH 4.5以下不能生长繁殖，故可利用控制氢离子浓度来抑制微生物生长，多用于各种蔬菜，如泡菜和酸菜等。
    \item \textbf{防腐剂：}在食品中可添加符合食品添加剂标准的防腐剂，来抑制或杀灭食品中引起腐败变质的微生物，如苯甲酸、山梨酸等。抗氧化剂用于防止油脂酸败。
\end{enumerate}
\end{keybox}

\subsection{低温保藏法}

\begin{keybox}
\textbf{低温保藏法：}
\begin{itemize}[leftmargin=*]
    \item \textbf{冷藏：}指在不冻结状态下的低温储藏，温度一般设定在$-1\sim10^\circ$C范围内。其原理是低温（10$^\circ$C以下）可抑制微生物的繁殖，降低食品中酶的活性，可有效延缓食品的变质。
    \item \textbf{冷冻储藏：}指在$-18^\circ$C以下保藏。在此温度下几乎所有的微生物不再繁殖，冷冻食品可以较长期地保藏。在工艺上，速冻（急冻）、缓慢解冻（缓化）有利于保持食品的品质。
\end{itemize}

\textbf{冰晶生成带：}
\begin{itemize}[leftmargin=*]
    \item $-1\sim-5^\circ$C为冰晶生成带，可导致食品中水结冰率为85\%，不利于保持食品品质。
    \item $-8\sim-120^\circ$C为冻结带，食品中水结冰率为90\%。速冻（急冻），迅速通过冰晶生成带，形成的冰晶数量多、体积小而不损伤食品组织细胞。
\end{itemize}
\end{keybox}

\subsection{高温杀菌保藏法}

\begin{keybox}
\textbf{高温杀菌保藏法：}

\textbf{D值（Decimal Reduction Time）*\difficulty：}在一定温度下，能杀死食品中某种细菌90\%的时间（min）。主要用于比较细菌的死亡速度。如金黄色葡萄球菌 D$_{63}$ = 7；肉毒梭菌芽孢 D$_{121}$ = 0.1$\sim$0.23。

\textbf{常用高温杀菌方法：}
\begin{enumerate}[leftmargin=*]
    \item \textbf{常压杀菌法：}温度控制在100$^\circ$C及以下，杀灭致病菌和繁殖型微生物，适用于液态食物。常用巴氏杀菌法处理食品（如牛乳、啤酒、酱油、葡萄酒等）。以牛乳为例，低温巴氏杀菌法温度63$^\circ$C，30分钟；高温短时巴氏杀菌法温度72$^\circ$C，15秒。
    \item \textbf{高压杀菌法：}温度100$\sim$121$^\circ$C（绝对压力为0.2 MPa），适用于肉类制品、罐头食品，可杀灭繁殖型和芽孢型细菌。
    \item \textbf{超高温瞬时杀菌（UHT）：}在封闭的系统中加热到120$^\circ$C以上，持续几秒钟后迅速冷却至常温的一种杀菌方法，用于牛乳。
    \item \textbf{微波杀菌：}目前多用915 MHz和2450 MHz两个频率。915 MHz可获得较大穿透深度，适用于加热含水量高、厚度或体积较大的食品；2450 MHz适用于含水量低的食品。

    特点：快速、节能、对食品品质影响小，能保留更多活性物质和营养成分。
\end{enumerate}
\end{keybox}

\subsection{食品的干燥脱水保藏}

\begin{keybox}
\textbf{食品脱水保藏的机制：}降低食品水分至15\%以下或Aw值在0.00$\sim$0.60之间，可抑制各种微生物生长，使食品在常温下长期保藏。

\textbf{常用方法：}日晒、阴干、喷雾干燥、减压蒸发、冷冻干燥等。

\textbf{酶性褐变：}是酚酶催化酚类物质形成醌及其聚合物的结果。

\textbf{非酶褐变（美拉德反应/Maillard反应）：}由蛋白质、肽和氨基酸等的氨基与糖及脂肪氧化的醛、酮等羰基所发生的反应，使食物呈现褐色和焦香味。影响因素：加热与给氧促进褐变；水分活性（Aw 0.6最适于褐变）；SO$_2$可与羰基结合而阻止褐变。
\end{keybox}

\subsection{食品辐照保藏}

\begin{keybox}
\textbf{辐照保藏：}

\textbf{辐照灭菌（radappertization）：}指使用波长小于200 nm的电磁波辐照处理食物，达到商业无菌要求的一种处理方式。用高剂量（10$\sim$15/50 kGy）杀灭食品中所有微生物。辐照可使食品中的微生物数量减少到零或有限个数，且不会因辐照而产生放射性。

\textbf{辐照消毒（radicidation）：}使用适当的中等剂量（5$\sim$10 kGy），杀死无芽孢致病菌，主要要求杀灭沙门菌以及各种微生物，辐照后的食品3天内不会发虫、霉变。

\textbf{辐照防腐（radurization）：}使用剂量1$\sim$5 kGy，主要杀灭腐败菌，延长新鲜食品的后熟期及保藏期。主要用于抑制发芽、延长后熟期、杀虫等。
\end{keybox}

% =====================================================================
%  第六节 农药和兽药残留
% =====================================================================
\section{农药和兽药残留}

\subsection{农药的定义与分类}

\begin{keybox}
\textbf{农药（pesticides）的定义：}指用于预防、控制危害农业、林业的病、虫、草、鼠和其他有害生物以及有目的地调节植物、昆虫生长的化学合成或者来源于生物、其他天然物质的一种物质或者几种物质的混合物及其制剂。

\textbf{农药分类（按五种方式）：}
\begin{enumerate}[leftmargin=*]
    \item \textbf{按用途分：}杀虫剂、杀鼠剂、杀菌剂、杀线虫剂、杀螺剂、除草剂、植物生长调节剂等。
    \item \textbf{按化学组成及结构分：}有机磷类、有机氯类、氨基甲酸酯类、拟除虫菊酯类、有机砷类、有机汞类等。
    \item \textbf{按成分和来源分：}无机类、有机类和生物农药。
    \item \textbf{按急性毒性大小分：}剧毒、高毒、中等毒和低毒农药。
    \item \textbf{按残留特性分：}高残留、中等残留和低残留农药。
\end{enumerate}
\end{keybox}

\subsection{兽药的定义}

\begin{defbox}
\textbf{兽药（veterinary drugs）：}指用于预防、治疗、诊断动物疾病或者有目的地调节动物生理机能的物质，包括药物饲料添加剂。主要包括血清制品、疫苗、诊断制品、微生态制品等。
\end{defbox}

\subsection{农药残留与最大残留限量}

\begin{keybox}
\textbf{农药残留（pesticide residues）：}指任何由于使用农药而在食品、农产品和动物饲料中出现的特定物质，包括被认为具有毒理学意义的农药衍生物，如农药转化物、代谢物、反应产物及杂质等。

\textbf{最大残留限量（Maximum Residue Limits, MRLs）：}指在食品或农产品内部或表面，按照良好农业生产规范（GAP）使用农药后，允许农药残留的最高浓度，以每千克食品或农产品中农药残留的毫克数（mg/kg）表示。
\end{keybox}

\subsection{农药污染食品的途径与预防}

\begin{defbox}
\textbf{农药污染食品的途径：}
\begin{enumerate}[leftmargin=*]
    \item 直接污染（施药）——最主要的污染途径；
    \item 农作物从污染环境中吸收；
    \item 通过食物链和生物富集作用污染食品；
    \item 储运过程中的污染；
    \item 其他来源的污染——食品加工、储藏、运输、销售过程中的污染。
\end{enumerate}

\textbf{预防控制措施：}
\begin{enumerate}[leftmargin=*]
    \item 加强农药管理；
    \item 安全合理使用农药；
    \item 制定和执行食品中农药残留限量标准；
    \item 研制和推广高效低毒低残留新品种。
\end{enumerate}
\end{defbox}

% [补充自往届笔记] — 农药分类毒性详解
\begin{tipbox}
\textbf{往届笔记补充——各类农药毒性特点：}

\textbf{有机氯农药（OCPs）：}化学性质极其稳定，高残留农药，脂溶性强，主要蓄积于脂肪组织，生物富集作用强。急性毒性为神经系统和肝脏毒性；慢性毒性影响神经系统、肝脏、血液系统、生殖系统，具有雌激素活性，可经胎盘和乳汁传递，有致畸性。已禁用。

\textbf{有机磷农药（OPPs）：}目前使用范围最广、用量最大的农药。怕光、怕热、怕水，易酶解，半衰期短，残留量低。急性毒性为抑制胆碱酯酶活性，中毒表现为神经功能紊乱，也可出现迟发性神经毒性。慢性毒性影响神经系统、血液系统、视觉。

\textbf{氨基甲酸酯类农药：}毒性低，易被分解，残留量低。毒性为胆碱酯酶抑制剂，但抑制强度较弱且可逆，故毒性比OPPs小，无迟发性神经毒性。可使染色体断裂，有致癌、致畸和致突变作用。

\textbf{拟除虫菊酯类农药：}毒性低，半衰期短，残留量低，但使用后易产生抗药性。急性毒性为神经系统毒性，可引起肌肉痉挛；慢性毒性有致癌性。皮肤接触有刺激和致敏作用，有致突变性和致畸性。
\end{tipbox}

% =====================================================================
%  第七节 $N$-亚硝基化合物污染
% =====================================================================
\section{$N$-亚硝基化合物污染}

\subsection{前体物与合成条件}

\begin{keybox}
\textbf{$N$-亚硝基化合物的前体物*：}
\begin{enumerate}[leftmargin=*]
    \item \textbf{硝酸盐（NO$_3^-$）和亚硝酸盐（NO$_2^-$）：}广泛存在于自然环境（土壤、水体）和食品中；
    \item \textbf{胺类物质：}广泛存在于环境和食品中（蛋白质分解产物）。
\end{enumerate}

\textbf{$N$-亚硝基化合物的体内合成条件*：}
\begin{itemize}[leftmargin=*]
    \item \textbf{pH $<$ 3的酸性环境}是该反应的最适条件。
    \item \textbf{胃是体内合成$N$-亚硝基化合物的最主要场所！}胃液pH约1$\sim$2，且硝酸盐、亚硝酸盐和胺类可随食物同时进入胃内。胃内还有Cl$^-$、SCN$^-$等催化剂。
\end{itemize}
\end{keybox}

\subsection{致癌作用特点}

\begin{keybox}
\textbf{$N$-亚硝基化合物的致癌作用特点*：}
\begin{enumerate}[leftmargin=*]
    \item \textbf{多动物种属诱癌：}对多种实验动物均有致癌性，至今未发现任何动物种属对$N$-亚硝基化合物的致癌作用有抵抗力；
    \item \textbf{多器官/组织靶向性：}对称性亚硝胺主要诱发肝癌，不对称性亚硝胺主要诱发食管癌，可通过血脑屏障和胎盘屏障，引起中枢神经系统肿瘤和经胎盘致癌；
    \item \textbf{多途径均诱癌：}经口、吸入、皮下和肌肉注射、皮肤接触等多种途径均可诱癌；
    \item \textbf{不同接触剂量均可致癌：}一次大剂量或长期小剂量给药均可致癌，有明显的剂量-效应关系；且各剂量组合有相加作用；
    \item \textbf{可以通过胎盘使子代致癌（经胎盘致癌作用）：}怀孕动物摄入后，可通过胎盘影响胎儿，出生后发生肿瘤。
\end{enumerate}
\end{keybox}

\subsection{预防措施}

\begin{keybox}
\textbf{$N$-亚硝基化合物危害的预防*：}
\begin{enumerate}[leftmargin=*]
    \item 防止食品被微生物污染；
    \item 改进食品加工工艺；
    \item 使用钼肥（Mo肥料）——钼是硝酸还原酶的组分，施钼可降低蔬菜中硝酸盐的积累；
    \item 阻断亚硝化反应——多食新鲜蔬菜和水果，增加维生素C、维生素E、酚类、黄酮类化合物的摄入（VC将NO$_2^-$还原为NO，从而阻止亚硝化反应）；
    \item 制定食品中限量标准并加强监测：水产品中$N$-亚硝基化合物限量4 $\mu$g/kg，肉制品中$N$-亚硝基化合物限量3 $\mu$g/kg。% [补充自往届笔记]
\end{enumerate}
\end{keybox}

% =====================================================================
%  第八节 多环芳烃（PAH）污染
% =====================================================================
\section{多环芳烃（PAH）污染}

\begin{keybox}
\textbf{多环芳烃（Polycyclic Aromatic Hydrocarbons, PAH）：}由两个或两个以上苯环融合而成的化合物。\textbf{苯并(a)芘（Benzo(a)pyrene, B(a)P）}是最重要、最具有代表性的PAH。

\textbf{B(a)P的特性：}
\begin{itemize}[leftmargin=*]
    \item 主要蓄积于脂肪组织；
    \item 致胃癌、致突变；
    \item 预防措施：防止污染；去除毒素（活性炭吸附）；制定食品中限量标准（熏烤肉 $\leq$ 5 $\mu$g/kg，植物油 $\leq$ 10 $\mu$g/kg）。% [补充自往届笔记]
\end{itemize}
\end{keybox}

\subsection{污染来源}

\begin{keybox}
\textbf{B(a)P污染食品的主要途径：}
\begin{enumerate}[leftmargin=*]
    \item \textbf{烘烤、烟熏过程中直接污染——最主要的来源！}熏烟中含有大量PAH，烟熏时间越长、温度越高，B(a)P含量越高。炭火烤肉/鱼，油脂滴落在炭火上，热解聚合生成PAH，随烟雾附在食品表面。
    \item 食品成分在高温下的热解/热聚合：脂肪在 $>$ 400$^\circ$C 可生成PAH。
    \item 农作物从环境中吸收：空气沉降的PAH被作物叶片和果实表面吸附。
    \item 食品加工过程中接触污染：如用沥青铺路晒粮（严重违规）。
\end{enumerate}
\end{keybox}

% =====================================================================
%  第九节 杂环胺（Heterocyclic Amines）污染
% =====================================================================
\section{杂环胺（Heterocyclic Amines, HCA）污染}

\begin{keybox}
\textbf{杂环胺污染食品的来源及影响因素：}
\begin{enumerate}[leftmargin=*]
    \item \textbf{烹调方式：}加热温度和时间——温度越高、时间越长、水分含量越少，杂环胺生成越多；
    \item \textbf{食物成分：}蛋白质含量——蛋白质含量越高，杂环胺生成越多，美拉德反应是前提。
\end{enumerate}

\textbf{致癌性：}诱发肝癌、致突变（形成DNA加合物）。

\textbf{预防措施：}
\begin{enumerate}[leftmargin=*]
    \item 改变不良烹调方式和饮食习惯；
    \item 多吃蔬菜水果——蔬菜、水果中的多酚、类黄酮等植物化学物可抑制杂环胺的致突变作用；
    \item 加强监测。
\end{enumerate}
\end{keybox}

% =====================================================================
%  第十节 氯丙醇、丙烯酰胺
% =====================================================================
\section{氯丙醇、丙烯酰胺}

\subsection{氯丙醇（Chloropropanols）}

\begin{defbox}
\begin{itemize}[leftmargin=*]
    \item \textbf{首次发现：}在酸水解植物蛋白液（HVP）中发现，用于酱油等调味品中。
    \item \textbf{代表性物质：}3-氯-1,2-丙二醇（3-MCPD）含量最高、最常见。
    \item \textbf{毒性：}3-MCPD具肾毒性和雄性生殖毒性。
    \item \textbf{存在形式：}游离态3-MCPD（酸水解HVP）以及脂肪酸酯形式（即3-MCPD酯，广泛存在于精炼油脂中）。
\end{itemize}
\end{defbox}

\subsection{丙烯酰胺（Acrylamide）}

\begin{defbox}
\begin{itemize}[leftmargin=*]
    \item \textbf{生成条件：}富含淀粉的食品在 $>$ 120$^\circ$C 高温加工（油炸、烘烤）过程中，由天冬酰胺（asparagine）和还原糖（葡萄糖、果糖）通过Maillard反应生成。
    \item \textbf{高含量食品：}薯条（French fries）、薯片（potato chips）、方便面、油条、速溶咖啡。土豆制品的丙烯酰胺含量约为谷物制品的\textbf{4倍}。
    \item \textbf{毒性和致癌性：}神经毒性（职业暴露），啮齿动物多器官致癌。IARC Group 2A（很可能对人类致癌）。
    \item \textbf{预防：}改进加工工艺（降低油炸温度和时间、避免过度金黄焦脆），原料选用低还原糖的马铃薯品种。
\end{itemize}
\end{defbox}

% =====================================================================
%  第十一节 食品容器、包装材料的卫生问题
% =====================================================================
\section{食品容器、包装材料的卫生问题}

\subsection{概念、范围与分类}

\begin{keybox}
\textbf{食品容器、包装材料的概念：}指在食品生产、加工、储存、运输、销售过程中，直接或间接接触食品的各种材料、制品及辅助材料。

\textbf{范围：}包括从原料到成品的全过程中所有接触食品的器具、管道、设备和包装物。

\textbf{分类及其基本卫生问题（掌握）：}
\begin{enumerate}[leftmargin=*]
    \item \textbf{塑料制品：}PET（聚对苯二甲酸乙二醇酯）、PC（聚碳酸酯）、PE（聚乙烯）、PP（聚丙烯）、PS（聚苯乙烯）、三聚氰胺-甲醛树脂（密胺）、脲醛树脂。基本卫生问题：单体残留（氯乙烯、苯乙烯、双酚A）、增塑剂迁移（邻苯二甲酸酯类——内分泌干扰物）、热稳定性不足。
    \item \textbf{橡胶制品：}天然橡胶、合成橡胶（奶嘴、密封圈）。基本卫生问题：硫化促进剂、抗氧化剂和填充剂的迁移。
    \item \textbf{纸、竹、木：}传统材料。基本卫生问题：漂白剂残留、荧光增白剂、印刷油墨迁移。
    \item \textbf{金属、搪瓷、陶瓷、玻璃：}基本卫生问题：金属离子（铅、镉、铬、镍、铝）的溶出，釉彩中的铅镉迁移。
\end{enumerate}
\end{keybox}

\subsection{影响食品受容器污染的因素}

\begin{defbox}
\begin{enumerate}[leftmargin=*]
    \item \textbf{接触时间和温度：}接触时间越长、温度越高，迁移量越大。热水或热食会显著加速塑料中有害物质溶出。
    \item \textbf{食品的性质：}酸性食品（如果汁、醋）、含酒精食品和油脂类食品对容器材料中有害物质的萃取能力更强。pH越低、脂肪含量越高，迁移量越大。
    \item \textbf{容器材料的性质：}材料的化学组成、加工工艺、添加剂种类和用量决定其卫生安全性。再生塑料的卫生风险较高。
    \item \textbf{接触面积与食品体积之比：}接触面积越大、食品体积越小，单位食品中迁移物的浓度越高。
    \item \textbf{材料的老化程度：}长期使用后材料表面出现裂纹和磨损，增加有害物质的释放。
\end{enumerate}
\end{defbox}

\subsection{主要卫生问题}

\begin{enumerate}[leftmargin=*]
    \item \textbf{有害物质向食品迁移（Migration）——最核心的卫生问题！}
    \begin{itemize}[leftmargin=*]
        \item 塑料中的增塑剂（如邻苯二甲酸酯类，phthalates——具有雌激素样活性，属于内分泌干扰物）
        \item 塑料中的稳定剂、抗氧化剂
        \item 密胺制品（仿瓷餐具）：游离甲醛毒性
        \item 铝制品：盛装酸性或碱性食品可溶解氧化铝膜，铝溶出（铝摄入过量与阿尔茨海默病的关联虽未最终定论，但已引起广泛关注）
    \end{itemize}
    \item 包装材料的回收和再生问题：回收塑料中的不明污染物转移入食品。
    \item 印刷油墨：油墨中的有害物质（如苯系物）可能迁移至食品接触面。
    \item 纳米包装材料：新型材料的安全性评价体系尚待完善。
\end{enumerate}

\subsection{相关卫生标准与指标}

\begin{defbox}
我国已制定GB 4806系列《食品安全国家标准 食品接触材料及制品》标准体系，涵盖：
\begin{itemize}[leftmargin=*]
    \item 各类食品接触材料的通用安全要求
    \item 特定迁移量（SML）和总迁移量（OML）限量
    \item 食品接触材料中添加剂的使用标准（GB 9685）
    \item 各类材料的具体卫生标准（如GB 4806.1--GB 4806.12）
\end{itemize}
\end{defbox}

% =====================================================================
%  复习题
% =====================================================================
\reviewcontent{
\section{复习题}

\begin{enumerate}[leftmargin=*, itemsep=8pt]
    \item \textbf{【名词解释】}食品污染；食品污染物；水分活性（Aw）；菌落总数；大肠菌群；MPN法；霉菌毒素；细菌菌相；霉菌菌相
    \item \textbf{【名词解释】}食品腐败变质；TVBN；K值；D值；食品保藏
    \item \textbf{【名词解释】}农药；农药残留；MRLs；兽药
    \item \textbf{【简答·重点】}简述食品污染的三大分类（生物性、化学性、物理性）及其具体内容。
    \item \textbf{【简答·重点】}简述水分活性（Aw）的定义公式，并说明Aw与微生物生长的关系（含各类微生物的最低Aw阈值及细菌芽孢形成的特殊要求）。
    \item \textbf{【简答·重点】}简述常见食品细菌（假单胞菌属、微球菌属、芽孢杆菌属、肠杆菌科）及其与腐败变质的关系。
    \item \textbf{【简答·重点】}简述细菌菌相的定义及食品卫生学意义。
    \item \textbf{【简答·重点】}简述菌落总数和大肠菌群的食品卫生学意义，并说明MPN法的适用条件及原理。
    \item \textbf{【简答·重点】}简述粪便污染指示菌应具备的条件。
    \item \textbf{【简答·重点】}简述霉菌产毒的条件（基质、水分、湿度、温度、通风）和特点（四项），以及霉菌毒素中毒的共同特点。
    \item \textbf{【简答·重点】}简述霉菌菌相及其食品卫生学意义。
    \item \textbf{【简答·重点】}简述黄曲霉毒素的化学结构、理化性质、毒性排序、产毒条件、体内代谢过程及毒性表现（急性、慢性、致癌性）。为何说黄曲霉毒素是已知最强的天然致癌物？
    \item \textbf{【简答·重点】}详述黄曲霉毒素预防措施，包括食物防霉和7种去毒方法。
    \item \textbf{【简答·重点】}简述镰刀菌毒素三大类（单端孢霉烯族化合物、玉米赤霉烯酮、伏马菌素）的来源及毒性特点。
    \item \textbf{【简答·重点】}简述食品腐败变质的概念、三大原因及各类营养素的腐败产物特征（蛋白质、脂肪、碳水化合物）。
    \item \textbf{【简答·重点】}详述食品腐败变质的鉴定指标（感官、物理、化学、微生物），含TVBN、三甲胺、组胺、K值、pH、过氧化值/酸价的具体卫生学意义。
    \item \textbf{【简答·重点】}简述腐败变质食品的处理原则（四种情况）。
    \item \textbf{【简答·重点】}详述食品保藏的各类方法和关键参数：化学保藏（盐腌浓度、酸渍pH）、低温保藏（冷藏与冷冻的温度范围、冰晶生成带）、高温杀菌（常压/高压/UHT参数、D值及实例、微波频率）、干燥脱水保藏（水分和Aw范围、褐变类型）、辐照保藏（三种剂量分级及用途）。
    \item \textbf{【简答】}简述农药的五大分类方式及各类农药的毒性特点。
    \item \textbf{【简答·重点】}简述$N$-亚硝基化合物的前体物、体内合成条件、致癌作用特点及至少四项预防措施。
    \item \textbf{【简答】}简述B(a)P污染食品的主要途径、蓄积特点和防控措施。
    \item \textbf{【简答·重点】}简述杂环胺污染食品的来源及影响因素、致癌性及预防措施。
    \item \textbf{【简答】}简述丙烯酰胺的生成条件、高含量食品及预防措施。
    \item \textbf{【简答】}简述氯丙醇的来源、代表物质及毒性。
    \item \textbf{【简答·重点】}简述食品容器包装材料的概念、分类及各类的基本卫生问题。
    \item \textbf{【简答】}简述影响食品受容器污染的因素（五项）。
    \item \textbf{【综合论述】}试从食品污染的三大来源出发，综合分析我国食品安全面临的主要化学性污染问题及其综合防控策略。
    \item \textbf{【综合论述】}某地区疑似发生黄曲霉毒素污染事件，请设计一份食品安全调查与防控方案。
\end{enumerate}

\subsection*{参考答案要点}

\begin{enumerate}[leftmargin=*, itemsep=6pt]
    \item \textbf{食品污染：}食品中含有对健康造成危害或影响食用及商品价值的生物性、化学性及物理性物质。\textbf{食品污染物：}在种植、养殖、生产、加工、储存、运输、销售、烹饪、食用等环节进入食品的有害因素。\textbf{Aw：}$\text{Aw}=P/P_0$，同温度下食品水分蒸汽压与纯水蒸汽压之比（0$\sim$1），反映可被微生物利用的实际含水量。\textbf{菌落总数：}单位质量/容积/表面积检样在规定条件下培养后生成的细菌菌落总数（CFU）。\textbf{大肠菌群：}能发酵乳糖、产酸产气、需氧/兼性厌氧的革兰阴性无芽孢杆菌，包含埃希菌属、柠檬酸杆菌属、肠杆菌属和克雷伯菌属。\textbf{MPN法：}最可能数法，基于泊松分布的间接计数法，用于大肠菌群数量低时的统计。\textbf{霉菌毒素：}霉菌污染食品后产生的有毒代谢产物。\textbf{细菌菌相：}共存于食品中的细菌种类及其相对数量构成，优势菌种决定腐败特征。\textbf{霉菌菌相：}食品中污染的真菌种类及其相对数量构成，田野霉占优势示粮食新鲜，毛霉/根霉/木霉检出示已开始霉变。

    \item \textbf{食品腐败变质：}在以微生物为主的各种环境因素作用下食品成分被分解、破坏，失去或降低食用价值的一切变化。\textbf{TVBN：}弱碱条件下水蒸气蒸馏出的总氮量，鱼肉类蛋白腐败鉴定的化学指标。\textbf{K值：}ATP分解的HxR$+$Hx占ATP系列分解产物的百分比，$K<20\%$新鲜，$K\geq40\%$开始腐败。\textbf{D值：}在一定温度下杀死食品中某种细菌90\%的时间（min），如金黄色葡萄球菌D$_{63}$=7，肉毒梭菌芽孢D$_{121}$=0.1$\sim$0.23。\textbf{食品保藏：}为防止食品腐败变质，对微生物进行杀灭或抑制其生长繁殖，通过改变温度、水分、氢离子浓度、渗透压及其他抑菌杀菌措施，延长食品可食用期限。

    \item \textbf{农药：}用于预防控制农林病、虫、草、鼠等有害生物及调节植物昆虫生长的化学合成或生物源物质及其制剂。\textbf{农药残留：}因使用农药在食品中出现的农药衍生物，包括转化物、代谢物、反应产物及杂质。\textbf{MRLs：}按GAP使用农药后在食品中允许的农药残留最高浓度（mg/kg）。\textbf{兽药：}用于预防、治疗、诊断动物疾病或调节动物生理机能的物质，包括药物饲料添加剂。

    \item \textbf{三大分类：}（1）生物性污染：微生物（细菌及毒素、霉菌及毒素、病毒）、寄生虫和昆虫，细菌及其毒素最常见最重要；（2）化学性污染：农药兽药残留、工业三废中汞镉铅砷等有害金属及有机污染物、食品接触材料迁移物、滥用食品添加剂、加工储存中产生的$N$-亚硝基化合物/PAH/杂环胺/丙烯酰胺、掺假制假物质（苏丹红、三聚氰胺）；（3）物理性污染：杂物污染（草籽、灰尘等）和放射性污染（核爆炸、核废料、核事故，水-水生生物-动植物-人转移）。

    \item \textbf{Aw与微生物生长关系：}Aw $= P/P_0$。（1）每种微生物有最低Aw要求，低于则生长受抑；（2）Aw $<$ 0.60绝大多数微生物无法生长；（3）细菌需Aw $>$ 0.9，酵母菌 $>$ 0.87，霉菌 $>$ 0.8；（4）同属不同种Aw要求不同；（5）细菌形成芽孢时需要比生长时更高的Aw值。

    \item \textbf{常见食品细菌：}假单胞菌属（需氧无芽孢杆菌，冷藏食品腐败主要细菌）；微球菌属（需氧无芽孢杆菌，常见于水、蔬菜腐败）；微球菌和葡萄球菌属（革兰阳性，产过氧化氢酶、色素，分解碳水化合物）；芽孢杆菌属与梭状芽孢杆菌属（革兰阳性芽孢杆菌，罐头食品常见腐败菌）；肠杆菌科（需氧无芽孢杆菌，水产、肉、蛋腐败有关，含大肠埃希菌——粪便污染指示菌，变形杆菌属——腐败菌代表）。

    \item \textbf{细菌菌相：}食品中细菌种类及相对数量构成，优势菌种（属、株）决定腐败特征。卫生学意义：（1）影响腐败变质特征，使食用价值降低甚至不能食用；（2）影响腐败变质速度与程度，腐败过程中可出现霉菌繁殖并产毒。

    \item \textbf{菌落总数：}（1）食品细菌污染程度及卫生状态标志；（2）预测食品耐储性——细菌越多腐败越快。\textbf{大肠菌群：}（1）粪便污染指示菌——典型大肠杆菌=近期污染，非典型=陈旧污染；（2）肠道致病菌指示菌——同源同存活时间；（3）检测加工者/设备/存储条件污染情况。\textbf{MPN法：}大肠菌群数量低时采用，基于泊松分布的间接计数法，GB/T 4789.3-2016，单位MPN/g或MPN/mL；数量高时用CFU平板计数法。

    \item \textbf{粪便污染指示菌条件：}（1）仅来自肠道；（2）肠道中数量多、易检出；（3）外环境能生存一定期间；（4）检验方法敏感简易。\textbf{肠球菌（粪链球菌）：}反映低温类食品粪便污染，对低温抵抗力比大肠菌群强。

    \item \textbf{霉菌产毒条件：}基质（营养丰富，天然食品更易繁殖）、水分（粮食17\%$\sim$18\%最适，Aw $<$ 0.7一般不能生长）、湿度（RH $<$ 70\%不能产毒）、温度（最适25$\sim$30$^\circ$C，0$^\circ$C以下或30$^\circ$C以上产毒能力下降）、通风（大部分需氧，毛霉等可厌氧，耐受高浓度CO$_2$）。\textbf{产毒特点：}（1）只限于少数产毒菌株，不同菌株产毒能力不同；（2）产毒能力可变可易变；（3）产毒无严格专一性；（4）产毒需要一定时间。\textbf{霉菌毒素中毒特点：}毒素稳定、经天然食品中毒、实质器官受损为主、无传染性和免疫性、季节性和地区性。

    \item \textbf{霉菌菌相：}田野霉（气连孢霉等）占优势=粮食新鲜；曲霉和青霉检出=储藏条件可能不良；毛霉、根霉、木霉检出=已开始霉变。

    \item \textbf{AF化学结构：}二呋喃氧杂萘邻酮。\textbf{毒性排序：}B1 $>$ M1 $>$ G1 $>$ B2 $>$ M2，监测以AFB1为指标。\textbf{理化性质：}耐热（268$\sim$269$^\circ$C才分解），脂溶性，碱性条件内酯环开环破坏，紫外线下产生荧光。\textbf{产毒条件：}最适温度25$\sim$33$^\circ$C，最适Aw 0.93$\sim$0.98。\textbf{代谢：}AFB1在肝脏CYP450催化下：AFB1 $\rightarrow$ AFM1（羟化，乳汁中）；AFB1 $\rightarrow$ AFQ1（羟化，解毒）；AFQ1 $\rightarrow$ AFH1（还原，解毒）；AFB1 $\rightarrow$ AFB1-8,9-环氧化物（终致癌物，与DNA形成加合物）。\textbf{毒性：}急性中毒性肝炎（黄疸、发热、腹水），慢性肝损伤（发育障碍、肝功降低），致癌性——肝癌，诱癌强度是二甲基亚硝胺75倍，多物种多器官致癌，灵长类可诱癌，与HBV协同致癌，IARC Group 1。AFM1在乳及乳制品中稳定，一般家庭加热及加工过程均对其水平无显著影响。

    \item \textbf{预防：}（1）食物防霉——控制水分在安全水分以下，低温保藏，通风；（2）去毒7法：挑选霉粒法、碾轧加工法、加水搓洗法、植物油加碱去毒（NaOH水解内酯环形成香豆素钠盐，碱炼后水洗）、物理吸附法（活性白陶土/活性炭，液体食品效果好，固体食品效果不明显）、氨气处理法（18kg压力/72$\sim$82$^\circ$C，去除98\%$\sim$100\%）、日光/紫外线照射法；（3）制定限量标准。

    \item \textbf{单端孢霉烯族化合物：}T-2毒素（ATA病原物），化学性质稳定，强细胞毒性/免疫抑制/致畸/潜在致癌性，污染麦类、玉米。\textbf{玉米赤霉烯酮（F-2毒素）：}雌激素样作用，生殖系统毒性，污染玉米、小麦、大麦。\textbf{伏马菌素：}FB1为主，干扰鞘脂类代谢，神经毒性（马脑白质软化症），肾脏毒性，促癌物诱发原发性肝癌，污染玉米。

    \item \textbf{腐败变质概念：}微生物为主各因素作用下，食品成分被分解破坏，失去或降低食用价值。\textbf{三大原因：}（1）微生物——细菌（分解蛋白质）、霉菌和酵母菌（分解碳水化合物、脂肪）；（2）食物本身——酶作用、营养成分和水分（高蛋白=蛋白质腐败，高碳水=产酸发酵，Aw越小越不易腐败）、pH（酸性/碱性不利细菌）；（3）环境——温度、湿度、氧气、光照。

    \textbf{蛋白质腐败：}蛋白质 $\rightarrow$ 多肽 $\rightarrow$ 氨基酸 $\rightarrow$ 脱羧/脱氨 $\rightarrow$ 酪胺、组胺、尸胺（恶臭）+ 硫醇（恶臭）+ 色胺、甲基色胺（粪臭）。\textbf{脂肪酸败：}水解+氧化，中性脂肪 $\rightarrow$ 甘油+脂肪酸 $\rightarrow$ 醛、酮、酸，不饱和脂肪酸氧化形成过氧化物再分解为醛酮酸，不饱和脂肪酸含量越高越易氧化，产生"哈喇味"。酸败早期即可出现酸度增高。鉴定：油脂气味、过氧化值上升（最早期指标）、酸价和羰基价阳性。\textbf{碳水化合物分解：}$\rightarrow$ 单糖、醇、醛、酮、CO$_2$和水，酸度升高，产生酸、醇、醛气味。

    \item \textbf{鉴定指标：}（1）感官——最敏感，色/香/味/形（不良气味、变色、软化、发黏）；（2）物理——油脂比重、折光率、黏度；（3）化学——TVBN（鱼肉类蛋白腐败鉴定，与腐败程度有明确对应关系）；三甲胺（鱼虾水产品新鲜程度）；组胺（$>$200 mg/100g可致过敏性食物中毒）；K值（$<$20\%新鲜，$\geq$40\%开始腐败，鉴定鱼类早期腐败）；pH（V字形变动）；过氧化值（脂肪酸败最早期指标）+ 酸价（脂肪酸败程度）；（4）微生物——菌落总数、大肠菌群、细菌菌相。

    \item \textbf{处理原则：}严重腐败变质——销毁或工业用；轻度腐败——加工处理降至安全限量以下，限期食用；已变质——禁止食用；局部变质——挑选去除变质部分，利用完好部位。

    \item \textbf{化学保藏：}盐腌——食盐浓度达10\%抑制大多数细菌，糖含量达60\%$\sim$65\%即可；酸渍——pH 4.5以下抑制微生物，用于各种蔬菜（泡菜、酸菜）；防腐剂——苯甲酸、山梨酸；抗氧化剂——防止油脂酸败。

    \textbf{低温保藏：}冷藏——$-1\sim10^\circ$C，低温抑制微生物繁殖和酶活性；冷冻——$\leq-18^\circ$C，几乎全部微生物不繁殖；急冻缓化有利于保持食品品质。冰晶生成带：$-1\sim-5^\circ$C（水结冰率85\%），$-8\sim-120^\circ$C（水结冰率90\%），速冻迅速通过冰晶生成带，冰晶数量多、体积小，不损伤组织细胞。

    \textbf{高温杀菌：}常压——100$^\circ$C以下，巴氏杀菌（低温63$^\circ$C/30min，高温短时72$^\circ$C/15s），处理牛乳/啤酒/酱油/葡萄酒；高压——100$\sim$121$^\circ$C（0.2 MPa），肉类制品/罐头，杀灭繁殖型和芽孢型；UHT——封闭系统加热至120$^\circ$C以上数秒后迅速冷却；微波——915 MHz（含水量高/体积大食品）、2450 MHz（含水量低食品），快速节能对品质影响小。

    \textbf{D值：}杀死90\%细菌所需时间（min）。金黄色葡萄球菌D$_{63}$=7，肉毒梭菌芽孢D$_{121}$=0.1$\sim$0.23。

    \textbf{干燥脱水：}水分降至15\%以下或Aw 0.00$\sim$0.60，抑制各种微生物生长。方法：日晒、阴干、喷雾干燥、减压蒸发、冷冻干燥。酶性褐变——酚酶催化酚类形成醌。非酶褐变/Maillard反应——氨基与羰基反应，Aw 0.6最适于褐变，加热与给氧促进褐变，SO$_2$可阻止褐变。

    \textbf{辐照保藏：}辐照灭菌——高剂量10$\sim$15/50 kGy，杀灭所有微生物达到商业无菌；辐照消毒——中等剂量5$\sim$10 kGy，杀灭无芽孢致病菌（沙门菌等），3天内不发虫霉变；辐照防腐——低剂量1$\sim$5 kGy，杀灭腐败菌，延长后熟期及保藏期（抑制发芽、杀虫）。

    \item \textbf{五大分类方式：}按用途（杀虫/杀鼠/杀菌剂等）、按化学组成及结构（有机磷/有机氯/氨基甲酸酯/拟除虫菊酯/有机砷/有机汞等）、按成分和来源（无机/有机/生物农药）、按急性毒性大小（剧毒/高毒/中等毒/低毒）、按残留特性（高残留/中等残留/低残留）。

    \item \textbf{$N$-亚硝基化合物：}前体物为硝酸盐/亚硝酸盐+胺类。合成条件为pH $<$ 3的酸性环境（胃是主要合成场所，胃液pH 1$\sim$2，Cl$^-$/SCN$^-$催化）。致癌特点：（1）多物种无抵抗；（2）多器官靶向（对称亚硝胺=肝癌，不对称=食管癌，可过血脑/胎盘屏障）；（3）多途径诱癌（经口/吸入/注射/皮肤）；（4）单次大剂量或长期小剂量均可致癌，有剂量-效应关系，各剂量组合相加；（5）经胎盘致子代肿瘤。预防：防微生物污染、改进加工工艺、施钼肥、多食新鲜蔬果（VC/VE/酚类/黄酮阻断亚硝化）、制定限量标准（水产品$\leq$4$\mu$g/kg，肉制品$\leq$3$\mu$g/kg）。

    \item \textbf{B(a)P：}主要蓄积于脂肪组织，致胃癌/致突变。污染途径：烘烤/烟熏直接污染（最主要，脂肪滴落炭火热解聚合）、高温热解热聚合（脂肪$>$400$^\circ$C）、农作物从环境吸收、加工接触。防控：防止污染、活性炭吸附去除、制定限量标准（熏烤肉$\leq$5$\mu$g/kg，植物油$\leq$10$\mu$g/kg）。

    \item \textbf{杂环胺：}来源及影响因素——烹调方式（温度越高/时间越长/水分越少生成越多）、食物成分（蛋白质越高生成越多，美拉德反应是前提）。致癌性——肝癌、致突变（DNA加合物）。预防——改变不良烹调方式和饮食习惯、多吃蔬菜水果、加强监测。

    \item \textbf{丙烯酰胺：}富淀粉食品 $>$ 120$^\circ$C高温加工时天冬酰胺+还原糖经Maillard反应生成。高含量食品：薯条、薯片、方便面、油条、速溶咖啡（土豆制品约为谷物制品4倍）。神经毒性，IARC Group 2A。预防：降低油炸温度和时间、选用低还原糖马铃薯品种。

    \item \textbf{氯丙醇：}首次发现于酸水解植物蛋白液（HVP），3-MCPD含量最高最常见，具肾毒性和雄性生殖毒性。游离态3-MCPD（酸水解HVP）及脂肪酸酯形式（3-MCPD酯，广泛存在于精炼油脂中）。

    \item \textbf{食品容器包装材料：}概念——直接或间接接触食品的各种材料、制品及辅助材料。分类：（1）塑料——单体残留（氯乙烯/苯乙烯/双酚A）、增塑剂迁移（邻苯二甲酸酯类，内分泌干扰物）；（2）橡胶——硫化促进剂/抗氧化剂/填充剂迁移；（3）纸竹木——漂白剂残留/荧光增白剂/印刷油墨迁移；（4）金属搪瓷陶瓷玻璃——铅镉铬镍铝溶出，釉彩铅镉迁移。

    \item \textbf{影响迁移的因素：}（1）接触时间和温度——越长/越高迁移量越大；（2）食品性质——酸性/含酒精/油脂类萃取能力更强，pH越低/脂肪越高迁移越大；（3）容器材料性质——化学组成/加工工艺/添加剂，再生塑料风险高；（4）接触面积与食品体积之比——面积越大/体积越小，迁移浓度越高；（5）材料老化程度——裂纹和磨损增加释放。

    \item \textbf{综合防控策略：}（要点）三大污染来源——生物性（微生物：强化HACCP/GMP、冷链管理；寄生虫：加强检验检疫）；化学性（农药：登记制度+安全间隔期+MRLs+高效低毒新品种；有害金属：消除工业污染源是根本+GB 2762限量+金属间拮抗营养策略；$N$-亚硝基化合物：防微生物污染+改进加工工艺+施钼肥+多食新鲜蔬果+限量标准；PAH：改进烟熏工艺+活性炭吸附+限量标准；杂环胺：合理烹调降低温度+多吃蔬菜水果）；物理性（放射性：核事故应急+食品监测；杂物：严格GMP）。综合防控：完善法规标准体系、建立全程可追溯制度、加强风险监测和风险评估、开展食品安全健康教育。

    \item \textbf{黄曲霉毒素污染调查与防控方案：}（要点）调查：采样检测可疑食品（花生/玉米/谷物/植物油等）中AFB1含量，现场调查储藏条件（温度/湿度/Aw），了解人员中毒情况（肝功能检测/流行病学调查）。防控：立即封存/销毁超标食品；查找污染源头（田间/储藏/加工环节）；防霉措施（控制水分至安全水分以下+低温保藏+通风）；已污染食品去毒（挑选霉粒/碾轧加工/碾轧加工/碱炼/吸附/氨气处理/紫外线照射）；严格执行限量标准（GB 2761）；加强健康教育（HBV疫苗接种+AF暴露风险告知，因AF与HBV有协同致癌作用）。
\end{enumerate}

}
\newpage

% =====================================================================
%  CHAPTER 9: 食品添加剂及其管理 (Food Additives and Management)
% =====================================================================
\chapter{食品添加剂及其管理}
\setcounter{chapter}{9}
\chapterintro{
\keyword{掌握}食品添加剂的定义、使用要求；掌握JECFA安全性评价分类；掌握各类食品添加剂（抗氧化剂、防腐剂、护色剂——亚硝酸盐*等）的作用。\keyword{熟悉}食品添加剂的管理制度。\keyword{了解}各类食品添加剂的常用品种和使用范围。
}

% =====================================================================
%  第一节 食品添加剂的定义与分类
% =====================================================================
\section{食品添加剂的定义与分类}

\subsection{定义}

\begin{keybox}
\textbf{食品添加剂的定义（掌握）——中国《食品安全法》：}
食品添加剂是指为改善食品品质和色、香、味，以及为防腐、保鲜和加工工艺的需要而加入食品中的\textbf{人工合成或者天然物质}。

\textbf{定义涵盖范围包括：}
\begin{itemize}[leftmargin=*]
    \item 食品用香料
    \item 胶基糖果中基础剂物质
    \item 食品工业用加工助剂
    \item 营养强化剂
\end{itemize}
\end{keybox}

\begin{tipbox}
\textbf{中国与CAC定义的差异：}
\begin{itemize}[leftmargin=*]
    \item \textbf{中国定义：}包含营养强化剂、食品用香料、加工助剂、胶基物质。
    \item \textbf{CAC（国际食品法典委员会）定义：}食品添加剂通常不包括污染物，也不包括为保持或提高营养价值而添加的物质（即营养强化剂不属于添加剂范畴）。
    \item 中国的定义范围更广，将营养强化剂也包括在食品添加剂中。
\end{itemize}
\end{tipbox}

\subsection{分类}

\textbf{中国GB 2760-2014《食品安全国家标准 食品添加剂使用标准》将食品添加剂按功能分为23类，主要包括：}

\begin{enumerate}[leftmargin=*]
    \item 酸度调节剂（Acidity Regulator）
    \item 抗结剂（Anticaking Agent）
    \item 消泡剂（Antifoaming Agent）
    \item 抗氧化剂（Antioxidant）
    \item 漂白剂（Bleaching Agent）
    \item 膨松剂（Bulking Agent）
    \item 着色剂（Color）
    \item 护色剂（Color Fixative）
    \item 乳化剂（Emulsifier）
    \item 酶制剂（Enzyme Preparation）
    \item 增味剂（Flavor Enhancer）
    \item 面粉处理剂（Flour Treatment Agent）
    \item 被膜剂（Coating Agent）
    \item 水分保持剂（Humectant）
    \item 营养强化剂（Nutrient Fortifier）
    \item 防腐剂（Preservative）
    \item 稳定剂和凝固剂（Stabilizer and Coagulator）
    \item 甜味剂（Sweetener）
    \item 增稠剂（Thickener）
    \item 食品用香料（Food Flavorings）
    \item 食品工业用加工助剂（Food Processing Aids）
    \item 胶基糖果中基础剂物质（Gum Base Materials）
    \item 其他（Others）
\end{enumerate}

% =====================================================================
%  第二节 食品添加剂的使用要求
% =====================================================================
\section{食品添加剂的使用要求}

\subsection{基本要求}

\begin{keybox}
\textbf{食品添加剂的使用原则（掌握）：}
\begin{enumerate}[leftmargin=*]
    \item \textbf{不应对人体产生任何健康危害。}
    \item \textbf{不应掩盖食品腐败变质。}
    \item \textbf{不应掩盖食品本身或加工过程中的质量缺陷（即不得以掺杂、掺假、伪造为目的使用食品添加剂！）。}
    \item \textbf{不应降低食品本身的营养价值。}
    \item \textbf{在达到预期效果的前提下，尽可能降低在食品中的使用量。}
\end{enumerate}
\end{keybox}

\begin{keybox}
\textbf{食品添加剂允许使用的情形（掌握）：}
\begin{enumerate}[leftmargin=*]
    \item 保持或提高食品本身的营养价值
    \item 作为某些特殊膳食用食品的必需配料或成分
    \item 提高食品的质量和稳定性，改进其感官特性
    \item 便于食品的生产、加工、包装、运输或者贮藏
\end{enumerate}
\end{keybox}

\subsection{带入原则（Carry-over Principle）}

\begin{defbox}
\textbf{带入原则（Carry-over Principle）：}某种食品添加剂不是直接加入到某食品中，而是通过含该添加剂的食品配料（或原料）带入到该食品中的情况。

\textbf{带入原则的适用条件：}
\begin{enumerate}[leftmargin=*]
    \item 根据GB 2760，食品配料中允许使用该添加剂
    \item 食品配料中该添加剂的用量不应超过允许的最大使用量
    \item 应在正常生产工艺条件下使用这些配料
    \item 由配料带入食品中的该添加剂的含量应明显低于直接将其添加到该食品中通常所需要的水平
\end{enumerate}
\end{defbox}

% =====================================================================
%  第三节 JECFA安全性评价分类
% =====================================================================
\section{JECFA对食品添加剂的安全性评价分类}

\begin{keybox}
\textbf{JECFA（FAO/WHO联合食品添加剂专家委员会）安全性评价分类（掌握）*：}
\end{keybox}

\begin{enumerate}[leftmargin=*]
    \item \textbf{第一类 GRAS（Generally Recognized As Safe，一般公认安全）：}
    \begin{itemize}[leftmargin=*]
        \item 以现有毒理学资料和化学、生物化学知识来看，按正常需要使用，对人体是安全的。
        \item 此类物质毒性极低、不规定ADI值。
        \item 可按照正常生产需要使用，不需限定最大用量。
        \item 示例：食盐、糖、柠檬酸、醋酸、碳酸氢钠等。
    \end{itemize}

    \item \textbf{第二类 A类：}
    \begin{itemize}[leftmargin=*]
        \item \textbf{A1类：}JECFA已制定正式ADI值（有完整的毒理学资料，安全性评价明确）。
        \item \textbf{A2类：}JECFA尚未制定正式ADI值（毒理学资料不完全），仅制定临时ADI值，在规定的期限内暂时允许使用。
    \end{itemize}

    \item \textbf{第三类 B类：}
    \begin{itemize}[leftmargin=*]
        \item JECFA曾进行过安全性评价，但因毒理学资料不足，\textbf{未能制定ADI值}。
        \item 需要进一步的毒理学数据补充后才能重新评价。
    \end{itemize}

    \item \textbf{第四类 C类：}
    \begin{itemize}[leftmargin=*]
        \item \textbf{C1类：}JECFA根据毒理学资料，认为在食品中使用是\textbf{不安全的}。
        \item \textbf{C2类：}JECFA根据毒理学资料，认为应\textbf{仅限于特殊用途食品或在特殊条件下}严格限制使用。
    \end{itemize}
\end{enumerate}

\begin{exambox}
\textbf{常见考试混淆点：}
\begin{itemize}[leftmargin=*]
    \item GRAS（第一类）= 无ADI（不需要限制），正常使用即可
    \item A1 = 有正式ADI；A2 = 仅临时ADI
    \item B类 = 毒理学资料不足，无ADI
    \item C1 = 不安全；C2 = 仅特殊情况限制使用
\end{itemize}
\end{exambox}

% =====================================================================
%  第四节 各类食品添加剂（常用品种）
% =====================================================================
\section{常用食品添加剂各论}

\subsection{酸度调节剂（Acidity Regulator）}

\begin{itemize}[leftmargin=*]
    \item \textbf{常用品种：}柠檬酸、乳酸、苹果酸、酒石酸、磷酸（酸化剂）；碳酸氢钠、碳酸钠、氢氧化钠（碱化剂）。
    \item \textbf{特点：}大多数酸度调节剂为机体正常代谢产物或参与正常代谢的物质，\textbf{毒性极低}，使用安全性高。
    \item \textbf{功能：}调节/维持食品酸度，改善风味，凝胶作用，防腐辅助（降低pH可抑制某些微生物）。
\end{itemize}

\subsection{抗氧化剂（Antioxidant）}

\begin{keybox}
\textbf{抗氧化剂的分类（掌握）：}
\begin{enumerate}[leftmargin=*]
    \item \textbf{脂溶性抗氧化剂：}
    \begin{itemize}[leftmargin=*]
        \item BHA（丁基羟基茴香醚）、BHT（二丁基羟基甲苯）、PG（没食子酸丙酯）、TBHQ（特丁基对苯二酚）。
        \item 主要用于油脂及含油食品，防止脂肪氧化酸败。
    \end{itemize}
    \item \textbf{水溶性抗氧化剂：}
    \begin{itemize}[leftmargin=*]
        \item 抗坏血酸（维生素C）及其钠盐、异抗坏血酸及其钠盐。
        \item 应用范围比脂溶性抗氧化剂更广。
    \end{itemize}
    \item \textbf{天然抗氧化剂：}
    \begin{itemize}[leftmargin=*]
        \item 茶多酚、维生素E（生育酚）、迷迭香提取物、竹叶提取物、甘草抗氧化物。
        \item 安全性优于合成抗氧化剂，日益受到重视。
    \end{itemize}
\end{enumerate}
\end{keybox}

\begin{tipbox}
\textbf{抗氧化剂的协同作用：}BHA与BHT混合使用时抗氧化效果优于单独使用。柠檬酸、磷酸等金属离子螯合剂可作为抗氧化剂的增效剂（Synergist），通过螯合催化氧化的金属离子来增强抗氧化效果。
\end{tipbox}

\subsection{着色剂（Color / Food Colorant）}

\begin{enumerate}[leftmargin=*]
    \item \textbf{天然色素：}
    \begin{itemize}[leftmargin=*]
        \item 辣椒红素（capsanthin）、姜黄素（curcumin）、红曲红（monascus red）、栀子黄、\(\beta\)-胡萝卜素、焦糖色。
        \item 优点：安全性较高，部分兼具营养功能（如\(\beta\)-胡萝卜素为维生素A前体）。
        \item 缺点：着色力较弱、稳定性较差（对光、热、pH敏感）、价格较高。
    \end{itemize}
    \item \textbf{合成色素：}
    \begin{itemize}[leftmargin=*]
        \item 柠檬黄（tartrazine）、苋菜红（amaranth）、亮蓝（brilliant blue）、日落黄、胭脂红、靛蓝。
        \item 优点：着色力强、色泽鲜艳、稳定性好、价格低廉。
        \item 缺点：安全性存疑（部分合成色素有致敏、行为影响等风险）。
    \end{itemize}
    \item \textbf{禁用色素（严禁用于食品！）：}
    \begin{itemize}[leftmargin=*]
        \item 苏丹红（Sudan Red）——致癌性，曾用于辣椒粉、鸭蛋（"红心鸭蛋"事件）
        \item 奶油黄（Butter Yellow）——致癌性
    \end{itemize}
\end{enumerate}

\subsection{防腐剂（Preservative）}

\begin{keybox}
\textbf{常用防腐剂分类（掌握）：}
\begin{enumerate}[leftmargin=*]
    \item \textbf{酸型防腐剂：}
    \begin{itemize}[leftmargin=*]
        \item 苯甲酸（Benzoic Acid）及其钠盐、山梨酸（Sorbic Acid）及其钾盐、丙酸（Propionic Acid）及其盐类。
        \item 有效形式为未解离的分子，因此\textbf{在酸性条件下（低pH）防腐效果更强}。在pH $>$ 5.5时效果显著下降。
        \item 山梨酸的安全性高于苯甲酸，是目前国际上应用最广泛的防腐剂之一。
    \end{itemize}
    \item \textbf{酯型防腐剂：}
    \begin{itemize}[leftmargin=*]
        \item 对羟基苯甲酸酯类（尼泊金酯类，Parabens）。
        \item 在较宽的pH范围内均有效（pH 4–8），不受酸性条件的限制。
        \item 缺点：水溶性差、具特殊气味。
    \end{itemize}
    \item \textbf{生物型防腐剂：}
    \begin{itemize}[leftmargin=*]
        \item \textbf{乳酸链球菌素（Nisin）：}由乳酸链球菌产生的一种多肽类抗菌物质。
        \item 特点：对G$^+$菌（特别是产芽孢菌如肉毒梭菌）有强烈抑制作用。对G$^-$菌、酵母、霉菌无效。
        \item 安全性高——可被人体消化道中的蛋白酶水解为氨基酸。
    \end{itemize}
\end{enumerate}
\end{keybox}

\subsection{护色剂（Color Fixative）— 亚硝酸盐}

\begin{keybox}
\textbf{护色剂（发色剂）——亚硝酸盐（Nitrite）*（掌握）：}
\end{keybox}

\begin{defbox}
\textbf{护色剂：}能使肉及肉制品呈良好色泽的物质。最常用的是\textbf{亚硝酸钠（NaNO$_2$）}和\textbf{硝酸钠（NaNO$_3$）}。
\end{defbox}

\begin{keybox}
\textbf{亚硝酸盐在肉制品中的三大作用*：}
\begin{enumerate}[leftmargin=*]
    \item \textbf{发色作用（护色）：}
    \begin{itemize}[leftmargin=*]
        \item 亚硝酸盐在肉中还原为NO $\rightarrow$ NO与肌红蛋白（Mb）反应生成亚硝基肌红蛋白（NOMb）$\rightarrow$ 加热后形成稳定的\textbf{亚硝基血色原}，赋予肉制品\textbf{鲜艳的桃红色}。
    \end{itemize}
    \item \textbf{抑菌作用（最重要的安全功能！）：}
    \begin{itemize}[leftmargin=*]
        \item 亚硝酸盐是目前已知\textbf{抑制肉毒梭菌（C.\ botulinum）最有效的物质！}
        \item 对肉毒梭菌芽孢的萌发和生长有强力抑制作用，是罐头类肉制品和腌制肉制品中\textbf{不可替代}的防腐成分。
    \end{itemize}
    \item \textbf{增强风味：}赋予腌腊肉制品特有的风味。
\end{enumerate}
\end{keybox}

\begin{exambox}
\textbf{亚硝酸盐危害的双面性（考试常见考点）：}
\begin{itemize}[leftmargin=*]
    \item \textbf{益处：}最佳的肉毒梭菌抑制剂 $+$ 良好的发色和风味作用。
    \item \textbf{风险：}与胺类物质在体内外结合生成$N$-亚硝基化合物（强烈致癌物）！尤其在胃的酸性条件下更易发生。
    \item \textbf{管理对策：}GB 2760规定亚硝酸盐在肉制品中的最大使用量（通常 $\leq$ 0.15 g/kg）及最终残留限量（$\leq$ 30 mg/kg），并要求与抗坏血酸或异抗坏血酸钠（VC可阻断亚硝化反应！）同时使用。
\end{itemize}
\end{exambox}

\subsection{甜味剂（Sweetener）}

\begin{enumerate}[leftmargin=*]
    \item \textbf{天然甜味剂：}
    \begin{itemize}[leftmargin=*]
        \item \textbf{糖醇类：}山梨糖醇（sorbitol）、木糖醇（xylitol）、麦芽糖醇（maltitol）、甘露糖醇（mannitol）。热量低、不引起血糖剧烈波动、不致龋齿。
        \item \textbf{非糖醇类：}甜菊糖苷（stevioside，从甜叶菊提取）、甘草酸。甜度极高（蔗糖的200–300倍）。
    \end{itemize}

    \item \textbf{人工合成甜味剂：}
    \begin{itemize}[leftmargin=*]
        \item \textbf{磺胺类：}甜蜜素（cyclamate）、糖精钠（saccharin）。历史悠久、价格低廉。甜蜜素在一些国家被禁用（膀胱致癌性存疑）。
        \item \textbf{二肽类：}阿斯巴甜（aspartame）——由天冬氨酸和苯丙氨酸组成。甜度为蔗糖的200倍！但苯丙酮尿症（PKU）患者禁用。热稳定性差，不宜用于高温加工食品。
        \item \textbf{蔗糖衍生物：}三氯蔗糖（sucralose）——以蔗糖为原料经氯代制成。甜度为蔗糖的600倍，不被人体代谢、无热量，耐高温。安赛蜜（acesulfame K）——甜度为蔗糖的200倍。
    \end{itemize}
\end{enumerate}

\begin{tipbox}
\textbf{阿斯巴甜的安全提醒：}必须在标签上注明"含苯丙氨酸"，以提醒苯丙酮尿症（PKU）患者避免食用。
\end{tipbox}

\subsection{增稠剂（Thickener）}

\begin{itemize}[leftmargin=*]
    \item \textbf{常用品种：}瓜尔胶、黄原胶、卡拉胶（carrageenan）、明胶、琼脂、羧甲基纤维素钠（CMC-Na）。
    \item \textbf{特征：}大多为大分子亲水性多糖类或多肽类，在肠道中不易消化吸收，大多数来自天然（植物、海藻、微生物发酵）。
    \item \textbf{功能：}提高食品黏度、改善口感和组织状态、稳定乳化体系、防止冰晶形成（冷冻食品中）。
\end{itemize}

\subsection{乳化剂（Emulsifier）}

\begin{itemize}[leftmargin=*]
    \item \textbf{常用品种：}卵磷脂（lecithin，从大豆或蛋黄提取）、单硬脂酸甘油酯（单甘酯）、聚甘油脂肪酸酯。
    \item \textbf{功能：}降低油-水界面张力，使不相溶的两相均匀分散，形成稳定的乳状液。面包中使脂肪均匀分布、延缓老化；冰淇淋中使脂肪均匀分散、改善口感。
\end{itemize}

\subsection{其他重要食品添加剂}

\begin{enumerate}[leftmargin=*]
    \item \textbf{漂白剂（Bleaching Agent）：}亚硫酸盐（还原性漂白）——用于干果、蜜饯、淀粉等。过量可破坏食品中维生素B1。GB 2760严格限制残留量（以SO$_2$计）。
    \item \textbf{膨松剂（Bulking Agent）：}碳酸氢钠、碳酸氢铵、明矾（硫酸铝钾——含铝，长期大量摄入有神经毒性风险）。
    \item \textbf{水分保持剂（Humectant）：}磷酸盐——提高肉制品保水性。
    \item \textbf{增味剂（Flavor Enhancer）：}味精（谷氨酸钠，MSG）、5$'$-呈味核苷酸二钠（IMP+GMP，与MSG有强烈协同增鲜效应）。
    \item \textbf{面粉处理剂：}过氧化苯甲酰（漂白面粉）——已在中国禁用（2011年起）。
\end{enumerate}

% =====================================================================
%  第五节 食品添加剂的管理制度
% =====================================================================
\section{食品添加剂的管理制度}

\begin{keybox}
\textbf{中国食品添加剂管理的法律依据：}
\begin{enumerate}[leftmargin=*]
    \item \textbf{《食品安全法》：}第39条规定国家对食品添加剂生产实行许可制度；第40条规定食品添加剂应当在技术上确有必要且经过风险评估证明安全可靠，方可列入允许使用的范围。
    \item \textbf{GB 2760《食品安全国家标准 食品添加剂使用标准》：}是食品添加剂使用的核心标准，规定了允许使用的添加剂品种、使用范围和最大使用量（或残留量）。
    \item \textbf{GB 14880《食品安全国家标准 食品营养强化剂使用标准》：}专门规定营养强化剂的使用范围和用量。
\end{enumerate}
\end{keybox}

\textbf{新食品添加剂的审批：}
\begin{itemize}[leftmargin=*]
    \item 审批机构：国家卫生健康委员会（NHCC，原卫生部）。
    \item 审批程序：申请人提交安全性评估资料、生产工艺、质量标准、使用范围等 $\rightarrow$ 食品安全风险评估中心评审 $\rightarrow$ NHCC批准公告 $\rightarrow$ 纳入GB 2760。
    \item 基本原则：技术上确有必要，且经风险评估证明安全可靠。
\end{itemize}

\textbf{食品添加剂的标识要求：}预包装食品的配料表中，食品添加剂应按GB 2760规定的通用名称标示；阿斯巴甜必须加注"含苯丙氨酸"字样。

% =====================================================================
%  复习题
% =====================================================================
\reviewcontent{
\section{复习题}

\subsection*{一、名词解释}

\begin{enumerate}[leftmargin=*, itemsep=6pt]
    \item \textbf{食品添加剂（中国《食品安全法》定义）、带入原则（Carry-over）}
    \item \textbf{ADI、GRAS}
    \item \textbf{护色剂、防腐剂、抗氧化剂}
\end{enumerate}

\subsection*{二、简答题}

\begin{enumerate}[leftmargin=*, itemsep=8pt]
    \item \textbf{【重点】}简述食品添加剂的定义及使用原则（基本要求）。
    \item \textbf{【重点】}简述JECFA对食品添加剂安全性评价的四类分级。
    \item \textbf{【重点】}简述亚硝酸盐在肉制品加工中的三大作用。为何亚硝酸盐虽然存在形成$N$-亚硝基化合物的风险，仍在肉制品加工中允许使用？如何控制这一风险？
    \item \textbf{【重点】}比较苯甲酸/山梨酸（酸型）与尼泊金酯类（酯型）防腐剂的作用特点差异。
    \item 简述抗氧化剂的分类，各举一例代表物质。
    \item 比较天然色素与合成色素的优缺点。
    \item 简述甜味剂的主要类型及阿斯巴甜使用的特殊注意事项。
    \item 简述食品添加剂"带入原则"的含义及其适用条件。
    \item 简述中国食品添加剂管理的主要法律法规和标准。
\end{enumerate}

\subsection*{三、综合论述题}

\begin{enumerate}[leftmargin=*, itemsep=8pt]
    \item \textbf{【综合】}结合GB 2760的使用原则，论述食品添加剂"安全性"与"技术必要性"的关系。试以亚硝酸盐为例，分析食品添加剂管理中"风险-获益"平衡的理念。
\end{enumerate}

\subsection*{参考答案要点}

\subsubsection*{名词解释（要点）}

\begin{enumerate}[leftmargin=*, itemsep=4pt]
    \item \textbf{食品添加剂：}为改善食品品质和色香味或防腐、保鲜和加工需要加入食品中的人工合成或天然物质。包括香料、胶基、加工助剂和营养强化剂。
    \item \textbf{带入原则：}添加剂非直接加入某食品，而是通过含此添加剂的配料带入。
    \item \textbf{ADI：}每日允许摄入量——终生每日摄入而无健康风险的剂量（mg/kg bw）。
    \item \textbf{GRAS：}第一类——一般公认安全，毒性极低，不规定ADI。
\end{enumerate}

\subsubsection*{简答题（要点）}

\begin{enumerate}[leftmargin=*, itemsep=6pt]
    \item \textbf{五项使用原则：}不产生健康危害、不掩盖腐败、不掩盖质量缺陷（不掺假）、不降低营养价值、最小使用量原则。
    \item \textbf{JECFA四类：}第一类GRAS（不设ADI）；第二类A（A1有正式ADI，A2有临时ADI）；第三类B（毒理资料不足，无ADI）；第四类C（C1不安全，C2仅特殊情况使用）。
    \item \textbf{亚硝酸盐三作用：}发色（形成亚硝基肌红蛋白）、抑制肉毒梭菌（最强抑菌剂）、增强风味。\textbf{风险-获益平衡：}其独特且不可替代的抑制肉毒梭菌的作用（可致死的食源性病原菌），使其在严格控制下利大于弊。\textbf{控制措施：}限量使用（$\leq$ 0.15 g/kg）、规定残留限量（$\leq$ 30 mg/kg）、与VC同时使用以阻断亚硝胺合成。
    \item \textbf{酸型防腐剂：}有效成分为未解离分子，低pH条件下效果好，pH $>$ 5.5效果剧降；\textbf{酯型（尼泊金酯）：}在pH 4–8均有效，不受酸性条件限制，但水溶性差。
    \item \textbf{抗氧化剂分类：}脂溶性（BHA/BHT/PG/TBHQ）、水溶性（抗坏血酸类）、天然（茶多酚/VE/迷迭香提取物）。
    \item \textbf{天然色素：}安全性高，但着色力弱、稳定性差、价贵。\textbf{合成色素：}着色力强、鲜艳、稳定、价廉，但安全性存疑。
    \item \textbf{甜味剂类型：}天然（糖醇类、甜菊糖苷）、人工合成（磺胺类—糖精、二肽类—阿斯巴甜、蔗糖衍生物—三氯蔗糖）。阿斯巴甜必须标注"含苯丙氨酸"以提醒PKU患者。
    \item \textbf{带入原则：}某食品中的添加剂由其原料带入，而非直接添加。需满足：原料允许使用、不超量、正常工艺、带入量显著低于直接添加水平。
    \item 法律：《食品安全法》第39/40条；核心标准：GB 2760（使用标准）和GB 14880（营养强化剂标准）；审批机构：国家卫健委（NHCC）。
\end{enumerate}

\newpage
}

% =====================================================================
%  CHAPTER 10: 各类食品卫生及其管理 (Food Hygiene by Category)
% =====================================================================
\chapter{各类食品卫生及其管理}
\setcounter{chapter}{10}
\chapterintro{
\keyword{掌握}畜禽肉卫生（肉的成熟过程）；掌握油脂酸败指标（AV/POV/CGV/MDA）。\keyword{熟悉}乳和蛋类卫生；保健食品/GMF/绿色食品/有机食品。
}

% =====================================================================
%  第一节 粮豆类卫生
% =====================================================================
\section{粮豆类卫生}

\begin{keybox}
\textbf{粮豆类的主要卫生问题：}
\begin{enumerate}[leftmargin=*]
    \item \textbf{微生物污染：}霉菌及其毒素是粮豆类最重要的卫生问题（黄曲霉毒素、镰刀菌毒素、赭曲霉毒素等）。
    \item \textbf{农药残留：}田间施药和储藏期熏蒸。
    \item \textbf{其他有害化学物质的污染：}工业废水灌溉、环境重金属本底值、有毒重金属（汞、镉、铅、砷）。
    \item \textbf{仓库害虫：}温度18--21$^\circ$C、相对湿度$>$65\%时害虫繁殖最活跃。常见害虫包括谷象、米象、谷蠹、螨类等，造成数量损失和质量损害。
    \item \textbf{其他问题：}自然陈化（储存过久风味营养降低）、有毒植物种子的污染（毒麦、麦仙翁籽、毛果洋茉莉等）、无机夹杂物（砂石、泥土）、掺杂、掺假、转基因。
\end{enumerate}

\textbf{粮豆类的卫生管理措施：}
\begin{enumerate}[leftmargin=*]
    \item \textbf{控制安全水分：}粮谷安全水分为12\%--14\%，豆类为10\%--13\%。将水分控制在安全水分以下是防霉最关键的措施。
    \item \textbf{控制储藏温度：}低温储藏（$\leq$ 10$^\circ$C可有效抑制霉菌生长）。
    \item \textbf{执行真菌毒素限量标准：}GB 2761-2017规定了谷物中AFB1、DON、ZEN等限量。
    \item \textbf{加强农药残留监管：}执行GB 2763-2021规定的MRLs。
    \item \textbf{防止仓库害虫：}清洁仓库、控制温湿度、合理使用熏蒸剂。
    \item \textbf{防止无机有害物质和有毒种子的污染。}
\end{enumerate}
\end{keybox}

% =====================================================================
%  第二节 蔬菜、水果类卫生
% =====================================================================
\section{蔬菜、水果类卫生}

\begin{keybox}
\textbf{蔬菜、水果的主要卫生问题：}
\begin{enumerate}[leftmargin=*]
    \item \textbf{细菌、霉菌及其毒素污染：}人畜粪便施肥可引入肠道致病菌和寄生虫卵。腐败菌可导致蔬菜水果腐烂溃烂。
    \item \textbf{寄生虫和虫卵的污染：}未经无害化处理的人畜粪便施肥，可引入蛔虫卵、钩虫卵等。
    \item \textbf{有害化学物污染：}农药残留（有机磷、拟除虫菊酯类等）、有害重金属（工业废水灌溉）、多环芳烃类化合物（大气沉降）、硝酸盐/亚硝酸盐（过量施用氮肥）。
\end{enumerate}

\textbf{农药、硝酸盐污染的预防：}
\begin{itemize}[leftmargin=*]
    \item 人畜粪便无害化处理后使用
    \item 生食前彻底清洗
    \item 剔除烂根残叶，防止肠道致病菌和寄生虫卵的污染
    \item 选用高效低毒低残留农药，制定和执行MRLs
    \item 慎重使用激素类农药
    \item 工业废水灌溉须经无害化处理，尽量采用地下灌溉
\end{itemize}

\textbf{蔬菜、水果贮藏的卫生要求：}
\begin{itemize}[leftmargin=*]
    \item 低温冷藏（一般蔬菜0--10$^\circ$C）
    \item 控制相对湿度
    \item 气调储藏（降低O$_2$、增加CO$_2$）
    \item 避免机械损伤
\end{itemize}
\end{keybox}

% =====================================================================
%  第三节 畜禽肉类及其制品卫生
% =====================================================================
\section{畜禽肉类及其制品卫生}

\subsection{肉的成熟过程}

\begin{keybox}
\textbf{肉的成熟过程（掌握）：}

\begin{enumerate}[leftmargin=*]
    \item \textbf{僵直（Rigor）：}宰后糖原和含磷有机化合物分解为乳酸和游离磷酸，pH降至5.4--6.7。pH=5.4时达到肌凝蛋白等电点，肌凝蛋白凝固，肌纤维硬化出现僵直。
    \item \textbf{后熟（After Ripening）：}糖原继续分解产乳酸，pH回升。肌肉结缔组织变软并具弹性，此时肉松软多汁、滋味鲜美。表面因蛋白凝固形成一层干膜，可阻止微生物侵入。
    \item \textbf{自溶（Autolysis）：}常温下存放过久，组织酶在无菌条件下继续活动，分解蛋白质、脂肪。蛋白质分解。
    \item \textbf{腐败（Putrefaction）：}自溶为细菌的入侵繁殖创造条件，细菌的酶使蛋白质、含氮物质分解，pH升高。产生H$_2$S、胺类、吲哚等恶臭物质，肉色变暗绿色（硫血红蛋白形成）。
\end{enumerate}

\textbf{鉴定指标：}感官（颜色、气味、弹性和黏度）+ TVBN。
\end{keybox}

% =====================================================================
%  第二节 油脂酸败及其指标
% =====================================================================
\section{油脂酸败及其指标}

\subsection{油脂酸败（Oil Rancidity）}

\begin{keybox}
\textbf{油脂酸败（Oil Rancidity）：}食用油脂在储藏过程中由于光、热、氧、水分和金属离子等的作用，发生一系列化学变化，使其感官性状变劣和产生不良气味的现象。

\textbf{原因：}
\begin{enumerate}[leftmargin=*]
    \item \textbf{生物学原因（水解）：}脂肪酶水解甘油三酯为甘油和游离脂肪酸。
    \item \textbf{化学原因（自动氧化——最主要原因）：}不饱和脂肪酸在O$_2$作用下经自由基链式反应，产生氢过氧化物，进一步分解为醛、酮、酸等。
\end{enumerate}
\end{keybox}

\subsection{油脂酸败的理化指标}

\begin{keybox}
\textbf{油脂酸败的理化指标（掌握*）：}

\begin{enumerate}[leftmargin=*]
    \item \textbf{酸价（Acid Value, AV）：}1g油脂中游离脂肪酸所需KOH的mg数。反映油脂水解程度。

    \item \textbf{过氧化值（Peroxide Value, POV）：}100g油脂中与KI反应析出碘的g数。POV $>$ 0.25g/100g为超标。反映油脂氧化初期产生的氢过氧化物含量——是脂肪酸败最早期、最灵敏的指标。

    \item \textbf{羰基价（Carbonyl Value, CGV）：}CGV $>$ 50meq/kg为酸败。反映油脂氧化分解产生的含羰基化合物（醛、酮等）的总量——与POV形成互补。POV反映初级氧化产物，羰基价反映次级氧化（分解）产物。

    \item \textbf{丙二醛（Malondialdehyde, MDA）：}MDA $>$ 0.25mg/100g为超标。常用于反映动物油脂酸败的程度，优点是可反映甘油三酯以外的其他物质（如磷脂）的氧化破坏程度。
\end{enumerate}
\end{keybox}

% =====================================================================
%  第三节 乳及乳制品卫生
% =====================================================================
\section{乳及乳制品卫生}

\begin{defbox}
\textbf{乳及乳制品卫生：}

\begin{itemize}[leftmargin=*]
    \item \textbf{巴氏杀菌乳：}LTLT（62.8$^\circ$C、30 min）或HTST（71.7$^\circ$C、15 s），杀灭所有致病菌，需冷链配送。
    \item \textbf{灭菌乳：}UHT（135--150$^\circ$C、2 s），达到商业无菌，常温保存。
    \item \textbf{发酵乳：}以乳为原料经乳酸菌发酵制成。
\end{itemize}

\textbf{主要卫生问题：}微生物污染（沙门菌、结核杆菌、布鲁氏菌、金葡菌肠毒素）、化学性污染（农药残留、兽药残留、AF M1）、掺伪。
\end{defbox}

% =====================================================================
%  第四节 蛋类卫生
% =====================================================================
\section{蛋类卫生}

\begin{tipbox}
\textbf{蛋类卫生：}鲜蛋的变质主要是微生物通过蛋壳气孔进入蛋内引起。沙门菌是禽蛋最主要的食源性致病菌。

\textbf{鲜蛋变质过程：}
\begin{enumerate}[leftmargin=*]
    \item \textbf{贴壳蛋：}蛋白变稀，蛋黄膜破裂，蛋黄贴于蛋壳内壁。
    \item \textbf{散黄蛋：}蛋黄膜完全破裂，蛋黄与蛋白混合。
    \item \textbf{浑汤蛋/黑腐蛋：}细菌大量繁殖，蛋白质彻底分解产生H$_2$S、胺类等，有恶臭。
\end{enumerate}
\end{tipbox}

% =====================================================================
%  第五节 水产品卫生
% =====================================================================
\section{水产品（鱼虾类）卫生}

\begin{defbox}
\textbf{水产品（鱼虾类）的主要卫生问题：}
\begin{enumerate}[leftmargin=*]
    \item \textbf{腐败变质（最主要问题）：}鱼类营养丰富、水分含量高（70\%--80\%）、污染的微生物多，且组织酶的活性高。与肉类相比，更易发生腐败变质。
    \item \textbf{有害金属富集：}甲基汞在鱼贝类中经食物链逐级富集；镉可富集于贝类体内。
    \item \textbf{病原微生物：}海产食品最易受到副溶血性弧菌（嗜盐菌）污染，是夏秋季节食物中毒的重要原因。淡水鱼虾挥发性盐基总氮 $\leq$ 20 mg/100g。
    \item \textbf{寄生虫感染：}我国常见鱼类寄生虫有华支睾吸虫（肝吸虫）、肺吸虫等。生食或半生食淡水鱼虾容易感染。
    \item \textbf{天然有毒品种：}河豚含河豚毒素（TTX）；含组胺高的青皮红肉鱼类（秋刀鱼、金枪鱼、鲐鱼）可致组胺中毒。高组胺鱼类组胺限值 $\leq$ 40 mg/100g。
\end{enumerate}

\textbf{鱼类保鲜的卫生要求：}
\begin{itemize}[leftmargin=*]
    \item 海水鱼虾挥发性盐基总氮 $\leq$ 30 mg/100g
    \item 淡水鱼虾挥发性盐基总氮 $\leq$ 20 mg/100g
    \item 有效的保鲜措施：低温（冷藏/冷冻）、盐腌、防止微生物污染和减少鱼体损伤
    \item \textbf{K值：}鱼类鲜度灵敏指标。K $<$ 20\%新鲜，K $>$ 40\%为初期腐败。
\end{itemize}
\end{defbox}

% =====================================================================
%  第七节 酒类卫生
% =====================================================================
\section{酒类卫生}

\begin{keybox}
\textbf{酒类的主要卫生问题：}

\textbf{蒸馏酒与配制酒中可能存在的有害物质：}
\begin{enumerate}[leftmargin=*]
    \item \textbf{甲醇：}主要来自制酒原辅料（薯干、马铃薯、水果、糠麸等）中的果胶。具有剧烈的神经毒性，主要侵害视神经，导致视网膜受损、视神经萎缩、视力减退和双目失明。长期少量摄入可导致慢性中毒，特征性临床表现为视野缩小、不可校正的视力减退。
    \item \textbf{醛类：}如乙醛、糠醛，为酒醅发酵过程中生成的中间产物。乙醛毒性大于甲醛。
    \item \textbf{杂醇油：}由蛋白质、氨基酸和糖类分解产生，包括丙醇、异丁醇、异戊醇等。毒性及麻醉力比乙醇强，使中枢神经系统充血，引起剧烈头痛。
    \item \textbf{氢氰酸：}来自木薯或其他含氰苷原料。中毒使机体陷于窒息状态，麻痹呼吸中枢及血管运动中枢。
    \item \textbf{铅：}来自蒸馏器、冷凝导管和储酒容器中的铅溶出。
    \item \textbf{锰：}高锰酸钾-活性炭脱臭除杂处理时引入。
\end{enumerate}

\textbf{发酵酒中的卫生问题：}展青霉素（水果原料霉变）、二氧化硫残留、亚硝胺、微生物污染（细菌、酵母菌）、添加剂使用问题。
\end{keybox}

% =====================================================================
%  第八节 冷饮食品卫生
% =====================================================================
\section{冷饮食品卫生}

\begin{keybox}
\textbf{冷饮食品的分类：}冰淇淋、雪糕、冰棍、食用冰、碳酸饮料、果汁饮料、乳酸菌饮料等。

\textbf{主要卫生问题：}
\begin{enumerate}[leftmargin=*]
    \item \textbf{微生物污染：}原料含菌量高、加工过程污染、杀菌不彻底、冷链不完善。致病菌包括沙门菌、金黄色葡萄球菌肠毒素等。
    \item \textbf{食品添加剂：}色素、香精、甜味剂、酸度调节剂等需符合GB 2760规定。
    \item \textbf{重金属污染：}铅、铜、砷等来自加工设备和原料。
\end{enumerate}
\end{keybox}

% =====================================================================
%  第九节 罐头食品卫生
% =====================================================================
\section{罐头食品卫生}

\begin{keybox}
\textbf{罐头食品的分类：}按内容物分为肉类罐头、禽类罐头、水产品罐头、水果蔬菜罐头等；按包装分为金属罐、玻璃罐、复合袋（软罐头）等。

\textbf{主要卫生问题：}
\begin{enumerate}[leftmargin=*]
    \item \textbf{罐头容器的卫生要求：}无渗漏、无锈蚀、内壁涂料完整无脱落。
    \item \textbf{原料微生物污染与灭菌：}必须达到商业无菌。D值、F值（通常12D）概念应用于灭菌条件设定。
    \item \textbf{重金属污染：}酸性内容物可腐蚀金属罐、产生氢气、溶出锡和铅。
\end{enumerate}

\textbf{罐头食品的卫生学鉴定及处理：}
\begin{enumerate}[leftmargin=*]
    \item \textbf{感官检查：}容器密封完好、无泄漏、无胖听；内容物具有该品种应有的色泽、气味、滋味、形态。
    \item \textbf{胖听（Swelling）：}由于罐内微生物活动或化学作用产生气体，形成正压，使一端或两端外凸。
    \begin{itemize}
        \item \textbf{化学性胖听：}金属罐受酸性内容物腐蚀产生氢气，叩击呈鼓音，穿洞有气体逸出但无腐败气味，一般不宜食用。
        \item \textbf{生物性胖听：}微生物活动产气所致，叩击呈浊音，穿洞有腐败臭味，\textbf{绝对不可食用！}
        \item \textbf{物理性胖听：}低温装罐、运输海拔变化等物理原因导致，内容物未变质，可食用。
    \end{itemize}
    \item \textbf{平酸腐败（Flat-sour Spoilage）：}由分解碳水化合物产酸不产气的平酸菌引起，内容物酸度增加而外观完全正常。
\end{enumerate}
\end{keybox}

% =====================================================================
%  第十节 调味品卫生
% =====================================================================
\section{调味品卫生}

\begin{keybox}
\textbf{调味品的定义和营养学意义：}以粮食、豆类、水果等为原料，经发酵或调配制成的具有调味功能的产品。主要提供咸、甜、酸、鲜、辣等基本味道，酱油和发酵酱提供少量蛋白质和B族维生素，食盐提供钠，食醋提供有机酸。

\textbf{酱油的主要卫生问题：}
\begin{itemize}[leftmargin=*]
    \item 原料霉变导致黄曲霉毒素污染
    \item 氯丙醇（3-MCPD）——酸水解植物蛋白液（HVP）配制酱油中的有害物
    \item 微生物污染（耐盐酵母菌、细菌等产生霉变和"生白"）
    \item 防腐剂使用
    \item 掺伪和假冒
\end{itemize}

\textbf{食醋的卫生问题：}
\begin{itemize}[leftmargin=*]
    \item 原料卫生
    \item 醋酸发酵过程的微生物管理
    \item 不应使用冰醋酸勾兑
    \item 昆虫（醋鳗、醋蝇）污染
\end{itemize}

\textbf{食盐的卫生问题：}
\begin{itemize}[leftmargin=*]
    \item 碘缺乏病——我国实行食盐加碘（GB 26878）
    \item 杂质控制（不溶性杂质、重金属）
    \item 亚铁氰化钾等抗结剂的使用
\end{itemize}
\end{keybox}

% =====================================================================
%  第十一节 特殊食品
% =====================================================================
\section{特殊食品}

\subsection{保健食品（Health Food）}

\begin{defbox}
\textbf{保健食品（Health Food）：}声称具有特定保健功能或以补充维生素、矿物质为目的的食品。即适用于特定人群食用，具有调节机体功能，不以治疗疾病为目的，且对人体不产生任何急性、亚急性或者慢性危害的食品。
\end{defbox}

\subsection{转基因食品（GMF）}

\begin{tipbox}
\textbf{转基因食品（Genetically Modified Food, GMF）：}利用基因工程技术改变基因组构成的生物所生产的食品。

\textbf{GMF分类：}分为I/II/III/IV类。安全性评价遵循实质等同性原则和个案评价原则。
\end{tipbox}

\subsection{无公害食品（Non-environmental Pollution Food）}

\begin{defbox}
\textbf{无公害食品（熟悉）：}产地环境、生产过程和产品质量符合国家有关标准和规范的要求，经认证合格获得认证证书并允许使用无公害农产品标志的、未经加工或者初加工的食用农产品。无公害食品是食品安全的基本门槛。
\end{defbox}

\subsection{绿色食品（Green Food）}

\begin{defbox}
\textbf{绿色食品（Green Food，熟悉）：}遵循可持续发展原则，按照绿色食品标准生产，经过专门机构认定，许可使用绿色食品标志的无污染、安全、优质、营养类食品。具备安全和营养的双重质量保证。
\begin{itemize}[leftmargin=*]
    \item \textbf{A级：}限量使用化学合成生产资料。
    \item \textbf{AA级：}完全按照有机生产方式生产（等同于有机食品）。
\end{itemize}
\end{defbox}

\subsection{有机食品（Organic Food）}

\begin{defbox}
\textbf{有机食品（Organic Food，熟悉）：}来自于有机农业生产体系，根据有机农业生产的规范生产加工，并经独立的认证机构认证的农产品及其加工产品。国际通行的最高级别认证，禁止使用任何化学合成的农药、化肥、激素、抗生素、转基因技术。
\end{defbox}

% =====================================================================
%  复习题
% =====================================================================
\reviewcontent{
\section{复习题}

\subsection*{一、名词解释}

\begin{enumerate}[leftmargin=*, itemsep=6pt]
    \item \textbf{僵直（Rigor）、后熟（After Ripening）、自溶（Autolysis）、腐败（Putrefaction）}
    \item \textbf{油脂酸败、酸价（AV）、过氧化值（POV）、羰基价（CGV）、丙二醛（MDA）}
    \item \textbf{巴氏杀菌乳、灭菌乳、K值}
    \item \textbf{保健食品、转基因食品（GMF）、绿色食品、有机食品}
\end{enumerate}

\subsection*{二、简答题}

\begin{enumerate}[leftmargin=*, itemsep=8pt]
    \item \textbf{【重点】}简述畜禽肉成熟的四个阶段及各阶段的特征。
    \item \textbf{【重点】}简述油脂酸败的定义、原因及四大理化指标（AV、POV、CGV、MDA）的卫生学意义。
    \item 简述乳及乳制品的种类和主要卫生问题。
    \item 简述鲜蛋变质的过程。
    \item 简述水产品的主要卫生问题及鲜度指标。
    \item 简述保健食品、转基因食品、绿色食品和有机食品的定义及区别。
\end{enumerate}

\subsection*{参考答案要点}

\subsubsection*{名词解释（要点）}

\begin{enumerate}[leftmargin=*, itemsep=4pt]
    \item \textbf{僵直：}宰后pH降至5.4--6.7，肌凝蛋白凝固，肌纤维硬化。\textbf{后熟：}pH回升，肉松软多汁滋味鲜美。\textbf{自溶：}组织酶分解蛋白质、脂肪。\textbf{腐败：}细菌作用，pH升高，产生H$_2$S、胺类、吲哚等恶臭物质。
    \item \textbf{油脂酸败：}油脂在光、热、氧、水分和金属离子等作用下发生感官变劣和产生不良气味的化学变化。\textbf{AV：}1g油脂中游离脂肪酸所需KOH的mg数。\textbf{POV：}100g油脂中与KI反应析出碘的g数，POV $>$ 0.25g/100g为超标，氧化早期最灵敏指标。\textbf{CGV：}CGV $>$ 50meq/kg为酸败，反映次级氧化产物。\textbf{MDA：}MDA $>$ 0.25mg/100g为超标。
    \item \textbf{巴氏杀菌乳：}LTLT或HTST消毒，杀灭所有致病菌。\textbf{灭菌乳：}UHT灭菌，商业无菌，常温保存。\textbf{K值：}鱼类鲜度指标，K $<$ 20\%新鲜，K $>$ 40\%腐败。
    \item \textbf{保健食品：}声称特定保健功能或以补充维生素、矿物质为目的的食品。\textbf{GMF：}利用基因工程技术改变基因组构成的生物所生产的食品，分为I/II/III/IV类。\textbf{绿色食品：}AA级和A级。\textbf{有机食品：}禁止化学合成农药、化肥、激素、抗生素、转基因技术的最高级别认证。
\end{enumerate}

\subsubsection*{简答题（要点）}

\begin{enumerate}[leftmargin=*, itemsep=6pt]
    \item \textbf{四阶段：}①僵直（pH 5.4--6.7，肌凝蛋白凝固，肌纤维硬化）；②后熟（pH回升，肉变松软多汁）；③自溶（蛋白质分解）；④腐败（pH升高，产生恶臭物质）。
    \item \textbf{原因：}生物学水解（脂肪酶）和化学自动氧化（自由基链式反应）。\textbf{AV：}反映水解程度。\textbf{POV：}氧化初期最灵敏指标，POV $>$ 0.25g/100g为超标。\textbf{CGV：}反映次级氧化产物，CGV $>$ 50meq/kg为酸败。\textbf{MDA：}反映动物油脂酸败程度，MDA $>$ 0.25mg/100g为超标。
    \item \textbf{种类：}巴氏杀菌乳、灭菌乳、发酵乳。\textbf{卫生问题：}微生物污染（沙门菌、结核杆菌、布鲁氏菌）、化学性污染（农药/兽药/AF M1）、掺伪。
    \item \textbf{变质过程：}贴壳蛋 $\rightarrow$ 散黄蛋 $\rightarrow$ 浑汤蛋/黑腐蛋。主要机制是微生物通过蛋壳气孔进入蛋内引起。
    \item \textbf{问题：}腐败变质（含水高、组织脆）、有害金属富集、病原微生物（副溶血性弧菌）、寄生虫（肝吸虫）、天然有毒（河豚）。\textbf{K值：}鱼类鲜度指标；\textbf{组胺中毒：}与青皮红肉鱼有关。
    \item \textbf{保健食品：}声称保健功能，不以治疗为目的。\textbf{GMF：}基因工程技术改造，分I/II/III/IV类。\textbf{绿色食品：}AA级（等同于有机）和A级（限量化学合成）。\textbf{有机食品：}最高认证，全禁化学合成物+转基因。
\end{enumerate}

\newpage
}

% =====================================================================
%  CHAPTER 11: 食源性疾病及其预防
% =====================================================================
\chapter{食源性疾病及其预防}
\setcounter{chapter}{11}
\chapterintro{
\keyword{掌握}食源性疾病的定义；掌握食物中毒的定义、特点；掌握细菌性食物中毒的发病原因、特点及常见类型（沙门氏菌、金黄色葡萄球菌、副溶血性弧菌、肉毒梭菌）；掌握河豚鱼中毒、毒蕈中毒；掌握亚硝酸盐食物中毒、有机磷农药中毒。\keyword{熟悉}食物过敏。
}

% =====================================================================
%  第一节：食源性疾病
% =====================================================================
\section{食源性疾病}

\begin{keybox}
\textbf{食源性疾病（foodborne disease）的定义：}
食品中致病因素进入人体引起的感染性、中毒性疾病。
\end{keybox}

食源性疾病包括食物中毒、食源性肠道传染病、食源性寄生虫病、人畜共患传染病、食物过敏、食物污染物引起的慢性中毒性疾病等。其中，食物中毒是最常见的一类。

% =====================================================================
%  第二节：食物中毒
% =====================================================================
\section{食物中毒}

\subsection{食物中毒的定义}

\begin{keybox}
\textbf{食物中毒（food poisoning）的定义：}
食用被有毒有害物质污染的食品或食用含有毒有害物质的食品后出现的急性、亚急性疾病。
\end{keybox}

\subsection{食物中毒的特点}

\begin{keybox}
\textbf{食物中毒的四大临床特点（掌握）：}
\begin{enumerate}[leftmargin=*]
    \item \textbf{发病与特定食物有关：}所有发病者均在近期内进食过同一种可疑食物，未进食该食物者不发病。停止食用该食物后，发病立即停止。
    \item \textbf{潜伏期短，发病曲线单峰型：}短时间内多人同时或相继发病，发病急剧，呈暴发性。发病曲线呈突然上升又迅速下降的单峰型，无传染病流行的余波（拖尾现象）。
    \item \textbf{临床表现相似：}以急性胃肠炎症状为主，最常见为消化道症状（腹痛、腹泻、恶心、呕吐）。
    \item \textbf{一般没有人与人间的直接传染：}停止进食致病食品后，发病立即下降，无二代病例出现。
\end{enumerate}
\end{keybox}

\subsection{食物中毒的分类}

食物中毒按致病因素可分为以下五类：
\begin{enumerate}[leftmargin=*]
    \item \textbf{细菌性食物中毒：}发病率最高、中毒人数最多
    \item \textbf{真菌毒素和霉变食品中毒}
    \item \textbf{动物性食物中毒：}如河豚鱼中毒
    \item \textbf{有毒植物食物中毒：}如毒蕈中毒
    \item \textbf{化学性食物中毒：}如亚硝酸盐中毒、有机磷农药中毒
\end{enumerate}

% =====================================================================
%  第三节：细菌性食物中毒
% =====================================================================
\section{细菌性食物中毒}

\begin{defbox}
\textbf{细菌性食物中毒：}摄入被致病菌或其毒素污染的食品后发生的急性或亚急性疾病。是最常见的食物中毒。
\end{defbox}

\subsection{发病原因}

\begin{keybox}
\textbf{细菌性食物中毒的三大发病原因（掌握）：}
\begin{enumerate}[leftmargin=*]
    \item \textbf{致病菌污染食品：}原料污染（生前感染或宰后污染）、加工过程污染、从业人员带菌污染。
    \item \textbf{食品的贮存方式不当：}在较高温度下存放，使致病菌大量生长繁殖和/或产生毒素。
    \item \textbf{食用前未彻底加热处理：}或生熟交叉污染，使食品中心温度未达到灭菌要求。
\end{enumerate}
\end{keybox}

\subsection{流行病学特点}

\begin{keybox}
\textbf{细菌性食物中毒的特点（掌握）：}
\begin{enumerate}[leftmargin=*]
    \item \textbf{潜伏期较短：}摄入致病菌或其毒素后，短时间内即可发病。
    \item \textbf{多呈暴发性：}短时间内多人同时或相继发病。
    \item \textbf{发病季节性明显：}夏秋季（5～10月）高发，此时气温高、湿度大，细菌易繁殖。
    \item \textbf{以胃肠道症状为主：}腹痛、腹泻、恶心、呕吐等急性胃肠炎表现。
    \item \textbf{病死率较低：}\keyword{肉毒梭菌中毒除外！}一般细菌性食物中毒如能及时治疗，预后良好。
\end{enumerate}
\end{keybox}

\subsection{常见类型}

\subsubsection{一、沙门氏菌食物中毒}

\begin{keybox}
\textbf{沙门氏菌食物中毒（掌握）：}
\begin{itemize}[leftmargin=*]
    \item \textbf{来源：}主要为动物性食品（畜禽肉、蛋类、乳类）。其中畜肉类及其制品居首位，其次为禽肉、蛋类、乳类。
    \item \textbf{病原特点：}沙门菌属（Salmonella），革兰阴性杆菌。\keyword{不分解蛋白质、不产生靛基质}（吲哚），因此\textbf{污染食品后无感官性状改变}——这是其最危险的特性。
    \item \textbf{中毒机制：}\keyword{活菌 + 内毒素}（感染型）。活菌随食物进入肠道后大量繁殖，侵入肠黏膜上皮细胞，引起炎症反应并释放内毒素。
    \item \textbf{临床特点：}\keyword{潜伏期12～36h}（一般12～14h）。\textbf{发热}、恶心呕吐、\textbf{腹痛腹泻（水样便或黄绿色稀便）}。
    \item \textbf{高发季节：}夏秋季（5～10月）占全年的80\%。
    \item \textbf{灭菌要求：}肉块中心温度$\ge$80$^\circ$C持续12分钟；鸡蛋煮沸8～10分钟。
\end{itemize}
\end{keybox}

\subsubsection{二、金黄色葡萄球菌食物中毒}

\begin{keybox}
\textbf{金黄色葡萄球菌食物中毒（掌握）：}
\begin{itemize}[leftmargin=*]
    \item \textbf{来源：}化脓性皮肤病或上呼吸道感染患者\textbf{对食品的污染}。人体鼻咽部带菌率20\%～50\%，皮肤化脓性病灶含大量致病性葡萄球菌。\keyword{带菌者污染食品是最重要的污染来源！}
    \item \textbf{中毒机制：}\keyword{肠毒素（毒素型）}。金黄色葡萄球菌在食品中繁殖时产生肠毒素，该肠毒素\textbf{耐热}——\keyword{100$^\circ$C加热2h才被破坏}，普通烹调方法难以彻底破坏。
    \item \textbf{临床特点：}\keyword{潜伏期1～6h}（极短，最短仅1h）。\textbf{剧烈反复呕吐}（最突出的症状！）、\textbf{上腹剧痛}、水样腹泻。\keyword{体温一般正常}。
    \item \textbf{预防：}\keyword{有局部化脓性感染或上呼吸道感染者应暂时调离食品工作岗位！}
\end{itemize}
\end{keybox}

\subsubsection{三、副溶血性弧菌食物中毒}

\begin{keybox}
\textbf{副溶血性弧菌食物中毒（掌握）：}
\begin{itemize}[leftmargin=*]
    \item \textbf{来源：}主要为\textbf{海产品}（鱼类、贝类、虾蟹等）。墨鱼的带菌率可达93\%。7～9月为高发季节。\keyword{嗜盐菌}——最适生长盐浓度3.5\% NaCl，淡水环境中迅速死亡。
    \item \textbf{中毒机制：}\keyword{活菌 + 溶血毒素}（混合型）。神奈川试验（Kanagawa Test）阳性——90\%的致病性菌株为神奈川现象阳性。
    \item \textbf{临床特点：}\keyword{潜伏期6～12h}。\textbf{上腹阵发性绞痛}（最突出的特点！）、\textbf{水样便或洗肉水样便}（血水样便）、\textbf{里急后重不明显}。发热较明显（37.5～39.5$^\circ$C）。
    \item \textbf{杀菌条件：}56$^\circ$C加热5分钟或90$^\circ$C加热1分钟即可杀灭；\keyword{含1\%食醋处理5分钟即可杀灭}——食用海鲜蘸醋对预防有意义。
\end{itemize}
\end{keybox}

\begin{exambox}
\textbf{副溶血性弧菌三大特征：}嗜盐性（3.5\% NaCl最适）、神奈川现象（90\%致病株阳性）、血水样便（洗肉水样便）。
\end{exambox}

\subsubsection{四、肉毒梭菌食物中毒}

\begin{keybox}
\textbf{肉毒梭菌食物中毒（掌握）}\difficulty：
\begin{itemize}[leftmargin=*]
    \item \textbf{来源：}家庭自制发酵食品（豆酱、臭豆腐、豆豉等）和罐头食品。中国以\textbf{植物性食品为主}（占91.48\%），新疆地区发病率最高。
    \item \textbf{中毒机制：}\keyword{外毒素——肉毒毒素（botulinum toxin）}，是\textbf{已知最剧烈的神经毒素}。毒素经小肠激活后，作用于神经肌肉接头，阻止乙酰胆碱释放，导致肌肉麻痹。人的致死剂量约10$^{-9}$mg/kg体重，毒性比KCN强10\,000倍。
    \item \textbf{临床特点：}\keyword{潜伏期12～48h}。以\textbf{运动神经麻痹症状为主}，胃肠道症状极少见。主要表现：\textbf{对称性颅神经损害}——\textbf{视力模糊、复视、眼睑下垂、吞咽困难、语言障碍、呼吸困难}（致死主要原因）。\textbf{病死率极高}。
    \item \textbf{芽胞耐热性：}\keyword{干热180$^\circ$C 5～15分钟，湿热100$^\circ$C 5小时，高压蒸汽灭菌121$^\circ$C 30分钟。}
    \item \textbf{治疗：}\keyword{尽早使用多价肉毒抗毒素血清}（A、B、E型），同时对症支持治疗。
    \item \textbf{婴儿肉毒中毒：}\keyword{蜂蜜是1岁以下婴儿肉毒中毒的嫌疑食品}，主要为E型。
    \item \textbf{预防：}家庭自制密封发酵食品为高危风险食品；罐头食品鼓盖/漏气绝对不要购买食用。
\end{itemize}
\end{keybox}

\begin{exambox}
\textbf{肉毒梭菌与其他细菌性食物中毒的根本区别：}胃肠道症状极少见，不以呕吐腹泻为主，主要表现为运动神经麻痹（对称性颅神经损害）。芽胞耐热性极强——湿热100$^\circ$C需5小时！中国以发酵豆制品为主要中毒食品，新疆高发！
\end{exambox}

\begin{tipbox}
\textbf{四种细菌性食物中毒对比速记表：}

\bgroup
\renewcommand{\arraystretch}{1.3}
\begin{tabularx}{\textwidth}{lXXXX}
\toprule
\textbf{病原体} & \textbf{来源} & \textbf{中毒机制} & \textbf{潜伏期} & \textbf{核心特征} \\
\midrule
沙门氏菌 & 动物性食品（肉、蛋、乳） & 活菌+内毒素（感染型） & 12～36h & 发热+水样便；食物无感官改变 \\
\addlinespace
金葡菌 & 带菌者（化脓/呼吸道感染） & 肠毒素（毒素型），耐热 & \textbf{1～6h} & 剧烈呕吐最突出；体温正常 \\
\addlinespace
副溶血性弧菌 & 海产品 & 活菌+溶血毒素（混合型） & 6～12h & 腹绞痛+洗肉水样便；嗜盐 \\
\addlinespace
肉毒梭菌 & 家庭发酵食品、罐头 & 肉毒毒素（神经毒素） & 12～48h & \textbf{运动神经麻痹}；病死率极高 \\
\bottomrule
\end{tabularx}
\egroup
\end{tipbox}

% =====================================================================
%  第四节：有毒动植物中毒
% =====================================================================
\section{有毒动植物中毒}

\subsection{河豚鱼中毒}

\begin{keybox}
\textbf{河豚鱼中毒（掌握）：}
\begin{itemize}[leftmargin=*]
    \item \textbf{有毒成分：}\keyword{河豚毒素（Tetrodotoxin, TTX）}，非蛋白质神经毒素。\textbf{耐热！煮沸、盐腌、日晒均不能破坏。}毒素分布：\keyword{卵巢毒性最强}，肝脏次之，血液、眼球、皮肤亦含毒素。新鲜洗净的肌肉基本无毒。每年\textbf{春季（2～5月）}为卵巢发育期，毒性最强。
    \item \textbf{中毒机制：}TTX\textbf{阻断Na$^+$通道}，阻止神经细胞膜钠离子内流，使神经传导阻滞。先麻痹神经末梢和感觉神经，后麻痹运动神经和中枢神经系统。
    \item \textbf{临床特点：}\keyword{潜伏期10min～3h}。发病急而剧烈，按顺序出现：
    \begin{enumerate}[leftmargin=*]
        \item \textbf{先感觉神经麻痹：}手指、\textbf{口唇、舌尖发麻}刺痛
        \item \textbf{后运动神经麻痹：}恶心、呕吐、四肢无力、肌肉麻痹、全身瘫痪
        \item \textbf{最后呼吸麻痹死亡：}呼吸衰竭是致死主要原因。通常死于发病后4～6小时，最快1.5小时。
    \end{enumerate}
    \item \textbf{治疗：}无特效解毒药！尽早洗胃 + 导泻 + 对症支持治疗（维持呼吸、升压等）。
    \item \textbf{预防：}\keyword{禁止销售河豚鱼。}严禁餐饮单位经营加工河豚鱼，加强宣传教育。
\end{itemize}
\end{keybox}

\subsection{毒蕈中毒}

\begin{keybox}
\textbf{毒蕈中毒（掌握）}\difficulty：
\begin{itemize}[leftmargin=*]
    \item \textbf{基本概况：}我国已知毒蕈约80余种，含剧毒可致死的不超过10种。不同类型毒素有\textbf{不同的临床表现}。
    \item \textbf{临床分型：}
    \begin{enumerate}[leftmargin=*]
        \item \textbf{胃肠炎型：}潜伏期短（10min～6h），恶心、呕吐、腹痛、腹泻。对症支持治疗，预后良好。
        \item \textbf{神经精神型：}毒蝇碱、裸盖菇素等。出汗、流涎、瞳孔缩小、幻视、精神错乱。\keyword{阿托品（Atropine）}拮抗治疗。
        \item \textbf{溶血型：}鹿花蕈毒素。溶血性贫血、黄疸、血红蛋白尿。\keyword{肾上腺皮质激素}治疗。
        \item \textbf{肝肾损害型（最严重！）：}毒肽（Phallotoxins）和毒伞肽（Amanitins）引起。\textbf{病死率最高。}潜伏期长（10～24h），病情发展可分为6期：
        \begin{enumerate}[leftmargin=*]
            \item 潜伏期（10～24h）
            \item 胃肠炎期
            \item \textbf{假愈期}（胃肠炎症状缓解但肝脏损害已开始——最易被忽视！）
            \item 内脏损害期（肝、肾、脑、心损害）
            \item 精神症状期（肝性昏迷、中毒性脑病）
            \item 恢复期
        \end{enumerate}
        抢救：\keyword{巯基解毒剂（二巯基丙磺酸钠、二巯基丁二酸钠等）是抢救关键！}
        \item 类光过敏型：类似日光性皮炎，身体暴露部位明显肿胀、疼痛，嘴唇肿胀外翻。
    \end{enumerate}
    \item \textbf{鹅膏菌中毒：}食用后\textbf{10小时以内彻底洗胃}是关键！
\end{itemize}
\end{keybox}

\begin{exambox}
\textbf{肝肾损害型毒蕈中毒的"假愈期"是重要考点：}胃肠炎症状暂时缓解后病情突然恶化，患者和家属误以为已康复，导致错过抢救时机。巯基解毒剂是关键抢救药物。
\end{exambox}

% =====================================================================
%  第五节：化学性食物中毒
% =====================================================================
\section{化学性食物中毒}

\subsection{亚硝酸盐食物中毒}

\begin{keybox}
\textbf{亚硝酸盐食物中毒（掌握）：}
\begin{itemize}[leftmargin=*]
    \item \textbf{中毒原因：}
    \begin{enumerate}[leftmargin=*]
        \item 误将亚硝酸盐当作食盐使用（最危险！建筑工地常见）
        \item 进食大量含硝酸盐/亚硝酸盐较高的蔬菜——新鲜腌制的蔬菜（"暴腌菜"），\keyword{腌制第4～8天亚硝酸盐含量达高峰}
        \item 肠源性青紫症（肠道功能紊乱时，肠道细菌将硝酸盐还原为亚硝酸盐）
        \item 某些井水含硝酸盐过高（苦井水）
    \end{enumerate}
    \item \textbf{中毒机制：}亚硝酸盐将\textbf{血红蛋白（Hb）中的Fe$^{2+}$氧化为Fe$^{3+}$}，形成\textbf{高铁血红蛋白（MetHb）}，失去携氧能力，导致组织缺氧，出现\textbf{肠源性青紫症（发绀）}。
    \item \textbf{临床特点：}\keyword{潜伏期10min～3h}。\textbf{口唇、指尖发绀}（肠源性青紫症），头晕、乏力、恶心、呕吐。严重者呼吸困难、昏迷、循环衰竭死亡。
    \item \textbf{特效解毒剂：}\keyword{美蓝（亚甲蓝，Methylene Blue）} + \textbf{维生素C}。小剂量美蓝（1～2 mg/kg）将MetHb还原为Hb；维生素C作为辅助还原剂。
    \item \textbf{安全食用腌菜：}\keyword{腌制20天以后亚硝酸盐基本消失}，方可安全食用。
\end{itemize}
\end{keybox}

\begin{tipbox}
\textbf{暴腌菜风险：}腌制第4～8天亚硝酸盐含量达峰值，此时食用中毒风险最大！腌制20天以上才安全。误将亚硝酸盐当食盐是化学性食物中毒的最常见原因。
\end{tipbox}

\subsection{有机磷农药中毒}

\begin{keybox}
\textbf{有机磷农药中毒（掌握）：}
\begin{itemize}[leftmargin=*]
    \item \textbf{中毒原因：}误食农药拌过的种子；误把有机磷农药当作食用油或酱油食用；食用喷洒农药后未经安全间隔期即采摘的蔬菜水果；误食被农药毒杀的家禽家畜。
    \item \textbf{中毒机制：}\keyword{抑制胆碱酯酶（ChE）活性}，使其失去催化水解乙酰胆碱（ACh）的能力 $\to$ 大量ACh在体内蓄积 $\to$ 胆碱能神经过度兴奋 $\to$ 出现毒蕈碱样症状、烟碱样症状和中枢神经系统症状。
    \item \textbf{临床特点：}\keyword{潜伏期30min～3h}。
    \begin{enumerate}[leftmargin=*]
        \item \textbf{毒蕈碱样症状（M样症状）：}\keyword{瞳孔缩小}、\keyword{流涎}、\keyword{大汗}、恶心呕吐、腹痛腹泻、支气管痉挛、呼吸困难。
        \item \textbf{烟碱样症状（N样症状）：}\keyword{肌束震颤}（从眼睑、面部小肌群开始，逐渐发展到全身）、肌肉无力甚至瘫痪。
        \item \textbf{中枢神经系统症状：}头晕、头痛、烦躁不安、抽搐、昏迷。严重者因呼吸中枢抑制和呼吸肌麻痹导致呼吸衰竭死亡。
    \end{enumerate}
    \item \textbf{全血胆碱酯酶活性分级：}
    \begin{itemize}[leftmargin=*]
        \item 轻度中毒：ChE活性降至正常值的50\%～70\%
        \item 中度中毒：ChE活性降至正常值的30\%～50\%
        \item 重度中毒：ChE活性降至正常值的30\%以下
    \end{itemize}
    \item \textbf{特效解毒剂：}\keyword{阿托品（Atropine） + 氯解磷定（Pralidoxime Chloride）/碘解磷定}。阿托品拮抗毒蕈碱样症状（对抗ACh蓄积），解磷定使被抑制的胆碱酯酶复活。
    \item \textbf{预防：}严格执行农药安全使用规定，遵守安全间隔期；加强农药保管；严禁食用被农药毒死的动物。
\end{itemize}
\end{keybox}

\begin{tipbox}
\textbf{有机磷中毒症状速记：}M样 = 毒蕈碱样 —— 流涎、大汗、瞳孔缩小（"水"的症状）；N样 = 烟碱样 —— 肌束震颤（肌肉症状）。特效解毒剂：阿托品（拮抗M样症状）+ 解磷定（复活ChE）。
\end{tipbox}

% =====================================================================
%  第六节：食物过敏
% =====================================================================
\section{食物过敏}

\begin{defbox}
\textbf{食物过敏（food allergy）的定义：}
机体免疫系统对进入体内的食物中的某些非感染性物质（过敏原）产生的异常免疫应答，导致生理功能紊乱和/或组织损伤。
\end{defbox}

\subsection{常见致敏食物}

\textbf{八大类主要致敏食物：}
\begin{enumerate}[leftmargin=*]
    \item 牛奶（Milk）
    \item 鸡蛋（Eggs）
    \item 花生（Peanuts）
    \item 大豆（Soy）
    \item 谷物（小麦等，Grains/Wheat）
    \item 鱼类（Fish）
    \item 甲壳类（Crustaceans）
    \item 坚果类（Tree Nuts）
\end{enumerate}

\subsection{食物过敏的特点}

\begin{enumerate}[leftmargin=*]
    \item 任何食物均可诱发过敏，但少数食物是主要过敏原
    \item 仅有部分食物成分具有变应原性，特定蛋白质为主要致敏成分
    \item 不同食物间可发生交叉过敏反应
    \item 食物加工可改变致敏性：加热可使部分过敏原变性降低致敏性，但某些热稳定型过敏原（如牛奶中$\beta$-乳球蛋白）不受加热影响
\end{enumerate}

\subsection{卫生假说（Hygiene Hypothesis）}

\begin{tipbox}
\textbf{卫生假说：}免疫系统在进化过程中为应对肠道寄生虫而设计，现代社会中缺乏肠道寄生虫感染，导致免疫系统过度活跃，对无害食物抗原产生异常免疫应答。这解释了工业化国家过敏性疾病高发的原因。
\end{tipbox}

% =====================================================================
%  复习题
% =====================================================================
\reviewcontent{
\section{复习题}

\subsection*{一、名词解释}

\begin{enumerate}[leftmargin=*, itemsep=6pt]
    \item 食源性疾病（foodborne disease）
    \item 食物过敏（food allergy）
    \item 食物中毒（food poisoning）
    \item 细菌性食物中毒
    \item 肠源性青紫症
    \item 高铁血红蛋白血症（Methemoglobinemia）
    \item 毒蕈中毒假愈期
\end{enumerate}

\subsection*{二、简答题}

\begin{enumerate}[leftmargin=*, itemsep=8pt]
    \item 简述食源性疾病和食物中毒的定义。
    \item 简述食物中毒的四大临床特点。
    \item 简述细菌性食物中毒的三大发病原因和五大特点。
    \item 比较沙门氏菌、金黄色葡萄球菌、副溶血性弧菌、肉毒梭菌四种食物中毒的来源、中毒机制、潜伏期和临床特点。
    \item 金黄色葡萄球菌肠毒素的耐热性如何？为什么带菌者管理对预防金葡菌食物中毒至关重要？
    \item 简述肉毒梭菌毒素的作用机制，并说明为什么肉毒梭菌中毒的病死率极高。
    \item 简述河豚毒素（TTX）的中毒机制和河豚鱼中毒的临床发展过程。
    \item 简述毒蕈中毒的临床分型，以及肝肾损害型的"假愈期"有何临床意义。
    \item 简述亚硝酸盐食物中毒的中毒机制、临床特点和特效解毒剂。
    \item 简述有机磷农药中毒的中毒机制和三大类临床表现。
    \item 简述有机磷农药中毒时全血胆碱酯酶活性的三级变化及相应的特效解毒剂。
    \item 简述食物过敏的八大类致敏食物及食物过敏的基本特点。
\end{enumerate}

\subsection*{三、案例分析}

\begin{enumerate}[leftmargin=*]
    \item \textbf{【案例分析1】}某工地食堂午餐后约2～4小时，50名工人陆续出现剧烈呕吐、上腹剧痛、水样腹泻，但体温均正常。调查发现午餐食用了凉拌熟肉。请回答：
    \begin{enumerate}[leftmargin=*]
        \item 最可能的中毒病原体是什么？
        \item 该病原体的中毒机制是什么？
        \item 肠毒素的耐热特点如何？
        \item 预防措施中最关键的一点是什么？
    \end{enumerate}

    \item \textbf{【案例分析2】}某家庭聚餐后次日，10人中8人出现视力模糊、复视、眼睑下垂、吞咽困难、发音不清，均无发热和明显胃肠道症状。调查发现食用了家庭自制的密封发酵豆酱。请回答：
    \begin{enumerate}[leftmargin=*]
        \item 最可能的诊断是什么？
        \item 该病原体产生的是何种毒素？其毒力如何？
        \item 应如何治疗？
        \item 如何开展预防？
    \end{enumerate}

    \item \textbf{【案例分析3】}某食堂午餐后1小时左右，多名就餐者出现口唇、指尖发绀，伴头晕、乏力、恶心。调查发现食用了腌制仅5天的暴腌菜和苦井水煮的饭菜。请回答：
    \begin{enumerate}[leftmargin=*]
        \item 最可能的中毒类型是什么？
        \item 其中毒机制是什么？
        \item 特效解毒剂是什么？
        \item 暴腌菜腌制多少天后亚硝酸盐基本消失？
    \end{enumerate}
\end{enumerate}

\subsection*{四、参考答案}

\subsubsection*{一、名词解释参考答案}
\begin{enumerate}[leftmargin=*, itemsep=4pt]
    \item \textbf{食源性疾病（foodborne disease）：}食品中致病因素进入人体引起的感染性、中毒性疾病。
    \item \textbf{食物过敏（food allergy）：}机体免疫系统对进入体内的食物中的某些非感染性物质（过敏原）产生的异常免疫应答，导致生理功能紊乱和/或组织损伤。
    \item \textbf{食物中毒（food poisoning）：}食用被有毒有害物质污染的食品或食用含有毒有害物质的食品后出现的急性、亚急性疾病。
    \item \textbf{细菌性食物中毒：}摄入被致病菌或其毒素污染的食品后发生的急性或亚急性疾病，是最常见的食物中毒。
    \item \textbf{肠源性青紫症：}亚硝酸盐将血红蛋白中的Fe$^{2+}$氧化为Fe$^{3+}$形成高铁血红蛋白，失去携氧能力，口唇、指尖出现发绀，因与肠道细菌将硝酸盐还原为亚硝酸盐有关，故称肠源性青紫症。
    \item \textbf{高铁血红蛋白血症（Methemoglobinemia）：}血红蛋白中的Fe$^{2+}$被氧化为Fe$^{3+}$形成高铁血红蛋白（MetHb），失去携氧能力导致的病理状态。
    \item \textbf{毒蕈中毒假愈期：}肝肾损害型毒蕈中毒病程中的特殊阶段，胃肠炎症状暂时缓解或消失，但肝脏损害已开始进行性发展，最容易被忽视而导致病情恶化。
\end{enumerate}

\subsubsection*{二、简答题参考答案}
\begin{enumerate}[leftmargin=*, itemsep=4pt]
    \item \textbf{食源性疾病和食物中毒的定义：}
    \begin{itemize}[leftmargin=*]
        \item 食源性疾病：食品中致病因素进入人体引起的感染性、中毒性疾病。
        \item 食物中毒：食用被有毒有害物质污染的食品或食用含有毒有害物质的食品后出现的急性、亚急性疾病。
    \end{itemize}

    \item \textbf{食物中毒的四大临床特点：}
    \begin{enumerate}[leftmargin=*]
        \item 发病与特定食物有关
        \item 潜伏期短，发病曲线单峰型
        \item 临床表现相似：最常见为消化道症状（腹痛、腹泻、恶心、呕吐）
        \item 一般没有人与人间的直接传染
    \end{enumerate}

    \item \textbf{细菌性食物中毒的发病原因和特点：}

    三大发病原因：①致病菌污染食品；②食品的贮存方式不当；③食用前未彻底加热处理。

    五大特点：①潜伏期较短；②多呈暴发性；③发病季节性明显（夏秋季5～10月高发）；④以胃肠道症状为主（腹痛、腹泻、恶心、呕吐）；⑤病死率较低（肉毒梭菌除外）。

    \item \textbf{四种细菌性食物中毒比较：}
    \begin{itemize}[leftmargin=*]
        \item \textbf{沙门氏菌：}来源为动物性食品（畜禽肉、蛋类、乳类）；中毒机制为活菌+内毒素（感染型）；潜伏期12～36h；临床特点为发热+水样便/黄绿色稀便，污染食品后无感官改变。
        \item \textbf{金黄色葡萄球菌：}来源为化脓性皮肤病或上呼吸道感染患者的污染；中毒机制为肠毒素（毒素型，耐热，100$^\circ$C 2h破坏）；潜伏期1～6h；临床特点为剧烈反复呕吐+上腹剧痛+水样腹泻，体温一般正常。
        \item \textbf{副溶血性弧菌：}来源为海产品（嗜盐菌）；中毒机制为活菌+溶血毒素（混合型）；潜伏期6～12h；临床特点为上腹阵发性绞痛+水样便或洗肉水样便，里急后重不明显。
        \item \textbf{肉毒梭菌：}来源为家庭自制发酵食品和罐头食品；中毒机制为外毒素（肉毒毒素，已知最剧烈的神经毒素）；潜伏期12～48h；临床特点为对称性颅神经损害（视力模糊、复视、眼睑下垂、吞咽困难、语言障碍、呼吸困难），病死率极高。
    \end{itemize}

    \item \textbf{金葡菌肠毒素耐热性及带菌者管理：}肠毒素耐热，100$^\circ$C加热2h才被破坏，普通烹调方法难以彻底破坏。带菌者（化脓性皮肤病、上呼吸道感染者）是食品污染的最重要来源，因此有局部化脓性感染或上呼吸道感染者应暂时调离食品工作岗位。

    \item \textbf{肉毒梭菌毒素作用机制及高病死率原因：}肉毒毒素作用于神经肌肉接头，阻止乙酰胆碱释放，导致肌肉麻痹。高病死率原因：(1)肉毒毒素是已知最剧烈的神经毒素，致死剂量极低（约10$^{-9}$mg/kg）；(2)主要致死原因为呼吸肌麻痹导致呼吸衰竭；(3)芽胞耐热性极强（湿热100$^\circ$C需5小时），普通加热无法杀灭。

    \item \textbf{河豚毒素（TTX）的中毒机制和临床发展过程：}中毒机制：TTX阻断Na$^+$通道，阻止神经细胞膜钠离子内流，使神经传导阻滞。临床发展过程：先感觉神经麻痹（口唇舌尖发麻）$\to$ 后运动神经麻痹（四肢无力、肌肉麻痹）$\to$ 最后呼吸麻痹死亡。

    \item \textbf{毒蕈中毒的临床分型及假愈期意义：}临床分型：胃肠炎型、神经精神型（阿托品治疗）、溶血型（肾上腺皮质激素）、肝肾损害型（巯基解毒剂，病死率最高）、类光过敏型。假愈期的意义：此期胃肠炎症状暂时缓解，但肝脏损害已开始进行性发展，需继续严密观察和积极治疗，不可因症状缓解而放松警惕。

    \item \textbf{亚硝酸盐食物中毒：}中毒机制：将Hb中的Fe$^{2+}$氧化为Fe$^{3+}$形成MetHb，失去携氧能力。临床特点：潜伏期10min～3h，口唇指尖发绀（肠源性青紫症），头晕乏力。特效解毒剂：美蓝（亚甲蓝）+ 维生素C。

    \item \textbf{有机磷农药中毒机制和临床表现：}中毒机制：抑制胆碱酯酶（ChE）活性，使ACh蓄积，胆碱能神经过度兴奋。三类临床表现：毒蕈碱样症状（瞳孔缩小、流涎、大汗）、烟碱样症状（肌束震颤）、中枢神经系统症状（头晕、昏迷）。

    \item \textbf{有机磷中毒ChE活性分级及特效解毒剂：}轻度：ChE 50\%～70\%；中度：ChE 30\%～50\%；重度：ChE $<$30\%。特效解毒剂：阿托品（拮抗毒蕈碱样症状）+ 氯解磷定/碘解磷定（复活胆碱酯酶）。

    \item \textbf{食物过敏的八大类致敏食物及特点：}八大类：牛奶、鸡蛋、花生、大豆、谷物（小麦等）、鱼类、甲壳类、坚果类。特点：任何食物均可诱发过敏；仅有部分成分具有变应原性；不同食物间可发生交叉过敏反应；食物加工可改变致敏性。
\end{enumerate}

\subsubsection*{三、案例分析参考答案}
\begin{enumerate}[leftmargin=*]
    \item \textbf{【案例分析1】}
    \begin{enumerate}[leftmargin=*]
        \item 最可能的病原体：金黄色葡萄球菌（肠毒素）。
        \item 中毒机制：毒素型——金葡菌在食品中预先产生肠毒素，毒素随食物进入体内引起中毒。
        \item 肠毒素耐热，100$^\circ$C加热2h才被破坏，普通烹调难以彻底破坏。
        \item 最关键预防措施：有局部化脓性感染或上呼吸道感染者应暂时调离食品工作岗位。
    \end{enumerate}

    \item \textbf{【案例分析2】}
    \begin{enumerate}[leftmargin=*]
        \item 最可能的诊断：肉毒梭菌毒素食物中毒。
        \item 肉毒毒素（外毒素），是已知最剧烈的神经毒素。毒性比KCN强10\,000倍，人的致死剂量约10$^{-9}$mg/kg。
        \item 尽早使用多价肉毒抗毒素血清（A、B、E型），同时对症支持治疗。
        \item 预防：家庭自制密封发酵食品为高危风险食品；罐头食品鼓盖/漏气绝对不要食用；蜂蜜不宜喂养1岁以下婴儿。
    \end{enumerate}

    \item \textbf{【案例分析3】}
    \begin{enumerate}[leftmargin=*]
        \item 最可能的诊断：亚硝酸盐食物中毒。
        \item 中毒机制：亚硝酸盐将Hb中的Fe$^{2+}$氧化为Fe$^{3+}$，形成MetHb，失去携氧能力，导致肠源性青紫症。
        \item 特效解毒剂：美蓝（亚甲蓝）+ 维生素C。
        \item 腌制20天以后亚硝酸盐基本消失。
    \end{enumerate}
\end{enumerate}

\newpage
}

% =====================================================================
%  CHAPTER 12: 食品安全性风险分析和控制
% =====================================================================
\chapter{食品安全性风险分析和控制}
\setcounter{chapter}{12}
\chapterintro{
\keyword{掌握}风险分析框架（危害识别、危害特征描述、暴露评估、风险特征描述）；掌握NOEL/NOAEL/ADI。\keyword{熟悉}营养毒理学。
}

% =====================================================================
%  第一节 食品安全性风险分析框架
% =====================================================================
\section{食品安全性风险分析框架}

\subsection{风险分析框架的四大步骤}

\begin{keybox}
\textbf{食品安全性风险分析框架（四大步骤*）：}

\begin{enumerate}[leftmargin=*]
    \item \textbf{危害识别（Hazard Identification）：}确定人体摄入某种食品中的危害因素与健康不良效应之间是否存在因果关系。信息来源：流行病学研究、动物毒理学试验、体外试验、构效关系分析。

    \item \textbf{危害特征描述（Hazard Characterization）：}对危害因素可能引起的不良健康效应进行定性和/或定量评价。核心任务：建立剂量-反应关系（Dose-Response Relationship），确定NOAEL或BMDL。

    \item \textbf{暴露评估（Exposure Assessment）：}对通过食品或其他途径摄入人体的生物、化学和物理因素的量进行评估。膳食暴露量 = $\sum$(食品中化学物浓度 $\times$ 食品消费量) / 体重。

    \item \textbf{风险特征描述（Risk Characterization）：}综合前三步的结果，分析和评估危害因素对人群健康产生不良效应的风险和严重程度。与健康指导值比较：暴露量 $<$ HBGV $\to$ 风险可接受。
\end{enumerate}
\end{keybox}

% =====================================================================
%  第二节 营养毒理学
% =====================================================================
\section{营养毒理学}

\begin{defbox}
\textbf{营养毒理学（Nutritional Toxicology）：}研究营养素过量摄入的毒副作用。UL（可耐受最高摄入量）是营养毒理学的重要指标。

\textbf{三大研究方向：}
\begin{enumerate}[leftmargin=*]
    \item 营养素过量摄入的不良反应及UL的制定。
    \item 营养素对外源化学物代谢和毒性的影响。
    \item 膳食中有毒物质对营养素代谢和营养过程的影响。
\end{enumerate}
\end{defbox}

% =====================================================================
%  第三节 安全性评价方法
% =====================================================================
\section{安全性评价方法}

\begin{defbox}
\textbf{安全性评价方法：}

\begin{enumerate}[leftmargin=*]
    \item \textbf{急性毒性试验（LD$_{50}$）：}了解受试物的急性毒性强度和性质，测定半数致死量。
    \item \textbf{30天喂养试验：}了解多次重复染毒条件下受试物的毒性。
    \item \textbf{90天喂养试验：}观察较长时间重复染毒的毒性效应，确定NOAEL。
    \item \textbf{传统致畸试验：}评价受试物是否具有致畸作用。
    \item \textbf{慢性毒性试验（包括致癌试验）：}观察长期（终生）染毒的慢性毒性和致癌性。
\end{enumerate}
\end{defbox}

% =====================================================================
%  第四节 关键毒理学指标
% =====================================================================
\section{关键毒理学指标}

\subsection{最大无作用剂量}

\begin{keybox}
\textbf{最大无作用剂量（MNL/NOEL/NOAEL）*：}在一定时间内，某种外源化学物按一定方式与机体接触，用现代检测方法和最灵敏的观察指标不能检出任何有害作用的最高剂量。

\begin{itemize}[leftmargin=*]
    \item \textbf{MNL（Maximum No-effect Level）：}最大无作用剂量。
    \item \textbf{NOEL（No Observed Effect Level）：}未观察到作用水平。
    \item \textbf{NOAEL（No Observed Adverse Effect Level）：}未观察到有害作用水平。
\end{itemize}
\end{keybox}

\subsection{每日允许摄入量}

\begin{keybox}
\textbf{每日允许摄入量（Acceptable Daily Intake, ADI）*：}人类每日摄入某物质直至终生，根据当时已知事实，不会产生可察觉健康危害的量。单位：mg/(kg bw$\cdot$d)。

\textbf{推导公式：}
\[
\text{ADI} = \frac{\text{NOAEL}}{\text{UFs}}
\]

其中UFs（Uncertainty Factors，不确定系数）常规取100倍（种间差异10 $\times$ 个体差异10）。
\end{keybox}

% =====================================================================
%  复习题
% =====================================================================
\reviewcontent{
\section{复习题}

\subsection*{一、名词解释}

\begin{enumerate}[leftmargin=*, itemsep=6pt]
    \item \textbf{危害识别（Hazard Identification）、危害特征描述（Hazard Characterization）}
    \item \textbf{暴露评估（Exposure Assessment）、风险特征描述（Risk Characterization）}
    \item \textbf{营养毒理学（Nutritional Toxicology）}
    \item \textbf{MNL、NOEL、NOAEL}
    \item \textbf{ADI（每日允许摄入量）}
    \item \textbf{UFs（不确定系数）}
\end{enumerate}

\subsection*{二、简答题}

\begin{enumerate}[leftmargin=*, itemsep=8pt]
    \item \textbf{【重点】}简述食品安全性风险分析框架的四大步骤及各步骤的主要任务。
    \item \textbf{【重点】}简述MNL/NOEL/NOAEL的定义及ADI的定义和推导公式。
    \item 简述营养毒理学的概念和主要研究内容。
    \item 简述食品安全性毒理学评价的主要方法。
\end{enumerate}

\subsection*{参考答案要点}

\subsubsection*{名词解释（要点）}

\begin{enumerate}[leftmargin=*, itemsep=4pt]
    \item \textbf{危害识别：}确定危害因素与健康不良效应之间是否存在因果关系。\textbf{危害特征描述：}对不良健康效应进行定性和/或定量评价，建立剂量-反应关系。
    \item \textbf{暴露评估：}对通过食品摄入人体的因素的量进行评估。\textbf{风险特征描述：}综合前三步，评估危害因素对人群健康产生不良效应的风险和严重程度。
    \item \textbf{营养毒理学：}研究营养素过量摄入的毒副作用。UL是营养毒理学的重要指标。
    \item \textbf{MNL/NOEL/NOAEL：}在一定时间内，某种外源化学物按一定方式与机体接触，用现代检测方法和最灵敏的观察指标不能检出任何有害作用的最高剂量。
    \item \textbf{ADI：}人类每日摄入某物质直至终生，根据当时已知事实，不会产生可察觉健康危害的量。ADI = NOAEL / UFs。
    \item \textbf{UFs：}不确定系数，常规取100倍（种间差异10 $\times$ 个体差异10）。
\end{enumerate}

\subsubsection*{简答题（要点）}

\begin{enumerate}[leftmargin=*, itemsep=6pt]
    \item \textbf{四大步骤：}①危害识别（确定因果关系）；②危害特征描述（定性和/或定量评价，建立剂量-反应关系）；③暴露评估（评估摄入量）；④风险特征描述（综合评估风险和严重程度）。
    \item \textbf{MNL/NOEL/NOAEL：}不能检出任何有害作用的最高剂量。\textbf{ADI：}人类每日摄入某物质直至终生不会产生可察觉健康危害的量。\textbf{公式：}ADI = NOAEL / UFs（常规100倍）。
    \item \textbf{营养毒理学：}研究营养素过量摄入的毒副作用。三大方向：①营养素过量不良反应及UL制定；②营养素对外源化学物代谢和毒性的影响；③膳食有毒物质对营养素代谢的影响。
    \item \textbf{方法：}急性毒性试验（LD$_{50}$）、30天喂养试验、90天喂养试验、传统致畸试验、慢性毒性试验（包括致癌试验）。
\end{enumerate}

\newpage
}

% =====================================================================
%  CHAPTER 13: 食品安全监督管理
% =====================================================================
\chapter{食品安全监督管理}
\setcounter{chapter}{13}
\chapterintro{
\keyword{掌握}GMP的概念；掌握HACCP的概念和7个基本原则。\keyword{熟悉}食品安全法。
}

% =====================================================================
%  第一节 GMP（良好生产规范）
% =====================================================================
\section{GMP（良好生产规范）}

\begin{keybox}
\textbf{GMP（Good Manufacturing Practice）：}良好生产规范，是保证食品具有高度安全性的良好生产管理体系。GMP是注重生产过程中产品质量与卫生安全的自主性管理制度。

\textbf{GMP的四大管理要素（4M）：}
\begin{enumerate}[leftmargin=*]
    \item \textbf{Man（人员）：}配备合适且能胜任的从业人员，健康检查合格，经过食品安全知识培训。
    \item \textbf{Material（物料）：}选用良好的原料，原料验收标准，供应商管理。
    \item \textbf{Machine（设备）：}采用标准规范的厂房设施和机械设备，厂房布局合理（人流物流分离），设备易于清洗消毒。
    \item \textbf{Method（方法）：}遵照既定的最优方法进行生产，标准操作规程（SOP），工艺参数控制。
\end{enumerate}

\textbf{GMP的基本精神：}
\begin{enumerate}[leftmargin=*]
    \item 降低生产过程中的人为错误。
    \item 防止食品在生产过程中遭到污染和品质劣变。
    \item 建立企业自主性质量保证体系。
\end{enumerate}
\end{keybox}

% =====================================================================
%  第二节 HACCP（危害分析与关键控制点）
% =====================================================================
\section{HACCP（危害分析与关键控制点）}

\subsection{HACCP概念}

\begin{keybox}
\textbf{HACCP（Hazard Analysis and Critical Control Point）：}危害分析与关键控制点。是"从农田到餐桌"全过程的食品安全预防体系。通过对食品生产全过程的危害分析，确定关键控制点并采取控制措施，将危害降低到可接受水平。

\textbf{HACCP建立的前提条件：}HACCP体系必须建立在GMP和SSOP的基础之上。
\end{keybox}

\subsection{HACCP的7个基本原则}

\begin{keybox}
\textbf{HACCP的7个基本原则（掌握*）：}

\begin{enumerate}[leftmargin=*]
    \item \textbf{进行危害分析（Hazard Analysis）：}对食品生产全过程进行系统的危害分析，识别所有潜在的生物性、化学性和物理性危害。对每个危害进行风险评估。

    \item \textbf{确定关键控制点（Critical Control Point, CCP）：}确定能够实施控制的、对防止和消除食品安全危害或将其降低到可接受水平所必需的某一步骤。CCP的三种控制形式：防止危害发生、消除危害、降低到可接受水平。

    \item \textbf{建立关键限值（Critical Limit, CL）：}在CCP上设定的最大值和/或最小值，用于将危害控制在一定范围内。如巴氏消毒最低72$^\circ$C、15秒。

    \item \textbf{建立CCP的监控系统（Monitoring）：}按照计划对CCP进行观察和测量，以判断CCP是否处于控制之中。监控四大要素：监控什么（What）、怎样监控（How）、监控频次（Frequency）、谁监控（Who）。

    \item \textbf{建立纠偏措施（Corrective Action）：}当监控发现CCP偏离关键限值时，立即采取的措施。包括隔离和评估受影响产品、纠正偏离原因、防止再次发生、记录全部纠偏活动。

    \item \textbf{建立验证程序（Verification）：}确认HACCP体系是否有效运行。验证的三个维度：适宜性、一致性、有效性。

    \item \textbf{建立记录保持程序（Record-keeping）：}所有HACCP相关活动必须有完整、准确的记录。文件包括：HACCP计划文本、危害分析表、CCP监控记录、纠偏记录、验证记录。
\end{enumerate}
\end{keybox}

% =====================================================================
%  第三节 食品安全法
% =====================================================================
\section{食品安全法}

\begin{defbox}
\textbf{食品安全法（2015年修订）：}确立了以食品安全风险监测和评估为基础的科学管理制度。2015年10月1日修订版正式实施，被称为"史上最严"食品安全法。

\textbf{四大原则：}预防为主、风险管理、全程控制、社会共治。

\textbf{核心制度：}食品安全风险监测制度、食品安全风险评估制度、食品安全标准制度、食品生产经营许可制度、食品召回制度、食品安全信息发布制度。
\end{defbox}

% =====================================================================
%  复习题
% =====================================================================
\reviewcontent{
\section{复习题}

\subsection*{一、名词解释}

\begin{enumerate}[leftmargin=*, itemsep=6pt]
    \item \textbf{GMP（Good Manufacturing Practice）}
    \item \textbf{HACCP（Hazard Analysis and Critical Control Point）}
    \item \textbf{CCP（关键控制点）}
    \item \textbf{CL（关键限值）}
\end{enumerate}

\subsection*{二、简答题}

\begin{enumerate}[leftmargin=*, itemsep=8pt]
    \item \textbf{【重点】}简述GMP的概念和四大管理要素（4M）。
    \item \textbf{【重点】}简述HACCP的定义及7个基本原则。
    \item \textbf{【重点】}简述HACCP的7个基本原则中每个原则的主要内容。
    \item 简述食品安全法（2015年修订）的基本特征和四大原则。
\end{enumerate}

\subsection*{参考答案要点}

\subsubsection*{名词解释（要点）}

\begin{enumerate}[leftmargin=*, itemsep=4pt]
    \item \textbf{GMP：}良好生产规范，是保证食品具有高度安全性的良好生产管理体系。
    \item \textbf{HACCP：}危害分析与关键控制点，是"从农田到餐桌"全过程的食品安全预防体系。
    \item \textbf{CCP：}关键控制点，能够实施控制的、对防止和消除食品安全危害或将其降低到可接受水平所必需的某一步骤。
    \item \textbf{CL：}关键限值，在CCP上设定的最大值和/或最小值，用于将危害控制在一定范围内。
\end{enumerate}

\subsubsection*{简答题（要点）}

\begin{enumerate}[leftmargin=*, itemsep=6pt]
    \item \textbf{GMP：}良好生产规范，保证食品具有高度安全性的良好生产管理体系。\textbf{4M：}Man（人员——合适且胜任的从业人员）、Material（物料——选用良好的原料）、Machine（设备——标准规范的厂房设施和机械设备）、Method（方法——遵照既定的最优方法进行生产）。
    \item \textbf{HACCP定义：}危害分析与关键控制点，"从农田到餐桌"全过程的食品安全预防体系。\textbf{7原则：}①进行危害分析；②确定关键控制点（CCP）；③建立关键限值；④建立CCP的监控系统；⑤建立纠偏措施；⑥建立验证程序；⑦建立记录保持程序。
    \item \textbf{7原则内容：}①危害分析——识别生物、化学、物理性危害；②确定CCP——防止危害发生、消除危害、降低到可接受水平；③建立CL——最大值和/或最小值；④监控——What/How/Frequency/Who；⑤纠偏——立即隔离评估、纠正原因、防止再发、记录；⑥验证——适宜性、一致性、有效性；⑦记录保持——完整准确的文档记录。
    \item \textbf{食品安全法（2015年修订）：}确立以食品安全风险监测和评估为基础的科学管理制度。\textbf{四大原则：}预防为主、风险管理、全程控制、社会共治。
\end{enumerate}

\newpage
}

% =====================================================================
%  CHAPTER 99: 英文缩写与专有名词汇总
% =====================================================================
\chapter{英文缩写与专有名词汇总}
\chapterintro{
本附录收录《营养与食品卫生学》课程中出现的全部英文缩写及专有名词，按类别分为九大类，以表格形式呈现，方便对照复习和查阅。
}

% ===== 一、国际组织与机构 =====
\section{一、国际组织与机构}

\begin{table}[H]
\centering
\begin{tabularx}{\textwidth}{lX}
\toprule
\textbf{缩写} & \textbf{全称与中文释义} \\
\midrule
WHO & World Health Organization \hfill 世界卫生组织 \\
FAO & Food and Agriculture Organization \hfill 联合国粮食及农业组织 \\
CAC & Codex Alimentarius Commission \hfill 国际食品法典委员会 \\
JECFA & Joint FAO/WHO Expert Committee on Food Additives \hfill FAO/WHO食品添加剂联合专家委员会 \\
IARC & International Agency for Research on Cancer \hfill 国际癌症研究中心 \\
EFSA & European Food Safety Authority \hfill 欧洲食品安全局 \\
WCRF & World Cancer Research Fund \hfill 世界癌症研究基金会 \\
AICR & American Institute for Cancer Research \hfill 美国癌症研究所 \\
WTO & World Trade Organization \hfill 世界贸易组织 \\
ISO & International Organization for Standardization \hfill 国际标准化组织 \\
\bottomrule
\end{tabularx}
\end{table}

% ===== 二、营养学基础 =====
\section{二、营养学基础}

\subsection*{（一）氨基酸与蛋白质}

\begin{table}[H]
\centering
\begin{tabularx}{\textwidth}{lX}
\toprule
\textbf{缩写} & \textbf{全称与中文释义} \\
\midrule
EAA & Essential Amino Acid \hfill 必需氨基酸 \\
Ile & Isoleucine \hfill 异亮氨酸 \\
Leu & Leucine \hfill 亮氨酸 \\
Lys & Lysine \hfill 赖氨酸 \\
Met & Methionine \hfill 蛋氨酸（甲硫氨酸） \\
Phe & Phenylalanine \hfill 苯丙氨酸 \\
Thr & Threonine \hfill 苏氨酸 \\
Trp & Tryptophan \hfill 色氨酸 \\
Val & Valine \hfill 缬氨酸 \\
His & Histidine \hfill 组氨酸（婴儿必需氨基酸） \\
Cys & Cysteine \hfill 半胱氨酸（条件必需氨基酸） \\
Tyr & Tyrosine \hfill 酪氨酸（条件必需氨基酸） \\
ONL & Obligatory Nitrogen Loss \hfill 必要氮损失 \\
PEM & Protein-Energy Malnutrition \hfill 蛋白质-能量营养不良 \\
AD & Apparent Digestibility \hfill 表观消化率 \\
TD & True Digestibility \hfill 真消化率 \\
BV & Biological Value \hfill 生物价 \\
NPU & Net Protein Utilization \hfill 蛋白质净利用率 \\
PER & Protein Efficiency Ratio \hfill 蛋白质功效比值 \\
AAS & Amino Acid Score \hfill 氨基酸评分 \\
PDCAAS & Protein Digestibility Corrected Amino Acid Score \hfill 经消化率修正的氨基酸评分 \\
DIAAS & Digestible Indispensable Amino Acid Score \hfill 可消化必需氨基酸评分 \\
LAA & Limiting Amino Acid \hfill 限制氨基酸 \\
\bottomrule
\end{tabularx}
\end{table}

\subsection*{（二）脂类}

\begin{table}[H]
\centering
\begin{tabularx}{\textwidth}{lX}
\toprule
\textbf{缩写} & \textbf{全称与中文释义} \\
\midrule
SFA & Saturated Fatty Acid \hfill 饱和脂肪酸 \\
MUFA & Monounsaturated Fatty Acid \hfill 单不饱和脂肪酸 \\
PUFA & Polyunsaturated Fatty Acid \hfill 多不饱和脂肪酸 \\
EFA & Essential Fatty Acid \hfill 必需脂肪酸 \\
DHA & Docosahexaenoic Acid \hfill 二十二碳六烯酸 \\
EPA & Eicosapentaenoic Acid \hfill 二十碳五烯酸 \\
AA & Arachidonic Acid \hfill 花生四烯酸 \\
ARA & Arachidonic Acid \hfill 花生四烯酸（同AA） \\
LA & Linoleic Acid \hfill 亚油酸 \\
ALA & $\alpha$-Linolenic Acid \hfill $\alpha$-亚麻酸 \\
MCFA & Medium Chain Fatty Acid \hfill 中链脂肪酸 \\
SCFA & Short Chain Fatty Acid \hfill 短链脂肪酸 \\
TFA & Trans Fatty Acid \hfill 反式脂肪酸 \\
LCT & Long Chain Triglyceride \hfill 长链甘油三酯 \\
MCT & Medium Chain Triglyceride \hfill 中链甘油三酯 \\
\bottomrule
\end{tabularx}
\end{table}

\subsection*{（三）碳水合物与能量}

\begin{table}[H]
\centering
\begin{tabularx}{\textwidth}{lX}
\toprule
\textbf{缩写} & \textbf{全称与中文释义} \\
\midrule
GI & Glycemic Index \hfill 血糖生成指数 \\
GL & Glycemic Load \hfill 血糖负荷 \\
BEE & Basal Energy Expenditure \hfill 基础代谢能量消耗 \\
BMR & Basal Metabolic Rate \hfill 基础代谢率 \\
PAL & Physical Activity Level \hfill 体力活动水平 \\
TEF & Thermic Effect of Food \hfill 食物热效应 \\
SDA & Specific Dynamic Action \hfill 食物特殊动力作用（即TEF） \\
RQ & Respiratory Quotient \hfill 呼吸商 \\
REE & Resting Energy Expenditure \hfill 静息能量消耗 \\
\bottomrule
\end{tabularx}
\end{table}

\subsection*{（四）维生素与辅酶}

\begin{table}[H]
\centering
\begin{tabularx}{\textwidth}{lX}
\toprule
\textbf{缩写} & \textbf{全称与中文释义} \\
\midrule
RE & Retinol Equivalent \hfill 视黄醇当量 \\
RAE & Retinol Activity Equivalent \hfill 视黄醇活性当量 \\
TPP & Thiamine Pyrophosphate \hfill 硫胺素焦磷酸酯（维生素B1活性形式） \\
TDP & Thiamine Diphosphate \hfill 硫胺素二磷酸（同TPP） \\
FMN & Flavin Mononucleotide \hfill 黄素单核苷酸 \\
FAD & Flavin Adenine Dinucleotide \hfill 黄素腺嘌呤二核苷酸 \\
NAD$^+$ & Nicotinamide Adenine Dinucleotide \hfill 烟酰胺腺嘌呤二核苷酸（辅酶I） \\
NADP$^+$ & Nicotinamide Adenine Dinucleotide Phosphate \hfill 烟酰胺腺嘌呤二核苷酸磷酸（辅酶II） \\
THF & Tetrahydrofolate \hfill 四氢叶酸 \\
CoA & Coenzyme A \hfill 辅酶A \\
ACP & Acyl Carrier Protein \hfill 酰基载体蛋白 \\
NE & Niacin Equivalent \hfill 烟酸当量 \\
DFE & Dietary Folate Equivalent \hfill 膳食叶酸当量 \\
$\alpha$-TE & $\alpha$-Tocopherol Equivalent \hfill $\alpha$-生育酚当量 \\
PLP & Pyridoxal Phosphate \hfill 磷酸吡哆醛（维生素B6活性形式） \\
MTHF & 5-Methyltetrahydrofolate \hfill 5-甲基四氢叶酸 \\
\bottomrule
\end{tabularx}
\end{table}

\subsection*{（五）矿物质}

\begin{table}[H]
\centering
\begin{tabularx}{\textwidth}{lX}
\toprule
\textbf{缩写} & \textbf{全称与中文释义} \\
\midrule
Ca & Calcium \hfill 钙 \\
Fe & Iron \hfill 铁 \\
Zn & Zinc \hfill 锌 \\
Se & Selenium \hfill 硒 \\
I & Iodine \hfill 碘 \\
Mg & Magnesium \hfill 镁 \\
Cu & Copper \hfill 铜 \\
Cr & Chromium \hfill 铬 \\
Mn & Manganese \hfill 锰 \\
F & Fluorine \hfill 氟 \\
P & Phosphorus \hfill 磷 \\
K & Potassium \hfill 钾 \\
Na & Sodium \hfill 钠 \\
Cl & Chlorine \hfill 氯 \\
\bottomrule
\end{tabularx}
\end{table}

% ===== 三、公共营养 =====
\section{三、公共营养}

\begin{table}[H]
\centering
\begin{tabularx}{\textwidth}{lX}
\toprule
\textbf{缩写} & \textbf{全称与中文释义} \\
\midrule
DRIs & Dietary Reference Intakes \hfill 膳食营养素参考摄入量 \\
EAR & Estimated Average Requirement \hfill 平均需要量 \\
RNI & Recommended Nutrient Intake \hfill 推荐摄入量 \\
AI & Adequate Intake \hfill 适宜摄入量 \\
UL & Tolerable Upper Intake Level \hfill 可耐受最高摄入量 \\
AMDR & Acceptable Macronutrient Distribution Range \hfill 宏量营养素可接受范围 \\
PI-NCD & Proposed Intake for Non-communicable Chronic Disease \hfill 预防非传染病慢性病的建议摄入量 \\
SPL & Specific Proposed Level \hfill 特定建议值 \\
BMI & Body Mass Index \hfill 体质指数 \\
NRV & Nutrient Reference Value \hfill 营养素参考值 \\
FFQ & Food Frequency Questionnaire \hfill 食物频率法 \\
NRS-2002 & Nutrition Risk Screening 2002 \hfill 营养风险筛查2002 \\
SGA & Subjective Global Assessment \hfill 主观全面评定 \\
MNA & Mini Nutritional Assessment \hfill 微型营养评定 \\
MUST & Malnutrition Universal Screening Tool \hfill 营养不良通用筛查工具 \\
DDP & Desirable Dietary Pattern \hfill 理想膳食模式 \\
DGE & German Nutrition Society \hfill 德国营养学会（DDP模型来源） \\
\bottomrule
\end{tabularx}
\end{table}

% ===== 四、临床营养 =====
\section{四、临床营养}

\begin{table}[H]
\centering
\begin{tabularx}{\textwidth}{lX}
\toprule
\textbf{缩写} & \textbf{全称与中文释义} \\
\midrule
EN & Enteral Nutrition \hfill 肠内营养 \\
PN & Parenteral Nutrition \hfill 肠外营养 \\
TPN & Total Parenteral Nutrition \hfill 全肠外营养 \\
PPN & Peripheral Parenteral Nutrition \hfill 周围静脉营养 \\
FSMP & Food for Special Medical Purposes \hfill 特殊医学用途配方食品 \\
DASH & Dietary Approaches to Stop Hypertension \hfill 终止高血压膳食疗法 \\
GDM & Gestational Diabetes Mellitus \hfill 妊娠期糖尿病 \\
DOHaD & Developmental Origins of Health and Disease \hfill 健康和疾病的发育起源（都哈理论） \\
T2DM & Type 2 Diabetes Mellitus \hfill 2型糖尿病 \\
CVD & Cardiovascular Disease \hfill 心血管疾病 \\
CHD & Coronary Heart Disease \hfill 冠心病 \\
NCD & Non-communicable Chronic Disease \hfill 非传染性慢性病 \\
\bottomrule
\end{tabularx}
\end{table}

% ===== 五、食品污染 =====
\section{五、食品污染}

\subsection*{（一）微生物与腐败}

\begin{table}[H]
\centering
\begin{tabularx}{\textwidth}{lX}
\toprule
\textbf{缩写} & \textbf{全称与中文释义} \\
\midrule
CFU & Colony Forming Unit \hfill 菌落形成单位 \\
MPN & Most Probable Number \hfill 大肠菌群最大可能数 \\
Aw & Water Activity \hfill 水分活性 \\
TVBN & Total Volatile Basic Nitrogen \hfill 挥发性盐基总氮 \\
TMA & Trimethylamine \hfill 三甲胺 \\
D值 & Decimal Reduction Time \hfill 十进制减菌时间（90\%杀灭时间） \\
F值 & Sterilization Time \hfill 杀菌值 \\
Z值 & Temperature Change for 10-fold D-value Change \hfill Z值 \\
UHT & Ultra High Temperature \hfill 超高温消毒法 \\
LTLT & Low Temperature Long Time \hfill 低温长时间巴氏消毒 \\
HTST & High Temperature Short Time \hfill 高温短时间巴氏消毒 \\
TTT & Time-Temperature-Tolerance \hfill 时间-温度-耐受原则 \\
MAP & Modified Atmosphere Packaging \hfill 气调包装 \\
\bottomrule
\end{tabularx}
\end{table}

\subsection*{（二）化学污染物与农药}

\begin{table}[H]
\centering
\begin{tabularx}{\textwidth}{lX}
\toprule
\textbf{缩写} & \textbf{全称与中文释义} \\
\midrule
MRLs & Maximum Residue Limits \hfill 最大残留限量 \\
ADI & Acceptable Daily Intake \hfill 每日允许摄入量 \\
BHA & Butylated Hydroxyanisole \hfill 丁基羟基茴香醚（抗氧化剂） \\
BHT & Butylated Hydroxytoluene \hfill 二丁基羟基甲苯（抗氧化剂） \\
PG & Propyl Gallate \hfill 没食子酸丙酯（抗氧化剂） \\
TBHQ & Tertiary Butylhydroquinone \hfill 特丁基对苯二酚（抗氧化剂） \\
DDT & Dichlorodiphenyltrichloroethane \hfill 二氯二苯三氯乙烷（有机氯农药） \\
BHC/666 & Benzene Hexachloride \hfill 六六六（有机氯农药） \\
PAH & Polycyclic Aromatic Hydrocarbon \hfill 多环芳烃 \\
B(a)P & Benzo(a)pyrene \hfill 苯并(a)芘 \\
NDMA & N-Nitrosodimethylamine \hfill 二甲基亚硝胺 \\
NDEA & N-Nitrosodiethylamine \hfill 二乙基亚硝胺 \\
3-MCPD & 3-Monochloropropane-1,2-diol \hfill 3-氯-1,2-丙二醇 \\
PCB & Polychlorinated Biphenyl \hfill 多氯联苯 \\
POPs & Persistent Organic Pollutants \hfill 持久性有机污染物 \\
\bottomrule
\end{tabularx}
\end{table}

\subsection*{（三）真菌毒素}

\begin{table}[H]
\centering
\begin{tabularx}{\textwidth}{lX}
\toprule
\textbf{缩写} & \textbf{全称与中文释义} \\
\midrule
AF & Aflatoxin \hfill 黄曲霉毒素 \\
AFB1 & Aflatoxin B1 \hfill 黄曲霉毒素B1（毒性和致癌性最强） \\
AFM1 & Aflatoxin M1 \hfill 黄曲霉毒素M1（乳中代谢产物） \\
OTA & Ochratoxin A \hfill 赭曲霉毒素A \\
DON & Deoxynivalenol \hfill 脱氧雪腐镰刀菌烯醇（呕吐毒素） \\
3-NPA & 3-Nitropropionic Acid \hfill 3-硝基丙酸（霉变甘蔗毒素） \\
ZEN/ZEA & Zearalenone \hfill 玉米赤霉烯酮 \\
T-2 & T-2 Toxin \hfill T-2毒素 \\
ATA & Alimentary Toxic Aleukia \hfill 食物中毒性白细胞缺乏症（T-2引起） \\
NIV & Nivalenol \hfill 雪腐镰刀菌烯醇 \\
\bottomrule
\end{tabularx}
\end{table}

\subsection*{（四）ATP分解产物}

\begin{table}[H]
\centering
\begin{tabularx}{\textwidth}{lX}
\toprule
\textbf{缩写} & \textbf{全称与中文释义} \\
\midrule
ATP & Adenosine Triphosphate \hfill 腺苷三磷酸 \\
ADP & Adenosine Diphosphate \hfill 腺苷二磷酸 \\
AMP & Adenosine Monophosphate \hfill 腺苷一磷酸 \\
IMP & Inosine Monophosphate \hfill 肌苷一磷酸（鲜味物质） \\
HxR & Inosine \hfill 肌苷/次黄嘌呤核苷 \\
Hx & Hypoxanthine \hfill 次黄嘌呤 \\
K值 & Freshness Index K value \hfill 鱼肉鲜度指标K值 \\
\bottomrule
\end{tabularx}
\end{table}

% ===== 六、食品安全风险分析与监管 =====
\section{六、食品安全风险分析与监管}

\begin{table}[H]
\centering
\begin{tabularx}{\textwidth}{lX}
\toprule
\textbf{缩写} & \textbf{全称与中文释义} \\
\midrule
HBGV & Health-Based Guidance Value \hfill 健康指导值 \\
POD & Point of Departure \hfill 出发点 \\
NOAEL & No Observed Adverse Effect Level \hfill 未观察到有害作用水平 \\
LOAEL & Lowest Observed Adverse Effect Level \hfill 观察到有害作用最低水平 \\
BMDL & Benchmark Dose Lower Confidence Limit \hfill 基准剂量下限 \\
UFs & Uncertainty Factors \hfill 不确定系数 \\
ARfD & Acute Reference Dose \hfill 急性参考剂量 \\
TDI & Tolerable Daily Intake \hfill 日耐受摄入量 \\
PTWI & Provisional Tolerable Weekly Intake \hfill 暂定每周耐受摄入量 \\
PTMI & Provisional Tolerable Monthly Intake \hfill 暂定每月耐受摄入量 \\
MOE & Margin of Exposure \hfill 暴露边界 \\
ALARA & As Low As Reasonably Achievable \hfill 合理可行的最低水平 \\
GMP & Good Manufacturing Practice \hfill 食品良好生产规范 \\
SSOP & Sanitation Standard Operating Procedure \hfill 卫生标准操作程序 \\
HACCP & Hazard Analysis and Critical Control Point \hfill 危害分析及关键控制点 \\
CCP & Critical Control Point \hfill 关键控制点 \\
CL & Critical Limit \hfill 关键限值 \\
OL & Operating Limit \hfill 操作限值 \\
HA & Hazard Analysis \hfill 危害分析 \\
CA & Corrective Action \hfill 纠偏程序 \\
V & Verification \hfill 验证程序 \\
R & Record-keeping \hfill 记录保存 \\
GRAS & Generally Recognized As Safe \hfill 公认安全 \\
SAMR & State Administration for Market Regulation \hfill 国家市场监督管理总局 \\
CFDA & China Food and Drug Administration \hfill 国家食品药品监督管理总局（前身） \\
GHP & Good Hygiene Practice \hfill 良好卫生规范 \\
TDS & Total Diet Study \hfill 总膳食研究 \\
\bottomrule
\end{tabularx}
\end{table}

% ===== 七、食源性疾病 =====
\section{七、食源性疾病}

\begin{table}[H]
\centering
\begin{tabularx}{\textwidth}{lX}
\toprule
\textbf{缩写} & \textbf{全称与中文释义} \\
\midrule
BSE & Bovine Spongiform Encephalopathy \hfill 疯牛病 \\
nvCJD & new variant Creutzfeldt-Jakob Disease \hfill 新型变异型克雅氏病 \\
PrP$^\text{C}$ & Cellular Prion Protein \hfill 细胞型朊蛋白（正常型） \\
PrP$^\text{Sc}$ & Scrapie Prion Protein \hfill 致病型朊病毒蛋白 \\
TTX & Tetrodotoxin \hfill 河豚毒素 \\
HCN & Hydrogen Cyanide \hfill 氢氰酸 \\
MNL & Maximum No-effect Level \hfill 最大无作用剂量 \\
3-NPA & 3-Nitropropionic Acid \hfill 3-硝基丙酸（霉变甘蔗） \\
FMD & Foot and Mouth Disease \hfill 口蹄疫 \\
HUS & Hemolytic Uremic Syndrome \hfill 溶血性尿毒综合征（E. coli O157引起） \\
ETEC & Enterotoxigenic E. coli \hfill 肠产毒性大肠杆菌 \\
EHEC & Enterohemorrhagic E. coli \hfill 肠出血性大肠杆菌 \\
VTEC & Vero-toxin producing E. coli \hfill 产Vero毒素大肠杆菌 \\
STEC & Shiga-toxin producing E. coli \hfill 产志贺毒素大肠杆菌 \\
\bottomrule
\end{tabularx}
\end{table}

% ===== 八、其他食品营养相关概念 =====
\section{八、食品营养相关其他概念}

\begin{table}[H]
\centering
\begin{tabularx}{\textwidth}{lX}
\toprule
\textbf{缩写} & \textbf{全称与中文释义} \\
\midrule
GMO & Genetically Modified Organism \hfill 转基因生物/转基因食品 \\
GSH-Px & Glutathione Peroxidase \hfill 谷胱甘肽过氧化物酶（含硒酶） \\
GTF & Glucose Tolerance Factor \hfill 葡萄糖耐量因子（铬的活性形式） \\
SOD & Superoxide Dismutase \hfill 超氧化物歧化酶 \\
IGF-1 & Insulin-like Growth Factor-1 \hfill 胰岛素样生长因子-1 \\
TNF-$\alpha$ & Tumor Necrosis Factor-$\alpha$ \hfill 肿瘤坏死因子-$\alpha$ \\
IL-6 & Interleukin-6 \hfill 白细胞介素-6 \\
CRP & C-Reactive Protein \hfill C反应蛋白 \\
HCG & Human Chorionic Gonadotropin \hfill 人绒毛膜促性腺激素 \\
HCS & Human Chorionic Somatomammotropin \hfill 人绒毛膜生长催乳素 \\
TSH & Thyroid Stimulating Hormone \hfill 促甲状腺激素 \\
T3 & Triiodothyronine \hfill 三碘甲状腺原氨酸 \\
T4 & Thyroxine \hfill 甲状腺素 \\
Hb & Hemoglobin \hfill 血红蛋白 \\
SF & Serum Ferritin \hfill 血清铁蛋白 \\
sTfR & Soluble Transferrin Receptor \hfill 血清运铁蛋白受体 \\
FEP & Free Erythrocyte Protoporphyrin \hfill 红细胞游离原卟啉 \\
TIBC & Total Iron Binding Capacity \hfill 总铁结合力 \\
LDL & Low Density Lipoprotein \hfill 低密度脂蛋白 \\
HDL & High Density Lipoprotein \hfill 高密度脂蛋白 \\
VLDL & Very Low Density Lipoprotein \hfill 极低密度脂蛋白 \\
IDL & Intermediate Density Lipoprotein \hfill 中间密度脂蛋白 \\
TC & Total Cholesterol \hfill 总胆固醇 \\
TG & Triglycerides \hfill 甘油三酯 \\
LDL-C & LDL Cholesterol \hfill 低密度脂蛋白胆固醇 \\
HDL-C & HDL Cholesterol \hfill 高密度脂蛋白胆固醇 \\
FFA & Free Fatty Acid \hfill 游离脂肪酸 \\
BP & Blood Pressure \hfill 血压 \\
SBP & Systolic Blood Pressure \hfill 收缩压 \\
DBP & Diastolic Blood Pressure \hfill 舒张压 \\
OGTT & Oral Glucose Tolerance Test \hfill 口服葡萄糖耐量试验 \\
HbA1c & Glycated Hemoglobin \hfill 糖化血红蛋白 \\
RDA & Recommended Dietary Allowance \hfill 推荐膳食供给量（旧DRIs体系） \\
RNI & Recommended Nutrient Intake \hfill 推荐摄入量（新DRIs体系） \\
\bottomrule
\end{tabularx}
\end{table}

% ===== 九、氨基酸三字母/单字母对照 =====
\section{九、常见氨基酸三字母与单字母缩写对照}

\begin{table}[H]
\centering
\begin{tabularx}{\textwidth}{llll}
\toprule
\textbf{中文名} & \textbf{英文名} & \textbf{三字母} & \textbf{单字母} \\
\midrule
甘氨酸 & Glycine & Gly & G \\
丙氨酸 & Alanine & Ala & A \\
缬氨酸* & Valine & Val & V \\
亮氨酸* & Leucine & Leu & L \\
异亮氨酸* & Isoleucine & Ile & I \\
苯丙氨酸* & Phenylalanine & Phe & F \\
色氨酸* & Tryptophan & Trp & W \\
蛋氨酸* & Methionine & Met & M \\
脯氨酸 & Proline & Pro & P \\
丝氨酸 & Serine & Ser & S \\
苏氨酸* & Threonine & Thr & T \\
半胱氨酸 & Cysteine & Cys & C \\
酪氨酸 & Tyrosine & Tyr & Y \\
天冬酰胺 & Asparagine & Asn & N \\
谷氨酰胺 & Glutamine & Gln & Q \\
天冬氨酸 & Aspartic acid & Asp & D \\
谷氨酸 & Glutamic acid & Glu & E \\
赖氨酸* & Lysine & Lys & K \\
精氨酸 & Arginine & Arg & R \\
组氨酸（婴儿*） & Histidine & His & H \\
\bottomrule
\multicolumn{4}{l}{\small * 标注为必需氨基酸（成人8种+婴儿组氨酸）} \\
\end{tabularx}
\end{table}

% ===== 十、常用希腊字母与符号 =====
\section{十、常用希腊字母与符号}

\begin{table}[H]
\centering
\begin{tabularx}{\textwidth}{llX}
\toprule
\textbf{符号} & \textbf{名称} & \textbf{常见用途} \\
\midrule
$\alpha$ & Alpha & $\alpha$-亚麻酸、$\alpha$-生育酚、$\alpha$-螺旋 \\
$\beta$ & Beta & $\beta$-胡萝卜素、$\beta$-折叠、$\beta$-氧化 \\
$\gamma$ & Gamma & $\gamma$-生育酚、$\gamma$-亚麻酸 \\
$\delta$ & Delta & $\delta$-生育酚 \\
$\Delta$ & Delta（大写） & 表示双键位置（如$\Delta$9-desaturase） \\
$\omega$ & Omega & $\omega$-3 / $\omega$-6脂肪酸 \\
$\mu$ & Mu & 微克（$\mu$g） \\
\bottomrule
\end{tabularx}
\end{table}

\newpage


\end{document}
